\documentclass[12pt]{report}

% ---------------------------------------------------------------
% 字型與中文排版
% ---------------------------------------------------------------
% 作者在 macOS 上用 Times New Roman／Songti TC／PingFang TC；
% 其他系統（例如 Linux）找不到這些字型時，改用 TeX Gyre Termes 與
% Noto Serif／Sans CJK TC，同一份原始碼兩邊都能編。
\usepackage{fontspec}
\usepackage{xeCJK}
\IfFontExistsTF{Times New Roman}
    {\setmainfont{Times New Roman}}
    {\setmainfont{TeX Gyre Termes}}
\IfFontExistsTF{Songti TC}
    {\setCJKmainfont{Songti TC}}
    {\setCJKmainfont{Noto Serif CJK TC}}
\IfFontExistsTF{PingFang TC}
    {\setCJKsansfont{PingFang TC}}
    {\setCJKsansfont{Noto Sans CJK TC}}
% 等寬字用 Latin Modern Mono。它不是系統字型，只在 TeX Live 樹裡；
% 用「檔名＋副檔名」讓 XeTeX 透過 kpathsea 找檔案，不寫死 TeX Live
% 的年份路徑。
\setmonofont{lmmono10-regular}[
    Extension = .otf,
    AutoFakeBold = 2.5,
]

% ---------------------------------------------------------------
% 本書排版套件（packages/ 目錄）
% ---------------------------------------------------------------
\usepackage{sleek}
% sleek-theorems.sty 不載入：本書沒有用到定理類環境，而它依賴的 thmtools 和 2026-11 以後的
% LaTeX 核心衝突（\newtheorem{lem}[thm] 會報「\c@lem already defined」），載入的話整本書編不過。
\usepackage{sleek-code}
\usepackage{sleek-diagrams}
\usepackage{sleek-recap}
% 注意：sleek-listings.sty 不載入——它是舊版 lstlistings/inconsolata 排版套件，
% 本書實際用的是 sleek-code.sty 的 minted 版 code/example 環境，兩者同時載入會讓
% minted 抓不到 \setmonofont 設定的等寬字（回退成 Latin Modern Roman），
% 且會讓程式碼裡的直引號被轉成全形彎引號。

% 章節標題中文化：「第 N 章」
\usepackage{titlesec}
\titleformat{\chapter}{\normalfont\huge\bfseries}{第~\thechapter~章}{1em}{}
\titlespacing*{\chapter}{0pt}{-20pt}{20pt}
% 頁首與目錄標題中文化（預設是 CHAPTER 1. 和 CONTENTS）
\renewcommand{\contentsname}{目錄}
% 附錄的頁首寫「附錄 A」：附錄開始時 \@chapapp 會被換掉，和這裡存下的原值比較就知道。
% （這幾行不能出現附錄指令的名字：Makefile 編單章時會刪掉含那個字的行。）
\makeatletter
\let\zh@chapapp\@chapapp
\renewcommand{\chaptermark}[1]{\markboth{\ifx\@chapapp\zh@chapapp 第~\thechapter~章\else 附錄~\thechapter\fi\quad #1}{}}
\makeatother

% 行內程式碼（\texttt）不能斷行，段落裡出現較長的識別字時容易凸出右邊界；
% 允許 TeX 多拉開一點字距來排版。
\setlength{\emergencystretch}{3em}

\title{C++ 程式設計：給學過 C 的人}
\author{林子祐}
\date{\today}

\begin{document}

\maketitle
\romantableofcontents

% ---------------------------------------------------------------
% 各章內容放在 "chapters/<單元>/chapter.tex"。
% 尚未動筆的章先註解掉，整本書才編得過。
% ---------------------------------------------------------------

\chapter{從 C 到 C++}
\label{ch:from-c-to-cpp}

學完 C，你已經會變數、迴圈、陣列、函式、指標和結構。 C++ 幾乎把這些
全部保留下來，再往上加了參考、類別、模板和一整套標準函式庫。 換句話說，
你不是從零開始學一個新語言，而是在熟悉的地基上多蓋幾層。

這一章只做最基本的轉換：寫出第一支 C++ 程式並編譯它，看懂到處出現的
\texttt{std::}，學會用 \texttt{cout} 輸出、用 \texttt{cin} 輸入，以及
解題時常見的輸入加速寫法。 最後一節整理 C 的寫法在 C++ 裡哪些照用、
哪些要改。

\section{第一支 C++ 程式}
\label{sec:first-program}

先看熟悉的 C 版本：

\begin{code}{c}{C 的 Hello}
#include <stdio.h>

int main(void) {
    printf("Hello, C!\n");
    return 0;
}
\end{code}

對應的 C++ 版本長這樣：

\begin{example}[label=ex:hello]{cpp}{(hello.cpp) 第一支 C++ 程式}
#include <iostream>

int main() {
    std::cout << "Hello, C++!" << '\n';
    std::cout << 1 + 2 << '\n';
    return 0;
}
\end{example}

\iosample{}{
Hello, C++! \\
3 \\
}

逐行對照：

\begin{itemize}
\item \texttt{\#include <iostream>}：C++ 的輸入輸出放在
\texttt{<iostream>}（input/output stream）這個標頭檔裡，地位相當於
C 的 \texttt{<stdio.h>}。 C++ 標準函式庫的標頭檔沒有 \texttt{.h}。
\item \texttt{int main()}：C++ 裡空括號就代表沒有參數，不必寫
\texttt{void}（寫了也合法）。
\item \texttt{std::cout}：標準輸出，也就是螢幕。 前面的 \texttt{std::}
下一節解釋。
\item \texttt{<<}：把右邊的東西送進 \texttt{cout}。 可以一直串下去，
一行送好幾個值。
\item \texttt{'\textbackslash n'}：換行字元，跟 C 一樣。
\end{itemize}

第二行輸出值得注意：\texttt{1 + 2} 印出 $3$，我們沒有寫任何
\texttt{\%d}。 \texttt{cout} 會依照值的型別決定怎麼印，這是它跟
\texttt{printf} 最大的不同，第~\ref{sec:cout}~節再細看。

\subsection{編譯與執行}

C++ 原始檔的副檔名慣例是 \texttt{.cpp}。 編譯器改用 \texttt{g++}，
用法跟 \texttt{gcc} 幾乎一樣：

\begin{code}{bash}{編譯並執行 hello.cpp}
g++ -std=c++17 -Wall hello.cpp -o hello
./hello
\end{code}

\texttt{-std=c++17} 指定用 C++17 標準編譯，\texttt{-Wall} 打開常用警告。
本書每一支完整範例都用這行指令編譯、實際跑過，印出的結果就是書上寫的。

Windows 上執行檔叫 \texttt{hello.exe}，在命令提示字元裡直接打
\texttt{hello} 執行。 用 IDE 的話，記得在設定裡把語言標準選成 C++17
或更新的版本。

\subsection{常見錯誤}

\begin{itemize}
\item \textbf{用 \texttt{gcc} 編譯 \texttt{.cpp}}：\texttt{gcc} 看到
\texttt{.cpp} 會用 C++ 模式編譯，但連結時不會自動加上 C++ 標準函式庫，
常見的結果是一長串 \texttt{undefined reference}。 編 C++ 一律用
\texttt{g++}。
\item \textbf{寫成 \texttt{<iostream.h>}}：這是很舊的寫法，現在的編譯器
找不到這個檔案。 C++ 標準標頭檔沒有 \texttt{.h}。
\item \textbf{忘了 \texttt{-std=c++17}}：本書後面會用到 C++17 才有的
寫法。 舊版編譯器的預設標準可能比較舊，明確寫出來最保險。
\end{itemize}

\section{namespace 與 \texttt{std::}}
\label{sec:namespace}

大型程式和函式庫裡名字很多，難免撞名：你寫了一個 \texttt{count}，某個
函式庫也有一個 \texttt{count}。 C 只能靠取名字時加前綴來避開，例如
\texttt{mylib\_count}。 C++ 提供\textbf{命名空間}（namespace），把名字
分裝在不同的空間裡，用 \texttt{空間名::名字} 指定要哪一個。
\texttt{::} 叫做\textbf{範圍解析運算子}（scope resolution operator）。

\begin{example}[label=ex:namespaces]{cpp}{(namespaces.cpp) 兩個空間裡的同名變數}
#include <iostream>

namespace school {
    int count = 30;
}

namespace club {
    int count = 12;
}

int main() {
    std::cout << school::count << ' ' << club::count << '\n';
    return 0;
}
\end{example}

\iosample{}{
30 12 \\
}

兩個 \texttt{count} 住在不同的空間，互不干擾。 C++ 標準函式庫的名字
幾乎都放在一個叫 \texttt{std}（standard）的命名空間裡，所以要寫
\texttt{std::cout}，意思是 \texttt{std} 裡的 \texttt{cout}。

\subsection{\texttt{using namespace std;}}

每次都寫 \texttt{std::} 有點囉嗦。 在檔案開頭加一行

\begin{code}{cpp}{把 std 的名字全部引進來}
using namespace std;
\end{code}

之後就能直接寫 \texttt{cout}，編譯器會自己到 \texttt{std} 裡找。 解題
和小程式常這樣寫，本書從下一節開始也這樣寫，讓範例短一點。

代價是 \texttt{std} 裡上千個名字全部倒進來，撞名的機會變大。 下面這段
程式自己定義了全域變數 \texttt{count}，而 \texttt{<algorithm>} 裡剛好
有一個函式 \texttt{std::count}：

\begin{code}{cpp}{using namespace std 造成撞名（錯誤示範）}
#include <algorithm>
#include <iostream>
using namespace std;

int count = 0;

int main() {
    count++;               // 編譯錯誤
    cout << count << '\n';
    return 0;
}
\end{code}

\texttt{g++} 會說：

\begin{code}{text}{編譯器的錯誤訊息}
error: reference to 'count' is ambiguous
\end{code}

ambiguous 是有歧義：編譯器不知道你要自己的 \texttt{count} 還是
\texttt{std::count}。 遇到這種錯誤，把自己的變數改個名字就好。 寫給別
人用的程式或多人合作的專案，通常不寫 \texttt{using namespace std;}，
老實寫 \texttt{std::}。 折衷的寫法是只引進用得到的名字，例如
\texttt{using std::cout;}。

\subsection{常見錯誤}

\begin{itemize}
\item \textbf{漏了 \texttt{std::}}：沒有寫 \texttt{using}，又直接寫
\texttt{cout}，\texttt{g++} 會報 \texttt{'cout' was not declared in
this scope; did you mean 'std::cout'?}。 照提示補上 \texttt{std::}，
或在開頭加 \texttt{using}。
\item \textbf{寫在標頭檔裡}：\texttt{using namespace std;} 若寫在
\texttt{.h} 標頭檔，每個 \texttt{\#include} 它的檔案都會被迫引進整個
\texttt{std}。 標頭檔裡一律寫完整的 \texttt{std::}。
\end{itemize}

\section{用 \texttt{cout} 輸出}
\label{sec:cout}

\begin{crecap}[printf 的格式字串]
\texttt{printf} 靠格式字串決定怎麼印每個引數：\texttt{\%d} 印
\texttt{int}、\texttt{\%lld} 印 \texttt{long long}、\texttt{\%f} 印
\texttt{double}、\texttt{\%c} 印字元、\texttt{\%s} 印字串，
\texttt{\%.2f} 印到小數點後兩位。 格式與型別對不上（例如用
\texttt{\%d} 印 \texttt{double}）是未定義行為，編譯器只會給警告，
程式照樣編得過、印出錯的東西。
\end{crecap}

\texttt{cout} 不用格式字串。 每個用 \texttt{<<} 送進去的值，編譯器都
知道它的型別，會自動選對應的印法，所以不會有格式與型別對不上的問題。

\begin{example}[label=ex:cout-basics]{cpp}{(cout\_basics.cpp) 不同型別的輸出與小數位數}
#include <iomanip>
#include <iostream>
using namespace std;

int main() {
    int n = 3;
    double avg = 247.0 / n;
    char grade = 'B';
    bool passed = avg >= 60;

    cout << "人數：" << n << '\n';
    cout << "平均：" << avg << '\n';
    cout << "等第：" << grade << '\n';
    cout << "及格：" << passed << '\n';

    cout << fixed << setprecision(2);
    cout << "平均：" << avg << '\n';
    cout << "再印一次：" << 1.0 / 3 << '\n';
    return 0;
}
\end{example}

\iosample{}{
人數：3 \\
平均：82.3333 \\
等第：B \\
及格：1 \\
平均：82.33 \\
再印一次：0.33 \\
}

幾個要注意的地方：

\begin{itemize}
\item \texttt{double} 預設印 $6$ 位\textbf{有效數字}，所以
$82.333\ldots$ 印成 \texttt{82.3333}，不是 C 的 \texttt{\%f} 那種固定
小數點後 $6$ 位。
\item \texttt{bool} 是 C++ 內建的型別，值是 \texttt{true} 或
\texttt{false}，用 \texttt{cout} 印出來是 \texttt{1} 或 \texttt{0}。
\item 要固定小數位數，先送 \texttt{fixed}，再送
\texttt{setprecision(位數)}，效果等於 \texttt{\%.2f}。
\texttt{setprecision} 在 \texttt{<iomanip>} 裡，要另外 include。
\item \texttt{fixed} 和 \texttt{setprecision} 一旦設定就一直有效。
範例最後一行沒有再設定，\texttt{1.0 / 3} 照樣印成兩位小數。
\end{itemize}

\subsection{\texttt{endl} 與 \texttt{'\textbackslash n'}}

很多教材用 \texttt{endl} 換行：

\begin{code}{cpp}{兩種換行寫法}
cout << "Hello" << endl;   // 換行，並且立刻清空輸出緩衝區
cout << "Hello" << '\n';   // 只換行
\end{code}

兩者都會換行，差別在 \texttt{endl} 每次還會\textbf{清空緩衝區}
（flush），把累積的輸出立刻寫出去。 偶爾印幾行看不出差別；在迴圈裡印
幾十萬行時，每行都 flush 會明顯拖慢程式。 本書一律用
\texttt{'\textbackslash n'}。

\subsection{常見錯誤}

\begin{itemize}
\item \textbf{字元做運算後變成整數}：\texttt{char c = 'A';} 之後
\texttt{cout << c + 1} 印出 \texttt{66}，不是 \texttt{B}。
\texttt{c + 1} 的型別已經提升成 \texttt{int}，\texttt{cout} 就照整數
印。 要印字元，寫 \texttt{(char)(c + 1)}。
\item \textbf{忘了加括號}：\texttt{<<} 的優先順序比比較運算子高，
\texttt{cout << a < b;} 會被解讀成 \texttt{(cout << a) < b}，編譯失敗。
要印比較結果，寫 \texttt{cout << (a < b);}。
\item \textbf{只寫 \texttt{setprecision} 沒寫 \texttt{fixed}}：這時
\texttt{setprecision(2)} 指的是\textbf{有效數字}兩位，$82.3333$ 會印成
\texttt{82}。 要固定小數點後幾位，兩個都要寫。
\end{itemize}

\section{用 \texttt{cin} 輸入}
\label{sec:cin}

\begin{crecap}[scanf 為什麼要 \&]
C 的函式參數是傳值：函式拿到的是引數的複本，改複本不會影響呼叫者的變
數。 \texttt{scanf} 要把讀到的值寫回你的變數，只能請你把變數的位址交
給它，它再沿著指標寫進去：

\begin{code}{c}{scanf 收的是位址}
int n;
scanf("%d", &n);   // 傳 n 的位址，scanf 透過指標寫入 n
\end{code}

漏了 \texttt{\&}，\texttt{scanf} 會把 \texttt{n} 裡的垃圾值當成位址去
寫，結果是未定義行為。 另外，\texttt{scanf} 的回傳值是成功讀到幾個
項目，讀到檔案結尾時回傳 \texttt{EOF}。
\end{crecap}

C++ 用 \texttt{cin}（標準輸入）和 \texttt{>>} 讀資料，箭頭方向是從
\texttt{cin} 流向變數：

\begin{code}{cpp}{cin 的基本用法}
int n;
cin >> n;          // 不用寫 &
int a, b;
cin >> a >> b;     // 一次讀兩個
\end{code}

\texttt{cin >> n} 不用寫 \texttt{\&n}，卻一樣能改到 \texttt{n}。 背後
的機制叫\textbf{參考}（reference），是第~2~章的主題；現在只要記得
\texttt{cin} 後面直接寫變數名稱。

\texttt{>>} 會先跳過空白、Tab 和換行，再讀一個值，所以數字之間用空白
還是換行隔開都可以，這點跟 \texttt{scanf("\%d")} 一樣。

\begin{example}[label=ex:cin-sum]{cpp}{(cin\_sum.cpp) 讀 n 個整數，輸出總和與最大值}
#include <iostream>
using namespace std;

int main() {
    int n;
    cin >> n;

    int sum = 0, mx = 0;
    for (int i = 0; i < n; i++) {
        int x;
        cin >> x;
        sum += x;
        if (i == 0 || x > mx) {
            mx = x;
        }
    }
    cout << sum << ' ' << mx << '\n';
    return 0;
}
\end{example}

\iosample{
5 \\
3 9 -2 7 4 \\
}{
21 9 \\
}

\subsection{讀到檔案結尾為止}

有些題目不先告訴你有幾筆資料，要一直讀到輸入結束。 C 的寫法是
\texttt{while (scanf("\%d", \&x) == 1)}。 C++ 的寫法更短：

\begin{example}[label=ex:cin-eof]{cpp}{(cin\_eof.cpp) 讀到輸入結束，輸出筆數與總和}
#include <iostream>
using namespace std;

int main() {
    int x;
    int cnt = 0;
    long long sum = 0;
    while (cin >> x) {
        cnt++;
        sum += x;
    }
    cout << cnt << ' ' << sum << '\n';
    return 0;
}
\end{example}

\iosample{
10 20 30 \\
40 \\
}{
4 100 \\
}

\texttt{cin >> x} 這個運算式本身可以當條件用：讀取成功就是真，讀到結尾
或讀取失敗就是假，迴圈自然結束。 在終端機手動輸入時，Linux 和 macOS
按 Ctrl+D、Windows 按 Ctrl+Z 再按 Enter，代表輸入結束。

\subsection{常見錯誤}

\begin{itemize}
\item \textbf{箭頭方向寫反}：輸出是 \texttt{cout <<}，輸入是
\texttt{cin >>}。 箭頭指向資料要去的地方。 寫成 \texttt{cin << x} 會
編譯失敗。
\item \textbf{照 C 的習慣寫 \texttt{cin >> \&x}}：\texttt{\&x} 是位址，
不能拿來存讀進來的值，編譯會失敗。 這類錯誤的訊息常常有幾十行，先看第一
個 \texttt{error:} 所在的行號，通常就是寫錯的那一行。
\item \textbf{輸入格式不符就繼續讀}：要讀整數卻碰到 \texttt{abc}，這次
讀取失敗，\texttt{cin} 進入失敗狀態。 呼叫 \texttt{cin.clear()} 清除
這個狀態之前，後面的 \texttt{>>} 都不會再讀任何東西。 用 \texttt{while (cin >> x)} 這種寫法，失敗時迴圈會停下來，
比較不會出事。
\end{itemize}

\section{輸入加速}
\label{sec:fast-io}

解題時資料量可能很大，例如一次讀進幾十萬個整數。 這時 \texttt{cin}
和 \texttt{cout} 的預設設定會拖慢速度，常見的做法是在 \texttt{main} 一
開頭加兩行：

\begin{code}{cpp}{輸入加速的兩行}
ios::sync_with_stdio(false);
cin.tie(nullptr);
\end{code}

這兩行各自關掉一項預設行為：

\begin{itemize}
\item \texttt{ios::sync\_with\_stdio(false)}：預設情況下，
\texttt{cin}／\texttt{cout} 和 C 的 \texttt{scanf}／\texttt{printf}
保持同步，混著用時順序不會亂。 同步有成本，這行把它關掉。
\item \texttt{cin.tie(nullptr)}：預設情況下 \texttt{cin} 綁在
\texttt{cout} 上，每次讀輸入前都會先 flush \texttt{cout}，確保提示文字
已經印出來。 這行解除綁定。 \texttt{nullptr} 是 C++ 的空指標，相當於
C 的 \texttt{NULL}，第~\ref{sec:c-in-cpp}~節會再提。
\end{itemize}

\begin{example}[label=ex:fastio]{cpp}{(fastio.cpp) 加速輸入輸出，逐行印出前綴和}
#include <iostream>
using namespace std;

int main() {
    ios::sync_with_stdio(false);
    cin.tie(nullptr);

    int n;
    cin >> n;
    long long prefix = 0;
    for (int i = 0; i < n; i++) {
        int x;
        cin >> x;
        prefix += x;
        cout << prefix << '\n';
    }
    return 0;
}
\end{example}

\iosample{
5 \\
3 9 -2 7 4 \\
}{
3 \\
12 \\
10 \\
17 \\
21 \\
}

輸出用的是 \texttt{'\textbackslash n'} 而不是 \texttt{endl}。 前面那兩
行省下來的時間，如果每一行都 \texttt{endl} 一次，很容易又被吃回去。

\subsection{常見錯誤}

關掉同步之後，\texttt{cout} 和 \texttt{printf} 各自有自己的緩衝區，
誰先寫到螢幕上就不再有保證。 下面這段程式在本書測試的環境下，
\texttt{printf} 那行跑到最後面：

\begin{code}{cpp}{關掉同步後混用 printf（錯誤示範）}
ios::sync_with_stdio(false);
cout << "第一行" << '\n';
printf("第二行\n");
cout << "第三行" << '\n';
\end{code}

\iosample{}{
第一行 \\
第三行 \\
第二行 \\
}

換一個環境，順序可能又不一樣。 規則很簡單：加了這兩行，整支程式就只
用 \texttt{cin}／\texttt{cout}，不要再混用 \texttt{scanf}／
\texttt{printf}。

另一個常見錯誤是把這兩行放在已經讀過輸入之後。 它們要在第一次輸入輸出
之前執行，所以放在 \texttt{main} 的最前面。

\section{C 的寫法在 C++ 裡}
\label{sec:c-in-cpp}

C++ 幾乎是 C 的超集合：你在 C 學過的語法，絕大部分在 C++ 裡照樣能
用。 這一節整理需要留意的幾個地方。

\subsection{C 的標頭檔}

C 標準函式庫在 C++ 裡都還在，只是標頭檔換個名字：去掉 \texttt{.h}，前
面加 \texttt{c}。

\begin{center}
\begin{tabular}{ll}
\toprule
C & C++ \\
\midrule
\texttt{<stdio.h>} & \texttt{<cstdio>} \\
\texttt{<stdlib.h>} & \texttt{<cstdlib>} \\
\texttt{<math.h>} & \texttt{<cmath>} \\
\texttt{<string.h>} & \texttt{<cstring>} \\
\texttt{<ctype.h>} & \texttt{<cctype>} \\
\bottomrule
\end{tabular}
\end{center}

\texttt{printf}、\texttt{sqrt}、\texttt{strlen} 這些函式的用法都跟 C
一樣。 寫成 \texttt{<stdio.h>} 通常也編得過，不過 C++ 程式慣用
\texttt{c} 開頭的版本。

\subsection{\texttt{bool} 是內建型別}

C 要 \texttt{\#include <stdbool.h>} 才有 \texttt{bool}、\texttt{true}、
\texttt{false}。 在 C++ 裡它們是語言本身的一部分，不用 include。

\subsection{空指標寫 \texttt{nullptr}}

\begin{crecap}[NULL 與 malloc]
指標不指向任何東西時設成 \texttt{NULL}，使用前先檢查。
\texttt{malloc} 向 heap 要一塊記憶體，回傳型別是 \texttt{void *}，用完
要 \texttt{free}：

\begin{code}{c}{C 的動態配置}
int *a = malloc(5 * sizeof(int));   // void * 自動轉成 int *
if (a != NULL) {
    /* 使用 a[0] .. a[4] */
    free(a);
}
\end{code}

C 允許 \texttt{void *} 自動轉成任何物件指標，所以不必轉型。
\end{crecap}

C++ 有兩處不同：

\begin{itemize}
\item 空指標建議寫 \texttt{nullptr}。 它是專門給指標用的空值，不會像
\texttt{NULL} 那樣在某些地方被當成整數 $0$。
\item C++ 不允許 \texttt{void *} 自動轉成 \texttt{int *}。 上面那行
\texttt{malloc} 原封不動放進 \texttt{.cpp}，\texttt{g++} 會報
\texttt{invalid conversion from 'void*' to 'int*'}，要自己加上
\texttt{(int *)} 轉型。
\end{itemize}

其實寫 C++ 時很少需要 \texttt{malloc}。 要一段長度可變的陣列，第~3~章
的 \texttt{vector} 會幫你管記憶體；真的要自己配置，C++ 有
\texttt{new}／\texttt{delete}，第~4~章再談。 下面的範例只是示範 C 的
寫法在 C++ 裡還能跑：

\begin{example}[label=ex:c-in-cpp]{cpp}{(c\_in\_cpp.cpp) 在 C++ 裡用 C 的函式庫}
#include <cmath>
#include <cstdio>
#include <cstdlib>
#include <iostream>
using namespace std;

bool isSquare(int n) {
    int r = (int)sqrt(n);
    return r * r == n;
}

int main() {
    int n = 5;
    int *a = (int *)malloc(n * sizeof(int));
    for (int i = 0; i < n; i++) {
        a[i] = i * i + 9;
    }
    for (int i = 0; i < n; i++) {
        printf("%d:%s ", a[i], isSquare(a[i]) ? "yes" : "no");
    }
    printf("\n");
    free(a);

    int *p = nullptr;
    if (p == nullptr) {
        cout << "p 沒有指向任何東西" << '\n';
    }
    return 0;
}
\end{example}

\iosample{}{
9:yes 10:no 13:no 18:no 25:yes \\
p 沒有指向任何東西 \\
}

這支程式沒有關掉同步，所以 \texttt{printf} 和 \texttt{cout} 混用，順序
仍然正確。

\subsection{常見錯誤}

\begin{itemize}
\item \textbf{把變數取名為 C++ 的關鍵字}：C++ 多了一些 C 沒有的關鍵字，
例如 \texttt{new}、\texttt{delete}、\texttt{class}、\texttt{this}、
\texttt{template}。 在 C 裡把變數取名為 \texttt{new} 沒問題，搬到 C++
就編不過，要改名。
\item \textbf{以為 C++ 會幫你檢查陣列越界}：C 的陣列在 C++ 裡還是原來
的陣列，越界一樣是未定義行為，編譯器和執行時都不一定會提醒你。 第~3~
章的 \texttt{vector} 提供有檢查的存取方式。
\end{itemize}

\section{本章練習}
\label{sec:ch1-exercises}

\begin{drills}[label=drills:ch1]
\item 預測下面程式的輸出：
\begin{code}{cpp}{練習：預測輸出}
#include <iomanip>
#include <iostream>
using namespace std;

int main() {
    cout << fixed << setprecision(1) << 3.14159 << '\n';
    cout << 2.0 << '\n';
    cout << 'A' + 1 << ' ' << (char)('A' + 1) << '\n';
    return 0;
}
\end{code}
\item 下面的程式編譯失敗，找出原因並用兩種不同的方法修好：
\begin{code}{cpp}{練習：找錯}
#include <iostream>

int main() {
    int n;
    cin >> n;
    cout << n * 2 << '\n';
    return 0;
}
\end{code}
\item 寫一支程式：讀入兩個整數 $a$、$b$（$b \neq 0$），分三行輸出
$a + b$、$a \times b$，以及 $a / b$ 的實數結果（小數點後三位）。 輸入
\texttt{7 2} 時應輸出 \texttt{9}、\texttt{14}、\texttt{3.500}。
\item 一支程式開頭加了 \texttt{ios::sync\_with\_stdio(false);}，後面
同時用了 \texttt{cout} 和 \texttt{printf}。 可能會出什麼問題？ 該怎麼
改？
\item 寫一支程式：不斷讀入整數直到輸入結束，輸出讀到幾個數，以及它們的
平均（小數點後兩位），兩者以空白隔開。 輸入 \texttt{3 4 8} 時應輸出
\texttt{3 5.00}。 讀不到任何數時輸出 \texttt{0}。
\item 為什麼大量輸出時建議用 \texttt{'\textbackslash n'} 而不是
\texttt{endl}？
\end{drills}

\addcontentsline{toc}{section}{練習解答}
\section*{練習解答}

\begin{enumerate}
\item 輸出三行：\texttt{3.1}、\texttt{2.0}、\texttt{66 B}。
\texttt{fixed} 和 \texttt{setprecision(1)} 設定後一直有效，所以
\texttt{2.0} 也印一位小數；\texttt{'A' + 1} 是 \texttt{int}，印出
$66$，轉回 \texttt{char} 才印 \texttt{B}。
\item \texttt{cin}、\texttt{cout} 在 \texttt{std} 裡，沒寫
\texttt{std::}，編譯器找不到。 改法一：寫成 \texttt{std::cin}、
\texttt{std::cout}。 改法二：在 \texttt{\#include} 下面加
\texttt{using namespace std;}（或 \texttt{using std::cin;} 與
\texttt{using std::cout;}）。
\item 一種寫法：
\begin{code}{cpp}{練習 3 參考解答}
#include <iomanip>
#include <iostream>
using namespace std;

int main() {
    int a, b;
    cin >> a >> b;
    cout << a + b << '\n';
    cout << a * b << '\n';
    cout << fixed << setprecision(3) << (double)a / b << '\n';
    return 0;
}
\end{code}
要先把 \texttt{a} 轉成 \texttt{double} 再除，否則 \texttt{7 / 2} 是整
數除法，得到 $3$。
\item 關掉同步後，\texttt{cout} 和 \texttt{printf} 各用各的緩衝區，輸出
順序可能跟程式碼順序不同。 改法是整支程式只用 \texttt{cout}（或者拿掉
那一行，保留同步）。
\item 一種寫法：
\begin{code}{cpp}{練習 5 參考解答}
#include <iomanip>
#include <iostream>
using namespace std;

int main() {
    ios::sync_with_stdio(false);
    cin.tie(nullptr);

    int x, cnt = 0;
    long long sum = 0;
    while (cin >> x) {
        cnt++;
        sum += x;
    }
    if (cnt == 0) {
        cout << 0 << '\n';
        return 0;
    }
    cout << cnt << ' ' << fixed << setprecision(2) << (double)sum / cnt << '\n';
    return 0;
}
\end{code}
\item \texttt{endl} 除了換行，每次還會 flush 緩衝區；印很多行時，每行都
flush 會拖慢程式。 \texttt{'\textbackslash n'} 只換行，讓輸出累積在緩衝
區裡一次寫出去。
\end{enumerate}

\chapter{參考與函式}
\label{ch:references-and-functions}

第~1~章留了一個問題：\texttt{cin >> n} 沒有寫 \texttt{\&n}，為什麼
改得到 \texttt{n}？ 答案是\textbf{參考}（reference）。 這一章先講參考
是什麼、跟指標差在哪裡，再看它怎麼用在函式的參數和回傳值上。 後半章
介紹 C++ 讓函式更好寫的三個功能：函式多載、預設引數和 \texttt{auto}。

\section{參考是變數的別名}
\label{sec:reference}

\begin{crecap}[用指標交換兩個變數]
C 的函式參數是傳值，函式裡改的是複本。 想讓函式改到呼叫者的變數，要
傳位址進去，函式再透過指標解參考（dereference）去改：

\begin{code}{c}{C 的 swap}
#include <stdio.h>

void swap(int *pa, int *pb) {
    int t = *pa;
    *pa = *pb;
    *pb = t;
}

int main(void) {
    int a = 3, b = 7;
    swap(&a, &b);          // 傳位址
    printf("%d %d\n", a, b);
    return 0;
}
\end{code}

輸出 \texttt{7 3}。 這個寫法有兩個麻煩：呼叫時每個引數都要記得加
\texttt{\&}，函式裡每次用到都要寫 \texttt{*}；而且指標可能是
\texttt{NULL}，嚴謹一點的函式還得先檢查。
\end{crecap}

C++ 的參考是一個已經存在的變數的\textbf{別名}（alias）：同一塊記憶體，
多一個名字。 宣告時在變數名稱前面加 \texttt{\&}：

\begin{example}[label=ex:ref-basics]{cpp}{(ref\_basics.cpp) 參考和原變數是同一個東西}
#include <iostream>
using namespace std;

int main() {
    int x = 10;
    int &r = x;    // r 是 x 的另一個名字

    r = 20;
    cout << x << '\n';
    x++;
    cout << r << '\n';
    cout << (&r == &x) << '\n';
    return 0;
}
\end{example}

\iosample{}{
20 \\
21 \\
1 \\
}

\texttt{int \&r = x;} 讀作「\texttt{r} 是一個 \texttt{int} 的參考，綁定
（bind）到 \texttt{x}」。 之後對 \texttt{r} 做的任何事，就是對
\texttt{x} 做：\texttt{r = 20} 改的是 \texttt{x}，\texttt{x++} 之後
讀 \texttt{r} 也看得到。 最後一行比較兩者的位址，結果相同，它們就是
同一個變數。

注意 \texttt{\&} 在這裡的意思。 寫在\textbf{宣告}裡、型別後面時，
\texttt{\&} 表示「參考」；寫在\textbf{運算式}裡、變數前面時，
\texttt{\&x} 還是 C 的取位址。 \texttt{*} 也是同樣的情況：宣告裡表示
指標，運算式裡表示解參考。

\subsection{參考和指標的差別}

參考在底層通常就是用位址實作的，但語言對它的限制比指標多：

\begin{center}
\begin{tabular}{lll}
\toprule
 & 指標 \texttt{int *p} & 參考 \texttt{int \&r} \\
\midrule
宣告時一定要初始化 & 否 & 是 \\
可以改成指向別的變數 & 可以 & 不行，綁定後不能換 \\
可以是空的 & 可以是 \texttt{nullptr} & 不行 \\
取值的寫法 & \texttt{*p} & 直接寫 \texttt{r} \\
\bottomrule
\end{tabular}
\end{center}

第二列最容易誤會。 對參考做指定，是把值寫進它綁定的那個變數，不是改綁
到別的變數：

\begin{example}[label=ex:ref-vs-ptr]{cpp}{(ref\_vs\_ptr.cpp) 指標可以改指向，參考不能改綁}
#include <iostream>
using namespace std;

int main() {
    int a = 1, b = 2;

    int *p = &a;
    p = &b;        // 指標可以改指向 b
    *p = 20;       // 改的是 b
    cout << a << ' ' << b << '\n';

    int &r = a;    // r 綁定 a，從此不會變
    r = b;         // 把 b 的值指定給 a，不是改綁到 b
    r = 10;        // 改的仍然是 a
    cout << a << ' ' << b << '\n';
    return 0;
}
\end{example}

\iosample{}{
1 20 \\
10 20 \\
}

\texttt{r = b} 執行完，\texttt{a} 變成 $20$；接著 \texttt{r = 10}
又把 \texttt{a} 改成 $10$，\texttt{b} 完全沒被動到。 C++ 根本沒有讓
參考改綁的語法。

這些限制換來的好處是：拿到一個參考，就可以確定它綁著一個真的變數，不
必像指標那樣先檢查是不是空的。 一般的原則是，能用參考就用參考；需要
「可能沒有」或「中途換目標」的時候才用指標。

\subsection{常見錯誤}

\begin{itemize}
\item \textbf{宣告參考卻沒有初始化}：\texttt{int \&r;} 編譯失敗，
\texttt{g++} 會說 \texttt{'r' declared as reference but not
initialized}。
\item \textbf{綁定到常數}：\texttt{int \&r = 5;} 也編不過。 \texttt{5}
不是一個變數，沒有一塊記憶體可以讓 \texttt{r} 當別名：
\begin{code}{text}{編譯器的錯誤訊息}
error: cannot bind non-const lvalue reference of type 'int&' to an rvalue of type 'int'
\end{code}
lvalue 大致是「有名字、有位址、可以放在等號左邊」的東西，例如變數；
rvalue 是 \texttt{5}、\texttt{a + b} 這種算完就丟的值。 一般的參考只能
綁定 lvalue。 第~\ref{sec:const-ref}~節的 \texttt{const} 參考例外。
\item \textbf{想用 \texttt{r = b} 改綁}：如範例~\ref{ex:ref-vs-ptr}，
這行是指定，不是改綁。
\end{itemize}

\section{傳參考}
\label{sec:pass-by-ref}

參考最常見的用途是當函式參數。 參數宣告成參考時，呼叫的那一刻參數就
綁定到傳進來的變數，函式裡改參數就是改呼叫者的變數。 這叫\textbf{傳參
考}（pass by reference）。

\begin{example}[label=ex:swap-ref]{cpp}{(swap\_ref.cpp) 用參考交換，以及用參考帶回兩個結果}
#include <iostream>
using namespace std;

void swapInt(int &a, int &b) {
    int t = a;
    a = b;
    b = t;
}

void divide(int a, int b, int &q, int &r) {
    q = a / b;
    r = a % b;
}

int main() {
    int x = 3, y = 7;
    swapInt(x, y);
    cout << x << ' ' << y << '\n';

    int q, r;
    divide(17, 5, q, r);
    cout << q << ' ' << r << '\n';
    return 0;
}
\end{example}

\iosample{}{
7 3 \\
3 2 \\
}

跟 C 的版本比，\texttt{swapInt} 裡沒有任何 \texttt{*}，呼叫時也不用
\texttt{\&}。 函式取名 \texttt{swapInt} 是為了避開標準函式庫裡現成的
\texttt{std::swap}（在 \texttt{<utility>} 裡），以後交換兩個變數直接用
它就好。

\texttt{divide} 示範另一種常見用法：C 的函式只能回傳一個值，要帶回兩
個結果時，把用來裝結果的變數用參考傳進去。 \texttt{cin >> n} 能改到
\texttt{n}，也是因為它收的是參考。

傳參考有一個缺點：只看呼叫端的 \texttt{swapInt(x, y)}，看不出
\texttt{x}、\texttt{y} 會被改。 所以函式名稱最好清楚表達它會改參數，
不打算改的參數就不要宣告成一般的參考。

\subsection{\texttt{const} 參考：不改，也不複製}
\label{sec:const-ref}

\begin{crecap}[傳結構的指標]
傳值會複製整個引數。 結構很大時，每次呼叫都複製一份很浪費，C 的慣例
是改傳指標，再加 \texttt{const} 表示函式不會修改它，例如
\texttt{double average(const struct Scores *s);}。 函式裡用 \texttt{s->n} 存取成員，呼叫時寫 \texttt{average(\&s)}。
\end{crecap}

C++ 的對應寫法是 \texttt{const} 參考，型別寫成 \texttt{const Scores \&}。
傳進去不會複製，函式裡也不能透過它修改。 順帶一提，C++ 裡用結構當型別
時不必再寫 \texttt{struct} 這個字，直接寫 \texttt{Scores} 就好。

\begin{example}[label=ex:const-ref]{cpp}{(const\_ref.cpp) 唯讀的參數用 const 參考，要修改的用參考}
#include <iostream>
using namespace std;

struct Scores {
    int n;
    int v[1000];
};

double average(const Scores &s) {
    long long sum = 0;
    for (int i = 0; i < s.n; i++) {
        sum += s.v[i];
    }
    return (double)sum / s.n;
}

void addBonus(Scores &s, int bonus) {
    for (int i = 0; i < s.n; i++) {
        s.v[i] += bonus;
    }
}

int main() {
    Scores s;
    s.n = 4;
    s.v[0] = 70; s.v[1] = 85; s.v[2] = 62; s.v[3] = 91;

    cout << sizeof(Scores) << '\n';
    cout << average(s) << '\n';
    addBonus(s, 5);
    cout << average(s) << '\n';
    return 0;
}
\end{example}

\iosample{}{
4004 \\
77 \\
82 \\
}

\texttt{Scores} 有 $4004$ 個位元組。 如果 \texttt{average} 寫成傳值
\texttt{average(Scores s)}，每呼叫一次就複製這 $4004$ 個位元組；寫成
\texttt{const Scores \&s}，只傳一個別名，不複製。 存取成員用
\texttt{.}，跟一般的結構變數一樣，不用 \texttt{->}。

參數要怎麼宣告，可以這樣決定：

\begin{itemize}
\item 函式要修改引數：用參考，例如 \texttt{Scores \&s}。
\item 不修改，而且型別很大（結構、第~3~章的 \texttt{string}、
\texttt{vector}）：用 \texttt{const} 參考。
\item 不修改，而且是 \texttt{int}、\texttt{double}、\texttt{char} 這類
小型別：直接傳值。 複製一個 \texttt{int} 的成本跟傳一個參考差不多，
傳值反而簡單。
\end{itemize}

\texttt{const} 參考還有一個特性：它可以綁定 rvalue。 所以
\texttt{const int \&r = 5;} 合法，宣告成 \texttt{const} 參考的參數也
能直接收 \texttt{5}、\texttt{a + b} 這種值。

\subsection{常見錯誤}

\begin{itemize}
\item \textbf{把常數傳給一般參考}：\texttt{swapInt(x, 3)} 編譯失敗，
錯誤訊息跟 \texttt{int \&r = 5;} 一樣是 \texttt{cannot bind non-const
lvalue reference}。 函式要改這個參數，傳常數進去本來就沒有意義。
\item \textbf{在 \texttt{const} 參考上修改}：在 \texttt{average} 裡寫
\texttt{s.n = 0;}，\texttt{g++} 會報 \texttt{assignment of member
'Scores::n' in read-only object}。 這正是 \texttt{const} 的用處：不小心改到不該改的東西時，在編譯階段就
發現。
\item \textbf{漏寫 \texttt{\&}}：寫成 \texttt{void addBonus(Scores s,
int bonus)} 照樣編得過，但改的是複本，呼叫完 \texttt{main} 裡的分數
沒有變。 函式該改的東西沒改到，先檢查參數有沒有 \texttt{\&}。
\end{itemize}

\section{回傳參考}
\label{sec:return-ref}

函式的回傳型別也可以是參考。 這時函式回傳的不是一個值，而是某個變數
本身，所以函式呼叫可以放在等號左邊：

\begin{example}[label=ex:return-ref]{cpp}{(return\_ref.cpp) 回傳較大的那一個變數}
#include <iostream>
using namespace std;

int &larger(int &a, int &b) {
    if (a > b) {
        return a;
    }
    return b;
}

int main() {
    int x = 3, y = 8;
    larger(x, y) = 0;      // y 比較大，y 變成 0
    cout << x << ' ' << y << '\n';
    larger(x, y) += 5;     // 現在 x 比較大，x 加 5
    cout << x << ' ' << y << '\n';
    return 0;
}
\end{example}

\iosample{}{
3 0 \\
8 0 \\
}

\texttt{larger(x, y)} 回傳的是 \texttt{y} 這個變數，所以
\texttt{larger(x, y) = 0} 把 \texttt{y} 設成 $0$。 第二次呼叫時
\texttt{x} 比較大，\texttt{+= 5} 就加在 \texttt{x} 上。

你其實已經用過回傳參考的函式：\texttt{cin >> a} 執行完回傳的是
\texttt{cin} 本身的參考，所以後面可以接著寫 \texttt{>> b}。
\texttt{cout <<} 能串下去也是同樣的道理。 第~6~章自己寫運算子時會再
看到這種寫法。

\subsection{常見錯誤}

回傳參考時，被回傳的那個變數在函式結束後必須還存在。 下面的函式回傳
區域變數的參考：

\begin{code}{cpp}{回傳區域變數的參考（錯誤示範）}
int &bad() {
    int x = 42;
    return x;
}
\end{code}

\texttt{x} 在 \texttt{bad} 結束時就消失了，呼叫者拿到的是一個綁著已經
不存在的變數的參考，稱為\textbf{懸置參考}（dangling reference）。 使用
它是未定義行為。 這段程式編得過，\texttt{g++ -Wall} 只給警告：

\begin{code}{text}{編譯器的警告}
warning: reference to local variable 'x' returned [-Wreturn-local-addr]
\end{code}

看到這個警告就要改，不能放著。 簡單的判斷方式：回傳的是參數本身（像
\texttt{larger} 的 \texttt{a}、\texttt{b} 都是參考參數）或全域變數，
就沒問題；回傳函式裡宣告的區域變數，就要改成傳值回傳。 變數什麼時候
出生、什麼時候消失，第~4~章會完整談。

\section{函式多載}
\label{sec:overloading}

\begin{crecap}[同一件事，不同名字]
C 不允許兩個函式同名。 所以 \texttt{<stdlib.h>} 裡取 \texttt{int}
絕對值的是 \texttt{abs}、取 \texttt{long} 的是 \texttt{labs}，
\texttt{<math.h>} 裡取 \texttt{double} 的是 \texttt{fabs}。 做的是
同一件事，只因為型別不同，就得記三個名字。
\end{crecap}

C++ 允許同一個名字有好幾個版本，只要\textbf{參數列}不同（參數的個數或
型別不同）。 這叫\textbf{函式多載}（function overloading）。 呼叫時，
編譯器看引數的型別挑出最合適的那一個。 C++ 的 \texttt{<cmath>} 就把
\texttt{std::abs} 多載成好幾個版本，\texttt{int}、\texttt{double} 都用
同一個名字。

\begin{example}[label=ex:overload]{cpp}{(overload.cpp) 同名函式，依引數型別選版本}
#include <iostream>
using namespace std;

void show(int x) {
    cout << "int: " << x << '\n';
}

void show(double x) {
    cout << "double: " << x << '\n';
}

void show(const char *s) {
    cout << "字串: " << s << '\n';
}

int main() {
    show(3);
    show(2.5);
    show("hi");
    show('A');     // char 提升成 int
    show(1.5f);    // float 提升成 double
    return 0;
}
\end{example}

\iosample{}{
int: 3 \\
double: 2.5 \\
字串: hi \\
int: 65 \\
double: 1.5 \\
}

前三個呼叫的引數型別跟某個版本完全一樣，直接選它。 後兩個沒有完全相
同的版本，編譯器就要比較每個版本需要什麼轉換。 規則的細節很多，掌握
大方向就夠了，由好到差大致是：

\begin{enumerate}
\item \textbf{完全相符}：型別一樣。
\item \textbf{提升}（promotion）：小的整數型別（\texttt{char}、
\texttt{short}、\texttt{bool}）變成 \texttt{int}，或 \texttt{float}
變成 \texttt{double}。
\item \textbf{一般轉換}：其他算術型別之間的轉換，例如 \texttt{int} 轉
\texttt{double}、\texttt{double} 轉 \texttt{int}、\texttt{long} 轉
\texttt{int}。
\end{enumerate}

\texttt{show('A')} 裡 \texttt{char} 提升成 \texttt{int} 比轉成
\texttt{double} 好，所以印出 \texttt{int: 65}。 \texttt{show(1.5f)}
同理選了 \texttt{double} 版本。 \texttt{cout <<} 能自動依型別印出不同
格式，靠的也是多載。

\subsection{歧義}

如果有兩個版本一樣好，編譯器不會替你挑，直接報錯。 常見的例子是兩個
引數各自偏向不同的版本：

\begin{code}{cpp}{呼叫有歧義（錯誤示範）}
int maxOf(int a, int b) { return a > b ? a : b; }
double maxOf(double a, double b) { return a > b ? a : b; }

maxOf(3, 2.5);   // 編譯錯誤
\end{code}

\begin{code}{text}{編譯器的錯誤訊息}
error: call of overloaded 'maxOf(int, double)' is ambiguous
\end{code}

選 \texttt{int} 版本，第二個引數要轉換；選 \texttt{double} 版本，第一
個引數要轉換。 兩邊各有一個轉換，分不出高下。 解決方法是讓引數型別一
致，例如寫成 \texttt{maxOf(3.0, 2.5)}。

\subsection{常見錯誤}

\begin{itemize}
\item \textbf{只有回傳型別不同}：\texttt{int half(int x)} 和
\texttt{double half(int x)} 不能同時存在，因為呼叫 \texttt{half(5)}
時編譯器無從分辨。 \texttt{g++} 會在第二個宣告報
\texttt{ambiguating new declaration of 'double half(int)'}。 多載只看參數列，不看回傳型別。
\item \textbf{引數型別哪一邊都不完全相符}：只有 \texttt{show(int)} 和
\texttt{show(double)} 兩個版本時，\texttt{show(5L)} 也有歧義：
\texttt{long} 轉 \texttt{int} 和 \texttt{long} 轉 \texttt{double} 都
是一般轉換。 傳字面常數時注意它的型別，或乾脆自己轉型。
\end{itemize}

\section{預設引數}
\label{sec:default-args}

有些參數大部分時候都傳同一個值。 C++ 可以在宣告參數時給它一個
\textbf{預設引數}（default argument），呼叫時沒傳，就用預設值。

\begin{example}[label=ex:default-args]{cpp}{(default\_args.cpp) 印一條橫線，字元與縮排可省略}
#include <iostream>
using namespace std;

void bar(int len, char ch = '#', int indent = 0);

int main() {
    bar(5);
    bar(3, '*');
    bar(4, '-', 2);
    return 0;
}

void bar(int len, char ch, int indent) {
    for (int i = 0; i < indent; i++) {
        cout << ' ';
    }
    for (int i = 0; i < len; i++) {
        cout << ch;
    }
    cout << '\n';
}
\end{example}

\iosample{}{
\#\#\#\#\# \\
*** \\
\ \ ---- \\
}

\texttt{bar(5)} 等於 \texttt{bar(5, '\#', 0)}，\texttt{bar(3, '*')}
等於 \texttt{bar(3, '*', 0)}。 規則有兩條：

\begin{itemize}
\item 預設引數只能從參數列的\textbf{最右邊}開始連續給。 呼叫時引數也是
從左邊依序對應，不能跳過中間的：想指定 \texttt{indent} 就一定要寫出
\texttt{ch}。
\item 預設值只寫一次。 有函式原型時寫在原型上，定義那裡不要再寫。 呼叫
端看得到原型，編譯器才知道可以省略哪些引數。
\end{itemize}

預設引數和多載有時能做同一件事：上面的 \texttt{bar} 也可以寫成三個多
載版本。 差別在於預設引數只寫一份函式本體，適合「同一件事，只是有些
參數常用固定值」的情況；各版本的做法真的不同時，才用多載。

\subsection{常見錯誤}

\begin{itemize}
\item \textbf{原型和定義都寫預設值}：
\begin{code}{text}{編譯器的錯誤訊息}
error: default argument given for parameter 2 of 'void bar(int, char, int)' [-fpermissive]
\end{code}
把定義那裡的 \texttt{= '\#'} 拿掉。
\item \textbf{預設引數後面接了沒有預設值的參數}：例如
\texttt{void bar(int len, char ch = '\#', int indent)}，\texttt{g++}
會報 \texttt{default argument missing for parameter 3}。
\item \textbf{和多載撞在一起}：同時有 \texttt{void f(int a, int b = 0)}
和 \texttt{void f(int a)}，呼叫 \texttt{f(1)} 時兩個都能用，編譯器報
\texttt{call of overloaded 'f(int)' is ambiguous}。 兩者選一種寫法就好。
\end{itemize}

\section{\texttt{auto}：讓編譯器推斷型別}
\label{sec:auto}

宣告變數時把型別寫成 \texttt{auto}，編譯器會依照初始值推斷型別：
\texttt{auto n = 10;} 的 \texttt{n} 是 \texttt{int}，\texttt{auto x
= 2.5;} 的 \texttt{x} 是 \texttt{double}。 推斷在編譯時就完成，
\texttt{n} 從此就是 \texttt{int}，不會因為之後指定別的值而變型別。

\begin{example}[label=ex:auto]{cpp}{(auto\_demo.cpp) auto 推斷出來的型別}
#include <iostream>
using namespace std;

int main() {
    auto half = 7 / 2;      // int
    auto ratio = 7 / 2.0;   // double
    cout << half << ' ' << ratio << '\n';

    int x = 3;
    int &r = x;
    auto a = r;             // int：複製一份
    auto &b = r;            // int&：綁定 x
    a = 10;
    cout << a << ' ' << x << '\n';
    b = 20;
    cout << x << '\n';
    return 0;
}
\end{example}

\iosample{}{
3 3.5 \\
10 3 \\
20 \\
}

這個範例說明兩件事：

\begin{itemize}
\item \texttt{auto} 推斷的是\textbf{初始值的型別}，不是你心裡想的型別。
\texttt{7 / 2} 是整數除法，所以 \texttt{half} 是 \texttt{int}，值是
$3$。
\item \texttt{auto} 會丟掉參考。 \texttt{r} 是 \texttt{int \&}，
但 \texttt{auto a = r} 推斷成 \texttt{int}，\texttt{a} 是一份複本，改
它不會影響 \texttt{x}。 函式回傳參考時也一樣：\texttt{auto m =
larger(x, y);} 拿到的是複本。 要參考就寫 \texttt{auto \&}；要唯讀又不複製，寫
\texttt{const auto \&}。
\end{itemize}

\subsection{什麼時候用}

\texttt{auto} 真正省事的地方，是型別名稱很長、寫出來也沒有多少資訊的
時候。 第~8、9~章的容器和迭代器，型別名稱常常長到一行寫不下，那時
\texttt{auto} 會很好用。 像 \texttt{int}、\texttt{double} 這種短型別，
寫出來讀者一眼就知道是什麼，寫 \texttt{auto} 反而要讓人回頭看初始值
才知道型別。 本書的原則是：型別很短或很重要時照寫，很長而且從右邊看
得出來時才用 \texttt{auto}。

\subsection{常見錯誤}

\begin{itemize}
\item \textbf{沒有初始值}：\texttt{auto x;} 沒東西可以推斷，編譯失敗，
\texttt{g++} 會報 \texttt{declaration of 'auto x' has no initializer}。
\item \textbf{初始值的型別不是你要的}：\texttt{auto sum = 0;} 之後一直
加小數，\texttt{sum} 是 \texttt{int}，小數部分每次都被截掉，而且
\texttt{-Wall -Wextra} 不會警告。 要 \texttt{double} 就寫
\texttt{0.0}，或乾脆寫出 \texttt{double}。
\item \textbf{以為 \texttt{auto s = "hi";} 是字串型別}：它是
\texttt{const char *}，就是 C 的字串指標。 C++ 的字串型別
\texttt{string} 第~3~章介紹。
\end{itemize}

\section{本章練習}
\label{sec:ch2-exercises}

\begin{drills}[label=drills:ch2]
\item 預測下面程式的輸出：
\begin{code}{cpp}{練習：預測輸出}
#include <iostream>
using namespace std;

void f(int a, int &b) {
    a++;
    b++;
}

int main() {
    int x = 1, y = 1;
    f(x, y);
    int &r = x;
    r = y;
    r++;
    cout << x << ' ' << y << ' ' << r << '\n';
    return 0;
}
\end{code}
\item 寫一個函式 \texttt{void sort3(int \&a, int \&b, int \&c)}，把三個
變數由小到大排好。 再寫 \texttt{main} 讀入三個整數，排序後輸出。 輸入
\texttt{9 2 5} 時應輸出 \texttt{2 5 9}。
\item 下面的程式想把較大的變數設成 $0$，編譯時 \texttt{g++ -Wall} 卻
給了警告，執行起來也有問題。 問題在哪裡？ 怎麼改？
\begin{code}{cpp}{練習：找錯}
#include <iostream>
using namespace std;

int &larger(int a, int b) {
    if (a > b) {
        return a;
    }
    return b;
}

int main() {
    int x = 3, y = 8;
    larger(x, y) = 0;
    cout << x << ' ' << y << '\n';
    return 0;
}
\end{code}
\item 預測下面程式的輸出。 如果在 \texttt{main} 裡再加一行
\texttt{p(5L);}，會發生什麼事？
\begin{code}{cpp}{練習：多載與預設引數}
#include <iostream>
using namespace std;

void p(int x) {
    cout << "A " << x << '\n';
}

void p(double x) {
    cout << "B " << x << '\n';
}

void q(int a, int b = 2, int c = 3) {
    cout << a + b + c << '\n';
}

int main() {
    p('0');
    p(2.5f);
    q(1);
    q(1, 0);
    return 0;
}
\end{code}
\item 寫出下面每個變數的型別：
\begin{code}{cpp}{練習：auto 推斷出什麼}
auto a = 5 / 2;
auto b = 5 / 2.0;
auto c = 'x';
auto d = "abc";
int k = 1;
int &rk = k;
auto e = rk;
auto &f = rk;
\end{code}
\end{drills}

\addcontentsline{toc}{section}{練習解答}
\section*{練習解答}

\begin{enumerate}
\item 輸出 \texttt{3 2 3}。 \texttt{a} 是傳值，\texttt{a++} 不影響
\texttt{x}；\texttt{b} 是參考，\texttt{y} 變成 $2$。 \texttt{r = y}
把 $2$ 寫進 \texttt{x}，\texttt{r++} 再把 \texttt{x} 加成 $3$。
\texttt{r} 就是 \texttt{x}，所以也印 $3$。
\item 一種寫法：
\begin{code}{cpp}{練習 2 參考解答}
#include <iostream>
#include <utility>
using namespace std;

void sort3(int &a, int &b, int &c) {
    if (a > b) swap(a, b);
    if (b > c) swap(b, c);
    if (a > b) swap(a, b);
}

int main() {
    int a, b, c;
    cin >> a >> b >> c;
    sort3(a, b, c);
    cout << a << ' ' << b << ' ' << c << '\n';
    return 0;
}
\end{code}
前兩次比較把最大的換到 \texttt{c}，第三次再排好 \texttt{a}、
\texttt{b}。 \texttt{swap} 是標準函式庫的 \texttt{std::swap}，在
\texttt{<utility>} 裡，也可以換成自己寫的 \texttt{swapInt}。
\item 參數 \texttt{a}、\texttt{b} 是傳值，是 \texttt{larger} 自己的區
域變數，函式結束就消失了。 回傳它們的參考是懸置參考，\texttt{g++} 的警
告是 \texttt{reference to local variable 'a' returned}。 就算程式沒有
當掉，改的也不是 \texttt{main} 的 \texttt{x}、\texttt{y}；在本書測試
的環境下，這支程式執行時直接當掉（Segmentation fault）。 改法是把參數
改成 \texttt{int \&a, int \&b}，回傳的就是 \texttt{main} 裡的變數，輸
出 \texttt{3 0}。
\item 輸出四行：\texttt{A 48}、\texttt{B 2.5}、\texttt{6}、
\texttt{4}。 \texttt{'0'} 是 \texttt{char}，提升成
\texttt{int}，值是字元 \texttt{0} 的編碼 $48$；\texttt{2.5f} 是
\texttt{float}，提升成 \texttt{double}。 \texttt{q(1)} 是
$1 + 2 + 3$，\texttt{q(1, 0)} 是 $1 + 0 + 3$。 加上 \texttt{p(5L)} 會
編譯失敗：\texttt{long} 轉 \texttt{int} 和轉 \texttt{double} 都是一般
轉換，編譯器報 \texttt{call of overloaded 'p(long int)' is
ambiguous}。
\item \texttt{a} 是 \texttt{int}（值 $2$），\texttt{b} 是
\texttt{double}，\texttt{c} 是 \texttt{char}，\texttt{d} 是
\texttt{const char *}，\texttt{e} 是 \texttt{int}（\texttt{auto} 丟掉
參考，\texttt{e} 是 \texttt{k} 的複本），\texttt{f} 是
\texttt{int \&}，綁定 \texttt{k}。
\end{enumerate}

\chapter{string 與 vector}
\label{ch:string-and-vector}

C 的字串是以 \texttt{'\textbackslash 0'} 結尾的 \texttt{char} 陣列，陣列
的長度在宣告時就定死了；要一段會變長的陣列，得自己 \texttt{malloc}、
\texttt{realloc}、\texttt{free}。 這兩件事在 C++ 裡分別交給標準函式庫的
\texttt{string} 和 \texttt{vector}：它們自己管記憶體，用起來像一般的變
數，可以指定、相加、比較，也可以用第~2~章的參考傳進函式。

這一章前半講 \texttt{string}：基本操作、找子字串與數字轉換，以及讀一整
行時最常踩的陷阱。 後半講 \texttt{vector}：怎麼建立和增減元素、怎麼走訪
和傳給函式，最後用 \texttt{vector<vector<int>>} 處理二維的格子。

\section{\texttt{string}：會自己管記憶體的字串}
\label{sec:string}

\begin{crecap}[char 陣列與 \texttt{'\textbackslash 0'}]
C 沒有字串型別。 字串只是一個約定：一串 \texttt{char}，最後放一個
\texttt{'\textbackslash 0'} 表示結束。 \texttt{char a[8] = "Ada";} 佔
$8$ 個位元組，用到前 $4$ 個（\texttt{A}、\texttt{d}、\texttt{a}、
\texttt{'\textbackslash 0'}）。 \texttt{<string.h>} 的函式都靠這個
\texttt{'\textbackslash 0'} 找結尾，也都不知道陣列實際有多大：

\begin{code}{c}{C 字串常見的三個坑}
char a[8] = "Ada";
strcat(a, " Lovelace");     // 需要 13 個位元組，a 只有 8 個
size_t n = strlen(a);       // 從頭數到 '\0'，每呼叫一次就數一次
if (a == "Ada") { /* ... */ }  // 比的是兩個位址，不是內容
\end{code}

\texttt{strcpy}、\texttt{strcat} 寫超過陣列大小是未定義行為，空間夠不夠
要自己算；\texttt{strlen} 每次都從頭數；比較內容要用
\texttt{strcmp(a, "Ada") == 0}，寫 \texttt{==} 比的是位址，\texttt{gcc
-Wall} 只會警告 \texttt{comparison with string literal results in
unspecified behavior}。
\end{crecap}

C++ 的 \texttt{string}（在 \texttt{<string>} 裡，全名
\texttt{std::string}）把字元和長度一起包起來，空間不夠時自己去要更大的
記憶體。 上面三件事在 C++ 裡都變成一般的運算子：

\begin{example}[label=ex:string-basics]{cpp}{(string\_basics.cpp) 建立、相加、比較字串}
#include <iostream>
#include <string>
using namespace std;

int main() {
    string first = "Ada";
    string last("Lovelace");
    string full = first + " " + last;
    cout << full << ' ' << full.size() << '\n';

    full += "!";
    full[0] = 'a';
    cout << full << ' ' << full.size() << '\n';

    string a = "apple", b = "banana";
    cout << (a == "apple") << ' ' << (a < b) << '\n';

    string line(10, '-');
    string empty;
    cout << line << ' ' << empty.size() << ' ' << empty.empty() << '\n';
    return 0;
}
\end{example}

\iosample{}{
Ada Lovelace 12 \\
ada Lovelace! 13 \\
1 1 \\
---------- 0 1 \\
}

逐項看：

\begin{itemize}
\item \textbf{建立}：\texttt{string first = "Ada";} 和 \texttt{string
last("Lovelace");} 都是用字串常數建立。 \texttt{string line(10, '-')}
是 $10$ 個 \texttt{'-'}。 什麼都不給的 \texttt{string empty;} 是空字串，
不是垃圾值。
\item \textbf{相加}：\texttt{+} 接起兩個字串，產生一個新字串；
\texttt{+=} 接在原字串後面。 不用先算空間夠不夠。
\item \textbf{長度}：\texttt{full.size()} 回傳字元數，不含結尾的
\texttt{'\textbackslash 0'}。 \texttt{string} 自己記著長度，呼叫
\texttt{size()} 不用從頭數。 \texttt{empty()} 判斷是不是空字串。
\item \textbf{下標}：\texttt{full[0]} 跟 C 的陣列一樣，取第 $0$ 個字元，
型別是 \texttt{char}，可以讀也可以改。
\item \textbf{比較}：\texttt{==}、\texttt{!=}、\texttt{<}、\texttt{>}
比的是內容，順序跟 \texttt{strcmp} 一樣是逐字元比的字典順序，所以
\texttt{"apple" < "banana"} 是真。 兩邊只要有一邊是 \texttt{string}，
另一邊可以是字串常數。
\end{itemize}

\texttt{s.size()} 這種寫法是呼叫 \texttt{s} 的\textbf{成員函式}
（member function），意思是「對 \texttt{s} 做 \texttt{size} 這件事」。
成員函式怎麼寫是第~5~章的內容，現在會呼叫就好。

\subsection{常見錯誤}

\begin{itemize}
\item \textbf{兩個字串常數相加}：
\texttt{string s = "Hello, " + "world";} 編譯失敗：
\begin{code}{text}{編譯器的錯誤訊息}
error: invalid operands of types 'const char [8]' and 'const char [6]' to binary 'operator+'
\end{code}
\texttt{"Hello, "} 和 \texttt{"world"} 都還是 C 的字串常數（\texttt{char}
陣列），不是 \texttt{string}，兩個陣列不能相加。 \texttt{+} 的兩邊至少要
有一個 \texttt{string}，例如 \texttt{string("Hello, ") + "world"}，或者
直接把兩個常數寫在一起。 範例裡的 \texttt{first + " " + last} 沒問題，
因為是從左邊算起，\texttt{first + " "} 已經是 \texttt{string}。
\item \textbf{用 \texttt{int} 跟 \texttt{size()} 比較}：
\texttt{for (int i = 0; i < s.size(); i++)} 編得過，但 \texttt{-Wall} 會
警告 \texttt{comparison of integer expressions of different signedness}：
\texttt{size()} 回傳的是無號整數，\texttt{i} 是有號的。 迴圈變數改用 \texttt{size\_t}，或寫成
\texttt{i < (int)s.size()}；第~\ref{sec:range-for}~節的寫法連下標都不用。
\item \textbf{以為 \texttt{s[i]} 會檢查範圍}：不會。 \texttt{s[100]}
越界跟 C 陣列一樣是未定義行為。 要檢查可以用 \texttt{s.at(100)}，
第~\ref{sec:vector}~節介紹 \texttt{vector} 時一起說明。
\end{itemize}

\section{找子字串與數字轉換}
\label{sec:string-ops}

處理字串常做的事是切出一段、找某個字元或子字串的位置，以及字串和數字
互轉。 C 要靠 \texttt{strncpy}、\texttt{strstr}、\texttt{atoi}、
\texttt{sprintf}；\texttt{string} 有對應的成員函式，\texttt{<string>} 裡
也有轉換函式：

\begin{example}[label=ex:string-ops]{cpp}{(string\_ops.cpp) 拆日期字串，年份加一}
#include <iostream>
#include <string>
using namespace std;

int main() {
    string date = "2026-11-15";
    string year = date.substr(0, 4);
    string month = date.substr(5, 2);
    cout << year << '/' << month << '/' << date.substr(8) << '\n';

    cout << date.find('-') << ' ' << date.find('-', 5) << '\n';
    cout << date.find("11") << '\n';
    if (date.find('x') == string::npos) {
        cout << "沒有 x" << '\n';
    }

    int y = stoi(year);
    string next = to_string(y + 1) + date.substr(4);
    cout << next << '\n';
    cout << stod("3.75") * 2 << '\n';
    return 0;
}
\end{example}

\iosample{}{
2026/11/15 \\
4 7 \\
5 \\
沒有 x \\
2027-11-15 \\
7.5 \\
}

\begin{itemize}
\item \texttt{s.substr(pos, len)}：從位置 \texttt{pos} 開始，取
\texttt{len} 個字元，回傳一個新字串。 省略 \texttt{len} 就取到結尾，所
以 \texttt{date.substr(8)} 是 \texttt{"15"}。
\item \texttt{s.find(x)}：從頭找 \texttt{x} 第一次出現的位置，\texttt{x}
可以是字元或字串。 \texttt{s.find(x, from)} 從位置 \texttt{from} 開始
找，所以 \texttt{date.find('-', 5)} 跳過第一個 \texttt{-}，找到第二個。
\item 找不到時回傳 \texttt{string::npos}。 \texttt{find} 的回傳型別是
\texttt{size\_t}，\texttt{npos} 是它能表示的最大值，不是任何合法的位置；
判斷有沒有找到，一律跟它比。
\item \texttt{stoi(s)}、\texttt{stod(s)}、\texttt{stoll(s)} 把字串轉成
\texttt{int}、\texttt{double}、\texttt{long long}，對應 C 的
\texttt{atoi}、\texttt{atof}、\texttt{atoll}。 \texttt{to\_string(x)}
反過來把數字轉成字串。
\end{itemize}

\subsection{常見錯誤}

\begin{itemize}
\item \textbf{把 \texttt{find} 的結果直接當條件}：
\texttt{if (s.find('a'))} 想問「有沒有 \texttt{a}」，但 \texttt{find}
回傳的是位置。 \texttt{s} 是 \texttt{"apple"} 時，\texttt{a} 在位置
$0$，條件是假；找不到時回傳 \texttt{npos}，條件反而是真。 要寫
\texttt{if (s.find('a') != string::npos)}。
\item \textbf{\texttt{stoi} 的字串不是數字}：\texttt{stoi("abc")} 不會
像 \texttt{atoi} 那樣回傳 $0$，而是讓程式中止：
\begin{code}{text}{執行時的訊息}
terminate called after throwing an instance of 'std::invalid_argument'
  what():  stoi
\end{code}
這是 C++ 的\textbf{例外}（exception）機制：函式遇到沒辦法處理的情況時
「丟出」一個例外，沒人接住，程式就結束。 例外怎麼接住本書不談，現在只
要知道 \texttt{stoi} 的輸入開頭必須是數字。 開頭是數字、後面接著別的字
元時，\texttt{stoi} 只轉前面的數字，\texttt{stoi("42abc")} 是 $42$。
\item \textbf{\texttt{substr} 的起點超出長度}：\texttt{date.substr(20)}
也會丟出例外（\texttt{std::out\_of\_range}），訊息裡寫著
\texttt{\_\_pos (which is 20) > this->size() (which is 10)}。
\texttt{len} 超過剩下的長度則沒關係，取到結尾為止。
\end{itemize}

\section{讀一整行：\texttt{getline}}
\label{sec:getline}

\texttt{cin >> s} 讀進一個 \texttt{string} 時，跟讀數字一樣先跳過空白，
再讀到下一個空白或換行為止，所以一次只讀一個「詞」。 輸入裡有空白（例如
姓名、一句話）時，要改用 \texttt{getline(cin, s)}：從目前的位置讀到行
尾，把換行字元丟掉，其餘全部放進 \texttt{s}。 它對應 C 的
\texttt{fgets}，但不用擔心緩衝區大小，結果裡也沒有 \texttt{'\textbackslash
n'}。

\begin{example}[label=ex:getline-demo]{cpp}{(getline\_demo.cpp) >> 讀一個詞，getline 讀到行尾}
#include <iostream>
#include <string>
using namespace std;

int main() {
    string first, rest, line;
    cin >> first;
    getline(cin, rest);
    getline(cin, line);
    cout << '[' << first << ']' << '\n';
    cout << '[' << rest << ']' << '\n';
    cout << '[' << line << ']' << '\n';
    return 0;
}
\end{example}

\iosample{
Ada Lovelace \\
Alan Mathison Turing \\
}{
{[}Ada] \\
{[} Lovelace] \\
{[}Alan Mathison Turing] \\
}

輸出加上方括號，才看得出空白在哪裡。 \texttt{cin >> first} 讀到
\texttt{Ada} 就停在空白前面；接下來的 \texttt{getline} 從那個空白開始讀，
所以 \texttt{rest} 是 \texttt{" Lovelace"}，前面帶一個空白。 第二個
\texttt{getline} 才讀到第二行。 重點是：\texttt{getline} 讀的是「目前這
一行剩下的部分」，不是「下一行」。

\subsection{\texttt{>>} 後面接 \texttt{getline} 的陷阱}

題目常常先給一個數字 $n$，再給 $n$ 行資料。 直接照順序寫：

\begin{code}{cpp}{先讀 n 再讀 n 行（錯誤示範）}
int n;
cin >> n;
for (int i = 0; i < n; i++) {
    string name;
    getline(cin, name);
    cout << i + 1 << ": [" << name << "]\n";
}
\end{code}

\iosample{
2 \\
Ada Lovelace \\
Alan Turing \\
}{
1: [] \\
2: [Ada Lovelace] \\
}

第一行的 \texttt{2} 後面還有一個換行。 \texttt{cin >> n} 讀完 \texttt{2}
就停下，換行留在輸入裡；第一次 \texttt{getline} 讀到的是「第一行剩下的
部分」，也就是空字串。 結果每一行都錯開一格，最後一行沒讀到。 C 的
\texttt{scanf("\%d")} 後面接 \texttt{fgets} 也有一模一樣的問題。

修法是在 \texttt{getline} 之前把留下來的換行吃掉，在 \texttt{cin >> n;}
後面加一行：

\begin{code}{cpp}{跳過 n 後面的換行}
cin >> ws;
\end{code}

加上這行，同樣的輸入就會印出 \texttt{1: [Ada Lovelace]} 和
\texttt{2: [Alan Turing]}。
\texttt{cin >> ws} 跳過所有空白字元（空白、Tab、換行），停在下一個不是空
白的字元前面。 它會連空行和下一行開頭的空白一起跳過；如果資料本身可能是
空行，改寫成一次 \texttt{getline(cin, rest);}，把 \texttt{n} 那一行剩下
的部分讀掉丟棄。

要一直讀到輸入結束，跟第~1~章的 \texttt{while (cin >> x)} 一樣，寫
\texttt{while (getline(cin, line))}。

\subsection{常見錯誤}

\begin{itemize}
\item \textbf{\texttt{>>} 和 \texttt{getline} 混用，沒處理換行}：如上面
的錯誤示範。 只用 \texttt{>>}，或只用 \texttt{getline}，都沒有這個問題；
兩種混用時，從 \texttt{>>} 換到 \texttt{getline} 的地方要先吃掉換行。
\item \textbf{中文字的 \texttt{size()}}：\texttt{string s = "王小明";}
的 \texttt{s.size()} 是 $9$，不是 $3$。 \texttt{string} 存的是位元組，\texttt{size()}
算的也是位元組；在 UTF-8 編碼下，一個中文字通常佔 $3$ 個位元組。 同樣
地，\texttt{s[0]} 只是第一個中文字的第一個位元組，不是一個完整的字。
\end{itemize}

\section{\texttt{vector}：會自己長大的陣列}
\label{sec:vector}

\begin{crecap}[固定大小的陣列與 realloc]
C 的陣列大小在宣告時決定，之後不能改。 資料筆數事先不知道時，要不就開一個
「應該夠大」的陣列（\texttt{int a[100000];}），要不就自己用
\texttt{malloc} 配置，滿了再 \texttt{realloc} 加大：

\begin{code}{c}{C 的動態陣列}
int cap = 4, n = 0;
int *a = malloc(cap * sizeof(int));
int x;
while (scanf("%d", &x) == 1) {
    if (n == cap) {                  // 滿了，容量加倍
        cap *= 2;
        int *tmp = realloc(a, cap * sizeof(int));
        if (tmp == NULL) { free(a); return 1; }
        a = tmp;
    }
    a[n++] = x;
}
/* 使用 a[0] .. a[n-1] */
free(a);
\end{code}

要記住目前有幾個（\texttt{n}）、空間有多大（\texttt{cap}），
\texttt{realloc} 要用暫存指標接，最後還不能忘了 \texttt{free}。
\end{crecap}

C++ 的 \texttt{vector}（在 \texttt{<vector>} 裡）就是把上面這套做好的
動態陣列。 宣告時在角括號裡寫元素型別：\texttt{vector<int>} 是一串
\texttt{int}，\texttt{vector<string>} 是一串 \texttt{string}。 這種寫法
叫\textbf{模板}（template），第~8~章再談；現在把 \texttt{vector<int>}
整個當成一個型別名稱就好。

\begin{example}[label=ex:vector-basics]{cpp}{(vector\_basics.cpp) 讀到輸入結束，倒著印出來}
#include <iostream>
#include <vector>
using namespace std;

int main() {
    vector<int> v;
    int x;
    while (cin >> x) {
        v.push_back(x);
    }
    cout << v.size() << '\n';

    for (int i = (int)v.size() - 1; i >= 0; i--) {
        cout << v[i] << ' ';
    }
    cout << '\n';

    v[0] = 100;
    cout << v.front() << ' ' << v.back() << '\n';
    v.pop_back();
    cout << v.size() << ' ' << v.back() << '\n';
    return 0;
}
\end{example}

\iosample{
3 9 -2 \\
7 4 \\
}{
5 \\
4 7 -2 9 3 \\
100 4 \\
4 7 \\
}

跟 C 的版本對照：

\begin{itemize}
\item \texttt{vector<int> v;} 一開始是空的，\texttt{size()} 是 $0$。
\item \texttt{v.push\_back(x)} 把 \texttt{x} 加在最後面。 空間不夠時
\texttt{vector} 自己配置更大的記憶體、把舊資料搬過去，就是 C 版本裡
\texttt{realloc} 那一段。
\item \texttt{v.size()} 是目前的元素個數，對應 C 版本的 \texttt{n}。
\item \texttt{v[i]} 跟陣列一樣用下標存取。 \texttt{v.front()}、
\texttt{v.back()} 是第一個和最後一個元素，\texttt{v.pop\_back()} 刪掉
最後一個；這三個都不能用在空的 \texttt{vector} 上。
\item 不用 \texttt{free}。 \texttt{v} 離開它的範圍時，記憶體會自動還回
去，第~4~章會說明這是怎麼做到的。
\end{itemize}

倒著走的迴圈寫成 \texttt{(int)v.size() - 1}，先轉成 \texttt{int} 再減。
\texttt{size()} 是無號整數，\texttt{v} 是空的時候 \texttt{v.size() - 1}
不是 $-1$，而是 \texttt{18446744073709551615}（在 \texttt{size\_t} 是
$64$ 位元的環境下），迴圈會從一個巨大的下標開始。

\subsection{建立時就給大小}

已經知道要幾個元素時，建立時直接給大小，之後用下標填值。 另外有兩個成員
函式可以整批改變內容：\texttt{resize} 改變長度，\texttt{assign} 整個換掉。
註解寫的是執行完那一行之後的內容：

\begin{code}{cpp}{建立 vector 的幾種寫法，resize 與 assign}
vector<int> a(3);          // 0 0 0
vector<int> b(3, 7);       // 7 7 7
vector<int> c = {3, 1, 4}; // 3 1 4
c.resize(5);               // c: 3 1 4 0 0（變長，補 0）
c.resize(2);               // c: 3 1（變短，後面的丟掉）
b.assign(4, 9);            // b: 9 9 9 9（整個換掉）
\end{code}

和 C 的區域陣列不同，\texttt{vector<int> a(3)} 的元素一定是 $0$，不會是
垃圾值。 \texttt{c.at(1)} 跟 \texttt{c[1]} 一樣取第 $1$ 個元素，差別在
越界時的反應。 \texttt{c[5]} 不檢查，跟 C 陣列一樣是未定義行為；
\texttt{c.at(5)} 會檢查，越界時丟出例外，程式中止並印出：

\begin{code}{text}{執行時的訊息}
terminate called after throwing an instance of 'std::out_of_range'
  what():  vector::_M_range_check: __n (which is 5) >= this->size() (which is 2)
\end{code}

\texttt{string} 也有 \texttt{at()}，行為相同。 \texttt{at()} 每次多做一次
檢查，所以解題時通常用 \texttt{[]}；程式結果怪怪的、懷疑有越界時，把
\texttt{[]} 暫時換成 \texttt{at()}，就能抓到是哪一個下標出錯。

\subsection{常見錯誤}

\begin{itemize}
\item \textbf{對空的 \texttt{vector} 用下標指定}：
\begin{code}{cpp}{還沒有元素就寫入（錯誤示範）}
vector<int> v;
for (int i = 0; i < 3; i++) {
    v[i] = i * 10;
}
\end{code}
\texttt{v} 的大小是 $0$，\texttt{v[0]} 根本不存在。 \texttt{[]} 不會幫你
新增元素，這是越界寫入；在本書測試的環境下，這段程式執行時直接當掉
（Segmentation fault）。 要嘛用 \texttt{push\_back}，要嘛建立時就給大小
\texttt{vector<int> v(3);}。
\item \textbf{給了大小又 \texttt{push\_back}}：\texttt{vector<int> w(3);}
之後 \texttt{w.push\_back(5);}，\texttt{w} 有 $4$ 個元素：
\texttt{0 0 0 5}。 \texttt{push\_back} 一定加在現有元素的後面。 兩種寫法
選一種：先給大小就用下標填，從空的開始就用 \texttt{push\_back}。
\end{itemize}

\section{走訪 vector 與傳給函式}
\label{sec:range-for}

用下標走訪要寫 \texttt{for (size\_t i = 0; i < v.size(); i++)}，有點長。
只是要從頭到尾看過每個元素時，C++ 有更短的\textbf{範圍 for 迴圈}
（range-based for），寫成 \texttt{for (int x : v) \{ ... \}}，讀作「對
\texttt{v} 裡的每一個 \texttt{x}」。 每一輪 \texttt{x} 依序是
\texttt{v[0]}、\texttt{v[1]}、……，不用下標，也不會越界。 這種迴圈也能
用在 \texttt{string} 上，每一輪拿到一個 \texttt{char}。

要注意的是 \texttt{x} 怎麼宣告。 \texttt{int x} 是\textbf{複本}，跟傳值
的參數一樣，改 \texttt{x} 不影響 \texttt{v}。 要修改元素，就用第~2~章
的參考：

\begin{center}
\begin{tabular}{lll}
\toprule
寫法 & \texttt{x} 是什麼 & 用在 \\
\midrule
\texttt{for (int x : v)} & 元素的複本 & 只讀，元素是 \texttt{int} 這類小型別 \\
\texttt{for (int \&x : v)} & 元素本身的參考 & 要修改元素 \\
\texttt{for (const string \&s : v)} & 唯讀的參考 & 只讀，元素是 \texttt{string} 這類大型別 \\
\bottomrule
\end{tabular}
\end{center}

選法跟第~2~章決定參數怎麼宣告完全一樣。 型別也可以寫 \texttt{auto}：
\texttt{for (auto \&x : v)}、\texttt{for (const auto \&s : names)}。

\texttt{vector} 和 \texttt{string} 傳給函式時，也照第~2~章的原則：只讀就
用 \texttt{const} 參考，要修改就用參考。 傳值會把整個 \texttt{vector} 複
製一份，元素一多就很浪費。

\begin{example}[label=ex:vector-pass]{cpp}{(vector\_pass.cpp) 範圍 for 迴圈與傳 vector 給函式}
#include <iostream>
#include <string>
#include <vector>
using namespace std;

void print(const vector<int> &v) {
    for (int x : v) {
        cout << x << ' ';
    }
    cout << '\n';
}

void addAll(vector<int> &v, int d) {
    for (int &x : v) {
        x += d;
    }
}

vector<int> evens(const vector<int> &v) {
    vector<int> result;
    for (int x : v) {
        if (x % 2 == 0) {
            result.push_back(x);
        }
    }
    return result;
}

int main() {
    vector<int> v = {3, 1, 4, 1, 5, 9, 2, 6};
    addAll(v, 1);
    print(v);
    print(evens(v));

    vector<string> names = {"Ada", "Alan", "Grace"};
    for (const string &name : names) {
        cout << name << ' ' << name.size() << '\n';
    }
    return 0;
}
\end{example}

\iosample{}{
4 2 5 2 6 10 3 7 \\
4 2 2 6 10 \\
Ada 3 \\
Alan 4 \\
Grace 5 \\
}

\begin{itemize}
\item \texttt{addAll} 要改呼叫者的 \texttt{vector}，參數是
\texttt{vector<int> \&}；迴圈也要改元素，寫 \texttt{int \&x}。 兩個
\texttt{\&} 缺一個，\texttt{main} 裡的 \texttt{v} 就不會變。
\item \texttt{evens} 回傳一個新的 \texttt{vector}。 C 的函式不能回傳陣
列，只能請呼叫者先準備好空間；C++ 可以直接回傳 \texttt{vector}。 回傳
函式裡的區域 \texttt{vector} 時，編譯器通常會把它直接交給呼叫者，不會再
把元素逐一複製一遍，放心這樣寫。
\item \texttt{names} 的元素是 \texttt{string}，用 \texttt{const string \&}
走訪，每一輪不用複製字串。
\end{itemize}

範圍 for 迴圈背後用的是\textbf{迭代器}（iterator），排序、搜尋這些現成
的演算法也是用迭代器來指定範圍，例如 \texttt{sort(v.begin(), v.end())}。
這些留到第~8、9~章。

\subsection{常見錯誤}

\begin{itemize}
\item \textbf{想修改卻用複本走訪}：\texttt{for (int x : v) \{ x *= 10; \}}
編得過，\texttt{-Wall -Wextra} 也不會警告，但改的是每一輪的複本，
\texttt{v} 完全沒變。 要改元素就寫 \texttt{int \&x}。
\item \textbf{函式參數漏了 \texttt{\&}}：\texttt{void addAll(vector<int>
v, int d)} 一樣編得過，改的是整份複本。 跟第~2~章的 \texttt{addBonus}
是同一個錯誤。
\item \textbf{走訪時增刪元素}：在 \texttt{for (int x : v)} 的迴圈裡對
\texttt{v} 做 \texttt{push\_back} 是未定義行為。 \texttt{push\_back} 可
能讓 \texttt{vector} 換一塊更大的記憶體、把元素搬過去，迴圈手上指向舊記
憶體的東西就失效了；之前用參考綁定的元素（例如 \texttt{int \&first =
v[0];}）也一樣。 這件事跟物件的生命週期有關，第~4~章再仔細談。 現在先記
住：要邊走邊加元素，改用下標迴圈，或者把要加的東西先放進另一個
\texttt{vector}。
\end{itemize}

\section{二維 vector}
\label{sec:vector-2d}

\texttt{vector} 的元素型別可以是任何型別，包括另一個 \texttt{vector}。
\texttt{vector<vector<int>>} 是「一串 \texttt{vector<int>}」，也就是一
個二維的格子：\texttt{g[i]} 是第 $i$ 列（一個 \texttt{vector<int>}），
\texttt{g[i][j]} 是第 $i$ 列第 $j$ 個元素，下標寫法跟 C 的二維陣列一樣。

建立一個 $r$ 列 $c$ 行、全部是 $0$ 的格子，寫成
\texttt{vector<vector<int>> g(r, vector<int>(c, 0));}。 外層是「$r$ 個元素，每個都是 \texttt{vector<int>(c, 0)}」，也就是
\texttt{vector<int> b(n, val)} 的寫法套了兩層。 \texttt{r}、\texttt{c}
可以是執行時才讀進來的變數。 列數是 \texttt{g.size()}，行數是
\texttt{g[0].size()}。

下面的範例是踩地雷的提示數字：讀入一張地圖，\texttt{*} 是地雷，\texttt{.}
是空地；在每塊空地上印出周圍八格裡有幾顆地雷。 地圖每一列是一個字串，
用 \texttt{vector<string>} 存最方便；計數用 \texttt{vector<vector<int>>}。

\begin{example}[label=ex:minesweeper]{cpp}{(minesweeper.cpp) 踩地雷的提示數字}
#include <iostream>
#include <string>
#include <vector>
using namespace std;

int main() {
    int r, c;
    cin >> r >> c;
    vector<string> g(r);
    for (int i = 0; i < r; i++) {
        cin >> g[i];
    }

    vector<vector<int>> cnt(r, vector<int>(c, 0));
    for (int i = 0; i < r; i++) {
        for (int j = 0; j < c; j++) {
            if (g[i][j] != '*') {
                continue;
            }
            for (int di = -1; di <= 1; di++) {
                for (int dj = -1; dj <= 1; dj++) {
                    int ni = i + di, nj = j + dj;
                    if (ni >= 0 && ni < r && nj >= 0 && nj < c) {
                        cnt[ni][nj]++;
                    }
                }
            }
        }
    }

    for (int i = 0; i < r; i++) {
        for (int j = 0; j < c; j++) {
            if (g[i][j] == '*') {
                cout << '*';
            } else {
                cout << cnt[i][j];
            }
        }
        cout << '\n';
    }
    return 0;
}
\end{example}

\iosample{
3 4 \\
*... \\
..*. \\
.... \\
}{
*211 \\
12*1 \\
0111 \\
}

做法是反過來數：每遇到一顆地雷，就把它周圍八格（以及它自己那格，反正
不會印出來）的計數各加 $1$。 \texttt{ni}、\texttt{nj} 超出邊界的格子跳
過。 \texttt{vector<string> g(r)} 先建立 $r$ 個空字串，再用 \texttt{>>}
逐列讀入；\texttt{g[i][j]} 是第 $i$ 個字串的第 $j$ 個字元。

傳給函式時，二維 \texttt{vector} 比 C 的二維陣列好寫。 C 的參數要寫成
\texttt{int a[][10]}，第二維的大小一定要是常數；\texttt{vector} 直接寫
\texttt{const vector<vector<int>> \&g}，大小由 \texttt{g.size()} 和
\texttt{g[0].size()} 取得。 每一列各自是一個 \texttt{vector}，所以長度
也可以不一樣，例如第 $i$ 列只放 $i + 1$ 個元素的三角形。

\subsection{常見錯誤}

\begin{itemize}
\item \textbf{把列數、行數直接寫成兩個數字}：\texttt{vector<vector<int>>
g(3, 4);} 編譯失敗，錯誤訊息的第一行是
\begin{code}{text}{編譯器的錯誤訊息}
error: no matching function for call to 'std::vector<std::vector<int> >::vector(int, int)'
\end{code}
後面還有好幾十行。 外層的第二個引數必須是一個 \texttt{vector<int>}，要
寫成 \texttt{g(3, vector<int>(4))}。
\item \textbf{只建立了外層}：\texttt{vector<vector<int>> g(r);} 有 $r$
列，但每一列都是空的 \texttt{vector}，\texttt{g[0][0]} 已經越界。 要嘛
建立時給每一列大小，要嘛之後對每一列 \texttt{push\_back}。
\item \textbf{列與行弄反}：\texttt{g(r, vector<int>(c))} 的外層是列。
寫成 \texttt{g(c, vector<int>(r))} 之後再用 \texttt{g[i][j]}（\texttt{i}
是列），列數和行數不同時就會越界。 下標也一樣：先列後行。
\end{itemize}

\section{本章練習}
\label{sec:ch3-exercises}

\begin{drills}[label=drills:ch3]
\item 預測下面程式的輸出：
\begin{code}{cpp}{練習：預測輸出}
#include <iostream>
#include <string>
using namespace std;

int main() {
    string s = "banana";
    cout << s.size() << ' ' << s.find("an") << ' ' << s.find("an", 2) << '\n';
    cout << s.substr(1, 3) << ' ' << s.substr(3) << '\n';
    string a = "10", b = "9";
    cout << (a < b) << ' ' << (stoi(a) < stoi(b)) << '\n';
    cout << to_string(12) + "3" << ' ' << stoi("12") + 3 << '\n';
    return 0;
}
\end{code}
\item 下面的程式要讀入 $n$ 個姓名（姓名中可能有空白），輸出最長的那一個。
輸入是 \texttt{3}、\texttt{Ada Lovelace}、\texttt{Alan Turing}、
\texttt{Barbara Liskov} 四行，程式卻輸出 \texttt{Ada Lovelace}。 問題在
哪裡？ 怎麼改？
\begin{code}{cpp}{練習：找錯}
#include <iostream>
#include <string>
using namespace std;

int main() {
    int n;
    cin >> n;
    string longest;
    for (int i = 0; i < n; i++) {
        string name;
        getline(cin, name);
        if (name.size() > longest.size()) {
            longest = name;
        }
    }
    cout << longest << '\n';
    return 0;
}
\end{code}
\item 預測下面程式的輸出。 要讓 \texttt{twice} 真的把 \texttt{v} 的每個
元素乘以 $2$，要改哪裡？
\begin{code}{cpp}{練習：誰被改到了}
#include <iostream>
#include <vector>
using namespace std;

void twice(vector<int> v) {
    for (int &x : v) {
        x *= 2;
    }
}

int main() {
    vector<int> v = {1, 2, 3};
    for (int x : v) {
        x += 1;
    }
    twice(v);
    for (int &x : v) {
        x -= 1;
    }
    for (int x : v) {
        cout << x << ' ';
    }
    cout << '\n';
    return 0;
}
\end{code}
\item 寫一支程式：讀入 $r$、$c$ 和一個 $r$ 列 $c$ 行的整數矩陣，輸出它
的轉置（$c$ 列 $r$ 行，原本的第 $i$ 列第 $j$ 行放到第 $j$ 列第 $i$ 行）。
同一列的數字以一個空白隔開。 輸入 \texttt{2 3}、\texttt{1 2 3}、
\texttt{4 5 6} 時應輸出 \texttt{1 4}、\texttt{2 5}、\texttt{3 6} 三行。
\end{drills}

\addcontentsline{toc}{section}{練習解答}
\section*{練習解答}

\begin{enumerate}
\item 輸出四行：\texttt{6 1 3}、\texttt{ana ana}、\texttt{1 0}、
\texttt{123 15}。 \texttt{"an"} 第一次出現在位置 $1$，從位置 $2$ 開始找
則是位置 $3$。 \texttt{substr(1, 3)} 從位置 $1$ 取三個字元，
\texttt{substr(3)} 取到結尾，剛好都是 \texttt{ana}。 字串比較是字典順序，
\texttt{'1'} 比 \texttt{'9'} 小，所以 \texttt{"10" < "9"} 是真；轉成整數
後 $10 < 9$ 是假。 \texttt{to\_string(12) + "3"} 是字串相接，
\texttt{stoi("12") + 3} 是整數相加。
\item \texttt{cin >> n} 讀完 \texttt{3} 後，換行留在輸入裡，第一次
\texttt{getline} 讀到空字串，之後每一行都錯開一格，\texttt{Barbara
Liskov} 沒有被讀到。 在 \texttt{cin >> n;} 後面加一行
\texttt{cin >> ws;}（或 \texttt{string rest; getline(cin, rest);}），
就會輸出 \texttt{Barbara Liskov}。
\item 輸出 \texttt{0 1 2}。 第一個迴圈的 \texttt{x} 是複本，\texttt{v}
沒變；\texttt{twice} 是傳值，乘以 $2$ 的是函式自己的複本；只有
\texttt{int \&x} 的那個迴圈真的改到 \texttt{v}，每個元素減 $1$。 要讓
\texttt{twice} 生效，把參數改成 \texttt{vector<int> \&v}，這時輸出變成
\texttt{1 3 5}。
\item 一種寫法：
\begin{code}{cpp}{練習 4 參考解答}
#include <iostream>
#include <vector>
using namespace std;

int main() {
    int r, c;
    cin >> r >> c;
    vector<vector<int>> a(r, vector<int>(c));
    for (int i = 0; i < r; i++) {
        for (int j = 0; j < c; j++) {
            cin >> a[i][j];
        }
    }

    for (int j = 0; j < c; j++) {
        for (int i = 0; i < r; i++) {
            cout << a[i][j] << (i + 1 < r ? ' ' : '\n');
        }
    }
    return 0;
}
\end{code}
不必真的建出轉置後的矩陣：轉置的第 $j$ 列就是原矩陣的第 $j$ 行，外層迴圈
走行、內層迴圈走列，照順序印出來就好。 每列最後一個數字後面接換行，其他
接空白，行尾就不會多出空白。
\end{enumerate}

\chapter{生命週期}
\label{ch:lifetime}

第~2~章的懸置參考和第~3~章「走訪時不要 \texttt{push\_back}」，都跟物件
什麼時候建立、什麼時候消失有關。 C 的變數消失時不會執行任何程式，記憶體
還給系統就結束了。 C++ 的物件消失時會執行一段程式，\texttt{string} 和
\texttt{vector} 用這段程式把記憶體還回去，所以不用替它們 \texttt{free}。

這一章先整理物件能活多久，再看物件建立和銷毀的順序，以及物件銷毀之後還
去用它會怎樣。 後半章講 \texttt{new}／\texttt{delete}，以及平常怎麼用
RAII 和智慧指標取代它們，最後談複製和移動。

\section{三種儲存期}
\label{sec:storage}

\begin{crecap}[區域變數、static 與 malloc]
C 的變數依照宣告的位置和方式，記憶體配置和釋放的時機不同：

\begin{code}{c}{C 的三種變數}
int total = 0;               // 全域：程式開始就有，程式結束才釋放

void f(void) {
    int x = 0;               // 區域：進入區塊時配置，離開區塊就釋放
    static int calls = 0;    // static 區域：只初始化一次，一直存在
    int *p = malloc(sizeof(int));   // heap：直到 free(p) 才釋放
    /* ... */
}
\end{code}

區域變數離開區塊就沒了，所以函式不能回傳區域變數的位址；static 區域變數
一直都在，回傳它的位址沒問題。 \texttt{malloc} 拿到的記憶體不會自己消失，
忘了 \texttt{free} 就是記憶體洩漏，\texttt{free} 之後還去用就是懸置指標。
\end{crecap}

C++ 沿用同一套分類，名稱叫\textbf{儲存期}（storage duration）：

\begin{center}
\begin{tabular}{llll}
\toprule
儲存期 & 例子 & 何時建立 & 何時銷毀 \\
\midrule
自動（automatic） & 區域變數、參數 & 執行到宣告 & 離開所在的區塊 \\
靜態（static） & 全域變數 & \texttt{main} 開始之前 & 程式結束 \\
 & \texttt{static} 區域變數 & 第一次執行到宣告 & 程式結束 \\
動態（dynamic） & \texttt{new} 出來的物件 & \texttt{new} & \texttt{delete} \\
\bottomrule
\end{tabular}
\end{center}

\texttt{new} 和 \texttt{delete} 是 C++ 版的 \texttt{malloc} 和
\texttt{free}，第~\ref{sec:new-delete}~節介紹。 先看前兩種：

\begin{example}[label=ex:storage]{cpp}{(storage.cpp) 自動與靜態儲存期}
#include <iostream>
using namespace std;

int nextTicket() {
    static int ticket = 100;   // static 區域變數：只初始化一次
    ticket++;
    return ticket;
}

int main() {
    for (int i = 0; i < 3; i++) {
        int step = 10 + i;     // 自動：每一輪重新建立
        cout << nextTicket() << ' ' << step << '\n';
    }
    return 0;
}
\end{example}

\iosample{}{
101 10 \\
102 11 \\
103 12 \\
}

\texttt{step} 每一輪都是新的，上一輪的值不會留下來；\texttt{ticket} 只在
第一次呼叫 \texttt{nextTicket} 時設成 $100$，之後每次呼叫都沿用同一個，
所以號碼一直往上加。

\subsection{範圍與生命週期不是同一件事}

這兩個詞要分開：

\begin{itemize}
\item \textbf{範圍}（scope）：名字在原始碼的哪一段看得到。 這是編譯時的
概念，編譯器用它決定一個名字指的是哪個變數。
\item \textbf{生命週期}（lifetime）：物件在執行時從建立到銷毀的那段期間。
物件只有在生命週期內可以使用。
\end{itemize}

自動變數的兩者剛好重疊：\texttt{step} 的名字只在迴圈本體裡看得到，物件也
只活一輪。 其他情況兩者不一致：

\begin{itemize}
\item \texttt{ticket} 的範圍只有 \texttt{nextTicket} 的本體，生命週期卻
是整個程式。 \texttt{main} 叫不出這個名字，但它一直都在。
\item 動態配置的物件根本沒有名字，只能透過指標找到。 指標變數離開範圍時，
物件不會跟著消失；沒有其他指標記得它，它就洩漏了。
\item 反過來，名字還在範圍內，它指向的物件卻可能已經死了。 指標
\texttt{p} 指著一個區塊裡的區域變數，區塊結束後 \texttt{p} 還能用，但
\texttt{*p} 已經沒有東西，這就是第~\ref{sec:dangling}~節的懸置指標。
\end{itemize}

\subsection{常見錯誤}

\begin{itemize}
\item \textbf{以為區域變數會保留上一次的值}：把 \texttt{nextTicket} 裡的
\texttt{static} 拿掉，每次呼叫都重新建立一個 \texttt{ticket}，三次都回傳
$101$。 要跨呼叫保留值，就要加 \texttt{static}。
\end{itemize}

\section{物件何時建立、何時銷毀}
\label{sec:ctor-dtor}

C++ 的結構可以寫兩個特別的函式，在物件建立和銷毀時自動執行：

\begin{itemize}
\item \textbf{建構子}（constructor）：名稱跟結構相同，沒有回傳型別。 物件
建立時自動呼叫，通常用來設定初始值。
\item \textbf{解構子}（destructor）：名稱是 \texttt{\textasciitilde} 加上
結構名稱，沒有參數也沒有回傳型別。 物件銷毀時自動呼叫，通常用來釋放資源。
\end{itemize}

寫法的細節是第~5~章的內容。 這裡只用它們在建立和銷毀時各印一行，看物件
什麼時候建立、什麼時候銷毀。

\begin{example}[label=ex:lifetime-order]{cpp}{(lifetime\_order.cpp) 會報告出生與死亡的 Tracer}
#include <iostream>
#include <string>
using namespace std;

struct Tracer {
    string name;
    Tracer(const string &n) {      // 建構子：物件建立時自動呼叫
        name = n;
        cout << "建立 " << name << '\n';
    }
    ~Tracer() {                    // 解構子：物件銷毀時自動呼叫
        cout << "銷毀 " << name << '\n';
    }
};

void f() {
    static Tracer s("static s");
}

int main() {
    cout << "main 開始" << '\n';
    Tracer a("a");
    {
        Tracer b("b");
        Tracer c("c");
        cout << "內層區塊結束" << '\n';
    }
    f();
    f();
    Tracer("暫時物件");
    cout << "main 結束" << '\n';
    return 0;
}
\end{example}

\iosample{}{
main 開始 \\
建立 a \\
建立 b \\
建立 c \\
內層區塊結束 \\
銷毀 c \\
銷毀 b \\
建立 static s \\
建立 暫時物件 \\
銷毀 暫時物件 \\
main 結束 \\
銷毀 a \\
銷毀 static s \\
}

\texttt{Tracer a("a");} 建立一個 \texttt{Tracer}，括號裡的字串交給建構子
的參數 \texttt{n}。 對照輸出逐段看：

\begin{itemize}
\item \textbf{區塊裡的物件}依宣告順序建立，離開區塊時以\textbf{相反順序}
銷毀：\texttt{b}、\texttt{c} 建立，\texttt{c}、\texttt{b} 銷毀。 後建立
的物件可能用到先建立的物件，所以先銷毀後建立的。
\item \textbf{\texttt{static} 區域物件}在第一次呼叫 \texttt{f} 時建立，
第二次呼叫不會再建立（一般的區域物件則是每次呼叫都建立、銷毀一次）。
\item \textbf{暫時物件}：\texttt{Tracer("暫時物件");} 建立一個沒有名字的
物件，它在這個運算式結束時（大致就是到分號為止）就銷毀。
\item \texttt{main} 的 \texttt{return} 之後，\texttt{a} 先銷毀，靜態儲存
期的 \texttt{s} 在程式結束時才銷毀。 全域物件則在 \texttt{main} 開始之前
就建立。
\end{itemize}

\texttt{string}、\texttt{vector} 也有解構子，負責把自己配置的記憶體還回
去。 第~3~章說 \texttt{vector} 離開範圍時記憶體會自動還回去，就是解構子
做的。

\subsection{常見錯誤}

\begin{itemize}
\item \textbf{替建構子寫回傳型別}：寫成 \texttt{void Tracer(const string
\&n)} 會編譯失敗，\texttt{g++} 報 \texttt{return type specification for
constructor invalid}。 建構子和解構子都不寫回傳型別，連 \texttt{void} 也
不寫。
\item \textbf{以為物件在最後一次使用後就銷毀}：\texttt{a} 在
\texttt{main} 開頭之後就沒再用到，仍然活到 \texttt{main} 結束。 自動物件
的銷毀時機只看它在哪個區塊，跟最後一次用到它的位置無關。
\end{itemize}

\section{懸置指標與懸置參考}
\label{sec:dangling}

物件的生命週期結束後，原本指向它的指標或綁定它的參考不會自動失效，位址
還在，只是那塊記憶體已經不屬於這個物件。 這樣的指標叫\textbf{懸置指標}
（dangling pointer），參考叫\textbf{懸置參考}（dangling reference）。
透過它們讀寫是未定義行為：可能印出舊值、垃圾，也可能當掉。 結果不固定，
所以這一節的錯誤都用 \textbf{AddressSanitizer}（簡稱 ASan）跑一次。 它是 \texttt{g++} 內建的
檢查工具，編譯時加上 \texttt{-fsanitize=address}，程式執行時一碰到已經
釋放或已經失效的記憶體，就停下來印出報告：

\begin{code}{bash}{用 ASan 編譯並執行}
g++ -std=c++17 -Wall -Wextra -g -fsanitize=address vec_dangling.cpp -o vec_dangling
./vec_dangling
\end{code}

\texttt{-g} 讓報告裡有行號。 加了 ASan 的程式會變慢、變大，只在開發和除錯
時用。 本書的 ASan 輸出取自 Linux 上的 \texttt{g++}，只節錄重點並刪掉目錄
路徑；位址和 \texttt{==} 之間的數字每次執行都不一樣。

\subsection{\texttt{vector} 換記憶體}

第~3~章提過，\texttt{push\_back} 可能讓 \texttt{vector} 換一塊更大的記憶
體。 \texttt{v.capacity()} 是目前這塊記憶體放得下幾個元素，
\texttt{size()} 超過它的時候就要換。

\begin{example}[label=ex:vec-dangling]{cpp}{(vec\_dangling.cpp) push\_back 之後，舊的參考懸置了}
#include <iostream>
#include <vector>
using namespace std;

int main() {
    vector<int> v = {10, 20, 30};
    int &first = v[0];
    int *last = &v[2];
    cout << first << ' ' << *last << ' ' << v.capacity() << '\n';

    v.push_back(40);         // 空間不夠，換一塊更大的記憶體
    cout << v.capacity() << '\n';
    cout << first << '\n';   // first 綁著舊記憶體：未定義行為
    return 0;
}
\end{example}

\begin{code}{text}{用 ASan 編譯後執行（節錄）}
10 30 3
6
=================================================================
==481641==ERROR: AddressSanitizer: heap-use-after-free on address 0x502000000010 at pc 0x57c85395f7e5 bp 0x7ffd17c330a0 sp 0x7ffd17c33090
READ of size 4 at 0x502000000010 thread T0
    #0 0x57c85395f7e4 in main vec_dangling.cpp:13

freed by thread T0 here:
    #6 0x57c85395fcff in std::vector<int, std::allocator<int> >::push_back(int&&) /usr/include/c++/13/bits/stl_vector.h:1299
    #7 0x57c85395f76b in main vec_dangling.cpp:11
\end{code}

報告從上往下看：\texttt{heap-use-after-free} 是讀了一塊已經釋放的 heap
記憶體，\texttt{READ of size 4} 是讀了一個 \texttt{int}，發生在第 $13$ 行
的 \texttt{cout << first}；\texttt{freed by} 說這塊記憶體是第 $11$ 行的
\texttt{push\_back} 釋放的（完整的報告還有一段
\texttt{previously allocated by}，說它是第 $6$ 行建立 \texttt{v} 時配置的）。
所以經過是：\texttt{v} 原本的容量是
$3$，第四個元素放不下，\texttt{push\_back} 配置一塊容量 $6$ 的新記憶體、
把元素搬過去、釋放舊的，\texttt{first} 和 \texttt{last} 還指著舊的那塊。
容量怎麼成長由實作決定，$3$ 和 $6$ 是本書測試環境的結果。 不管在哪個環境，
\texttt{size()} 超過 \texttt{capacity()} 時，所有指向元素的參考、指標和
迭代器（第~9~章）都會失效。

修法有三種：存下標而不是參考（\texttt{v[0]} 每次都重新算位址）；在
\texttt{push\_back} 之後才取參考；或者事先用 \texttt{v.reserve(n)} 把容量
預留到至少 $n$，\texttt{size()} 不超過 $n$ 之前就不會換記憶體。

\subsection{離開區塊的區域變數}

第~\ref{sec:storage}~節提過的情況，寫成程式是這樣：

\begin{code}{cpp}{指標比它指的變數活得久（錯誤示範）}
int *p = nullptr;
{
    int x = 42;
    p = &x;
}                       // x 的生命週期在這裡結束
cout << *p << '\n';     // p 還在範圍內，*p 是未定義行為
\end{code}

ASan 回報 \texttt{stack-use-after-scope}（範圍結束後還在用 stack 上的變數），
報告裡會指出讀到的是 \texttt{'x'}。

第~2~章的「回傳區域變數的參考」是同一類錯誤，只是區塊換成了函式。
\texttt{g++ -Wall} 會先警告 \texttt{reference to local variable 'x'
returned}。 在本書測試的環境下，\texttt{g++} 讓這種函式直接回傳空位址，
ASan 報的是 \texttt{SEGV on unknown address 0x000000000000}。 修法都一樣：
改成傳值回傳。

\subsection{暫時物件的 \texttt{c\_str()}}

要把 \texttt{string} 交給只收 C 字串的函式（例如 \texttt{printf}、
\texttt{fopen}），要呼叫 \texttt{s.c\_str()}，它回傳一個
\texttt{const char *}，指向 \texttt{s} 內部的字元。 這個指標只在
\texttt{s} 還活著、而且沒被修改時有效：

\begin{code}{cpp}{指向暫時字串的指標（錯誤示範）}
string first = "Grace", last = "Hopper";
const char *p = (first + " Brewster " + last).c_str();
printf("%s\n", p);      // 暫時字串已經銷毀
\end{code}

\texttt{first + " Brewster " + last} 是一個暫時物件，第 $2$ 行的分號一到
就銷毀了，\texttt{p} 隨即懸置。 ASan 報 \texttt{heap-use-after-free}，
\texttt{freed by} 指向宣告 \texttt{p} 的那一行。 這段程式 \texttt{-Wall -Wextra} 不會
警告。

暫時物件活到運算式結束，所以在同一個運算式裡用完就沒問題：
\texttt{printf("\%s\textbackslash n", (first + " Brewster " +
last).c\_str());}。 或者先存成變數 \texttt{string full = ...;}，再用
\texttt{full.c\_str()}。

\subsection{常見錯誤}

\begin{itemize}
\item \textbf{程式沒當掉就以為沒事}：未定義行為不保證當掉。 不開 ASan 時，
範例~\ref{ex:vec-dangling} 在本書測試的環境下照樣跑完，最後一行印出一個
每次執行都不同的數字；換個環境，它也可能剛好印出 \texttt{10}，看起來沒有
問題。 懷疑有懸置的時候，加 \texttt{-fsanitize=address} 跑一次。
\item \textbf{把 \texttt{c\_str()} 的結果存起來，之後又修改字串}：
\texttt{const char *p = s.c\_str();} 之後做 \texttt{s += "..."}，
\texttt{s} 可能換了一塊記憶體，\texttt{p} 就懸置了。 要用 C 字串時再呼叫
\texttt{c\_str()}，不要把結果存著。
\end{itemize}

\section{\texttt{new} 與 \texttt{delete}}
\label{sec:new-delete}

\begin{crecap}[malloc 與 free 的規則]
\texttt{malloc(n)} 向 heap 要 $n$ 個位元組，內容是垃圾值，失敗時回傳
\texttt{NULL}；\texttt{free(p)} 把它還回去。 要遵守三條規則：每塊記憶體剛好
\texttt{free} 一次；\texttt{free} 之後不能再用；最後一個指向它的指標消失
之前要 \texttt{free}。 違反第一條是雙重釋放
（double free），第二條是懸置指標，第三條是記憶體洩漏（memory leak）。
\end{crecap}

C++ 動態配置用的是運算子 \texttt{new} 和 \texttt{delete}。 單一物件用
\texttt{new}／\texttt{delete}，例如 \texttt{int *p = new int(42);} 配置一個
初始值 $42$ 的 \texttt{int}，用完 \texttt{delete p;}；陣列用
\texttt{new[]}／\texttt{delete[]}。 下面的 \texttt{Tracer} 跟
範例~\ref{ex:lifetime-order} 的一樣，只是寫得緊湊一點：

\begin{example}[label=ex:new-delete]{cpp}{(new\_delete.cpp) 配置單一物件、陣列與結構}
#include <iostream>
#include <string>
using namespace std;

struct Tracer {
    string name;
    Tracer(const string &n) { name = n; cout << "建立 " << name << '\n'; }
    ~Tracer() { cout << "銷毀 " << name << '\n'; }
};

int main() {
    int n = 5;
    int *a = new int[n];           // n 個 int 的陣列
    for (int i = 0; i < n; i++) {
        a[i] = i * i;
    }
    cout << a[n - 1] << '\n';
    delete[] a;                    // 陣列要用 delete[]

    Tracer *t = new Tracer("heap 上的 t");
    cout << t->name << '\n';
    delete t;                      // 先呼叫解構子，再釋放記憶體
    return 0;
}
\end{example}

\iosample{}{
16 \\
建立 heap 上的 t \\
heap 上的 t \\
銷毀 heap 上的 t \\
}

跟 \texttt{malloc}／\texttt{free} 比較：

\begin{center}
\begin{tabular}{lll}
\toprule
 & \texttt{malloc}／\texttt{free} & \texttt{new}／\texttt{delete} \\
\midrule
要多少 & 位元組數 & 型別（和個數） \\
回傳型別 & \texttt{void *}，C++ 裡要轉型 & 對應型別的指標 \\
建構子、解構子 & 不呼叫 & 呼叫 \\
失敗時 & 回傳 \texttt{NULL} & 丟出例外 \texttt{std::bad\_alloc} \\
陣列 & 同一組函式 & \texttt{new[]} 配 \texttt{delete[]} \\
\bottomrule
\end{tabular}
\end{center}

差別最大的是第三列。 \texttt{Tracer} 裡有一個 \texttt{string}，用
\texttt{malloc} 拿到的只是一塊沒有初始化的位元組，裡面的 \texttt{string}
從來沒有被建構過，拿來用是未定義行為。 \texttt{new} 先配置記憶體再呼叫建
構子，\texttt{delete} 先呼叫解構子再釋放記憶體。 所以
C++ 的程式裡，含有 \texttt{string}、\texttt{vector} 這類成員的結構不能用
\texttt{malloc}。

另外，\texttt{new int[n]} 跟 \texttt{malloc} 一樣不會把元素設成 $0$；要
全部是 $0$ 就寫 \texttt{new int[n]()}。

\subsection{常見錯誤}

\texttt{malloc} 的三種錯誤 \texttt{new} 都會犯，另外還多一種：配錯對。
這些錯誤都編得過，下面附上 ASan 的報告。

\begin{itemize}
\item \textbf{\texttt{new[]} 配 \texttt{delete}}：\texttt{int *a = new
int[5];} 之後寫 \texttt{delete a;}，是未定義行為。 \texttt{g++ -Wall}
這次有警告（\texttt{called on pointer returned from a mismatched
allocation function [-Wmismatched-new-delete]}），ASan 執行時則報
\texttt{alloc-dealloc-mismatch (operator new [] vs operator delete)}。
規則是成對：\texttt{new} 配 \texttt{delete}、\texttt{new[]} 配
\texttt{delete[]}、\texttt{malloc} 配 \texttt{free}，不能交叉。
\item \textbf{重複 \texttt{delete}}：同一個指標 \texttt{delete} 兩次，ASan
報 \texttt{attempting double-free}，並列出兩次釋放各在哪一行。
\texttt{delete nullptr} 什麼都不做，所以
\texttt{delete} 之後把指標設成 \texttt{nullptr} 能擋掉一部分這類錯誤。
\item \textbf{忘了 \texttt{delete}}：最常見的情況是函式有好幾個出口，其中
一個忘了釋放：
\begin{code}{cpp}{(leak.cpp) 提早 return 造成洩漏（錯誤示範，main 省略）}
#include <iostream>
using namespace std;

// 0*0 + 1*1 + ... + (n-1)*(n-1)，超過 100 就回傳 -1
int sumSquares(int n) {
    int *a = new int[n];
    for (int i = 0; i < n; i++) {
        a[i] = i * i;
    }
    int s = 0;
    for (int i = 0; i < n; i++) {
        if (s + a[i] > 100) {
            return -1;           // 提早離開，忘了 delete[]
        }
        s += a[i];
    }
    delete[] a;
    return s;
}
\end{code}
\texttt{main} 裡只有一行 \texttt{cout << sumSquares(5) << ' ' <<
sumSquares(10) << '\textbackslash n';}，印出 \texttt{30 -1}，看不出問題。 ASan 會在程式結束時檢查還有哪些記
憶體沒還：
\begin{code}{text}{ASan 的洩漏報告（節錄）}
==482383==ERROR: LeakSanitizer: detected memory leaks

Direct leak of 40 byte(s) in 1 object(s) allocated from:
    #0 0x70254f4fe6c8 in operator new[](unsigned long) ../../../../src/libsanitizer/asan/asan_new_delete.cpp:98
    #1 0x5c27037fc2d7 in sumSquares(int) leak.cpp:6
\end{code}
$40$ 個位元組就是 \texttt{sumSquares(10)} 那次的 $10$ 個 \texttt{int}。
洩漏檢查（LeakSanitizer）在 Linux 上預設開啟；macOS 內建的編譯器通常不支
援這一項，其他檢查照樣可用。
\end{itemize}

\section{RAII 與智慧指標}
\label{sec:raii}

上一節的洩漏，修法是不要自己 \texttt{new}：把 \texttt{leak.cpp} 第
$6$ 行的 \texttt{int *a = new int[n];} 換成 \texttt{vector<int> a(n);}
（記得 \texttt{\#include <vector>}），刪掉 \texttt{delete[] a;} 那一行，
其他不用改。

\texttt{a} 現在是自動變數，函式不管從哪一個 \texttt{return} 離開，都會執
行 \texttt{a} 的解構子，把記憶體還回去。 用 ASan 跑同樣的 \texttt{main}，
輸出還是 \texttt{30 -1}，沒有報告。

這個做法叫 \textbf{RAII}（Resource Acquisition Is Initialization，取得
資源就是初始化），意思是：

\begin{itemize}
\item 資源（記憶體、開啟的檔案……）交給一個物件管。 物件的建構子取得資源，
解構子釋放資源。
\item 這個物件放在自動變數裡。 區塊結束時解構子執行，資源跟著釋放，不必
自己記得釋放。
\end{itemize}

\texttt{string} 和 \texttt{vector} 都是這樣寫的，所以本書到目前為止沒有
寫過 \texttt{delete}。 第~5~章會自己寫一個這樣的類別。

\subsection{\texttt{unique\_ptr} 與 \texttt{shared\_ptr}}

有時候需要的是 heap 上的單一物件，例如要讓物件活得比建立它的函式更久，
或者第~7~章用指標處理不同型別的物件。 這時用 \texttt{<memory>} 裡的
\textbf{智慧指標}（smart pointer），也就是套用 RAII 的指標：

\begin{itemize}
\item \texttt{unique\_ptr<T>}：\textbf{獨占}一個物件。 同一時間只有一個
\texttt{unique\_ptr} 擁有它，\texttt{unique\_ptr} 銷毀時物件跟著
\texttt{delete}。 用 \texttt{make\_unique<T>(建構子的引數)} 建立。
\item \texttt{shared\_ptr<T>}：好幾個 \texttt{shared\_ptr} \textbf{共同
擁有}一個物件。 它們一起記著目前有幾個擁有者（\texttt{use\_count()}），
最後一個擁有者銷毀或放手時，物件才 \texttt{delete}。 用
\texttt{make\_shared<T>(...)} 建立。
\end{itemize}

兩者用起來都像一般指標：\texttt{*p} 取物件，\texttt{p->name} 取成員，可
以跟 \texttt{nullptr} 比較。

下面的範例沿用同一個 \texttt{Tracer}：

\begin{example}[label=ex:smart-ptr]{cpp}{(smart\_ptr.cpp) unique\_ptr 交出所有權，shared\_ptr 計算擁有者}
#include <iostream>
#include <memory>
#include <string>
using namespace std;

struct Tracer {
    string name;
    Tracer(const string &n) { name = n; cout << "建立 " << name << '\n'; }
    ~Tracer() { cout << "銷毀 " << name << '\n'; }
};

unique_ptr<Tracer> makeTracer(const string &name) {
    unique_ptr<Tracer> p = make_unique<Tracer>(name);
    cout << "makeTracer 結束" << '\n';
    return p;                      // 所有權交給呼叫者
}

int main() {
    unique_ptr<Tracer> u = makeTracer("u");

    shared_ptr<Tracer> s1 = make_shared<Tracer>("共用");
    {
        shared_ptr<Tracer> s2 = s1;  // 兩個 shared_ptr 指向同一個物件
        cout << s1.use_count() << '\n';
    }
    cout << s1.use_count() << '\n';
    s1.reset();                      // 最後一個擁有者放手，物件銷毀
    cout << "main 結束" << '\n';
    return 0;
}
\end{example}

\iosample{}{
建立 u \\
makeTracer 結束 \\
建立 共用 \\
2 \\
1 \\
銷毀 共用 \\
main 結束 \\
銷毀 u \\
}

\begin{itemize}
\item \texttt{makeTracer} 結束時，區域變數 \texttt{p} 銷毀了，物件卻沒有
銷毀：回傳時所有權已經交給 \texttt{main} 的 \texttt{u}。 物件最後在
\texttt{main} 結束、\texttt{u} 銷毀時才跟著銷毀。
\item \texttt{s2 = s1} 之後有兩個擁有者；內層區塊結束，\texttt{s2} 銷毀，
剩一個。 \texttt{s1.reset()} 讓 \texttt{s1} 放手，擁有者變成 $0$，物件
立刻銷毀，所以 \texttt{銷毀 共用} 出現在 \texttt{main 結束} 之前。
\end{itemize}

怎麼選：先用一般變數、\texttt{string}、\texttt{vector}，大部分情況這樣
就夠了；需要 heap 上的單一物件，用 \texttt{unique\_ptr}；好幾個地方共同擁
有一個物件、不確定誰最後用完時，才用 \texttt{shared\_ptr}，它要多維護一
個計數，成本比 \texttt{unique\_ptr} 高。 照這個順序選，一般的程式碼裡就
不需要寫 \texttt{new} 和 \texttt{delete}。

\subsection{常見錯誤}

\begin{itemize}
\item \textbf{複製 \texttt{unique\_ptr}}：\texttt{unique\_ptr<int> q =
p;} 編譯失敗，訊息的開頭是 \texttt{use of deleted function
'std::unique\_ptr<\_Tp, \_Dp>::unique\_ptr(const
std::unique\_ptr<\_Tp, \_Dp>\&)'}，意思是 \texttt{unique\_ptr} 不能複
製。 要把所有權交給
\texttt{q}，寫 \texttt{unique\_ptr<int> q = std::move(p);}，下一節說明。
\item \textbf{自己 \texttt{delete} 智慧指標管的物件}：\texttt{p.get()} 會
回傳裡面的原始指標，給只收一般指標的函式用。 對它 \texttt{delete}，之後
\texttt{p} 銷毀時又 \texttt{delete} 一次，就是雙重釋放。
\end{itemize}

\section{複製與移動}
\label{sec:copy-move}

\texttt{vector<int> b = a;} 是\textbf{複製}（copy）：\texttt{b} 配置一塊
自己的記憶體，把 \texttt{a} 的元素逐一複製過去。 之後兩者各自獨立，改一
個不影響另一個。 代價跟元素個數成正比，\texttt{a} 有一百萬個元素，就要
配置約 $4$ MB 的記憶體、複製一百萬次。 這是第~2、3~章一直建議用
\texttt{const} 參考傳 \texttt{vector} 的原因。

如果 \texttt{a} 之後不會再用，可以不複製，直接把 \texttt{a} 的那塊記憶
體交給 \texttt{b}。 這叫\textbf{移動}（move），寫成
\texttt{std::move(a)}（在 \texttt{<utility>} 裡，習慣上連 \texttt{std::} 一起寫）：

\begin{example}[label=ex:copy-move]{cpp}{(copy\_move.cpp) 複製與移動 vector}
#include <iostream>
#include <utility>
#include <vector>
using namespace std;

int main() {
    vector<int> a(1000000, 7);
    vector<int> b = a;               // 複製：b 有自己的一百萬個 int
    b[0] = 1;
    cout << a[0] << ' ' << b[0] << '\n';

    const int *before = a.data();    // a 的元素放在哪裡
    vector<int> c = std::move(a);    // 移動：c 接手 a 的記憶體
    cout << c.size() << ' ' << (c.data() == before) << '\n';
    return 0;
}
\end{example}

\iosample{}{
7 1 \\
1000000 1 \\
}

\begin{itemize}
\item \texttt{v.data()} 回傳元素所在記憶體的位址。 移動之後
\texttt{c.data()} 跟原本的 \texttt{a.data()} 相同，表示 \texttt{c} 拿走
的就是 \texttt{a} 那塊記憶體，一個元素都沒有複製。 不管有幾個元素，移動
一個 \texttt{vector} 只要搬幾個指標。
\item \texttt{std::move} 本身不搬任何東西，只是把 \texttt{a} 標成「可以
被搬走」；接手記憶體的是建構 \texttt{c} 的那一步。 \texttt{string}、
\texttt{vector}、\texttt{unique\_ptr} 都會看這個標記。 第~5~章會看到，
成員是 \texttt{vector} 或 \texttt{unique\_ptr} 的類別，不用另外寫什麼就能
移動。
\item 被移走的物件（moved-from）處於\textbf{有效但未指定}（valid but
unspecified）的狀態：它還是一個合法的物件，可以重新指定、可以正常銷毀，
但內容沒有保證，除非該型別另外說明（\texttt{unique\_ptr} 就保證被移走後
是 \texttt{nullptr}）。 所以範例移動之後就不再讀 \texttt{a}；要再用，先重新指定，例如
\texttt{a = \{1, 2, 3\};}。 \texttt{string} 也一樣：
\texttt{names.push\_back(std::move(s));} 把 \texttt{s} 的字元搬進
\texttt{vector}，之後不要再讀 \texttt{s}。
\end{itemize}

很多時候移動是自動發生的，不必寫 \texttt{std::move}。 第~2~章說過，
\texttt{a + b} 這種算完就丟的值是 rvalue；暫時物件反正馬上要銷毀，編譯器
會直接移動它。 \texttt{names.push\_back(first + last)} 就會移動那個相加
出來的暫時字串。 函式回傳區域變數時也一樣，編譯器會直接把它交給呼叫者，
這就是第~3~章說回傳 \texttt{vector} 不必擔心複製的原因。

\subsection{常見錯誤}

\begin{itemize}
\item \textbf{移動之後繼續使用}：上面的 \texttt{push\_back} 之後印
\texttt{s}，編譯器不會阻止，但印出什麼沒有保證。 移走之後，除非重
新指定，否則不要再讀它。
\item \textbf{對 \texttt{const} 物件 \texttt{std::move}}：例如
\texttt{const string s = "...";} 之後寫 \texttt{string t =
std::move(s);}，編得過，也沒有
警告，但 \texttt{const} 物件不能被改，所以實際上做的是複製。 想移動的
物件不要宣告成 \texttt{const}。
\item \textbf{回傳時多寫 \texttt{std::move}}：\texttt{return
std::move(result);} 反而讓編譯器沒辦法直接把 \texttt{result} 交給呼叫者，
\texttt{g++ -Wall} 會警告 \texttt{moving a local object in a return
statement prevents copy elision}。 回傳區域變數直接寫
\texttt{return result;}。
\end{itemize}

\section{本章練習}
\label{sec:ch4-exercises}

\begin{drills}[label=drills:ch4]
\item 預測下面程式的輸出。 \texttt{T} 跟本章的 \texttt{Tracer} 一樣，只是
建立時印 \texttt{+} 和名字，銷毀時印 \texttt{-} 和名字，後面各接一個空白。
\begin{code}{cpp}{練習：建立與銷毀的順序}
void f(int k) {
    T x("x" + to_string(k));
    if (k > 1) {
        T y("y");
        return;
    }
    T z("z");
}

int main() {
    T a("a");
    f(1);
    cout << '\n';
    f(2);
    cout << '\n';
    T("tmp");
    cout << '\n';
    return 0;
}
\end{code}
\item 下面的程式用 \texttt{g++ -Wall -Wextra} 編譯沒有任何警告，但它有錯。
錯在哪裡？ 用 ASan 執行會看到什麼？ 怎麼改？
\begin{code}{cpp}{練習：找錯}
#include <cstdio>
#include <string>
using namespace std;

const char *greeting(const string &name) {
    string s = "Hello, " + name + "!";
    return s.c_str();
}

int main() {
    printf("%s\n", greeting("Barbara Liskov"));
    return 0;
}
\end{code}
\item 下面的程式要輸出陣列裡的最大值。 它確實印出 \texttt{9}，但接著在本書測試的環境下
印出 \texttt{free(): double free detected in tcache 2} 並異常結束。 裡
面有三個跟 \texttt{new}／\texttt{delete} 有關的錯誤。 找出來，再用 \texttt{vector}
改寫，讓程式裡沒有任何 \texttt{new} 和 \texttt{delete}。
\begin{code}{cpp}{練習：三個錯誤}
#include <iostream>
using namespace std;

int main() {
    int *a = new int[4]{3, 9, 2, 9};
    int *best = new int(a[0]);
    for (int i = 1; i < 4; i++) {
        if (a[i] > *best) {
            best = new int(a[i]);
        }
    }
    cout << *best << '\n';
    delete a;
    delete best;
    delete best;
    return 0;
}
\end{code}
\item 預測下面程式的輸出。 如果在最後再印一次 \texttt{b.size()}，印出的值
有保證嗎？
\begin{code}{cpp}{練習：複製、移動與智慧指標}
#include <iostream>
#include <memory>
#include <string>
#include <utility>
using namespace std;

int main() {
    string a = "abc";
    string b = a;
    string c = std::move(b);
    unique_ptr<int> p = make_unique<int>(8);
    unique_ptr<int> q = std::move(p);
    cout << a.size() << ' ' << c.size() << '\n';
    cout << (p == nullptr) << ' ' << *q << '\n';
    return 0;
}
\end{code}
\end{drills}

\addcontentsline{toc}{section}{練習解答}
\section*{練習解答}

\begin{enumerate}
\item 輸出四行：\texttt{+a +x1 +z -z -x1}、\texttt{+x2 +y -y -x2}、
\texttt{+tmp -tmp}、\texttt{-a}（每個名字後面都有一個空白；\texttt{-a} 是
\texttt{main} 結束後才印的，後面沒有換行）。 \texttt{f(1)} 依序建立 \texttt{x1}、\texttt{z}，結束時反過來銷毀。
\texttt{f(2)} 從 \texttt{if} 裡提早 \texttt{return}：先離開 \texttt{if}
的區塊，銷毀 \texttt{y}，再離開函式，銷毀 \texttt{x2}；\texttt{z} 根本沒有
建立。 \texttt{tmp} 在同一個敘述裡建立又銷毀；\texttt{a} 在最後的換行之後才銷毀。
\item \texttt{s} 是 \texttt{greeting} 的區域變數，函式結束時銷毀，
\texttt{c\_str()} 回傳的指標指向它已經釋放的記憶體，\texttt{printf} 讀的
是懸置指標。 用 ASan 執行會看到 \texttt{heap-use-after-free}。 改法是讓
\texttt{greeting} 回傳 \texttt{string}，呼叫端寫
\texttt{printf("\%s\textbackslash n", greeting("Barbara
Liskov").c\_str());}，輸出 \texttt{Hello, Barbara Liskov!}。
暫時字串活到分號，\texttt{printf} 用完時它還在。
\item 三個錯誤：(1) 迴圈裡每找到更大的數就 \texttt{new} 一個新的
\texttt{int}，舊的那個沒有 \texttt{delete}，洩漏；(2) \texttt{a} 是
\texttt{new[]} 配置的，要用 \texttt{delete[]}，寫成 \texttt{delete} 是
未定義行為；(3) \texttt{best} 被 \texttt{delete} 兩次。 改寫時陣列換成
\texttt{vector<int> a = \{3, 9, 2, 9\};}，最大值用一個 \texttt{int best}
記就夠了，迴圈裡寫 \texttt{if (a[i] > best) best = a[i];}，整支程式不需要
任何 \texttt{new} 或 \texttt{delete}。
\item 輸出兩行：\texttt{3 3} 和 \texttt{1 8}。 \texttt{a} 沒有被動過，還是
$3$ 個字元；\texttt{c} 拿到 \texttt{b} 從 \texttt{a} 複製來的 \texttt{"abc"}；
\texttt{p} 被移走之後保證是 \texttt{nullptr}，印 \texttt{1}；\texttt{*q}
是 $8$。 \texttt{b} 是被移走的 \texttt{string}，處於有效但未指定的狀態，
所以 \texttt{b.size()} 沒有保證；要用 \texttt{b}，先重新指定或呼叫
\texttt{b.clear()}。
\end{enumerate}

\chapter{類別}
\label{ch:classes}

第~3~章用過 \texttt{s.size()}、\texttt{v.push\_back(x)} 這種寫法，第~4~章
的 \texttt{Tracer} 寫過建構子和解構子。 這一章把這些整理起來，自己寫
\textbf{類別}（class）：把資料和操作它的函式放在一起，並且規定外面的程式
只能透過這些函式改資料。

前三節從 C 的 \texttt{struct} 出發，依序加上成員函式、\texttt{private}
和建構子。 第四節寫一個自己管理記憶體的類別，順便看它被複製時會出什麼
問題。 最後兩節講 \texttt{this}、\texttt{const} 成員函式和
\texttt{static} 成員。

\section{從 \texttt{struct} 到成員函式}
\label{sec:member-func}

\begin{crecap}[struct 與傳結構指標的函式]
C 把相關的資料放進一個 \texttt{struct}，操作它的函式另外寫，第一個參數傳
結構的指標：

\begin{code}{c}{C 的銀行帳戶}
struct Account {
    char owner[32];
    int balance;
};

void deposit(struct Account *a, int amount) {
    a->balance += amount;
}

int withdraw(struct Account *a, int amount) {
    if (amount > a->balance) {
        return 0;
    }
    a->balance -= amount;
    return 1;
}

/* main 裡 */
deposit(&acc, 100);
withdraw(&acc, 30);
acc.balance = -500;        // 沒有人擋得住
\end{code}

\texttt{withdraw} 檢查了餘額夠不夠，但 \texttt{main} 可以直接改
\texttt{acc.balance}，函式裡的檢查等於白做。 \texttt{deposit}、
\texttt{withdraw} 和 \texttt{struct Account} 是一組的，也只靠命名和參數
型別看得出來。
\end{crecap}

C++ 可以把函式寫進結構裡，叫做\textbf{成員函式}（member function）。 呼
叫時寫 \texttt{物件.函式(引數)}，就像第~3~章的 \texttt{s.size()}：

\begin{example}[label=ex:member-func]{cpp}{(member\_func.cpp) 把函式寫進結構}
#include <iostream>
#include <string>
using namespace std;

struct Account {
    string owner;
    int balance;

    void deposit(int amount) {        // 寫在結構裡面
        balance += amount;
    }
    bool withdraw(int amount);        // 只宣告，定義寫在外面
    void print() {
        cout << owner << ": " << balance << '\n';
    }
};

bool Account::withdraw(int amount) {
    if (amount > balance) {
        return false;
    }
    balance -= amount;
    return true;
}

int main() {
    Account a;
    a.owner = "Ada";
    a.balance = 0;
    a.deposit(100);
    cout << a.withdraw(30) << ' ' << a.withdraw(500) << '\n';
    a.print();

    Account b = {"Alan", 50};
    b.deposit(5);
    b.print();
    return 0;
}
\end{example}

\iosample{}{
1 0 \\
Ada: 70 \\
Alan: 55 \\
}

\begin{itemize}
\item \texttt{a.deposit(100)} 對應 C 的 \texttt{deposit(\&acc, 100)}。
呼叫的物件寫在點的前面，不再當成參數傳。
\item 成員函式裡直接寫 \texttt{balance}，指的是「呼叫這個函式的那個物件」
的 \texttt{balance}。 \texttt{a.deposit(100)} 改的是 \texttt{a.balance}，
\texttt{b.deposit(5)} 改的是 \texttt{b.balance}。 不用寫
\texttt{a->balance}。
\item 成員函式可以直接寫在結構裡，也可以在結構裡只寫宣告，定義寫在外面。
寫在外面時，名稱前面要加 \texttt{Account::}，表示這是 \texttt{Account}
的 \texttt{withdraw}。 短的函式通常寫在裡面，長的寫在外面。
\item \texttt{Account b = \{"Alan", 50\};} 依成員的順序給初值，跟 C 一樣。
下一節把成員改成 \texttt{private} 之後就不能這樣寫了，要改用建構子。
\item 成員函式不佔物件的空間，\texttt{sizeof(Account)} 只算 \texttt{owner}
和 \texttt{balance}。
\end{itemize}

\subsection{常見錯誤}

\begin{itemize}
\item \textbf{結構最後的分號}：\texttt{struct Account \{ ... \}} 的右大括號
後面要接分號，跟 C 一樣。 漏了的話，\texttt{g++} 報 \texttt{expected ';'
after struct definition}。
\item \textbf{呼叫成員函式沒有寫物件}：在 \texttt{main} 裡寫
\texttt{deposit(100);}，編譯器報 \texttt{'deposit' was not declared in this
scope}。 \texttt{deposit} 是 \texttt{Account} 的成員，不是全域函式，要寫
\texttt{a.deposit(100)}。
\item \textbf{在外面定義時漏了 \texttt{Account::}}：寫成
\texttt{bool withdraw(int amount) \{ ... \}}，這就是一個跟
\texttt{Account} 無關的普通函式，裡面的 \texttt{balance} 找不到，報
\texttt{'balance' was not declared in this scope}。
\end{itemize}

\section{存取控制：\texttt{public} 與 \texttt{private}}
\label{sec:access}

上一節的 \texttt{Account} 還是誰都能改 \texttt{balance}。 C++ 可以把成員
標成 \texttt{private}（私有），只有這個類別自己的成員函式能存取；標成
\texttt{public}（公開）的成員，外面的程式才能用。

\begin{example}[label=ex:access]{cpp}{(access.cpp) balance 改成 private}
#include <iostream>
using namespace std;

class Account {
public:
    void deposit(int amount) {
        if (amount > 0) {
            balance += amount;
        }
    }
    bool withdraw(int amount) {
        if (amount <= 0 || amount > balance) {
            return false;
        }
        balance -= amount;
        return true;
    }
    int getBalance() {
        return balance;
    }

private:
    int balance = 0;    // 不變量：balance >= 0
};

int main() {
    Account a;
    a.deposit(100);
    a.deposit(-20);                  // 不合理的金額，忽略
    cout << a.withdraw(30) << ' ' << a.withdraw(500) << '\n';
    cout << a.getBalance() << '\n';
    // a.balance = -500;             // 編譯失敗：balance 是 private
    return 0;
}
\end{example}

\iosample{}{
1 0 \\
70 \\
}

\begin{itemize}
\item \texttt{public:} 和 \texttt{private:} 是標籤，從標籤開始到下一個標籤
（或類別結尾）的成員都套用它。 習慣上 \texttt{public} 寫前面，使用這個類
別的人只需要看這一段。
\item \texttt{main} 裡寫 \texttt{a.balance = -500;}，編譯失敗：
\begin{code}{text}{編譯器的錯誤訊息}
error: 'int Account::balance' is private within this context
\end{code}
\item 外面要知道餘額，就透過 \texttt{getBalance()}。 這種只回傳成員值的函
式叫 \textbf{getter}。
\item \texttt{int balance = 0;} 是\textbf{預設成員初始值}：每個
\texttt{Account} 建立時 \texttt{balance} 都從 $0$ 開始。 沒有寫的話，
\texttt{Account a;} 的 \texttt{balance} 跟 C 的區域變數一樣是垃圾值。
\end{itemize}

\texttt{class} 和 \texttt{struct} 只差一件事：沒有寫標籤的成員，
\texttt{struct} 預設是 \texttt{public}，\texttt{class} 預設是
\texttt{private}。 其他功能完全相同，\texttt{struct} 也可以有
\texttt{private} 成員和建構子。 習慣上，單純放資料、所有成員都公開的用
\texttt{struct}（例如第~4~章的 \texttt{Tracer}），有 \texttt{private} 成員
的用 \texttt{class}。

\subsection{不變量}

範例裡的註解寫了「\texttt{balance >= 0}」。 這種對這個類別的每個物件隨時
都成立的條件，叫\textbf{不變量}（invariant）。 \texttt{balance} 是
\texttt{private}，能改它的只有 \texttt{deposit} 和 \texttt{withdraw}，兩
個函式都會檢查金額，所以不管外面怎麼呼叫，餘額都不會變成負的。 要確認這件
事，只要檢查這兩個函式；C 回顧框裡的版本則要檢查每一個碰到 \texttt{acc}
的地方。

\subsection{常見錯誤}

\begin{itemize}
\item \textbf{\texttt{class} 忘了寫 \texttt{public:}}：\texttt{class} 的
成員預設是 \texttt{private}，沒寫標籤的話 \texttt{a.deposit(100)} 編譯失
敗，訊息是 \texttt{'void Account::deposit(int)' is private within this
context}。 把 \texttt{struct} 改成 \texttt{class} 時常常漏掉。
\item \textbf{每個成員都配一組 getter 和 setter}：加一個
\texttt{void setBalance(int x) \{ balance = x; \}}，外面又可以把餘額設成
$-500$，\texttt{private} 就沒有用了。 只提供有意義的操作（存款、提款），
不要提供「把成員改成任意值」的函式。
\end{itemize}

\section{建構子}
\label{sec:ctor}

成員是 \texttt{private} 之後，外面不能再用 \texttt{\{"Alan", 50\}} 給初值，
要靠\textbf{建構子}（constructor）。 第~4~章提過，建構子的名稱跟類別相同，
沒有回傳型別，物件建立時自動呼叫。 一個類別可以有好幾個建構子，跟第~2~章的
函式多載一樣，依引數決定呼叫哪一個。

下面的 \texttt{Time} 表示一天裡的時刻，不變量是
$0 \le \texttt{hour} < 24$、$0 \le \texttt{minute} < 60$：

\begin{example}[label=ex:time-ctor]{cpp}{(time\_ctor.cpp) 預設建構子與帶參數的建構子}
#include <iostream>
using namespace std;

class Time {
public:
    Time() : hour(0), minute(0) {}            // 預設建構子
    Time(int h, int m) : hour(h), minute(m) {  // 帶參數的建構子
        normalize();
    }
    void add(int minutes) {
        minute += minutes;
        normalize();
    }
    void print() {
        cout << (hour < 10 ? "0" : "") << hour << ':'
             << (minute < 10 ? "0" : "") << minute << '\n';
    }

private:
    int hour, minute;    // 不變量：0 <= hour < 24，0 <= minute < 60
    void normalize() {
        hour = (hour + minute / 60) % 24;
        minute %= 60;
    }
};

int main() {
    Time a;
    Time b(9, 45);
    Time c(23, 130);
    a.print();
    b.print();
    c.print();

    b.add(30);
    b.print();
    return 0;
}
\end{example}

\iosample{}{
00:00 \\
09:45 \\
01:10 \\
10:15 \\
}

\begin{itemize}
\item 沒有參數的 \texttt{Time()} 叫\textbf{預設建構子}（default
constructor）。 \texttt{Time a;} 會用到它，\texttt{vector<Time> v(2);}
建立兩個元素時也是。
\item \texttt{Time b(9, 45);} 呼叫 \texttt{Time(int h, int m)}。 也可以寫
成 \texttt{Time b\{9, 45\};}，效果一樣。
\item \texttt{Time c(23, 130);} 給的分鐘超過 $60$，建構子呼叫
\texttt{normalize()} 調整成 \texttt{01:10}。 建構子結束時不變量就成立，之
後只有 \texttt{add} 會改時間，\texttt{add} 也會呼叫 \texttt{normalize()}。
（這裡假設傳進來的數字都不是負的。）
\item \texttt{normalize} 放在 \texttt{private}。 它只是內部用的工具，外面
不需要呼叫，成員函式也可以是 \texttt{private}。
\end{itemize}

\subsection{成員初始化列表}

建構子參數列後面、本體前面的 \texttt{: hour(h), minute(m)} 叫\textbf{成員
初始化列表}（member initializer list）。 物件的成員在建構子本體執行
\emph{之前}就建立好了，初始化列表決定它們用什麼值建立。 也可以寫成在本
體裡寫 \texttt{hour = h; minute = m;}。 對 \texttt{int} 來說，兩種寫法
結果相同。 但有兩種成員一定要用初始化列表：
\texttt{const} 成員和參考成員。 它們建立之後就不能再指定，只能在建立的
那一刻給值。

\begin{example}[label=ex:init-list]{cpp}{(init\_list.cpp) const 成員與參考成員}
#include <iostream>
#include <string>
#include <vector>
using namespace std;

class Recorder {
public:
    Recorder(int n, vector<string> &out) : id(n), log(out) {}
    void record(const string &msg) {
        log.push_back(to_string(id) + ": " + msg);
    }

private:
    const int id;              // 建立後不能改
    vector<string> &log;       // 綁著呼叫者的 vector
};

int main() {
    vector<string> log;
    Recorder r1(1, log), r2(2, log);
    r1.record("start");
    r2.record("load");
    r1.record("stop");
    for (const string &s : log) {
        cout << s << '\n';
    }
    return 0;
}
\end{example}

\iosample{}{
1: start \\
2: load \\
1: stop \\
}

\texttt{r1}、\texttt{r2} 的 \texttt{log} 都綁著 \texttt{main} 的同一個
\texttt{log}，所以三筆紀錄依序進到同一個 \texttt{vector}。 要是把建構子
改成在本體裡寫 \texttt{id = n; log = out;}，編譯失敗：

\begin{code}{text}{編譯器的錯誤訊息（節錄）}
error: uninitialized const member in 'const int' [-fpermissive]
error: uninitialized reference member in 'class std::vector<std::__cxx11::basic_string<char> >&' [-fpermissive]
error: assignment of read-only member 'Recorder::id'
\end{code}

進到本體時，\texttt{id} 和 \texttt{log} 已經建立了，卻沒有給值；本體裡的
\texttt{id = n} 是指定，\texttt{const} 不能指定。 參考也一樣，第~2~章說
過參考建立時就要綁定，之後的 \texttt{=} 是改它綁著的那個變數，不是重新綁定。

\texttt{string}、\texttt{vector} 這類成員用初始化列表也比較好：在本體裡
指定，等於先建一個空的，再指定一次。 所以習慣上一律用初始化列表，本體只
放建構時要做的其他事（例如 \texttt{normalize()}）。

\subsection{常見錯誤}

\begin{itemize}
\item \textbf{寫了建構子，預設建構子就沒了}：一個建構子都沒寫的時候，編譯器
會自動產生一個什麼都不做的預設建構子。 只要自己寫了任何一個建構子，編譯
器就不再產生。 把範例裡的 \texttt{Time()} 刪掉，\texttt{Time a;} 和
\texttt{vector<Time> v(2);} 都編譯失敗：
\begin{code}{text}{編譯器的錯誤訊息}
error: no matching function for call to 'Time::Time()'
\end{code}
需要的話就自己寫一個預設建構子。
\item \textbf{\texttt{Time a();}}：想用預設建構子建立 \texttt{a}，多寫了
括號。 這行會被當成「宣告一個函式 \texttt{a}，沒有參數，回傳
\texttt{Time}」。 \texttt{g++} 會警告 \texttt{empty parentheses were
disambiguated as a function declaration}，接著 \texttt{a.print()} 報
\texttt{request for member 'print' in 'a', which is of non-class type
'Time()'}。 用預設建構子時不加括號，寫 \texttt{Time a;}，或者寫
\texttt{Time a\{\};}。
\item \textbf{以為成員依初始化列表的順序建立}：成員一律依照在類別裡
\emph{宣告}的順序建立，跟初始化列表怎麼排無關。
\begin{code}{cpp}{cells 比 rows、cols 先建立（錯誤示範）}
class Grid {
public:
    Grid(int r, int c) : rows(r), cols(c), cells(rows * cols, 0) {}
private:
    vector<int> cells;
    int rows, cols;
};
\end{code}
\texttt{cells} 宣告在前，最先建立，這時 \texttt{rows}、\texttt{cols} 還
沒有值。 \texttt{-Wall} 會警告 \texttt{'Grid::cols' will be initialized
after}。 改法是把初始化列表的順序排成跟宣告一樣，或者直接用參數寫
\texttt{cells(r * c, 0)}。
\end{itemize}

\section{解構子：自己管理資源的類別}
\label{sec:dtor}

\textbf{解構子}（destructor）的名稱是 \texttt{\textasciitilde} 加上類別
名稱，沒有參數，物件銷毀時自動呼叫。 第~4~章說過 RAII 的做法：建構子取得
資源，解構子釋放資源。 這一節照這個做法寫一個簡化版的 \texttt{vector}：
建立時給長度，元素初始值是 $0$。 為了看出配置和釋放的時機，建構子和解構子
各印一行。

\begin{example}[label=ex:int-array]{cpp}{(int\_array.cpp) 建構子 new[]，解構子 delete[]}
#include <iostream>
using namespace std;

class IntArray {
public:
    IntArray(int n) : len(n), data(new int[n]()) {
        cout << "配置 " << len << " 個 int" << '\n';
    }
    ~IntArray() {
        delete[] data;
        cout << "釋放 " << len << " 個 int" << '\n';
    }
    int size() {
        return len;
    }
    int get(int i) {
        return data[i];
    }
    void set(int i, int x) {
        if (i >= 0 && i < len) {
            data[i] = x;
        }
    }

private:
    int len;
    int *data;
};

int sumSquares(int n) {
    IntArray a(n);
    for (int i = 0; i < n; i++) {
        a.set(i, i * i);
    }
    int s = 0;
    for (int i = 0; i < a.size(); i++) {
        s += a.get(i);
    }
    return s;
}

int main() {
    cout << sumSquares(4) << '\n';
    {
        IntArray b(3);
        b.set(0, 7);
        b.set(5, 1);             // 超出範圍，set 不做事
        cout << b.get(0) << ' ' << b.size() << '\n';
    }
    cout << "main 結束" << '\n';
    return 0;
}
\end{example}

\iosample{}{
配置 4 個 int \\
釋放 4 個 int \\
14 \\
配置 3 個 int \\
7 3 \\
釋放 3 個 int \\
main 結束 \\
}

\begin{itemize}
\item 用 \texttt{IntArray} 的人不會寫到 \texttt{new} 或 \texttt{delete}。
\texttt{sumSquares} 結束時 \texttt{a} 銷毀，解構子把記憶體還回去，所以
「釋放 4 個 int」印在 \texttt{14} 之前：\texttt{return s;} 先算好回傳值，
接著銷毀區域變數，最後 \texttt{main} 才印出結果。
\item \texttt{data} 是 \texttt{private}，外面拿不到這個指標，也就沒辦法
\texttt{delete} 它，或在物件銷毀後還留著它。
\item \texttt{set} 會檢查下標，越界時什麼都不做；\texttt{get} 沒有檢查，
跟 \texttt{vector} 的 \texttt{[]} 一樣。
\item 第~4~章的 \texttt{leak.cpp} 有好幾個 \texttt{return}，其中一個漏了
\texttt{delete[]}。 寫成類別之後，釋放只寫在解構子裡這一個地方。
\end{itemize}

\subsection{複製會出問題}

\texttt{IntArray} 還有一個漏洞。 把範例~\ref{ex:int-array} 的
\texttt{main} 換成下面這段，複製一個 \texttt{IntArray}：

\begin{code}{cpp}{複製擁有原始指標的物件（錯誤示範）}
int main() {
    IntArray a(3);
    IntArray b = a;          // 複製：b.data 和 a.data 是同一個位址
    b.set(0, 99);
    cout << a.get(0) << '\n';
    return 0;
}                            // b、a 各 delete[] 一次同一塊記憶體
\end{code}

\begin{code}{text}{執行結果（本書測試的環境）}
配置 3 個 int
99
釋放 3 個 int
free(): double free detected in tcache 2
\end{code}

\texttt{IntArray b = a;} 用的是編譯器自動產生的\textbf{複製建構子}（copy
constructor），它把成員一個一個複製過去。 \texttt{data} 是指標，複製的是
位址，不是陣列，所以 \texttt{a} 和 \texttt{b} 共用同一塊記憶體：改
\texttt{b} 的元素，\texttt{a.get(0)} 也變成 $99$。 配置只發生一次，
\texttt{main} 結束時卻有兩個解構子：\texttt{b} 先釋放，\texttt{a} 再釋放
同一塊，程式在這裡中止。 用 ASan 跑，報的是 \texttt{attempting
double-free}。

已經存在的兩個物件之間寫 \texttt{b = a;}，用的是自動產生的\textbf{複製指定
運算子}（copy assignment operator），一樣只複製指標，問題相同。

要讓複製正確，可以自己寫複製建構子和複製指定運算子，配置一塊新的陣列再把
元素抄過去。 一個類別需要自己寫解構子時，通常也需要自己寫這兩個，這叫
\textbf{三法則}（rule of three）；再加上第~4~章的移動（移動建構子和移動
指定運算子）共五個，叫\textbf{五法則}（rule of five）。 本書不細談這些
函式的寫法，建議的做法是成員不要用原始指標，改用 \texttt{vector}
或 \texttt{unique\_ptr}。

\begin{code}{cpp}{(copy\_fixed.cpp) 成員改用 vector（節錄）}
class IntArray {
public:
    IntArray(int n) : data(n) {}
    int size() {
        return (int)data.size();
    }
    // get 不變；set 裡的 len 改成 size()

private:
    vector<int> data;        // 不必寫解構子
};
\end{code}

用同樣的 \texttt{main}，最後一行改成印 \texttt{a.get(0)} 和
\texttt{b.get(0)}，輸出 \texttt{0 99}。 自動產生的複製建構子還是逐一複製
成員，但複製 \texttt{vector} 會複製所有元素，所以 \texttt{b} 有自己的一份。
銷毀時兩個 \texttt{vector} 各自釋放自己的記憶體，\texttt{IntArray} 不用寫
解構子。 自動產生的移動也會移動 \texttt{data}，所以這個類別可以直接用
\texttt{std::move}。

成員改用 \texttt{unique\_ptr} 的話，因為 \texttt{unique\_ptr} 不能複製，
整個類別也不能複製，\texttt{IntArray b = a;} 編譯失敗（\texttt{use of
deleted function}），但可以移動。 編譯失敗比執行時重複釋放容易處理。

\subsection{常見錯誤}

\begin{itemize}
\item \textbf{傳值傳進函式}：\texttt{void print(IntArray a)} 每次呼叫都會
複製一個 \texttt{IntArray}。 用原始指標的版本，函式結束時複本的解構子就把
陣列釋放了，呼叫端的物件從此懸置。 傳類別物件時照第~2~章的原則用
\texttt{const} 參考，就不會複製（\texttt{const} 參考要搭配
第~\ref{sec:this-const}~節的 \texttt{const} 成員函式）。
\item \textbf{解構子用錯 \texttt{delete}}：建構子用 \texttt{new int[n]}，
解構子就要用 \texttt{delete[]}。 寫成 \texttt{delete data;} 是第~4~章說
的配錯對。
\item \textbf{以為解構子要自己呼叫}：解構子在物件銷毀時自動執行，不要寫
\texttt{a.\textasciitilde IntArray();}。 自己呼叫一次，離開區塊時又自動
呼叫一次，又是重複釋放。
\end{itemize}

\section{\texttt{this} 與 \texttt{const} 成員函式}
\label{sec:this-const}

成員函式裡寫 \texttt{balance}，怎麼知道是哪個物件的？ 每個成員函式都有一
個看不到的參數 \texttt{this}，是指向呼叫物件的指標。 \texttt{a.deposit(100)}
呼叫時，\texttt{this} 就是 \texttt{\&a}，跟 C 版本的
\texttt{deposit(\&acc, 100)} 做的事一樣。 函式裡的 \texttt{balance} 就是
\texttt{this->balance} 的簡寫。

平常不用寫 \texttt{this}，有兩種情況會用到：

\begin{itemize}
\item 參數名稱跟成員相同時，\texttt{balance} 指的是參數，要寫
\texttt{this->balance} 才是成員。 初始化列表不受影響：
\texttt{owner(owner)} 括號外的是成員，括號內的是參數。
\item 回傳 \texttt{*this}，也就是呼叫的物件本身。 回傳型別寫成參考，呼叫
就能串起來，例如 \texttt{a.deposit(100).deposit(50)}。 第~2~章的
\texttt{cout <<} 能串下去也是這個道理。
\end{itemize}

\subsection{\texttt{const} 成員函式}

第~2~章說過，只讀的參數用 \texttt{const} 參考。 對類別物件這樣做，會遇到
一個問題：

\begin{code}{cpp}{透過 const 參考呼叫 getBalance（編譯失敗）}
void report(const Account &a) {
    cout << a.getBalance() << '\n';
}
\end{code}

\begin{code}{text}{編譯器的錯誤訊息}
error: passing 'const Account' as 'this' argument discards qualifiers [-fpermissive]
\end{code}

\texttt{a} 是 \texttt{const}，但編譯器不知道 \texttt{getBalance} 會不會
改它。 訊息裡的 \texttt{this argument} 指的就是 \texttt{this} 這個看不到
的參數：\texttt{a} 是 \texttt{const}，\texttt{getBalance} 的
\texttt{this} 卻不是指向 \texttt{const} 的指標，傳進去會丟掉
\texttt{const}。

在參數列後面加 \texttt{const}，宣告這個函式不會修改物件，就是\textbf{const
成員函式}。 這時 \texttt{this} 是指向 \texttt{const} 的指標，
\texttt{const} 物件和 \texttt{const} 參考都能呼叫它：

\begin{example}[label=ex:this-const]{cpp}{(this\_const.cpp) this、串接呼叫與 const 成員函式}
#include <iostream>
#include <string>
using namespace std;

class Account {
public:
    Account(const string &owner) : owner(owner), balance(0) {}

    Account &deposit(int amount) {
        if (amount > 0) {
            this->balance += amount;   // 等於 balance += amount
        }
        return *this;                  // 回傳呼叫它的物件本身
    }
    int getBalance() const {
        return balance;
    }
    const string &getOwner() const {
        return owner;
    }

private:
    string owner;
    int balance;
};

void report(const Account &a) {
    cout << a.getOwner() << ' ' << a.getBalance() << '\n';
}

int main() {
    Account a("Ada");
    a.deposit(100).deposit(50);
    report(a);

    const Account b("Alan");
    report(b);
    // b.deposit(10);                  // 編譯失敗：b 是 const
    return 0;
}
\end{example}

\iosample{}{
Ada 150 \\
Alan 0 \\
}

\begin{itemize}
\item \texttt{deposit} 回傳 \texttt{Account \&}，\texttt{a.deposit(100)}
的結果就是 \texttt{a} 本身，可以接著 \texttt{.deposit(50)}。 回傳型別如
果寫成 \texttt{Account}（沒有 \texttt{\&}），第二次 \texttt{deposit} 存進
的是一個暫時的複本，\texttt{a} 只多了 $100$。
\item \texttt{getBalance} 和 \texttt{getOwner} 不改物件，宣告成
\texttt{const}。 \texttt{report} 收 \texttt{const} 參考，只能呼叫這兩個。
\item \texttt{getOwner} 回傳 \texttt{const string \&}，不複製字串，呼叫端
也不能透過它改 \texttt{owner}。
\item \texttt{const Account b} 是 \texttt{const} 物件，建立之後只能呼叫
\texttt{const} 成員函式，\texttt{b.deposit(10)} 會出現跟上面一樣的錯誤訊
息。
\end{itemize}

不改物件的成員函式都要加上 \texttt{const}。 漏加的話，自己寫的
時候不會發現，等到別人用 \texttt{const} 參考傳你的物件才會編譯失敗。

\subsection{常見錯誤}

\begin{itemize}
\item \textbf{參數名稱跟成員相同}：
\begin{code}{cpp}{參數蓋掉成員（錯誤示範）}
void deposit(int balance) {
    balance += balance;     // 兩個都是參數
}
\end{code}
函式裡的 \texttt{balance} 都是參數，成員沒有改到。 \texttt{-Wall -Wextra}
不會警告，要另外加 \texttt{-Wshadow} 才會看到 \texttt{declaration of
'balance' shadows a member of 'Account'}。 參數換個名字，或者寫
\texttt{this->balance += balance;}。
\item \textbf{在 \texttt{const} 成員函式裡改成員}：在 \texttt{getBalance()
const} 裡寫 \texttt{balance = 0;}，報 \texttt{assignment of member
'Account::balance' in read-only object}。 要改成員的函式就不要宣告成
\texttt{const}。
\item \textbf{\texttt{const} 函式回傳成員的一般參考}：
\texttt{string \&getOwner() const \{ return owner; \}} 編譯失敗，訊息是
\texttt{binding reference of type 'std::string\&' ... discards
qualifiers}。 \texttt{const} 函式裡 \texttt{owner} 是 \texttt{const}，不
能綁到非 \texttt{const} 參考上；回傳型別要寫 \texttt{const string \&}。
\item \textbf{\texttt{const} 寫錯位置}：\texttt{const int getBalance()} 的
\texttt{const} 修飾的是回傳值，這不是 \texttt{const} 成員函式。
\texttt{const} 要寫在參數列的右括號後面。
\end{itemize}

\section{\texttt{static} 成員}
\label{sec:static-member}

一般的成員是每個物件各有一份。 加上 \texttt{static} 的成員則屬於整個類別，
所有物件共用一份，不建立物件也存在。 下面的 \texttt{Ticket} 用
\texttt{static} 成員發號碼，並記錄目前有幾張票還在：

\begin{example}[label=ex:static-member]{cpp}{(static\_member.cpp) 所有物件共用的計數}
#include <iostream>
using namespace std;

class Ticket {
public:
    Ticket() : id(nextId) {
        nextId++;
        alive++;
    }
    ~Ticket() {
        alive--;
    }
    int getId() const {
        return id;
    }
    static int aliveCount() {
        return alive;
    }

private:
    int id;                        // 每張票各有一個
    inline static int nextId = 1;  // 整個類別共用一個
    inline static int alive = 0;
};

int main() {
    Ticket a, b;
    cout << a.getId() << ' ' << b.getId() << ' ' << Ticket::aliveCount() << '\n';
    {
        Ticket c;
        cout << c.getId() << ' ' << Ticket::aliveCount() << '\n';
    }
    Ticket d;
    cout << d.getId() << ' ' << Ticket::aliveCount() << '\n';
    return 0;
}
\end{example}

\iosample{}{
1 2 2 \\
3 3 \\
4 3 \\
}

\begin{itemize}
\item \texttt{nextId} 和 \texttt{alive} 只有一份。 每建立一張票，建構子
把目前的 \texttt{nextId} 給這張票當 \texttt{id}，再把 \texttt{nextId} 加
一；\texttt{c} 銷毀時解構子把 \texttt{alive} 減一，所以 \texttt{d} 拿到
$4$，但還在的票是 $3$ 張。 這跟第~4~章 \texttt{nextTicket} 裡的
\texttt{static} 區域變數是同一個想法，只是放在類別裡。
\item \texttt{static} 成員變數要在類別裡寫 \texttt{inline static}，才能直接
給初始值（C++17 起）。
\item \texttt{aliveCount} 是 \texttt{static} 成員函式，用
\texttt{Ticket::aliveCount()} 呼叫，不需要物件。 它沒有 \texttt{this}，所
以只能用 \texttt{static} 成員，不能用 \texttt{id}。
\end{itemize}

\texttt{static} 成員適合放跟個別物件無關、整個類別共用的東西，例如計數、
共用的設定值。 不要拿它來放每個物件應該各有一份的資料。

\subsection{常見錯誤}

\begin{itemize}
\item \textbf{漏了 \texttt{inline}}：在類別裡寫
\texttt{static int alive = 0;}，編譯失敗：\texttt{ISO C++ forbids in-class
initialization of non-const static member 'Ticket::alive'}。 C++17 以前的
寫法是在類別裡只寫 \texttt{static int alive;}，再到類別外面寫一行定義
\texttt{int Ticket::alive = 0;}。 只寫了類別裡那行、忘了外面那行，編譯會過，
連結時報 \texttt{undefined reference to `Ticket::alive'}。
\item \textbf{\texttt{static} 成員函式用了一般成員}：在 \texttt{aliveCount}
裡寫 \texttt{return id;}，報 \texttt{invalid use of member 'Ticket::id' in
static member function}。 沒有物件，就不知道是哪一張票的 \texttt{id}。
\end{itemize}

\section{本章練習}
\label{sec:ch5-exercises}

\begin{drills}[label=drills:ch5]
\item 預測下面程式的輸出。
\begin{code}{cpp}{練習：建構子、串接與解構子}
#include <iostream>
using namespace std;

class Box {
public:
    Box() : w(1), h(1) { cout << "A "; }
    Box(int s) : w(s), h(s) { cout << "B "; }
    Box(int w, int h) : w(w), h(h) { cout << "C "; }
    ~Box() { cout << "~" << w * h << ' '; }
    Box &grow(int d) {
        w += d;
        h += d;
        return *this;
    }
    int area() const { return w * h; }

private:
    int w, h;
};

int main() {
    Box a;
    Box b(4);
    Box c(2, 5);
    cout << '\n';
    cout << a.grow(1).grow(1).area() << ' ' << c.area() << '\n';
    return 0;
}
\end{code}
\item 下面的程式編譯失敗，\texttt{g++} 報了五個錯誤。 找出原因並改正，讓 \texttt{show(p)} 印出
\texttt{4 6}，\texttt{q} 是原點 $(0, 0)$。
\begin{code}{cpp}{練習：找錯}
#include <iostream>
using namespace std;

class Point {
public:
    Point(int x, int y) : x(x), y(y) {}
    int getX() { return x; }
    void move(int dx, int dy) const {
        x += dx;
        y += dy;
    }

private:
    int x, y;
};

void show(const Point &p) {
    cout << p.getX() << ' ' << p.y << '\n';
}

int main() {
    Point p(1, 2);
    p.move(3, 4);
    show(p);
    Point q;
    return 0;
}
\end{code}
\item 寫一個類別 \texttt{Stats}，記錄加進來的整數，提供這些成員函式：
\texttt{add(x)} 加入一個整數；\texttt{count()}、\texttt{min()}、
\texttt{max()} 回傳個數、最小值、最大值；\texttt{mean()} 回傳平均
（\texttt{double}）。 不要把所有數字存起來，只存需要的量。 \texttt{main}
讀整數直到輸入結束，輸出個數、最小值、最大值、平均，以一個空白隔開；沒有
任何輸入時輸出 \texttt{沒有資料}。 輸入 \texttt{7 -2 10 5} 時應輸出
\texttt{4 -2 10 5}。 哪些函式應該是 \texttt{const}？
\item 寫一個 \texttt{Account}，建構子收擁有者的名字。 每個帳戶有一個編號，
從 $1$ 開始依建立順序編，建立後不能改；另外用 \texttt{static} 成員記錄
總共開過幾個帳戶（銷毀的也算）。 提供 \texttt{getId()}、
\texttt{getOwner()} 和 \texttt{static} 的 \texttt{openedCount()}。 寫完之
後試試 \texttt{a = b;}，會發生什麼事？
\end{drills}

\addcontentsline{toc}{section}{練習解答}
\section*{練習解答}

\begin{enumerate}
\item 輸出三行：\texttt{A B C }、\texttt{9 10}、
\texttt{\textasciitilde 10 \textasciitilde 16 \textasciitilde 9 }（最後一
行沒有換行）。 \texttt{a}、\texttt{b}、\texttt{c} 依序呼叫三個不同的建構
子。 \texttt{grow} 回傳 \texttt{*this} 的參考，兩次 \texttt{grow(1)} 都
作用在 \texttt{a} 上，邊長從 $1$ 變成 $3$，面積 $9$。 \texttt{main} 結束
時以相反順序銷毀：\texttt{c}（$2 \times 5$）、\texttt{b}（$4 \times 4$）、
\texttt{a}（$3 \times 3$）。
\item 五個錯誤訊息來自四個問題：
(1) \texttt{move} 會改 \texttt{x}、\texttt{y}，不能是 \texttt{const}，報
兩次 \texttt{assignment of member ... in read-only object}；
(2) \texttt{getX} 沒有 \texttt{const}，\texttt{show} 透過 \texttt{const}
參考呼叫它，報 \texttt{passing 'const Point' as 'this' argument discards
qualifiers}；
(3) \texttt{p.y} 是 \texttt{private}，\texttt{show} 不是成員函式，不能存
取；
(4) 寫了帶參數的建構子，就沒有預設建構子，\texttt{Point q;} 報
\texttt{no matching function for call to 'Point::Point()'}。
改法：拿掉 \texttt{move} 的 \texttt{const}；\texttt{getX} 加上
\texttt{const}；新增一個 \texttt{int getY() const \{ return y; \}}，
\texttt{show} 改用 \texttt{p.getY()}，不要把 \texttt{y} 改成
\texttt{public}；加一個預設建構子 \texttt{Point() : x(0), y(0) \{\}}。
改完後 \texttt{show(p)} 印出 \texttt{4 6}，\texttt{show(q)} 印出
\texttt{0 0}。
\item 一種寫法：
\begin{code}{cpp}{練習 3 參考解答}
#include <iostream>
using namespace std;

class Stats {
public:
    void add(int x) {
        if (n == 0 || x < lo) {
            lo = x;
        }
        if (n == 0 || x > hi) {
            hi = x;
        }
        sum += x;
        n++;
    }
    int count() const { return n; }
    int min() const { return lo; }
    int max() const { return hi; }
    double mean() const { return n == 0 ? 0 : (double)sum / n; }

private:
    int n = 0;
    long long sum = 0;
    int lo = 0, hi = 0;
};

int main() {
    Stats s;
    int x;
    while (cin >> x) {
        s.add(x);
    }
    if (s.count() == 0) {
        cout << "沒有資料" << '\n';
        return 0;
    }
    cout << s.count() << ' ' << s.min() << ' ' << s.max() << ' ' << s.mean() << '\n';
    return 0;
}
\end{code}
只要存個數、總和、目前的最小值和最大值。 除了 \texttt{add}，其他四個函式
都不改物件，宣告成 \texttt{const}。 \texttt{n == 0} 時 \texttt{lo}、
\texttt{hi} 沒有意義，所以第一個數字一定要直接存進去；\texttt{main} 也先
檢查 \texttt{count()} 再印。 總和用 \texttt{long long}，避免很多個大數字
相加時溢位。
\item 一種寫法：
\begin{code}{cpp}{練習 4 參考解答}
#include <iostream>
#include <string>
using namespace std;

class Account {
public:
    Account(const string &owner) : id(opened + 1), owner(owner) {
        opened++;
    }
    int getId() const { return id; }
    const string &getOwner() const { return owner; }
    static int openedCount() { return opened; }

private:
    const int id;
    string owner;
    inline static int opened = 0;
};

int main() {
    Account a("Ada");
    Account b("Alan");
    cout << a.getId() << ' ' << a.getOwner() << '\n';
    cout << b.getId() << ' ' << b.getOwner() << '\n';
    {
        Account c("Grace");
        cout << c.getId() << ' ' << Account::openedCount() << '\n';
    }
    cout << Account::openedCount() << '\n';
    return 0;
}
\end{code}
輸出 \texttt{1 Ada}、\texttt{2 Alan}、\texttt{3 3}、\texttt{3}。
\texttt{id} 是 \texttt{const} 成員，只能在初始化列表給值。
\texttt{opened} 只增不減，所以 \texttt{c} 銷毀後還是 $3$。
加上 \texttt{a = b;} 會編譯失敗：
\begin{code}{text}{編譯器的錯誤訊息}
error: use of deleted function 'Account& Account::operator=(const Account&)'
error: non-static const member 'const int Account::id', cannot use default assignment operator
\end{code}
自動產生的複製指定運算子要逐一指定成員，\texttt{const} 的 \texttt{id} 不
能指定，所以編譯器不產生它。 這也符合需求：兩個帳戶的編號本來就不該變成
一樣。 \texttt{Account c = a;} 這種複製建構還是編得過，\texttt{c} 會拿到跟
\texttt{a} 一樣的編號；要連這個也禁止，得自己處理複製建構子，本書不細談。
\end{enumerate}

\chapter{運算子多載與複製}
\label{ch:operators-copy}

第~3~章的 \texttt{string} 可以寫 \texttt{a + b}、\texttt{a == b}、
\texttt{s[0]}，\texttt{cout << x} 什麼型別都能印。 這些運算子都是函式，
只是呼叫的寫法不一樣。 這一章先替自己的類別寫運算子：算術、比較、輸出和
下標。

後半章處理第~5~章留下的問題：擁有原始指標的 \texttt{IntArray} 一複製就
重複釋放。 我們會自己寫複製建構子和複製指定運算子，再看移動，最後整理成
三條規則：三法則、五法則和零法則。

\section{運算子也是函式}
\label{sec:op-basics}

\begin{crecap}[對結構做運算的函式]
C 的運算子只能用在內建型別上。 自己定義的結構要相加，只能寫一個函式：

\begin{code}{c}{C 的分數}
struct Frac {
    int num, den;
};

struct Frac frac_add(struct Frac a, struct Frac b) {
    struct Frac r = {a.num * b.den + b.num * a.den, a.den * b.den};
    return r;          /* 約分省略 */
}

struct Frac frac_mul(struct Frac a, struct Frac b) {
    struct Frac r = {a.num * b.num, a.den * b.den};
    return r;
}

/* main 裡 */
struct Frac d = frac_add(frac_mul(a, b), c);   /* a * b + c */
\end{code}

這樣寫沒有錯，但式子一長，就得對著括號慢慢拆，才看得出算的是
\texttt{a * b + c}。
\end{crecap}

C++ 可以替自己的類別定義運算子，叫\textbf{運算子多載}（operator
overloading）。 運算子就是一個名稱特別的函式：\texttt{operator+} 是
\texttt{+}，\texttt{operator+=} 是 \texttt{+=}。 寫了 \texttt{operator+}
之後，\texttt{a + b} 就等於呼叫 \texttt{operator+(a, b)}。

下面的 \texttt{Fraction} 表示一個分數。 不變量是分母大於 $0$，而且一定是
約分過的。 約分用 \texttt{<numeric>} 裡的 \texttt{gcd}（最大公因數，
C++17 起才有）：

\begin{example}[label=ex:fraction]{cpp}{(fraction.cpp) 分數的加法與乘法}
#include <iostream>
#include <numeric>
using namespace std;

class Fraction {
public:
    Fraction(int n = 0, int d = 1) : num(n), den(d) {
        normalize();
    }
    Fraction &operator+=(const Fraction &r) {
        num = num * r.den + r.num * den;
        den = den * r.den;
        normalize();
        return *this;
    }
    Fraction &operator*=(const Fraction &r) {
        num = num * r.num;
        den = den * r.den;
        normalize();
        return *this;
    }
    int getNum() const { return num; }
    int getDen() const { return den; }

private:
    int num, den;    // 不變量：den > 0，而且 num/den 已經約分
    void normalize() {
        if (den < 0) {
            num = -num;
            den = -den;
        }
        int g = gcd(num, den);
        num /= g;
        den /= g;
    }
};

Fraction operator+(Fraction a, const Fraction &b) {
    a += b;
    return a;
}
Fraction operator*(Fraction a, const Fraction &b) {
    a *= b;
    return a;
}

void show(const Fraction &f) {
    cout << f.getNum() << '/' << f.getDen() << '\n';
}

int main() {
    Fraction a(1, 2), b(2, -6);
    show(b);
    show(a + b);
    show(a * b + Fraction(1, 4));
    show(1 + b);

    a += b;
    a *= 3;
    show(a);
    return 0;
}
\end{example}

\iosample{}{
-1/3 \\
1/6 \\
1/12 \\
2/3 \\
1/2 \\
}

\begin{itemize}
\item \texttt{+=}、\texttt{*=} 寫成成員函式。 \texttt{a += b}
等於 \texttt{a.operator+=(b)}，左邊的 \texttt{a} 就是 \texttt{*this}，右
邊的 \texttt{b} 是參數 \texttt{r}。 它們會改 \texttt{a}，所以不是
\texttt{const}；回傳 \texttt{*this} 的參考，跟第~5~章的
\texttt{deposit} 一樣，內建型別的 \texttt{x += 1} 本身也是 \texttt{x}。
\item \texttt{+}、\texttt{*} 寫成類別外面的普通函式，用對應的
複合指定來做。 參數 \texttt{a} 是傳值，函式拿到的是左邊運算元的複本，在複
本上做 \texttt{+=}，再回傳複本。 \texttt{a + b} 不會改到 \texttt{a}，跟
\texttt{int} 的 \texttt{+} 一樣。
\item 加法的公式只寫在 \texttt{operator+=} 裡一個地方。 先寫複合指定，再用
它寫 \texttt{+}，是常見的寫法。
\item \texttt{Fraction(1, 4)} 建立一個暫時物件，跟 \texttt{a * b} 的結果
相加。 \texttt{a * b + Fraction(1, 4)} 照一般的優先順序，先乘再加，就是
C 回顧框裡想寫的那種式子。 減法照同樣的方式寫 \texttt{-=} 和 \texttt{-}。
\item 建構子的兩個參數都有預設引數，所以 \texttt{Fraction(1)} 是
$1/1$。 編譯器需要時會用這個建構子把 \texttt{int} 轉成
\texttt{Fraction}：\texttt{1 + b} 先把 \texttt{1} 轉成 \texttt{Fraction(1)}，
再呼叫 \texttt{operator+}；\texttt{a *= 3} 也一樣。
\item 輸出暫時用 \texttt{show} 函式，下一節改成用 \texttt{cout <<}。
\end{itemize}

\subsection{成員還是非成員}

運算子可以寫成成員函式，也可以寫成類別外面的函式。 寫成成員
\texttt{Fraction operator+(const Fraction \&r) const} 的話，\texttt{a + b}
等於 \texttt{a.operator+(b)}；寫成非成員，等於 \texttt{operator+(a, b)}。

寫成成員時，左邊的運算元一定要是 \texttt{Fraction}，編譯器不會為了呼叫成
員函式去轉換它。 \texttt{1 + b} 的左邊是 \texttt{int}，\texttt{int} 沒有
成員函式，所以找不到。 寫成非成員時，兩個運算元都是普通參數，兩邊都能轉
換，\texttt{b + 1} 和 \texttt{1 + b} 都可以用。 一般的選法是：

\begin{itemize}
\item 會改左邊運算元的（\texttt{+=}、\texttt{-=} 這類）寫成成員。
\item 兩邊地位相同、不改運算元的（\texttt{+}、\texttt{-}、\texttt{==}、
\texttt{<}）寫成非成員。
\item \texttt{=}、\texttt{[]} 規定只能寫成成員（第~\ref{sec:subscript}~節）。
\end{itemize}

運算子多載也有限制：不能發明新的運算子（例如 \texttt{**}），不能改變優先
順序和運算元的個數，運算元至少要有一個是自己定義的型別（不能重新定義
\texttt{int + int}）。 \texttt{.}、\texttt{::}、\texttt{?:}、
\texttt{sizeof} 不能多載。

\subsection{常見錯誤}

\begin{itemize}
\item \textbf{把 \texttt{+} 寫成成員}：拿掉範例裡類別外面的
\texttt{operator+}，改成成員函式
\texttt{Fraction operator+(const Fraction \&r) const}，\texttt{a + b} 還是
可以用，\texttt{1 + b} 卻編譯失敗：
\begin{code}{text}{編譯器的錯誤訊息}
error: no match for 'operator+' (operand types are 'int' and 'Fraction')
\end{code}
\item \textbf{參數漏了 \texttt{const}}：\texttt{operator+=} 的參數寫成
\texttt{Fraction \&r}，兩個地方會編譯失敗。 \texttt{operator+} 裡的
\texttt{a += b} 報 \texttt{binding reference of type 'Fraction\&' to
'const Fraction' discards qualifiers}，因為 \texttt{b} 是 \texttt{const}
參考；另外寫一行 \texttt{a += 3}，報 \texttt{cannot bind non-const lvalue reference of
type 'Fraction\&' to an rvalue of type 'Fraction'}，因為 \texttt{3} 轉出來
的 \texttt{Fraction} 是暫時物件（第~2~章）。 不改的參數一律用
\texttt{const} 參考。
\item \textbf{\texttt{+} 改到了左邊}：成員版的 \texttt{operator+} 寫成
\texttt{*this += r; return *this;}（而且沒有標 \texttt{const}），編譯會過，
但 \texttt{c = a + b;} 之後 \texttt{a} 也變了。 \texttt{+} 要在複本上做。
\end{itemize}

\section{比較與輸出}
\label{sec:compare-output}

\texttt{==} 和 \texttt{<} 回傳 \texttt{bool}。 \texttt{Fraction} 一定是約
分過的，所以兩個分數相等，就是分子和分母都相等；不維持這個不變量的話，
$1/2$ 和 $2/4$ 就會被判成不相等。 比較大小時，因為分母都是正的，
$a/b < c/d$ 等於 $a \times d < c \times b$。

有了 \texttt{==} 和 \texttt{<}，其他四個都可以用它們寫出來，不必各寫一套
邏輯。 \texttt{sort} 預設用 \texttt{<} 比較元素，所以有了 \texttt{<}，
\texttt{vector<Fraction>} 就能直接排序。 \texttt{sort} 在
\texttt{<algorithm>} 裡，第~9~章再仔細講，這裡只用最基本的
\texttt{sort(v.begin(), v.end())}。

\begin{example}[label=ex:fraction-cmp]{cpp}{(fraction\_cmp.cpp) 比較運算子與 \texttt{<<}}
#include <algorithm>
#include <iostream>
#include <numeric>
#include <vector>
using namespace std;

class Fraction {
public:
    Fraction(int n = 0, int d = 1) : num(n), den(d) {
        normalize();
    }
    // += *= 和上一個範例相同，這裡省略
    int getNum() const { return num; }
    int getDen() const { return den; }

    friend ostream &operator<<(ostream &os, const Fraction &f) {
        os << f.num;
        if (f.den != 1) {
            os << '/' << f.den;
        }
        return os;
    }

private:
    int num, den;    // 不變量：den > 0，而且 num/den 已經約分
    void normalize() {
        if (den < 0) {
            num = -num;
            den = -den;
        }
        int g = gcd(num, den);
        num /= g;
        den /= g;
    }
};

bool operator==(const Fraction &a, const Fraction &b) {
    return a.getNum() == b.getNum() && a.getDen() == b.getDen();
}
bool operator!=(const Fraction &a, const Fraction &b) { return !(a == b); }
bool operator<(const Fraction &a, const Fraction &b) {
    return a.getNum() * b.getDen() < b.getNum() * a.getDen();
}
bool operator>(const Fraction &a, const Fraction &b) { return b < a; }
bool operator<=(const Fraction &a, const Fraction &b) { return !(b < a); }
bool operator>=(const Fraction &a, const Fraction &b) { return !(a < b); }

int main() {
    Fraction a(2, 4), b(1, 2), c(3, -4);
    cout << a << ' ' << c << ' ' << Fraction(6, 3) << '\n';
    cout << (a == b) << ' ' << (a != b) << ' ' << (c < a) << ' ' << (a >= 1) << '\n';

    vector<Fraction> v = {Fraction(2, 3), c, Fraction(1, 6), 1, b};
    sort(v.begin(), v.end());
    for (const Fraction &f : v) {
        cout << f << ' ';
    }
    cout << '\n';
    return 0;
}
\end{example}

\iosample{}{
1/2 -3/4 2 \\
1 0 1 0 \\
-3/4 1/6 1/2 2/3 1 \\
}

\begin{itemize}
\item \texttt{a >= 1} 的右邊是 \texttt{int}，跟上一節的 \texttt{1 + b} 一
樣先轉成 \texttt{Fraction}。 比較運算子寫成非成員，兩邊都能轉換。
\item \texttt{a > b} 就是 \texttt{b < a}；\texttt{a <= b} 是「\texttt{b < a}
不成立」。 六個運算子裡真正在比的只有 \texttt{==} 和 \texttt{<}，其他四
個都是一行。
\item 初始化 \texttt{v} 時，\texttt{1} 也轉成了 \texttt{Fraction}。 排序
後由小到大印出，\texttt{2/1} 這種分母是 $1$ 的印成整數。
\end{itemize}

C++20 加了一個 \texttt{<=>} 運算子，寫一個就能讓編譯器補出其他比較。 本
書用 C++17，所以照上面這樣寫齊。

\subsection{輸出運算子 \texttt{<<}}

\texttt{cout << f} 的左邊是 \texttt{cout}，型別是 \texttt{ostream}（在
\texttt{<iostream>} 裡）。 寫成成員的話，左邊運算元得是
\texttt{Fraction}，所以 \texttt{<<} 只能寫成非成員：

\begin{itemize}
\item 第一個參數是 \texttt{ostream \&}。 寫入資料會改變串流，所以不能是
\texttt{const}；串流也不能複製，所以不能傳值。
\item 回傳同一個 \texttt{ostream \&}。 \texttt{cout << a << ' ' << c} 是
\texttt{((cout << a) << ' ') << c}，每一段都要回傳 \texttt{cout}，下一個
\texttt{<<} 才接得上。 第~2~章說 \texttt{cout <<} 能串下去，就是這樣做的。
\item 寫入時用參數 \texttt{os}，不要寫死 \texttt{cout}。 這樣同一個函式也
能寫到其他串流，例如 \texttt{cerr}。
\end{itemize}

非成員函式不能存取 \texttt{private} 成員。 \texttt{Fraction} 有 getter，
所以 \texttt{<<} 可以跟比較運算子一樣寫在外面，用 \texttt{getNum()}、
\texttt{getDen()}。 範例示範另一種寫法：在類別裡把它宣告成
\textbf{friend}（友元）。 \texttt{friend} 函式不是成員，沒有
\texttt{this}，呼叫方式跟寫在外面完全一樣，但可以存取 \texttt{private}
成員。 已經有適當的 getter 就用 getter；沒有、又不想為了輸出多開
getter 時，才用 \texttt{friend}。

\subsection{常見錯誤}

\begin{itemize}
\item \textbf{把 \texttt{<<} 寫成成員}：寫成
\texttt{ostream \&operator<<(ostream \&os) const}，編譯器把它當成
\texttt{f << cout} 的意思，\texttt{cout << a} 報 \texttt{no match for
'operator<<' (operand types are 'std::ostream' \{aka
'std::basic\_ostream<char>'\} and 'Fraction')}。 後面還有幾百行，列出標準函
式庫裡每一個 \texttt{operator<<} 為什麼不合用。 看第一行就夠了：
\texttt{operand types} 告訴你兩邊的型別。
\item \textbf{回傳 \texttt{ostream} 而不是 \texttt{ostream \&}}：回傳型別
少了 \texttt{\&}，就要複製串流，報 \texttt{use of deleted function}，後
面接著的是 \texttt{basic\_ostream} 的複製建構子。 跟第~4~章的
\texttt{unique\_ptr} 一樣，\texttt{deleted function} 表示這個型別禁止複製。
\item \textbf{沒有 \texttt{<} 就排序}：拿掉 \texttt{operator<}，錯誤的位置在
標準函式庫的 \texttt{predefined\_ops.h}，訊息是
\texttt{no match for 'operator<' (operand types are 'Fraction' and
'Fraction')}。 錯誤指著標準函式庫的檔案時，多半是自己的型別少了它需要
的運算子。
\end{itemize}

\section{下標運算子}
\label{sec:subscript}

第~5~章的 \texttt{IntArray} 用 \texttt{get(i)} 和 \texttt{set(i, x)} 存
取元素。 \texttt{vector} 用的是 \texttt{v[i]}，讀和寫都是它。 下標運算子
\texttt{operator[]} 只能寫成成員，參數是下標，回傳元素的參考：

\begin{example}[label=ex:subscript]{cpp}{(subscript.cpp) const 與非 const 兩個版本的 []}
#include <iostream>
using namespace std;

class IntArray {
public:
    IntArray(int n) : len(n), data(new int[n]()) {}
    ~IntArray() { delete[] data; }
    int size() const { return len; }

    int &operator[](int i) { return data[i]; }             // 可以改
    const int &operator[](int i) const { return data[i]; } // 只能讀

private:
    int len;
    int *data;
};

int sum(const IntArray &a) {
    int s = 0;
    for (int i = 0; i < a.size(); i++) {
        s += a[i];                       // 呼叫 const 版本
    }
    return s;
}

int main() {
    IntArray a(5);
    for (int i = 0; i < a.size(); i++) {
        a[i] = i * i;                    // 呼叫非 const 版本
    }
    a[4] += 100;
    cout << a[2] << ' ' << a[4] << ' ' << sum(a) << '\n';
    return 0;
}
\end{example}

\iosample{}{
4 116 130 \\
}

\begin{itemize}
\item \texttt{a[i]} 等於 \texttt{a.operator[](i)}。 回傳 \texttt{int \&}，
所以 \texttt{a[i] = i * i} 和 \texttt{a[4] += 100} 改的是陣列裡的元素，
跟第~2~章回傳參考的 \texttt{larger} 一樣。
\item 兩個版本只差在 \texttt{const}。 \texttt{a} 不是 \texttt{const}，用
第一個；\texttt{sum} 裡的 \texttt{a} 是 \texttt{const} 參考，只能呼叫
\texttt{const} 成員函式，用第二個。 第二個回傳 \texttt{const int \&}，透
過 \texttt{const} 參考拿到的元素也只能讀。 元素是 \texttt{int} 這種小型別
時，\texttt{const} 版本也可以直接回傳 \texttt{int}。
\item 跟 \texttt{vector} 的 \texttt{[]} 一樣，不檢查下標。 要檢查可以另外
寫一個 \texttt{at}。
\item 這個 \texttt{IntArray} 一複製還是會重複釋放，下一節處理。
\end{itemize}

\subsection{讓運算子不出人意料}

讀程式的人看到 \texttt{a + b}，會用他對 \texttt{int} 的 \texttt{+} 的理
解去讀。 自己寫的運算子要跟這個理解一致：

\begin{itemize}
\item \textbf{意思要明顯}。 \texttt{Fraction} 的 \texttt{+} 就是分數加
法，沒有疑問。 二維向量 \texttt{Vector2} 的 \texttt{+} 是向量加法，但
\texttt{p * q} 是內積還是逐項相乘？ 讀的人猜不出來，就不要用運算子，取名
\texttt{dot(p, q)}。 向量乘上一個數字倒是意思明顯，\texttt{v * 2} 可以用運
算子寫（本章練習第~3~題）。
\item \textbf{不改不該改的東西}。 \texttt{+}、\texttt{==}、\texttt{<} 不
改運算元，\texttt{+=} 只改左邊。
\item \textbf{成組的運算子要一致}。 有 \texttt{+} 就有 \texttt{+=}，而且
\texttt{a = a + b} 跟 \texttt{a += b} 結果相同；有 \texttt{==} 就有
\texttt{!=}；\texttt{<}、\texttt{>}、\texttt{<=}、\texttt{>=} 彼此不矛盾。
用 \texttt{+=} 寫 \texttt{+}、用 \texttt{<} 寫其他三個，就是為了這個。
\item \textbf{\texttt{<} 要是真的「小於」}。 \texttt{sort} 假設
\texttt{a < a} 是假，\texttt{a < b} 和 \texttt{b < c} 成立時
\texttt{a < c} 也成立。 寫成 \texttt{<=} 或其他不滿足這些條件的比較，
\texttt{sort} 的行為是未定義的，甚至可能越界。
\item \textbf{不要多載 \texttt{\&\&}、\texttt{||} 和逗號}。 內建的
\texttt{\&\&} 左邊是假就不算右邊，自己多載的版本是普通函式，兩邊都會先
算，讀的人以為還有短路求值。
\end{itemize}

\subsection{常見錯誤}

\begin{itemize}
\item \textbf{只寫非 \texttt{const} 版本}：刪掉範例裡的 \texttt{const}
版本，\texttt{sum} 裡的 \texttt{a[i]} 報第~5~章看過的
\texttt{passing 'const IntArray' as 'this' argument discards
qualifiers}。
\item \textbf{非 \texttt{const} 版本回傳 \texttt{int}}：少了 \texttt{\&}，
回傳的是元素的複本，\texttt{a[i] = i * i} 報 \texttt{lvalue required as
left operand of assignment}，意思是等號左邊不是一個可以寫入的東西。
\item \textbf{\texttt{const} 版本回傳 \texttt{int \&}}：寫成
\texttt{int \&operator[](int i) const}，編譯會過，\texttt{-Wall -Wextra}
也沒有警告，結果 \texttt{sum} 裡可以寫 \texttt{a[i] = 0} 改掉元素。
\texttt{const} 成員函式裡，\texttt{data} 這個指標本身不能改，它指向的
\texttt{int} 卻不受影響，所以編譯器不會擋。 要靠自己把回傳型別寫成
\texttt{const int \&}。
\end{itemize}

\section{複製建構子與複製指定運算子}
\label{sec:copy}

\begin{crecap}[結構指定只複製成員]
C 的結構可以用 \texttt{=} 指定，效果是把每個成員的位元組照抄一份。 成員是
指標時，抄的是位址：

\begin{code}{c}{複製含指標的結構}
struct Array {
    int len;
    int *data;
};

struct Array a = {3, malloc(3 * sizeof(int))};
a.data[0] = 7;
struct Array b = a;     /* 複製 len 和 data 這兩個成員 */
b.data[0] = 99;
printf("%d\n", a.data[0]);   /* 99 */
free(a.data);           /* 只能 free 一次，b.data 也跟著失效 */
\end{code}

\texttt{a} 和 \texttt{b} 共用同一塊記憶體。 要各自一份，C 只能自己寫一個
複製函式，配置新的記憶體再 \texttt{memcpy}，而且每次複製都要記得呼叫它。
\end{crecap}

C++ 的類別也是這樣：沒有自己寫的話，編譯器會產生一個\textbf{複製建構子}
和一個\textbf{複製指定運算子}，逐一複製每個成員。 成員是 \texttt{int}、
\texttt{string}、\texttt{vector} 時，逐一複製就是正確的複製；成員是原始指
標時，就是第~5~章看到的重複釋放。

C++ 跟 C 不同的地方是，可以自己寫這兩個函式，之後每一次複製都自動用自己的
版本，不用記得呼叫什麼：

\begin{itemize}
\item 複製建構子 \texttt{IntArray(const IntArray \&other)}：用一個已經存在
的物件\emph{建立}新物件時呼叫，例如 \texttt{IntArray b = a;}、
\texttt{IntArray b(a);}，以及傳值傳進函式。
\item 複製指定運算子 \texttt{IntArray \&operator=(const IntArray \&other)}：
兩個\emph{已經存在}的物件之間寫 \texttt{c = a;} 時呼叫。
\end{itemize}

\texttt{IntArray b = a;} 雖然有等號，但 \texttt{b} 是這一行才建立的，用的
是複製建構子。 判斷的方法是看左邊的物件之前存不存在。

\begin{example}[label=ex:copy-ctor]{cpp}{(copy\_ctor.cpp) 自己寫的複製建構子與複製指定運算子}
#include <iostream>
using namespace std;

class IntArray {
public:
    IntArray(int n) : len(n), data(new int[n]()) {}

    IntArray(const IntArray &other)                 // 複製建構子
        : len(other.len), data(new int[other.len]) {
        for (int i = 0; i < len; i++) {
            data[i] = other.data[i];
        }
        cout << "複製建構" << '\n';
    }

    IntArray &operator=(const IntArray &other) {    // 複製指定運算子
        cout << "複製指定" << '\n';
        if (this == &other) {                       // a = a 什麼都不做
            return *this;
        }
        int *fresh = new int[other.len];
        for (int i = 0; i < other.len; i++) {
            fresh[i] = other.data[i];
        }
        delete[] data;                              // 舊的陣列不要了
        data = fresh;
        len = other.len;
        return *this;
    }

    ~IntArray() { delete[] data; }

    int size() const { return len; }
    int &operator[](int i) { return data[i]; }

private:
    int len;
    int *data;
};

int main() {
    IntArray a(3);
    a[0] = 7;
    IntArray b = a;          // 建立 b：複製建構子
    b[0] = 99;
    cout << a[0] << ' ' << b[0] << '\n';

    IntArray c(5);
    c = a;                   // c 已經存在：複製指定運算子
    c = c;                   // 自己指定給自己
    cout << c.size() << ' ' << c[0] << '\n';
    return 0;
}
\end{example}

\iosample{}{
複製建構 \\
7 99 \\
複製指定 \\
複製指定 \\
3 7 \\
}

\begin{itemize}
\item 複製建構子配置一塊新的陣列，把元素一個一個抄過去。 這叫\textbf{深
複製}（deep copy）；只抄指標的叫\textbf{淺複製}（shallow copy）。 改
\texttt{b[0]} 不再影響 \texttt{a}，三個物件各自釋放自己的陣列，用 ASan 跑
也沒有錯誤。
\item 複製建構子的參數一定要是參考。 寫成 \texttt{IntArray(IntArray
other)}，傳值本身就要先複製一次，會變成無窮遞迴，所以編譯器直接拒絕：
\texttt{invalid constructor; you probably meant 'IntArray (const
IntArray\&)'}。
\item 複製指定比複製建構多了兩件事。 \texttt{c} 原本有一個長度 $5$ 的陣
列，換成新的之前要先 \texttt{delete[]}，不然就洩漏了。 另外要處理
\texttt{c = c} 這種自己指定給自己的情況，下面說明。
\item 複製指定回傳 \texttt{*this} 的參考，跟 \texttt{+=} 一樣。 內建型別
可以寫 \texttt{x = y = 0}，自己的類別也照這個習慣。
\end{itemize}

\subsection{自我指定}

範例最後的 \texttt{c = c} 看起來沒有人會寫，但 \texttt{v[i] = v[j]} 在 \texttt{i}
和 \texttt{j} 相等時就是這種情況。 假如 \texttt{operator=} 照直覺寫成三
步：先 \texttt{delete[] data;}，再 \texttt{data = new int[other.len];}，
最後用迴圈把 \texttt{other.data[i]} 抄進 \texttt{data[i]}。

\texttt{other} 就是 \texttt{*this} 時，\texttt{other.data} 和
\texttt{data} 是同一個成員。 第一步把原本的陣列釋放了，第二步
\texttt{data} 換成一塊還沒有初始化的新陣列，第三步把它自己抄給自己，原本的元素就
這樣不見了。 \texttt{a[0] = 7; a = a;} 之後印 \texttt{a[0]}，在本書測
試的環境下每次執行印出不同的數字，不是 \texttt{7}。

範例的寫法用兩個方法避開這個問題。 一是開頭檢查 \texttt{this == \&other}，
是同一個物件就直接回傳。 二是先配置新陣列、抄完，最後才
\texttt{delete[]} 舊的；就算沒有那個檢查，自我指定也只是多抄一次。 這個
順序還有一個好處：\texttt{new} 配置失敗時會丟出例外（第~4~章），這時舊的
陣列還在，\texttt{c} 保持原樣。

\subsection{複製後交換}

複製指定還有一種更短的寫法，叫\textbf{複製後交換}（copy-and-swap）：

\begin{code}{cpp}{用複製後交換寫 operator=}
IntArray &operator=(IntArray other) {   // 傳值：other 是一份複本
    swap(len, other.len);
    swap(data, other.data);
    return *this;
}                                       // other 帶著舊的陣列銷毀
\end{code}

參數改成傳值，呼叫時就先用複製建構子做出一份複本 \texttt{other}。 接著把
自己和 \texttt{other} 的成員互換（\texttt{std::swap}，在
\texttt{<utility>} 裡）：\texttt{*this} 拿到新複製的陣列，\texttt{other}
拿到舊的。 函式結束時 \texttt{other} 銷毀，它的解構子替我們釋放舊的陣列。

複製的工作交給複製建構子，釋放交給解構子，\texttt{operator=} 自己不碰
\texttt{new} 和 \texttt{delete}。 自我指定也不用特別處理：\texttt{other}
是另一份複本，跟自己互換不會出問題。 代價是自我指定時也會複製一次，通常
不在意。 下一節會看到，這個寫法同時也能當移動指定。

\subsection{常見錯誤}

\begin{itemize}
\item \textbf{只寫了複製建構子}：忘了寫 \texttt{operator=}，編譯器產生的
複製指定還是只抄指標，\texttt{c = a;} 照樣重複釋放。 \texttt{g++ -Wextra}
在用到它的地方會警告 \texttt{implicitly-declared ... operator=(const
...\&)' is deprecated [-Wdeprecated-copy]}，原因寫在下一行：類別有自己寫的
複製建構子（本章練習第~4~題）。 看到這個警告，就要把 \texttt{operator=}
補上。
\item \textbf{複製指定沒有釋放舊的陣列}：直接寫 \texttt{data = new
int[other.len];}，原本的陣列就沒有指標記得了，ASan 在程式結束時報
\texttt{detected memory leaks}。
\end{itemize}

\section{移動，以及三法則、五法則、零法則}
\label{sec:rule-of-zero}

第~4~章說過，\texttt{vector} 被 \texttt{std::move} 之後，新的
\texttt{vector} 直接接手舊的那塊記憶體，一個元素都不用複製。 自己的類別
要做到這件事，要寫\textbf{移動建構子}（move constructor）。 移動指定運
算子可以不另外寫：上一節複製後交換的 \texttt{operator=} 參數是傳值，
\texttt{other} 由引數建立，引數是 rvalue 時用的就是移動建構子。 下面把它
們放在一起，複製和移動時各印一行：

\begin{example}[label=ex:move]{cpp}{(move.cpp) 移動建構子與複製後交換}
#include <iostream>
#include <utility>
using namespace std;

class IntArray {
public:
    IntArray(int n) : len(n), data(new int[n]()) {}

    IntArray(const IntArray &other)                 // 複製建構子
        : len(other.len), data(new int[other.len]) {
        for (int i = 0; i < len; i++) {
            data[i] = other.data[i];
        }
        cout << "複製 " << len << " 個" << '\n';
    }

    IntArray(IntArray &&other) noexcept             // 移動建構子
        : len(other.len), data(other.data) {
        other.len = 0;
        other.data = nullptr;
        cout << "移動" << '\n';
    }

    IntArray &operator=(IntArray other) {           // 複製指定兼移動指定
        swap(len, other.len);
        swap(data, other.data);
        return *this;
    }                                               // other 帶著舊陣列銷毀

    ~IntArray() { delete[] data; }

    int size() const { return len; }
    int &operator[](int i) { return data[i]; }

private:
    int len;
    int *data;
};

int main() {
    IntArray a(4);
    a[3] = 9;
    IntArray b = a;                  // a 之後還要用：複製
    IntArray c = std::move(a);       // a 不再用了：移動
    cout << a.size() << ' ' << c.size() << ' ' << c[3] << '\n';

    b = c;                           // 參數 other 由 c 複製而來
    b = IntArray(2);                 // 參數 other 直接拿暫時物件
    a = std::move(c);                // 參數 other 由 c 移動而來
    cout << a.size() << ' ' << b.size() << ' ' << b[1] << '\n';
    return 0;
}
\end{example}

\iosample{}{
複製 4 個 \\
移動 \\
0 4 9 \\
複製 4 個 \\
移動 \\
4 2 0 \\
}

先看移動建構子：

\begin{itemize}
\item 參數型別 \texttt{IntArray \&\&} 是 \textbf{rvalue 參考}，只能綁定
rvalue：暫時物件，或者用 \texttt{std::move} 標成「可以被搬走」的物件。
複製建構子的參數 \texttt{const IntArray \&} 什麼都能綁，所以兩個都能用時，
編譯器對 rvalue 選移動建構子。
\item 移動就是把指標拿過來，再把 \texttt{other} 的指標設成
\texttt{nullptr}。 少了這一步，兩個物件又會指著同一個陣列，\texttt{other}
銷毀時就把它釋放了。 \texttt{delete[] nullptr} 什麼都不做，所以被移走的
物件可以正常銷毀。
\item \texttt{noexcept} 表示這個函式不會丟出例外。 移動只是搬幾個成員，本
來就不會失敗，寫上它，\texttt{vector} 擴充容量時才會用移動（見本節的常見
錯誤）。
\end{itemize}

再對照輸出：

\begin{itemize}
\item \texttt{IntArray c = std::move(a);} 用移動建構子。 \texttt{a} 被移
走之後，\texttt{a.size()} 是 $0$。 這是我們自己的移動建構子保證的；一般
的型別被移走後是有效但未指定的狀態（第~4~章），不要依賴它的內容。
\item 後面三行都呼叫同一個 \texttt{operator=}，差別在參數 \texttt{other}
怎麼建立：\texttt{b = c} 的 \texttt{c} 是 lvalue，用複製建構子；
\texttt{IntArray(2)} 是暫時物件，直接成為 \texttt{other}，什麼都沒
印；\texttt{std::move(c)} 用移動建構子。
\item 用 ASan 跑，沒有重複釋放也沒有洩漏。
\end{itemize}

\subsection{編譯器會產生哪些函式}

建構子以外，下面五個函式叫\textbf{特殊成員函式}。 沒有自己寫的時候，編譯
器會在需要時產生，做法都是對每個成員做同一件事：

\begin{center}
\begin{tabular}{lll}
\toprule
函式 & 宣告 & 編譯器產生的版本 \\
\midrule
解構子 & \texttt{\textasciitilde IntArray()} & 銷毀每個成員 \\
複製建構子 & \texttt{IntArray(const IntArray \&)} & 逐一複製成員 \\
複製指定 & \texttt{IntArray \&operator=(const IntArray \&)} & 逐一指定成員 \\
移動建構子 & \texttt{IntArray(IntArray \&\&)} & 逐一移動成員 \\
移動指定 & \texttt{IntArray \&operator=(IntArray \&\&)} & 逐一移動指定成員 \\
\bottomrule
\end{tabular}
\end{center}

要注意的規則是：只要自己寫了解構子、複製建構子或複製指定其中任何一個，編
譯器就\emph{不}產生移動建構子和移動指定。 這時 \texttt{std::move(a)} 不
會出錯，而是改用複製。 例如範例~\ref{ex:copy-ctor} 的 \texttt{IntArray}
沒有移動建構子，寫 \texttt{IntArray c = std::move(a);} 會印出「複製建構」，
之後 \texttt{a.size()} 還是 $3$。 結果正確，只是慢。

第~5~章那個只寫了解構子的 \texttt{IntArray} 也是這樣：沒有移動，複製又是
編譯器產生的淺複製，所以 \texttt{std::move} 也救不了它。

不想要某個函式時，在宣告後面寫 \texttt{= delete}；想明確表示「用編譯器產生
的版本」時，寫 \texttt{= default}。 第~5~章練習第~4~題的 \texttt{Account}
每個帳戶有自己的編號，不應該被複製，可以這樣禁止：

\begin{code}{cpp}{禁止複製 Account}
class Account {
public:
    Account(const string &owner) : id(opened + 1), owner(owner) {
        opened++;
    }
    Account(const Account &) = delete;
    Account &operator=(const Account &) = delete;
    // ...
};
\end{code}

這時 \texttt{Account c = a;} 報 \texttt{use of deleted function
'Account::Account(const Account\&)'}，跟複製 \texttt{unique\_ptr} 時看到的一樣。

\subsection{三法則、五法則、零法則}

\begin{itemize}
\item \textbf{三法則}（rule of three）：解構子、複製建構子、複製指定，自己
寫了其中一個，通常三個都要寫。 需要自己寫解構子，表示類別直接管著某個資源
（例如 \texttt{new[]} 來的陣列），編譯器產生的複製只會複製指標。
\item \textbf{五法則}（rule of five）：再加上移動建構子和移動指定。 寫了
三法則的那三個，編譯器就不產生移動，想要移動得自己寫（或者像範例那樣，
用傳值的 \texttt{operator=} 一次處理兩種指定）。
\item \textbf{零法則}（rule of zero）：一個都不寫。 資源交給已經寫好這五
個函式的型別來管，例如 \texttt{string}、\texttt{vector}、
\texttt{unique\_ptr}，類別本身只放這些成員，五個函式都用編譯器產生的。
\end{itemize}

本書的建議是零法則。 範例~\ref{ex:move} 的 \texttt{IntArray} 寫了四十多
行，才做到 \texttt{vector<int>} 本來就會的事，其中任何一行寫錯都是重複釋放
或洩漏。 三法則、五法則是在你真的要自己寫一個管理資源的型別時才用得上，
例如自己實作一個容器；一般的類別用零法則：

\begin{example}[label=ex:rule-of-zero]{cpp}{(rule\_of\_zero.cpp) 成員管好自己，類別什麼都不用寫}
#include <iostream>
#include <string>
#include <utility>
#include <vector>
using namespace std;

class Scores {
public:
    Scores(const string &name) : name(name) {}
    void add(int x) { data.push_back(x); }
    int &operator[](int i) { return data[i]; }

    friend ostream &operator<<(ostream &os, const Scores &s) {
        os << s.name << ':';
        for (int x : s.data) {
            os << ' ' << x;
        }
        return os;
    }

private:
    string name;
    vector<int> data;        // 沒有寫解構子、複製和移動，全部由編譯器產生
};

int main() {
    Scores a("Ada");
    a.add(90);
    a.add(75);

    Scores b = a;            // 複製：b 有自己的 name 和 data
    b[1] = 100;
    cout << a << '\n' << b << '\n';

    a = b;                   // 複製指定
    vector<Scores> all;
    all.push_back(std::move(b));   // 移動：不複製元素
    all.push_back(Scores("Alan"));
    all[1].add(60);
    cout << a << '\n' << all[0] << '\n' << all[1] << '\n';
    return 0;
}
\end{example}

\iosample{}{
Ada: 90 75 \\
Ada: 90 100 \\
Ada: 90 100 \\
Ada: 90 100 \\
Alan: 60 \\
}

編譯器產生的複製建構子逐一複製 \texttt{name} 和 \texttt{data}，而
\texttt{string}、\texttt{vector} 的複製本來就是深複製，所以 \texttt{b}
有自己的一份。 移動、指定、銷毀也都是交給成員自己處理。 成員是
\texttt{unique\_ptr} 的類別也一樣不用寫：\texttt{unique\_ptr} 不能複製，類
別也就自動不能複製，但可以移動。

\subsection{常見錯誤}

\begin{itemize}
\item \textbf{移動建構子沒有寫 \texttt{noexcept}}：\texttt{vector} 擴充容量
時要把舊元素搬到新的記憶體。 搬到一半如果有例外，已經移動過去的元素回不
來，所以 \texttt{vector} 只在移動建構子標了 \texttt{noexcept} 時才用移
動，否則改用複製。 拿範例~\ref{ex:move} 的 \texttt{IntArray} 做三次
\texttt{v.push\_back(IntArray(i))}（\texttt{i} 是 $1$、$2$、$3$），容量依
序變成 $1$、$2$、$4$。 有 \texttt{noexcept} 時印出六次「移動」：三次是把
暫時物件移進 \texttt{vector}，三次是擴充時搬舊元素。 拿掉
\texttt{noexcept} 後，擴充時搬舊元素的那三次變成「複製 1 個」、「複製 1
個」、「複製 2 個」。
\item \textbf{寫了解構子，以為還有移動}：如上面所說，這時編譯器不產生移
動，\texttt{std::move} 會默默改用複製。 只寫了解構子的類別，複製還是淺複
製。
\item \textbf{零法則的類別裡又寫了空的解構子}：在 \texttt{Scores} 裡加一行
\texttt{\textasciitilde Scores() \{\}}，程式還是對的，但編譯器就不產生移動
了，\texttt{all.push\_back(std::move(b))} 變成複製整個 \texttt{vector}。
不需要做事的解構子就不要寫。
\end{itemize}

\section{本章練習}
\label{sec:ch6-exercises}

\begin{drills}[label=drills:ch6]
\item 預測下面程式的輸出。
\begin{code}{cpp}{練習：哪一行呼叫哪一個函式}
#include <iostream>
using namespace std;

struct T {
    int v;
    T(int x) : v(x) { cout << "C" << v << ' '; }
    T(const T &o) : v(o.v) { cout << "K" << v << ' '; }
    T &operator=(const T &o) {
        v = o.v;
        cout << "A" << v << ' ';
        return *this;
    }
    ~T() { cout << "D" << v << ' '; }
};

void f(T t) { cout << "f" << t.v << ' '; }
void g(const T &t) { cout << "g" << t.v << ' '; }

int main() {
    T a(1);
    T b = a;
    T c(2);
    c = b;
    cout << '\n';
    f(c);
    g(c);
    cout << '\n';
    c = T(3);
    cout << '\n';
    return 0;
}
\end{code}
\item 下面的程式編譯失敗。 找出所有問題並改正，讓它輸出
\texttt{0.05 1.50 2.75 4.30}。 其中有一個問題，\texttt{g++} 這次沒有報錯。
\begin{code}{cpp}{練習：找錯}
#include <algorithm>
#include <iostream>
#include <vector>
using namespace std;

class Money {
public:
    Money(int cents = 0) : cents(cents) {}
    Money &operator+=(Money &r) {
        cents += r.cents;
        return *this;
    }
    bool operator<(const Money &r) {
        return cents < r.cents;
    }
    ostream &operator<<(ostream &os) const {
        return os << cents / 100 << '.' << cents % 100 / 10 << cents % 10;
    }

private:
    int cents;
};

Money operator+(Money a, const Money &b) {
    a += b;
    return a;
}

int main() {
    vector<Money> v = {Money(275), Money(150), Money(5)};
    sort(v.begin(), v.end());
    Money total;
    for (const Money &m : v) {
        total += m;
        cout << m << ' ';
    }
    cout << total << '\n';
    return 0;
}
\end{code}
\item 寫一個二維向量類別 \texttt{Vector2}，成員是兩個 \texttt{double}
\texttt{x}、\texttt{y}。 提供 \texttt{+}、\texttt{+=}、向量乘以數字（\texttt{v * 2} 和
\texttt{2 * v} 都要能用）、\texttt{==}、輸出成 \texttt{(x, y)} 的 \texttt{<<}，內積用一
個普通函式 \texttt{dot}。 用下面的 \texttt{main} 測試：
\begin{code}{cpp}{練習 3 的 main}
int main() {
    Vector2 p(1, 2), q(3, -1);
    cout << p + q << ' ' << 2 * p << ' ' << q * 0.5 << '\n';
    cout << dot(p, q) << ' ' << (p + q == Vector2(4, 1)) << '\n';
    p += q;
    p *= 3;
    cout << p << '\n';
    return 0;
}
\end{code}
應該輸出 \texttt{(4, 1) (2, 4) (1.5, -0.5)}、\texttt{1 1}、
\texttt{(12, 3)} 三行。 \texttt{p *= 3} 表示你還需要哪一個運算子？
\item 下面的 \texttt{Text} 用 C 的字串函式管理一塊 \texttt{char} 陣列。
用 \texttt{g++ -Wall -Wextra} 編譯會出現一個警告；在終端機執行，印出
\texttt{hello jello hello} 之後，接著印出
\texttt{free(): double free detected in tcache 2} 並異常結束。 說明原因，
再用兩種方法改正：(a) 照三法則補上缺的函式；(b) 照零法則改用
\texttt{string}。
\begin{code}{cpp}{(text.cpp) 練習：少了一個函式}
#include <cstring>
#include <iostream>
using namespace std;

class Text {
public:
    Text(const char *s) : len(strlen(s)), buf(new char[len + 1]) {
        strcpy(buf, s);
    }
    Text(const Text &other) : len(other.len), buf(new char[len + 1]) {
        strcpy(buf, other.buf);
    }
    ~Text() { delete[] buf; }

    char &operator[](int i) { return buf[i]; }
    friend ostream &operator<<(ostream &os, const Text &t) {
        return os << t.buf;
    }

private:
    int len;
    char *buf;
};

int main() {
    Text a("hello");
    Text b = a;
    b[0] = 'j';
    Text c("hi");
    c = a;
    cout << a << ' ' << b << ' ' << c << '\n';
    return 0;
}
\end{code}
\item 下面每個類別只列出成員變數。 判斷每一個應該照零法則什麼都不寫、照三
法則（或五法則）自己寫，還是應該用 \texttt{= delete} 禁止複製：
\begin{enumerate}
\item \texttt{string name; vector<int> scores;}
\item \texttt{int *buf; int len;}，\texttt{buf} 指向建構子 \texttt{new[]}
來的陣列。
\item \texttt{unique\_ptr<int[]> buf; int len;}
\item \texttt{FILE *fp;}，建構子 \texttt{fopen}、解構子 \texttt{fclose}。
\item \texttt{vector<string> \&log;}，跟第~5~章的 \texttt{Recorder} 一樣綁
著呼叫者的 \texttt{vector}。
\end{enumerate}
\end{drills}

\addcontentsline{toc}{section}{練習解答}
\section*{練習解答}

\begin{enumerate}
\item 輸出四行：\texttt{C1 K1 C2 A1 }、\texttt{K1 f1 D1 g1 }、
\texttt{C3 A3 D3 }、\texttt{D3 D1 D1 }（最後一行沒有換行）。
\texttt{T b = a;} 是建立 \texttt{b}，用複製建構子（\texttt{K}）；
\texttt{c = b;} 時 \texttt{c} 已經存在，用複製指定（\texttt{A}）。
\texttt{f(c)} 傳值，參數 \texttt{t} 由複製建構子建立，\texttt{f} 結束時
\texttt{t} 銷毀，印 \texttt{D1}；\texttt{g(c)} 傳 \texttt{const} 參考，不
複製。 \texttt{c = T(3)} 先建立暫時物件（\texttt{C3}），指定給
\texttt{c}（\texttt{A3}），分號一到暫時物件就銷毀（\texttt{D3}）。 這個
\texttt{T} 沒有移動指定，因為自己寫了複製指定和解構子，所以暫時物件也是用
複製指定。 最後 \texttt{c}、\texttt{b}、\texttt{a} 以相反順序銷毀，
\texttt{c} 的值已經是 $3$。
\item 三個問題：
(1) \texttt{operator+=} 的參數要是 \texttt{const Money \&}。
\texttt{operator+} 裡的 \texttt{b} 和迴圈裡的 \texttt{m} 都是
\texttt{const} 參考，兩處都報 \texttt{binding reference of type
'Money\&' to 'const Money' discards qualifiers}。
(2) \texttt{<<} 寫成了成員，\texttt{cout << m} 和 \texttt{cout << total}
報 \texttt{no match for 'operator<<'}。 改成
\texttt{friend ostream \&operator<<(ostream \&os, const Money \&m)}，函式
裡的 \texttt{cents} 都改成 \texttt{m.cents}。
(3) \texttt{operator<} 不改物件，卻沒有標 \texttt{const}。 這次
\texttt{sort} 剛好只用非 \texttt{const} 的元素比較，所以編譯得過；兩個
\texttt{const Money} 比較時就會失敗。 最好照第~\ref{sec:compare-output}~節
寫成非成員 \texttt{bool operator<(const Money \&a, const Money \&b)}，再替
\texttt{Money} 加一個 \texttt{getCents() const}。
改完後輸出 \texttt{0.05 1.50 2.75 4.30}。 \texttt{\%} 那一串是在補足小數點
後兩位：$5$ 分印成 \texttt{0.05}，不是 \texttt{0.5}。
\item 一種寫法：
\begin{code}{cpp}{練習 3 參考解答（main 同題目）}
#include <iostream>
using namespace std;

class Vector2 {
public:
    Vector2(double x = 0, double y = 0) : x(x), y(y) {}
    double getX() const { return x; }
    double getY() const { return y; }

    Vector2 &operator+=(const Vector2 &r) {
        x += r.x;
        y += r.y;
        return *this;
    }
    Vector2 &operator*=(double k) {
        x *= k;
        y *= k;
        return *this;
    }

private:
    double x, y;
};

Vector2 operator+(Vector2 a, const Vector2 &b) { return a += b; }
Vector2 operator*(Vector2 v, double k) { return v *= k; }
Vector2 operator*(double k, Vector2 v) { return v *= k; }

bool operator==(const Vector2 &a, const Vector2 &b) {
    return a.getX() == b.getX() && a.getY() == b.getY();
}

double dot(const Vector2 &a, const Vector2 &b) {
    return a.getX() * b.getX() + a.getY() * b.getY();
}

ostream &operator<<(ostream &os, const Vector2 &v) {
    return os << '(' << v.getX() << ", " << v.getY() << ')';
}
\end{code}
\texttt{p *= 3} 需要 \texttt{operator*=}，\texttt{v * k} 和 \texttt{k * v}
都用它來寫。 \texttt{return a += b;} 是把 \texttt{a += b; return a;} 寫成
一行，\texttt{+=} 回傳的就是 \texttt{a}。 有 getter，所以 \texttt{<<} 不
需要 \texttt{friend}。 成員只有兩個 \texttt{double}，照零法則，複製和解構
都不用寫。 要注意 \texttt{==} 直接比 \texttt{double}，經過計算的小數可能
有誤差（例如 \texttt{0.1 + 0.2} 不等於 \texttt{0.3}），實際使用時常改成
「差距小於某個很小的數」。
\item 警告是 \texttt{implicitly-declared 'constexpr Text\& Text::operator=(const
Text\&)' is deprecated [-Wdeprecated-copy]}，下一行說明原因是
\texttt{Text} 有自己寫的複製建構子。 \texttt{Text b = a;} 用的是自己寫的
複製建構子，\texttt{b} 有自己的陣列，所以 \texttt{b[0] = 'j'} 不影響
\texttt{a}。 \texttt{c = a;} 用的卻是編譯器產生的複製指定，只複製了指標：
\texttt{c.buf} 和 \texttt{a.buf} 指向同一塊，\texttt{c} 原本的
\texttt{"hi"} 沒有人釋放。 \texttt{main} 結束時 \texttt{c}、\texttt{a} 先
後 \texttt{delete[]} 同一塊，就是重複釋放。 用 ASan 跑，報告是
\texttt{attempting double-free}，兩次釋放都在 \texttt{Text::\textasciitilde
Text()}，配置的位置是 \texttt{Text::Text(char const*)}。

(a) 補上複製指定，用複製後交換（記得 \texttt{\#include <utility>}）：
\begin{code}{cpp}{練習 4 (a)：補上 operator=}
Text &operator=(Text other) {
    swap(len, other.len);
    swap(buf, other.buf);
    return *this;
}
\end{code}

(b) 成員改成 \texttt{string str;}，建構子寫成
\texttt{Text(const char *s) : str(s) \{\}}，\texttt{operator[]} 回傳 \texttt{str[i]}，
\texttt{<<} 印 \texttt{t.str}；刪掉複製建構子、解構子，也不用
\texttt{<cstring>}。

兩種改法都輸出 \texttt{hello jello hello}，用 ASan 跑沒有錯誤。 (b) 少了
十幾行，而且多了移動：(a) 寫了解構子和複製建構子，編譯器不產生移動建構
子，要移動得自己再寫。
\item
(a) 零法則。 \texttt{string} 和 \texttt{vector} 會自己複製、移動、釋放。
(b) 三法則（最好是五法則），或者改用 \texttt{vector<int>} 變成 (a) 的情況。
這就是本章的 \texttt{IntArray}。
(c) 零法則。 編譯器產生的版本會讓類別不能複製、可以移動，通常正好是要的。
要能複製的話，改用 \texttt{vector}。
(d) 要寫解構子，所以照三法則處理。 開啟的檔案沒有好的「複製」方法，通常把
複製建構子和複製指定 \texttt{= delete}，需要的話再自己寫移動（移動時把來源
的 \texttt{fp} 設成 \texttt{nullptr}，解構子裡檢查不是 \texttt{nullptr} 才
\texttt{fclose}）。
(e) 什麼都不用寫，編譯器產生的複製建構子會讓新物件綁著同一個
\texttt{vector}，這通常就是要的。 但參考成員不能重新綁定，所以編譯器不會
產生複製指定，\texttt{r1 = r2;} 編譯失敗。
\end{enumerate}

\chapter{繼承與多型}
\label{ch:inheritance}

第~5、6~章的類別各自獨立：\texttt{Account} 是帳戶，\texttt{Fraction} 是
分數，彼此沒有關係。 實際的程式常有一群「同一種東西的不同版本」：圓和矩
形都是形狀，汽車和機車都是車。 它們有共同的部分（名字、都能算面積），也有
各自不同的部分（半徑、寬和高、算面積的公式）。

這一章先用\textbf{繼承}（inheritance）表達「圓是一種形狀」，看衍生類別
怎麼建立和銷毀。 接著用\textbf{虛擬函式}讓同一行 \texttt{s.area()} 依物件
實際的型別呼叫不同的函式，這叫\textbf{多型}（polymorphism）。 後半章講抽
象類別、虛擬解構子，以及怎麼把不同型別的物件放進同一個 \texttt{vector}。
最後一節講什麼時候不該用繼承。

\section{衍生類別}
\label{sec:derived}

\begin{crecap}[型別標記與共同開頭]
C 要表達「圓和矩形都是形狀」，常見的寫法是在結構裡放一個\textbf{型別標
記}，用 \texttt{switch} 分開處理：

\begin{code}{c}{用型別標記區分形狀}
enum Kind { CIRCLE, RECT };

struct Shape {
    enum Kind kind;        /* 型別標記 */
    const char *name;
    double a, b;           /* 圓：a 是半徑；矩形：a、b 是寬和高 */
};

double area(const struct Shape *s) {
    switch (s->kind) {
    case CIRCLE: return 3.14159 * s->a * s->a;
    case RECT:   return s->a * s->b;
    }
    return 0;
}
\end{code}

\texttt{a}、\texttt{b} 代表什麼，只能靠註解約定。 多一種形狀，就要找出每
一個對 \texttt{kind} 做 \texttt{switch} 的函式，一個一個補上。

另一種寫法是把共同的部分做成一個結構，放在每種形狀的第一個成員，再把「算
面積」做成函式指標：

\begin{code}{c}{共同開頭加函式指標}
struct Shape {
    const char *name;
    double (*area)(const struct Shape *self);
};

struct Circle {
    struct Shape base;     /* 共同的部分放在第一個成員 */
    double r;
};

/* c 是一個 struct Circle，建立時把 c.base.area 設成 circle_area */
struct Shape *p = (struct Shape *)&c;   /* 第一個成員的位址就是 c 的位址 */
printf("%s %g\n", p->name, p->area(p)); /* 圓 3.14159 */
\end{code}

\texttt{circle\_area} 收 \texttt{struct Shape *}，自己轉回
\texttt{struct Circle *} 再算。 拿到 \texttt{struct Shape *} 的程式不用知
道它是哪一種形狀，呼叫 \texttt{p->area(p)} 就會執行對應的函式。 代價是到
處都要強制轉型，函式指標也要自己填，填錯了編譯器不會發現。
\end{crecap}

C++ 的繼承做的是第二種寫法的前半：把共同的部分寫成一個類別，其他類別「繼
承」它。 被繼承的叫\textbf{基底類別}（base class），繼承的叫\textbf{衍生
類別}（derived class）。 函式指標那一半是第~\ref{sec:virtual}~節的虛擬
函式。

\begin{example}[label=ex:derived]{cpp}{(derived.cpp) Circle 和 Rect 都繼承 Shape}
#include <iostream>
#include <string>
using namespace std;

class Shape {
public:
    Shape(const string &name) : name(name) {}
    const string &getName() const { return name; }

private:
    string name;
};

class Circle : public Shape {
public:
    Circle(double r) : Shape("圓"), r(r) {}
    double area() const { return 3.14159 * r * r; }

private:
    double r;
};

class Rect : public Shape {
public:
    Rect(double w, double h) : Shape("矩形"), w(w), h(h) {}
    double area() const { return w * h; }

private:
    double w, h;
};

int main() {
    Circle c(1);
    Rect r(3, 4);
    cout << c.getName() << ' ' << c.area() << '\n';
    cout << r.getName() << ' ' << r.area() << '\n';

    const Shape &s = r;        // Rect 也是一個 Shape
    cout << s.getName() << '\n';
    // s.area();               // 編譯失敗：Shape 沒有 area
    return 0;
}
\end{example}

\iosample{}{
圓 3.14159 \\
矩形 12 \\
矩形 \\
}

\begin{itemize}
\item \texttt{class Circle : public Shape} 表示 \texttt{Circle} 繼承
\texttt{Shape}。 每個 \texttt{Circle} 物件裡都有一份 \texttt{Shape}（跟
C 的 \texttt{base} 成員一樣），再加上自己的 \texttt{r}。
\texttt{Shape} 的 \texttt{public} 成員函式也變成 \texttt{Circle} 的，所以
可以寫 \texttt{c.getName()}。
\item 建構子不會繼承。 \texttt{Circle} 的建構子在初始化列表寫
\texttt{Shape("圓")}，用 \texttt{Shape} 的建構子建立裡面那一份
\texttt{Shape}，再建立自己的成員 \texttt{r}。
\item \texttt{name} 是 \texttt{Shape} 的 \texttt{private} 成員。
\texttt{Circle} 的物件裡有它，但 \texttt{Circle} 的成員函式不能直接用，只
能透過 \texttt{getName()}。 \texttt{private} 對衍生類別也一樣有效，
\texttt{Shape} 的不變量還是只要檢查 \texttt{Shape} 自己的函式。
\item \texttt{const Shape \&s = r;} 把 \texttt{Rect} 綁到 \texttt{Shape}
的參考上，同樣可以寫 \texttt{Shape *p = \&r;}。 這就是 C 回顧框裡的
\texttt{(struct Shape *)\&c}，只是不用轉型，編譯器知道
\texttt{Rect} 是一種 \texttt{Shape}。 反過來不行：\texttt{Shape} 不一定是
\texttt{Rect}。
\item 透過 \texttt{s} 只能用 \texttt{Shape} 有的成員。 \texttt{s.area()}
報 \texttt{'const class Shape' has no member named 'area'}。 讓
\texttt{s.area()} 能用、而且呼叫到 \texttt{Rect} 的 \texttt{area}，是
第~\ref{sec:virtual}~節的內容。
\end{itemize}

\texttt{public} 繼承的意思是「\texttt{Circle} 是一種 \texttt{Shape}」
（is-a）：任何需要 \texttt{Shape} 的地方，都可以交給它一個
\texttt{Circle}。 決定要不要用繼承時，先問這句話說不說得通。
第~\ref{sec:composition}~節會再回來談這件事。

\subsection{\texttt{protected}}

衍生類別有時候需要碰基底類別的內部，但又不想開放給所有人。 這時用第三種
存取標籤 \texttt{protected}（保護）：\texttt{protected} 成員對外面的程式跟
\texttt{private} 一樣，衍生類別的成員函式卻可以用。 例如在 \texttt{Shape}
加一段 \texttt{protected: void setName(const string \&n) \{ name = n; \}}，
\texttt{Circle} 的成員函式就可以呼叫 \texttt{setName("單位圓")}，
\texttt{main} 裡的 \texttt{c.setName(...)} 則編譯失敗。

\texttt{protected} 也可以用在成員變數上，例如把 \texttt{name} 直接標成
\texttt{protected}。 一般不建議這樣做：任何人寫一個衍生類別，就能任意改
\texttt{name}，要確認不變量就得檢查所有衍生類別。 成員變數維持
\texttt{private}，需要時開 \texttt{protected} 的成員函式。

\subsection{常見錯誤}

\begin{itemize}
\item \textbf{漏寫 \texttt{public}}：寫成 \texttt{class Circle : Shape}，
\texttt{class} 預設的繼承方式是 \texttt{private}，\texttt{Shape} 的成員在
外面全都不能用了。 \texttt{c.getName()} 報兩行：
\begin{code}{text}{編譯器的錯誤訊息}
error: 'const std::string& Shape::getName() const' is inaccessible within this context
error: 'Shape' is not an accessible base of 'Circle'
\end{code}
\texttt{private} 繼承很少用到。 本書所有的繼承都是 \texttt{public}，記得
每次都寫出來。
\item \textbf{衍生類別直接用基底的 \texttt{private} 成員}：在
\texttt{Circle} 裡寫 \texttt{void rename() \{ name = "圓形"; \}}，報
\texttt{'std::string Shape::name' is private within this context}。 改用
\texttt{Shape} 提供的函式。
\item \textbf{沒有呼叫基底的建構子}：把 \texttt{Circle} 的初始化列表寫成
\texttt{: r(r)}，編譯器會改用 \texttt{Shape} 的預設建構子，但
\texttt{Shape} 沒有，報 \texttt{no matching function for call to
'Shape::Shape()'}。 寫成 \texttt{: name("圓"), r(r)} 也不行，衍生類別的
初始化列表只能寫基底類別和自己的成員，報
\texttt{class 'Circle' does not have any field named 'name'}。
\end{itemize}

\section{建構與解構的順序}
\label{sec:inherit-order}

第~4~章看過，區塊裡的物件依宣告順序建立，以相反順序銷毀；第~5~章看過，成
員在建構子本體執行之前就建立好了。 有了繼承，衍生類別的物件裡有三個部分：
基底類別那一份、自己的成員、建構子本體。 下面的 \texttt{Tracer} 跟第~4~章
的一樣，建立和銷毀時各印一行：

\begin{example}[label=ex:inherit-order]{cpp}{(inherit\_order.cpp) 基底、成員、本體的順序}
#include <iostream>
#include <string>
using namespace std;

struct Tracer {
    string name;
    Tracer(const string &n) : name(n) { cout << "建立 " << name << '\n'; }
    ~Tracer() { cout << "銷毀 " << name << '\n'; }
};

class Base {
public:
    Base() : part("Base 的成員") { cout << "Base 建構子本體" << '\n'; }
    ~Base() { cout << "Base 解構子本體" << '\n'; }

private:
    Tracer part;
};

class Derived : public Base {
public:
    Derived() : part("Derived 的成員") { cout << "Derived 建構子本體" << '\n'; }
    ~Derived() { cout << "Derived 解構子本體" << '\n'; }

private:
    Tracer part;
};

int main() {
    {
        Derived d;
        cout << "--- d 建好了 ---" << '\n';
    }
    cout << "main 結束" << '\n';
    return 0;
}
\end{example}

\iosample{}{
建立 Base 的成員 \\
Base 建構子本體 \\
建立 Derived 的成員 \\
Derived 建構子本體 \\
--- d 建好了 --- \\
Derived 解構子本體 \\
銷毀 Derived 的成員 \\
Base 解構子本體 \\
銷毀 Base 的成員 \\
main 結束 \\
}

建立 \texttt{d} 時，先建立基底類別那一份（\texttt{Base} 自己也是先成員、
後本體），再依宣告順序建立 \texttt{Derived} 自己的成員，最後執行
\texttt{Derived} 建構子的本體。

銷毀時剛好反過來：先執行 \texttt{Derived} 解構子的本體，銷毀
\texttt{Derived} 的成員，最後才輪到 \texttt{Base} 那一份。 道理跟第~4~章一
樣：後建立的可能用到先建立的。 \texttt{Derived} 的建構子本體執行時，基底
和成員都已經建好，可以放心使用；\texttt{Derived} 的解構子本體執行時，它們
也都還在。

衍生類別的解構子不用、也不能自己呼叫基底的解構子，編譯器會在最後自動呼叫。
建構子則要在初始化列表選一個基底的建構子（上一節的 \texttt{Shape("圓")}），
沒寫的話就用基底的預設建構子；範例裡的 \texttt{Base} 只有預設建構子，所以
不用寫。

\subsection{常見錯誤}

\begin{itemize}
\item \textbf{以為初始化列表的順序決定建立順序}：把 \texttt{Derived} 的初
始化列表寫成 \texttt{: part("Derived 的成員"), Base()}，執行結果完全相同，
基底類別一定最先建立。 \texttt{-Wall} 會警告 \texttt{'Derived::part' will
be initialized after}，下一行是 \texttt{base 'Base'}。 跟第~5~章的成員
順序一樣，初始化列表照「基底、成員宣告順序」來排。
\item \textbf{在解構子裡自己呼叫基底的解構子}：在 \texttt{\textasciitilde
Derived()} 裡寫 \texttt{Base::\textasciitilde Base();}，\texttt{Base} 的解
構子會執行兩次。 \texttt{Base} 管著資源的話，就是第~5~章的重複釋放。
\end{itemize}

\section{虛擬函式}
\label{sec:virtual}

範例~\ref{ex:derived} 裡，透過 \texttt{Shape} 的參考不能呼叫
\texttt{area}，因為 \texttt{Shape} 沒有這個函式。 我們想要的是 C 回顧框
裡函式指標的效果：寫一個只認得 \texttt{Shape} 的函式，傳進圓就用圓的公式，
傳進矩形就用矩形的公式。

做法是在基底類別宣告 \texttt{area}，前面加上 \texttt{virtual}，成為
\textbf{虛擬函式}（virtual function）。 衍生類別寫一個同名、參數和
\texttt{const} 都一樣的函式，叫\textbf{覆寫}（override）。 下面的範例另外
放了一個沒有 \texttt{virtual} 的 \texttt{hello}，用來對照：

\begin{example}[label=ex:virtual]{cpp}{(virtual.cpp) 虛擬函式與一般函式}
#include <iostream>
#include <string>
using namespace std;

class Shape {
public:
    Shape(const string &name) : name(name) {}
    const string &getName() const { return name; }
    virtual double area() const { return 0; }                 // 虛擬函式
    void hello() const { cout << "Shape 的 hello" << '\n'; }  // 一般函式

private:
    string name;
};

class Circle : public Shape {
public:
    Circle(double r) : Shape("圓"), r(r) {}
    double area() const override { return 3.14159 * r * r; }
    void hello() const { cout << "Circle 的 hello" << '\n'; }

private:
    double r;
};

class Rect : public Shape {
public:
    Rect(double w, double h) : Shape("矩形"), w(w), h(h) {}
    double area() const override { return w * h; }

private:
    double w, h;
};

void report(const Shape &s) {
    cout << s.getName() << ' ' << s.area() << '\n';
}

int main() {
    Circle c(1);
    Rect r(3, 4);
    report(c);
    report(r);

    Shape *p = &c;
    p->hello();                // 一般函式：看 p 的型別
    c.hello();
    p = &r;
    cout << p->area() << '\n'; // 虛擬函式：看 p 指著的物件
    return 0;
}
\end{example}

\iosample{}{
圓 3.14159 \\
矩形 12 \\
Shape 的 hello \\
Circle 的 hello \\
12 \\
}

\texttt{report} 只知道參數是 \texttt{Shape}，\texttt{s.area()} 卻分別呼叫
了 \texttt{Circle} 和 \texttt{Rect} 的版本。 這裡要分兩種型別：

\begin{itemize}
\item \textbf{靜態型別}：變數宣告的型別，編譯時就知道。 \texttt{s} 的靜態
型別是 \texttt{const Shape \&}，\texttt{p} 是 \texttt{Shape *}。
\item \textbf{動態型別}：參考綁著的、指標指著的物件實際的型別，執行時才知
道。 第一次呼叫 \texttt{report} 時 \texttt{s} 的動態型別是
\texttt{Circle}，第二次是 \texttt{Rect}。
\end{itemize}

一般函式看靜態型別，所以 \texttt{p->hello()} 呼叫 \texttt{Shape} 的
\texttt{hello}，就算 \texttt{p} 指著 \texttt{Circle}。 虛擬函式透過參考或
指標呼叫時看動態型別，執行時才決定呼叫哪一個，這叫\textbf{動態繫結}
（dynamic binding）。 同一行 \texttt{s.area()} 依物件不同而有不同的行為，
就是多型。

\texttt{Circle} 的 \texttt{hello} 跟 \texttt{Shape} 的同名，但
\texttt{Shape} 的不是虛擬函式，所以這只是另一個函式，把基底的
\texttt{hello} 遮住了。 \texttt{c.hello()} 用的是 \texttt{Circle} 的，透過
\texttt{Shape *} 用的是 \texttt{Shape} 的，同一個物件兩種行為，容易搞混。
不要在衍生類別重新定義基底的一般函式。

\subsection{\texttt{override} 與 \texttt{final}}

衍生類別的 \texttt{area} 後面寫了 \texttt{override}。 它不改變程式的行
為，只是告訴編譯器「這個函式要覆寫基底的虛擬函式」。 如果基底沒有一樣的
虛擬函式可以覆寫，編譯器就報錯。 例如 \texttt{Circle} 的 \texttt{area} 漏
了 \texttt{const}：

\begin{code}{text}{編譯器的錯誤訊息}
error: 'double Circle::area()' marked 'override', but does not override
\end{code}

沒有 \texttt{const} 的 \texttt{area()} 跟有 \texttt{const} 的是兩個不同的
函式，沒寫 \texttt{override} 的話就會編譯成功，結果是 \texttt{report(c)}
印出 \texttt{圓 0}（見本節的常見錯誤）。 覆寫虛擬函式時一律寫
\texttt{override}。 基底類別裡的那一個寫 \texttt{virtual}；衍生類別裡寫了
\texttt{override} 就不用再寫 \texttt{virtual}，覆寫的函式自動也是虛擬的。

\texttt{final} 禁止再往下覆寫或繼承：函式後面寫 \texttt{final}，衍生類別
就不能再覆寫它；寫成 \texttt{class Circle final : public Shape}，就不能再
有類別繼承 \texttt{Circle}，否則報 \texttt{cannot derive from 'final' base
'Circle'}。 平常不常用。

虛擬函式背後的做法跟 C 回顧框的函式指標差不多：每個有虛擬函式的類別有一
張函式指標表，每個物件裡藏著一個指向自己類別那張表的指標，呼叫
\texttt{s.area()} 時先找到表，再從表裡拿出 \texttt{area} 的位址。 代價是
每個物件多一個指標（本書測試的環境下，只有一個 \texttt{string} 成員的類
別，加了虛擬函式後 \texttt{sizeof} 從 $32$ 變成 $40$），每次呼叫多查一次
表，一般的程式不用在意。 表由編譯器填，不會填錯，也不用強制轉型。

\subsection{常見錯誤}

\begin{itemize}
\item \textbf{忘了 \texttt{virtual}}：基底的 \texttt{area} 沒有寫
\texttt{virtual}，衍生類別也沒寫 \texttt{override}，程式照樣編譯成功。
\texttt{report} 和 \texttt{p->area()} 都用 \texttt{Shape} 的版本，三行都印
出 $0$。 衍生類別有寫 \texttt{override} 的話，編譯器會報上面的
\texttt{marked 'override', but does not override}，馬上就發現了。
\item \textbf{參數或 \texttt{const} 不一樣}：\texttt{Circle} 寫成
\texttt{double area()}，沒有 \texttt{const}，也沒有 \texttt{override}。 這
是一個新的函式，沒有覆寫任何東西，\texttt{report(c)} 印出 \texttt{圓 0}。
\texttt{g++ -Wall} 會警告 \texttt{'virtual double Shape::area() const' was
hidden [-Woverloaded-virtual=]}，但警告很容易被忽略，寫上
\texttt{override} 讓它變成錯誤。
\end{itemize}

\section{抽象類別與虛擬解構子}
\label{sec:abstract}

範例~\ref{ex:virtual} 的 \texttt{Shape::area} 回傳 $0$，其實是一個沒有意
義的值：「形狀」本身沒有面積，只有具體的圓或矩形才有。 這種基底類別無法
提供的函式，可以宣告成\textbf{純虛擬函式}（pure virtual function），在宣
告後面寫 \texttt{= 0}，不寫本體。

有純虛擬函式的類別叫\textbf{抽象類別}（abstract class），不能直接建立物
件，只能拿來當基底類別。 衍生類別要覆寫所有純虛擬函式，才能建立物件：

\begin{example}[label=ex:abstract]{cpp}{(abstract.cpp) 抽象類別 Shape}
#include <iostream>
#include <string>
using namespace std;

class Shape {                                   // 抽象類別
public:
    Shape(const string &name) : name(name) {}
    virtual ~Shape() = default;                 // 虛擬解構子，下面說明
    const string &getName() const { return name; }
    virtual double area() const = 0;            // 純虛擬函式
    virtual double perimeter() const = 0;

private:
    string name;
};

class Circle : public Shape {
public:
    Circle(double r) : Shape("圓"), r(r) {}
    double area() const override { return 3.14159 * r * r; }
    double perimeter() const override { return 2 * 3.14159 * r; }

private:
    double r;
};

class Rect : public Shape {
public:
    Rect(double w, double h) : Shape("矩形"), w(w), h(h) {}
    double area() const override { return w * h; }
    double perimeter() const override { return 2 * (w + h); }

private:
    double w, h;
};

void report(const Shape &s) {
    cout << s.getName() << " 面積 " << s.area() << " 周長 " << s.perimeter() << '\n';
}

int main() {
    Circle c(1);
    Rect r(3, 4);
    report(c);
    report(r);
    // Shape s("形狀");                         // 編譯失敗：Shape 是抽象類別
    return 0;
}
\end{example}

\iosample{}{
圓 面積 3.14159 周長 6.28318 \\
矩形 面積 12 周長 14 \\
}

\begin{itemize}
\item \texttt{Shape} 規定了「每種形狀都要能算面積和周長」，但不說怎麼算。
\texttt{report} 只依賴這個規定，之後多寫一個三角形，\texttt{report} 一行都
不用改。 C 回顧框的 \texttt{switch} 版本則要每個函式都補一個
\texttt{case}。
\item 拿掉最後一行的註解，\texttt{Shape s("形狀");} 編譯失敗，訊息還列出
了是哪些函式讓它變成抽象的：
\begin{code}{text}{編譯器的錯誤訊息}
error: cannot declare variable 's' to be of abstract type 'Shape'
note:   because the following virtual functions are pure within 'Shape':
note:     'virtual double Shape::area() const'
note:     'virtual double Shape::perimeter() const'
\end{code}
\end{itemize}

只有純虛擬函式（加上虛擬解構子）、沒有成員變數的抽象類別，常叫做
\textbf{介面}（interface），它只規定衍生類別要提供哪些函式。

\subsection{虛擬解構子}

\texttt{Shape} 裡有一行 \texttt{virtual \textasciitilde Shape() = default;}。
\texttt{= default} 是第~6~章說的「用編譯器產生的版本」，這裡要的是前面那個
\texttt{virtual}。 原因是物件常常透過基底類別的指標銷毀，例如
\texttt{unique\_ptr<Shape>} 指著一個 \texttt{Circle}。 \texttt{unique\_ptr}
銷毀時對它的指標做 \texttt{delete}，這個指標的型別是 \texttt{Shape *}。
解構子不是虛擬函式的話，\texttt{delete} 就看靜態型別，只呼叫
\texttt{Shape} 的解構子。 下面的 \texttt{Derived} 有一個 \texttt{vector}
成員，基底 \texttt{Base} 的解構子沒有 \texttt{virtual}：

\begin{example}[label=ex:vdtor]{cpp}{(vdtor.cpp) 基底的解構子不是 virtual（錯誤示範）}
#include <iostream>
#include <memory>
#include <vector>
using namespace std;

class Base {
public:
    ~Base() { cout << "~Base" << '\n'; }        // 沒有 virtual
    virtual void work() const { cout << "Base::work" << '\n'; }
};

class Derived : public Base {
public:
    Derived() : buf(1000, 7) {}
    ~Derived() { cout << "~Derived" << '\n'; }
    void work() const override { cout << "Derived::work " << buf.size() << '\n'; }

private:
    vector<int> buf;
};

int main() {
    unique_ptr<Base> p = make_unique<Derived>();
    p->work();
    return 0;
}                                               // p 銷毀：delete 一個 Base *
\end{example}

\iosample{}{
Derived::work 1000 \\
\textasciitilde Base \\
}

\texttt{g++ -Wall -Wextra} 沒有任何警告，執行也沒有當掉，但
\texttt{\textasciitilde Derived} 沒有印出來：\texttt{Derived} 的解構子沒有
執行，\texttt{buf} 也就沒有銷毀。 用 ASan 跑：

\begin{code}{text}{用 ASan 編譯後執行（節錄）}
Derived::work 1000
~Base
==2041472==ERROR: AddressSanitizer: new-delete-type-mismatch on 0x503000000040 in thread T0:
  object passed to delete has wrong type:
  size of the allocated type:   32 bytes;
  size of the deallocated type: 8 bytes.
    #3 0x5cb786796532 in main vdtor.cpp:26

allocated by thread T0 here:
    #2 0x5cb786796492 in main vdtor.cpp:23
\end{code}

\texttt{new-delete-type-mismatch} 是 \texttt{delete} 的型別跟
\texttt{new} 的不一樣：第 $23$ 行配置的是 $32$ 個位元組的 \texttt{Derived}，
第 $26$ 行（\texttt{p} 銷毀）卻當成 $8$ 個位元組的 \texttt{Base} 釋放。
照標準，透過基底指標 \texttt{delete} 衍生類別的物件，而基底的解構子不是虛
擬的，是未定義行為。 這個例子實際的後果是 \texttt{buf} 那 $4000$ 個位元組
洩漏了（關掉這項檢查再跑，LeakSanitizer 報的就是
\texttt{Direct leak of 4000 byte(s)}）。

把第 $8$ 行改成 \texttt{virtual \textasciitilde Base() \{ ... \}}，
\texttt{delete} 就看動態型別，先呼叫 \texttt{\textasciitilde Derived}，再
呼叫 \texttt{\textasciitilde Base}（上一節的順序），輸出多一行
\texttt{\textasciitilde Derived}，ASan 也沒有報告。 規則很簡單：\textbf{有虛
擬函式的基底類別，解構子一定要是 \texttt{virtual}}。 不需要做事的話，寫
\texttt{virtual \textasciitilde Shape() = default;} 就好。

第~6~章說過，自己宣告了解構子，編譯器就不產生移動。 虛擬解構子也是，但通
常沒有關係：多型的物件透過指標使用，不會整個複製或移動（原因見下一節）。

\subsection{常見錯誤}

\begin{itemize}
\item \textbf{漏了一個純虛擬函式}：\texttt{Rect} 沒有寫 \texttt{perimeter}，
它就還是抽象類別，\texttt{Rect r(3, 4);} 報 \texttt{cannot declare variable
'r' to be of abstract type 'Rect'}，\texttt{note} 列出還沒實作的
\texttt{'virtual double Shape::perimeter() const'}。
\item \textbf{以為 \texttt{unique\_ptr} 會處理好}：上面的範例用的就是
\texttt{unique\_ptr}，照樣只呼叫 \texttt{\textasciitilde Base}。
\texttt{unique\_ptr} 保證會 \texttt{delete}，但 \texttt{delete} 呼叫哪個解
構子，要看基底的解構子是不是虛擬的。 如果是自己寫 \texttt{Base *p = new
Derived; delete p;}，\texttt{-Wall} 會警告 \texttt{deleting object of
polymorphic class type 'Base' which has non-virtual destructor might cause
undefined behavior}；包在 \texttt{unique\_ptr} 裡反而看不到這個警告。
\item \textbf{在建構子或解構子裡呼叫虛擬函式}：\texttt{Base} 的建構子執行
時，\texttt{Derived} 的部分還沒建立（第~\ref{sec:inherit-order}~節），所以
這時呼叫虛擬函式，用的是 \texttt{Base} 的版本，不會呼叫到
\texttt{Derived} 的。 解構子也一樣。 呼叫的是純虛擬函式的話，結果是未定義
行為。 不要在建構子和解構子裡依賴虛擬函式（本章練習第~1~題）。
\end{itemize}

\section{用 \texttt{vector<unique\_ptr<Base>>} 存放多型物件}
\label{sec:poly-container}

C 回顧框的 \texttt{switch} 版本可以把所有形狀放進一個
\texttt{struct Shape} 陣列。 用繼承之後，圓和矩形是不同的型別，大小也不一
樣，不能直接放進同一個 \texttt{vector}。 能放進去的是指向它們的指標：
\texttt{vector<unique\_ptr<Shape>>}。 每個元素都是 \texttt{unique\_ptr<Shape>}，
指著的物件可以是任何一種形狀。

下面的程式讀入一串形狀，算出總面積和最大的那一個。 \texttt{Shape}、
\texttt{Circle}、\texttt{Rect} 跟範例~\ref{ex:abstract} 相同：

\begin{example}[label=ex:poly-vector]{cpp}{(shapes.cpp) 把不同的形狀放進同一個 vector}
#include <iostream>
#include <memory>
#include <string>
#include <vector>
using namespace std;

// Shape、Circle、Rect 同範例 abstract.cpp，這裡省略

unique_ptr<Shape> makeShape(const string &kind) {
    if (kind == "circle") {
        double r;
        cin >> r;
        return make_unique<Circle>(r);
    }
    if (kind == "rect") {
        double w, h;
        cin >> w >> h;
        return make_unique<Rect>(w, h);
    }
    return nullptr;
}

int main() {
    vector<unique_ptr<Shape>> shapes;
    string kind;
    while (cin >> kind) {
        unique_ptr<Shape> s = makeShape(kind);
        if (s == nullptr) {
            break;                              // 不認得的種類就停止讀取
        }
        shapes.push_back(std::move(s));
    }

    double total = 0;
    const Shape *largest = nullptr;             // 只用來看，不擁有
    for (const unique_ptr<Shape> &s : shapes) {
        cout << s->getName() << ' ' << s->area() << '\n';
        total += s->area();
        if (largest == nullptr || s->area() > largest->area()) {
            largest = s.get();
        }
    }
    cout << "總面積 " << total << '\n';
    if (largest != nullptr) {
        cout << "最大的是" << largest->getName() << '\n';
    }
    return 0;
}                                               // shapes 銷毀，每個形狀跟著銷毀
\end{example}

\iosample{
circle 1 \\
rect 3 4 \\
circle 2 \\
}{
圓 3.14159 \\
矩形 12 \\
圓 12.5664 \\
總面積 27.708 \\
最大的是圓 \\
}

\begin{itemize}
\item \texttt{makeShape} 依讀到的種類建立不同的物件，回傳型別都是
\texttt{unique\_ptr<Shape>}。 \texttt{make\_unique<Circle>(r)} 建立的是
\texttt{unique\_ptr<Circle>}，可以自動轉成 \texttt{unique\_ptr<Shape>}，就
像 \texttt{Circle *} 可以轉成 \texttt{Shape *}。 這種「依資料決定建立哪個
衍生類別」的函式，常叫做\textbf{工廠函式}（factory function）。
\item \texttt{unique\_ptr} 不能複製，所以放進 \texttt{vector} 時要
\texttt{std::move}（第~4~章），走訪時用 \texttt{const} 參考。
\texttt{s->area()} 透過指標呼叫虛擬函式，依每個元素實際的型別呼叫。
\item \texttt{largest} 是一般指標，用 \texttt{s.get()} 取得。 它只是記住
「最大的是哪一個」，不擁有物件，也不負責 \texttt{delete}；物件仍然由
\texttt{vector} 裡的 \texttt{unique\_ptr} 擁有。 \texttt{shapes} 還在，
\texttt{largest} 就有效。
\item \texttt{shapes} 銷毀時，每個 \texttt{unique\_ptr} 透過
\texttt{Shape *} 刪除自己的物件；\texttt{Shape} 有虛擬解構子，所以都正確
銷毀，用 ASan 跑沒有報告。
\end{itemize}

多一種形狀，只要寫一個新的衍生類別，再在 \texttt{makeShape} 加一個
\texttt{if}；\texttt{main} 裡的迴圈不用改。

\subsection{物件切割}

如果不用指標，直接寫 \texttt{vector<Shape>} 會怎樣？ \texttt{Shape} 是抽
象類別的話，編譯器不讓你這樣做。 換成範例~\ref{ex:virtual} 那個不是抽象的
\texttt{Shape}（\texttt{area} 回傳 $0$），把它的 \texttt{main} 換成下面這
段（記得 \texttt{\#include <vector>}）：

\begin{code}{cpp}{(slicing.cpp) 用傳值的方式存 Circle 和 Rect（錯誤示範）}
int main() {
    Circle c(1);
    Shape s = c;               // 只複製 c 裡 Shape 的那一部分
    cout << s.getName() << ' ' << s.area() << '\n';

    vector<Shape> v;
    v.push_back(c);
    v.push_back(Rect(3, 4));
    for (const Shape &x : v) {
        report(x);
    }
    return 0;
}
\end{code}

編譯沒有警告，執行結果是 \texttt{圓 0}、\texttt{圓 0}、\texttt{矩形 0} 三
行，面積全部變成 $0$。 \texttt{Shape s = c;} 建立的是一個
\texttt{Shape} 物件，用 \texttt{Shape} 的複製建構子，只複製 \texttt{c} 裡
\texttt{Shape} 的那一部分（\texttt{name}），\texttt{r} 被切掉了。 \texttt{s}
從頭到尾就是一個 \texttt{Shape}，呼叫的當然是 \texttt{Shape::area}。 這叫
\textbf{物件切割}（object slicing）。 \texttt{vector<Shape>} 的每個元素也
都是 \texttt{Shape}，\texttt{push\_back(c)} 一樣只複製 \texttt{Shape} 的部
分。 傳值的參數 \texttt{void report(Shape s)} 也會切割。

要保留動態型別，就不能複製出新的物件，只能透過參考或指標使用原本的物件：
參數用 \texttt{const Shape \&}，容器用 \texttt{vector<unique\_ptr<Shape>>}。
這也是基底類別最好寫成抽象類別的一個理由：\texttt{Shape} 是抽象的時候，
上面的 \texttt{Shape s = c;} 直接編譯失敗（\texttt{cannot declare variable
's' to be of abstract type 'Shape'}），傳值的參數也是
（\texttt{cannot declare parameter 's' to be of abstract type 'Shape'}）。

\subsection{\texttt{dynamic\_cast}}

偶爾需要知道一個 \texttt{Shape} 實際上是不是 \texttt{Circle}，可以寫
\texttt{const Circle *c = dynamic\_cast<const Circle *>(s.get());}。
\texttt{s} 指著的是 \texttt{Circle} 時，\texttt{c} 指向它，可以用
\texttt{Circle} 自己的成員；不是的話，\texttt{c} 是 \texttt{nullptr}。
\texttt{dynamic\_cast} 只能用在有虛擬函式的類別上。 程式裡常常需要
\texttt{dynamic\_cast}，通常表示基底類別少了一個虛擬函式：與其問「你是不是
圓」再分別處理，不如在 \texttt{Shape} 加一個虛擬函式，讓每種形狀自己處理。

\subsection{常見錯誤}

\begin{itemize}
\item \textbf{\texttt{vector<Shape *>} 加上 \texttt{new}}：用一般指標也能做到
多型，但每個元素都要記得 \texttt{delete}，中途 \texttt{return} 或出例外就洩
漏，這是第~4~章的問題。 用 \texttt{unique\_ptr}。
\item \textbf{把 \texttt{s.get()} 交給另一個 \texttt{unique\_ptr}}：
\texttt{unique\_ptr<Shape> q(shapes[0].get());} 讓兩個
\texttt{unique\_ptr} 都以為自己擁有同一個物件，最後 \texttt{delete} 兩次。
\texttt{get()} 取出來的指標只能拿來看。
\end{itemize}

\section{什麼時候不要用繼承}
\label{sec:composition}

繼承容易被拿來「重複使用程式碼」：某個類別已經有我要的函式，繼承它就全部
有了。 下面是一個常見的例子，想要一個堆疊（只能從頂端放入和取出），於是
直接繼承 \texttt{vector<int>}。 對照的 \texttt{Stack} 則是把
\texttt{vector} 放成成員，這叫\textbf{組合}（composition）：

\begin{example}[label=ex:stack]{cpp}{(stack.cpp) 繼承 vector 與把 vector 當成員}
#include <iostream>
#include <vector>
using namespace std;

class BadStack : public vector<int> {   // 繼承 vector，不建議
public:
    void push(int x) { push_back(x); }
    void pop() { pop_back(); }
    int top() const { return back(); }
};

class Stack {                           // vector 當成員
public:
    void push(int x) { data.push_back(x); }
    void pop() { data.pop_back(); }
    int top() const { return data.back(); }
    int size() const { return (int)data.size(); }

private:
    vector<int> data;
};

int main() {
    BadStack b;
    b.push(1);
    b.push(2);
    b[0] = 99;                  // 連堆疊底部都改得到
    b.push_back(3);
    cout << b.top() << ' ' << b[0] << ' ' << b.size() << '\n';

    Stack s;
    s.push(1);
    s.push(2);
    s.push(3);
    s.pop();
    // s[0] = 99;               // 編譯失敗
    cout << s.top() << ' ' << s.size() << '\n';
    return 0;
}
\end{example}

\iosample{}{
3 99 3 \\
2 2 \\
}

\begin{itemize}
\item \texttt{BadStack} 是一種 \texttt{vector}，所以 \texttt{vector} 的
\texttt{public} 成員函式它全都有：\texttt{b[0] = 99} 改了堆疊最底下的元素，
\texttt{push\_back} 繞過了 \texttt{push}。 「只能從頂端存取」這條規則，類
別本身擋不住。
\item \texttt{Stack} 有一個 \texttt{vector}（has-a），外面只看得到
\texttt{Stack} 自己開放的四個函式。 \texttt{s[0] = 99} 報 \texttt{no match
for 'operator[]' (operand types are 'Stack' and 'int')}。 之後想把
\texttt{vector} 換成別的容器，用 \texttt{Stack} 的程式也不用改。
\item 標準函式庫的容器沒有虛擬解構子，本來就不是設計來當基底類別的。
\end{itemize}

判斷的方法是第~\ref{sec:derived}~節那句話：「\texttt{A} 是一種
\texttt{B}」說得通，而且任何需要 \texttt{B} 的地方交給它一個 \texttt{A} 都
不會出問題，才用 \texttt{public} 繼承。 只是「\texttt{A} 用到
\texttt{B}」或「\texttt{A} 有一個 \texttt{B}」，就把 \texttt{B} 放成成員。
堆疊不是一種 \texttt{vector}，任何 \texttt{vector} 能做的事堆疊都不該做。

\subsection{型別標記加 \texttt{switch} 有時候比較好}

本章開頭 C 回顧框的第一種寫法，在 C++ 裡也不算錯。 型別標記改用
\texttt{enum class Kind \{ Circle, Rect, Triangle \};}，\texttt{switch} 裡
沒有 \texttt{default}，只寫了 \texttt{Circle} 和 \texttt{Rect} 兩個
\texttt{case} 的話，\texttt{g++ -Wall} 會警告 \texttt{enumeration value
'Triangle' not handled in switch [-Wswitch]}。 多一種形狀時，編譯器會指出
每一個要補的 \texttt{switch}。 兩種寫法各有適合的情況：

\begin{itemize}
\item 種類會一直增加，對每個物件做的操作比較固定（例如外掛、遊戲裡的各種
角色）：用虛擬函式。 多一種，只要多寫一個類別。
\item 種類固定只有幾種，操作卻會一直增加（例如一個小計算機的運算子種類、
一份資料只有三種紀錄格式）：用型別標記加 \texttt{switch}。 多一種操作，只
要多寫一個函式，不用改每個類別。
\end{itemize}

整體的建議是：先想清楚是不是真的需要多型，需要的話保持繼承階層又淺又
小，一個抽象的基底類別加上幾個衍生類別，就能處理大部分情況。 只是想重複
使用程式碼，用組合。

\subsection{常見錯誤}

\begin{itemize}
\item \textbf{為了少寫幾行而繼承}：\texttt{class Student : public
vector<int>}，因為學生「有一些分數」。 學生不是一種 \texttt{vector}，分
數應該是成員。
\item \textbf{很深的繼承階層}：\texttt{Shape} 下面有 \texttt{Polygon}，
下面有 \texttt{Quadrilateral}，下面有 \texttt{Rect}，再下面有
\texttt{Square}。 讀一個函式要往上翻好幾層才知道它從哪裡來，改基底的一行
可能影響所有衍生類別。 層數多到要畫圖才看得懂時，通常可以用組合拆開。
\end{itemize}

\section{本章練習}
\label{sec:ch7-exercises}

\begin{drills}[label=drills:ch7]
\item 預測下面程式的輸出。
\begin{code}{cpp}{練習：建構、虛擬函式與切割}
#include <iostream>
using namespace std;

struct A {
    A() { cout << "A "; hi(); }
    virtual ~A() { cout << "~A "; }
    virtual void hi() const { cout << "A.hi "; }
    void name() const { cout << "A.name "; }
};

struct B : A {
    B() { cout << "B "; hi(); }
    ~B() override { cout << "~B "; }
    void hi() const override { cout << "B.hi "; }
    void name() const { cout << "B.name "; }
};

void call(const A &a) {
    a.hi();
    a.name();
}

int main() {
    B b;
    cout << '\n';
    call(b);
    b.name();
    cout << '\n';
    A a = b;
    a.hi();
    cout << '\n';
    return 0;
}
\end{code}
\item 下面的程式編譯失敗，\texttt{g++} 報了四個錯誤。 找出原因並改正，讓
它輸出 \texttt{Lucky: 汪} 和 \texttt{Momo: 喵} 兩行，而且 \texttt{Cat} 仍
然可以改名。 改到編譯成功之後，還有一個問題是 \texttt{g++} 和 ASan 都不會
告訴你的，也一起改掉。
\begin{code}{cpp}{練習：找錯}
#include <iostream>
#include <memory>
#include <string>
#include <vector>
using namespace std;

class Animal {
public:
    Animal(const string &name) : name(name) {}
    virtual string sound() const = 0;
    void speak() const { cout << name << ": " << sound() << '\n'; }

private:
    string name;
};

class Dog : Animal {
public:
    Dog(const string &name) : Animal(name) {}
    string sound() override { return "汪"; }
};

class Cat : public Animal {
public:
    Cat(const string &name) : Animal(name) {}
    string sound() const override { return "喵"; }
    void rename(const string &n) { name = n; }
};

int main() {
    vector<unique_ptr<Animal>> pets;
    pets.push_back(make_unique<Dog>("Lucky"));
    auto cat = make_unique<Cat>("Mimi");
    cat->rename("Momo");
    pets.push_back(std::move(cat));
    for (const unique_ptr<Animal> &p : pets) {
        p->speak();
    }
    return 0;
}
\end{code}
\item 停車場的收費方式：汽車每小時 $40$ 元；機車每小時 $20$ 元，但最多收
$100$ 元。 寫一個抽象類別 \texttt{Vehicle}，建構子收停車的時數，提供
\texttt{getHours()}，以及純虛擬函式 \texttt{kind()}（回傳 \texttt{"汽車"}
或 \texttt{"機車"}）和 \texttt{fee()}；再寫兩個衍生類別 \texttt{Car} 和
\texttt{Motorcycle}。 \texttt{main} 每行讀入種類（\texttt{car} 或
\texttt{moto}）和時數直到輸入結束，不認得的種類略過；全部讀完後，依輸入順
序印出每一輛的種類、時數和費用，最後印出合計。 輸入
\begin{code}{text}{練習 3 的輸入}
car 3
moto 2
moto 8
bus 5
car 1
\end{code}
應該輸出 \texttt{汽車 3 小時 120}、\texttt{機車 2 小時 40}、
\texttt{機車 8 小時 100}、\texttt{汽車 1 小時 40}、\texttt{合計 300} 五行。
\item 下面的程式用 \texttt{g++ -Wall -Wextra} 編譯沒有警告，輸出
\texttt{Ada 990} 和 \texttt{Alan 4990}。 \texttt{Alan} 的是儲蓄帳戶，月底
應該加 $1\%$ 利息、不扣手續費，變成 $5050$ 才對。 說明原因並改正。
\begin{code}{cpp}{練習：月底結算}
#include <iostream>
#include <string>
#include <vector>
using namespace std;

class Account {
public:
    Account(const string &owner, int balance) : owner(owner), balance(balance) {}
    virtual ~Account() = default;
    void deposit(int amount) { balance += amount; }
    int getBalance() const { return balance; }
    virtual void monthEnd() { balance -= 10; }   // 每月手續費 10 元
    void print() const { cout << owner << ' ' << balance << '\n'; }

private:
    string owner;
    int balance;
};

class Savings : public Account {
public:
    Savings(const string &owner, int balance) : Account(owner, balance) {}
    void monthEnd() override { deposit(getBalance() / 100); }   // 利息 1%，免手續費
};

int main() {
    vector<Account> accounts;
    accounts.push_back(Account("Ada", 1000));
    accounts.push_back(Savings("Alan", 5000));
    for (Account &a : accounts) {
        a.monthEnd();
    }
    for (const Account &a : accounts) {
        a.print();
    }
    return 0;
}
\end{code}
\item 下面每一組，判斷應該用 \texttt{public} 繼承、組合（放成成員），還是
型別標記加 \texttt{switch}，並說明理由。
\begin{enumerate}
\item \texttt{Car} 和 \texttt{Engine}（引擎）。
\item \texttt{Student} 和 \texttt{Person}，程式裡有些函式收
\texttt{const Person \&}，要能把學生傳進去。
\item 一個可以改寬和高的 \texttt{Rect}（有 \texttt{setWidth}、
\texttt{setHeight}），想再寫一個 \texttt{Square}。
\item \texttt{Queue}（佇列，只能從尾端放入、從前端取出）和
\texttt{vector<int>}。
\item 一個記帳程式的每筆紀錄只有「收入」和「支出」兩種，之後會一直加新的
報表功能（月報、分類統計、匯出）。
\end{enumerate}
\end{drills}

\addcontentsline{toc}{section}{練習解答}
\section*{練習解答}

\begin{enumerate}
\item 輸出四行：\texttt{A A.hi B B.hi }、\texttt{B.hi A.name B.name }、
\texttt{A.hi }、\texttt{\textasciitilde A \textasciitilde B \textasciitilde A }
（最後一行沒有換行）。 建立 \texttt{b} 時先建立 \texttt{A} 的部分，這時
\texttt{B} 的部分還不存在，所以 \texttt{A} 建構子裡的 \texttt{hi()} 呼叫的
是 \texttt{A.hi}；接著 \texttt{B} 的建構子印 \texttt{B}，這時呼叫
\texttt{hi()} 才是 \texttt{B.hi}。 \texttt{call(b)} 裡 \texttt{a.hi()} 是
虛擬函式，看動態型別，印 \texttt{B.hi}；\texttt{a.name()} 不是虛擬函式，
看靜態型別 \texttt{A}，印 \texttt{A.name}。 \texttt{b.name()} 的靜態型別
是 \texttt{B}，印 \texttt{B.name}。 \texttt{A a = b;} 是物件切割，用編譯器
產生的複製建構子（不印東西），\texttt{a} 就是一個 \texttt{A}，所以
\texttt{a.hi()} 印 \texttt{A.hi}。 \texttt{main} 結束時先銷毀後建立的
\texttt{a}（\texttt{\textasciitilde A}），再銷毀 \texttt{b}
（\texttt{\textasciitilde B} 然後 \texttt{\textasciitilde A}）。
\item 四個錯誤訊息來自三個問題：
(1) \texttt{class Dog : Animal} 少了 \texttt{public}，是 \texttt{private}
繼承，\texttt{unique\_ptr<Dog>} 不能轉成 \texttt{unique\_ptr<Animal>}。 這
個錯誤報在 \texttt{pets.push\_back(...)} 那一行，訊息是
\texttt{no matching function for call to ...push\_back(...)}，後面的
\texttt{note} 很長，看不太出原因。 改成 \texttt{class Dog : public Animal}。
(2) \texttt{Dog::sound} 少了 \texttt{const}，報 \texttt{marked 'override',
but does not override}。 寫了 \texttt{override}，這個錯誤才被抓到。 因為
沒有覆寫到，\texttt{Dog} 還是抽象類別，所以第四個錯誤
\texttt{invalid new-expression of abstract class type 'Dog'}（位置在
\texttt{unique\_ptr.h}，是 \texttt{make\_unique<Dog>} 裡的 \texttt{new}）
也是同一個原因。
(3) \texttt{Cat::rename} 直接改 \texttt{Animal} 的 \texttt{private} 成員
\texttt{name}。 在 \texttt{Animal} 加一個 \texttt{protected} 的
\texttt{void setName(const string \&n) \{ name = n; \}}，\texttt{rename}
改成呼叫 \texttt{setName(n)}。 不要把 \texttt{name} 改成
\texttt{protected} 或 \texttt{public}。
還有一個問題：\texttt{Animal} 有虛擬函式，物件透過
\texttt{unique\_ptr<Animal>} 銷毀，解構子卻不是虛擬的。 這個程式裡
\texttt{Dog}、\texttt{Cat} 沒有自己的成員，大小跟 \texttt{Animal} 一樣，所以
ASan 剛好沒有報告，但這仍然是未定義行為；之後替 \texttt{Cat} 加一個
\texttt{string} 成員，它就會洩漏。 在 \texttt{Animal} 加上
\texttt{virtual \textasciitilde Animal() = default;}。
\item 一種寫法：
\begin{code}{cpp}{練習 3 參考解答（類別部分）}
class Vehicle {
public:
    Vehicle(int hours) : hours(hours) {}
    virtual ~Vehicle() = default;
    int getHours() const { return hours; }
    virtual string kind() const = 0;
    virtual int fee() const = 0;

private:
    int hours;
};

class Car : public Vehicle {
public:
    Car(int hours) : Vehicle(hours) {}
    string kind() const override { return "汽車"; }
    int fee() const override { return 40 * getHours(); }
};

class Motorcycle : public Vehicle {
public:
    Motorcycle(int hours) : Vehicle(hours) {}
    string kind() const override { return "機車"; }
    int fee() const override {
        int f = 20 * getHours();
        return f > 100 ? 100 : f;
    }
};
\end{code}
時數是每輛車都有的，放在 \texttt{Vehicle}，\texttt{private} 加上 getter；
收費方式每種不同，寫成純虛擬函式。 \texttt{Vehicle} 有虛擬函式，又透過
\texttt{unique\_ptr<Vehicle>} 銷毀，所以要有虛擬解構子。 \texttt{main} 照範
例~\ref{ex:poly-vector} 的寫法：寫一個工廠函式 \texttt{makeVehicle(type,
hours)}，\texttt{"car"} 回傳 \texttt{make\_unique<Car>(hours)}，
\texttt{"moto"} 回傳 \texttt{make\_unique<Motorcycle>(hours)}，其他回傳
\texttt{nullptr}；\texttt{while (cin >> type >> hours)} 讀入，不是
\texttt{nullptr} 才 \texttt{push\_back(std::move(v))} 進
\texttt{vector<unique\_ptr<Vehicle>> lot}；最後走訪 \texttt{lot}，印出
\texttt{v->kind()}、\texttt{v->getHours()}、\texttt{v->fee()} 並累加合計。

\item \texttt{vector<Account>} 的元素都是 \texttt{Account}。
\texttt{push\_back(Savings("Alan", 5000))} 只把 \texttt{Savings} 裡
\texttt{Account} 的部分複製進去，這是物件切割，元素已經不是
\texttt{Savings} 了，所以 \texttt{monthEnd} 用的是 \texttt{Account} 的版本，
扣了 $10$ 元手續費。 \texttt{Account} 不是抽象類別，所以編譯器不會擋。
改法是元素改存指標：
\begin{code}{cpp}{練習 4：改用 vector<unique\_ptr<Account>>}
vector<unique_ptr<Account>> accounts;
accounts.push_back(make_unique<Account>("Ada", 1000));
accounts.push_back(make_unique<Savings>("Alan", 5000));
for (const unique_ptr<Account> &a : accounts) {
    a->monthEnd();
}
for (const unique_ptr<Account> &a : accounts) {
    a->print();
}
\end{code}
記得 \texttt{\#include <memory>}。 改完輸出 \texttt{Ada 990} 和
\texttt{Alan 5050}。 第一個迴圈的 \texttt{a} 是 \texttt{const} 參考，但
\texttt{const} 的是 \texttt{unique\_ptr} 本身（不能改指向誰），不是它指著
的 \texttt{Account}，所以還是可以呼叫會改餘額的 \texttt{monthEnd}。
\item
(a) 組合。 汽車有一個引擎，不是一種引擎，\texttt{Engine} 放成
\texttt{Car} 的成員。
(b) \texttt{public} 繼承。 學生是一種人，需要 \texttt{Person} 的地方交給它
學生也說得通，正好是繼承要表達的關係。
(c) 不要讓 \texttt{Square} 繼承這個 \texttt{Rect}。 正方形在數學上是一種矩
形，但這個 \texttt{Rect} 可以只改寬不改高：一個函式收 \texttt{Rect \&}，
呼叫 \texttt{setWidth(5)} 之後預期高度不變，傳進正方形就出問題（要嘛高度
跟著變，要嘛它不再是正方形）。 「任何需要 \texttt{Rect} 的地方都能交給它
一個 \texttt{Square}」不成立。 可以讓 \texttt{Square} 和 \texttt{Rect} 都
直接繼承 \texttt{Shape}，或者讓形狀建立後不能改大小。
(d) 組合，跟 \texttt{Stack} 一樣。 佇列不該讓人用 \texttt{[]} 改中間的元
素。
(e) 型別標記加 \texttt{switch}。 種類固定只有兩種，會增加的是操作；每多一
個報表就是多一個函式，用 \texttt{enum class} 的話，漏了哪一種
\texttt{-Wall} 也會提醒。
\end{enumerate}

\chapter{模板與 STL 容器}
\label{ch:templates-stl}

第~3~章說 \texttt{vector<int>} 的角括號是「模板」，當時只要把
\texttt{vector<int>} 整個當成一個型別名稱。 這一章先自己寫模板：同一份
函式或類別，交給編譯器替 \texttt{int}、\texttt{double}、\texttt{string}
各產生一份。 看懂模板之後，標準函式庫的宣告和那些很長的錯誤訊息也就讀得
懂了。

後半章介紹解題和資料結構最常用的容器：\texttt{pair}、\texttt{set} 與
\texttt{map}、\texttt{unordered\_set} 與 \texttt{unordered\_map}、
\texttt{stack}、\texttt{queue}、\texttt{deque}，以及
\texttt{priority\_queue}。 每一種都說明它解決什麼問題、各個操作要花多少
時間，以及常見的錯誤。 這些容器在標準函式庫裡合稱 STL（Standard
Template Library）。 走訪一律用範圍 \texttt{for}，迭代器和演算法留到
第~9~章。

\section{函式模板}
\label{sec:func-template}

\begin{crecap}[巨集與 \texttt{void *}]
C 要寫一個「什麼型別都能用」的函式，常見的做法有兩種。 第一種是巨集：

\begin{code}{c}{用巨集寫 MAX}
#define MAX(a, b) ((a) > (b) ? (a) : (b))

int i = 5, j = 3;
int m = MAX(i++, j);       // m 是 6，i 變成 7
\end{code}

巨集只是文字替換，展開後是 \texttt{((i++) > (j) ? (i++) : (j))}，
\texttt{i++} 執行了兩次，所以 \texttt{m} 是 $6$、\texttt{i} 是 $7$，不是
預期的 $5$ 和 $6$。 巨集也不檢查型別，引數是什麼都照樣展開。

第二種是 \texttt{void *}，例如 \texttt{qsort}。 比較函式收兩個
\texttt{const void *}，自己轉回真正的型別：

\begin{code}{c}{qsort 的比較函式用錯型別}
int cmpInt(const void *a, const void *b) {
    int x = *(const int *)a, y = *(const int *)b;
    return (x > y) - (x < y);
}

double d[] = {2.5, 1.5, 3.5, 0.5};
qsort(d, 4, sizeof(double), cmpInt);   // 比較函式是給 int 用的
\end{code}

\texttt{gcc -Wall -Wextra} 沒有任何警告。 \texttt{cmpInt} 把每個
\texttt{double} 的前 $4$ 個位元組當成 \texttt{int} 來比，在本書測試的環境
下，排完還是 \texttt{2.5 1.5 3.5 0.5}，完全沒排好。 轉成 \texttt{void *}
之後，編譯器就不知道原本是什麼型別了。
\end{crecap}

C++ 的做法是\textbf{函式模板}（function template）：把型別也當成一個參
數，寫一份程式，讓編譯器依呼叫時的引數型別產生對應的函式。

\begin{example}[label=ex:max-template]{cpp}{(max\_template.cpp) 函式模板 maxOf}
#include <iostream>
#include <string>
using namespace std;

template <typename T>
T maxOf(T a, T b) {
    return a > b ? a : b;
}

int main() {
    int i = 5, j = 3;
    int m = maxOf(i++, j);
    cout << m << ' ' << i << '\n';

    cout << maxOf(2.5, 1.5) << '\n';
    string a = "apple", b = "banana";
    cout << maxOf(a, b) << '\n';
    cout << maxOf<double>(3, 2.5) << '\n';
    return 0;
}
\end{example}

\iosample{}{
5 6 \\
2.5 \\
banana \\
3 \\
}

\begin{itemize}
\item \texttt{template <typename T>} 宣告一個型別參數 \texttt{T}，下面的函式
裡 \texttt{T} 就當成一個型別來用。 \texttt{typename} 也可以寫成
\texttt{class}，意思完全一樣，標準函式庫的宣告裡兩種都會看到。
\item 呼叫 \texttt{maxOf(i++, j)} 時，編譯器從引數推斷出 \texttt{T} 是
\texttt{int}，產生一個 \texttt{int maxOf(int a, int b)}。 這個步驟叫
\textbf{實例化}（instantiation）。 \texttt{maxOf(2.5, 1.5)} 和
\texttt{maxOf(a, b)} 又各產生一份 \texttt{double} 和 \texttt{string} 的
版本。 跟第~2~章的多載一樣，最後是幾個不同的函式，只是不用自己寫。
\item \texttt{maxOf} 是真的函式，\texttt{i++} 只執行一次，輸出
\texttt{5 6}。 每份實例化出來的函式也都照一般函式檢查型別，不會發生
\texttt{qsort} 那種把 \texttt{double} 當 \texttt{int} 比的事。
\item \texttt{maxOf<double>(3, 2.5)} 在函式名稱後面直接寫出 \texttt{T}，
不用推斷。 \texttt{3} 照一般的規則轉成 \texttt{double}。
\end{itemize}

模板對 \texttt{T} 的要求，就是函式本體裡用到的操作。 \texttt{maxOf} 用了
\texttt{>}，\texttt{T} 就要能用 \texttt{>} 比較。 第~6~章替 \texttt{Fraction} 寫了比較運算子
之後 \texttt{sort} 就能排序，也是這個道理：\texttt{sort} 是一個函式模板，
只要求元素能用 \texttt{<} 比較。 它不需要 \texttt{qsort} 的
\texttt{void *} 和比較函式，比較時的型別在編譯時就確定了。

標準函式庫已經有 \texttt{maxOf} 的正式版本：\texttt{<algorithm>} 裡的
\texttt{max(a, b)} 和 \texttt{min(a, b)}，本章後面的範例會直接用。

\subsection{常見錯誤}

\begin{itemize}
\item \textbf{兩個引數推斷出不同的型別}：\texttt{maxOf(3, 2.5)} 編譯失
敗。 第一個引數說 \texttt{T} 是 \texttt{int}，第二個說是 \texttt{double}，
編譯器不會替你挑：
\begin{code}{text}{編譯器的錯誤訊息（節錄）}
error: no matching function for call to 'maxOf(int, double)'
note:   deduced conflicting types for parameter 'T' ('int' and 'double')
\end{code}
寫成 \texttt{maxOf<double>(3, 2.5)} 或 \texttt{maxOf(3.0, 2.5)}。
\texttt{std::max(3, 2.5)} 也是同樣的錯誤。
\item \textbf{型別不支援函式裡用到的操作}：對兩個沒有 \texttt{>} 的
\texttt{Point} 呼叫 \texttt{maxOf}，錯誤訊息分成兩段：
\begin{code}{text}{編譯器的錯誤訊息（節錄）}
In instantiation of 'T maxOf(T, T) [with T = Point]':
  required from here
error: no match for 'operator>' (operand types are 'Point' and 'Point')
\end{code}
\texttt{In instantiation of} 說的是「替 \texttt{T = Point} 實例化時出了
錯」，\texttt{required from here} 指出是哪一行呼叫造成的，最後一行才是真
正的原因。 錯誤位置指著模板裡面，問題卻多半在呼叫端傳進去的型別。
\item \textbf{傳兩個字串常數}：\texttt{maxOf("apple", "banana")} 的
\texttt{T} 推斷成 \texttt{const char *}，\texttt{>} 比的是兩個位址，不是
內容（跟第~3~章 C 回顧框裡 \texttt{a == "Ada"} 的問題一樣）。 編譯沒有警告，在本
書測試的環境下印出 \texttt{apple}。 先轉成 \texttt{string} 再比。
\item \textbf{把模板的定義放在另一個 \texttt{.cpp} 檔}：照 C 的習慣，標頭
檔只放 \texttt{maxOf} 的宣告，定義寫在 \texttt{maxof.cpp}，分開編譯後連結
失敗，報 \texttt{undefined reference to `int maxOf<int>(int, int)'}。 編譯
\texttt{main.cpp} 時看不到模板的本體，就沒辦法實例化。 模板要整個寫在標
頭檔裡。
\end{itemize}

\section{類別模板與 \texttt{pair}}
\label{sec:class-template}

類別也可以有型別參數。 第~7~章的 \texttt{Stack} 只能放 \texttt{int}，把
\texttt{int} 換成 \texttt{T}，就是一個什麼都能放的堆疊：

\begin{example}[label=ex:stack-template]{cpp}{(stack\_template.cpp) 類別模板 Stack<T>}
#include <iostream>
#include <string>
#include <vector>
using namespace std;

template <typename T>
class Stack {
public:
    void push(const T &x) { data.push_back(x); }
    void pop() { data.pop_back(); }
    const T &top() const { return data.back(); }
    bool empty() const { return data.empty(); }
    int size() const { return (int)data.size(); }

private:
    vector<T> data;
};

int main() {
    Stack<int> s;
    s.push(1);
    s.push(2);
    s.push(3);
    s.pop();
    cout << s.top() << ' ' << s.size() << '\n';

    Stack<string> words;
    words.push("hello");
    words.push("world");
    while (!words.empty()) {
        cout << words.top() << '\n';
        words.pop();
    }
    return 0;
}
\end{example}

\iosample{}{
2 2 \\
world \\
hello \\
}

\begin{itemize}
\item 類別前面加 \texttt{template <typename T>}，類別裡就可以用
\texttt{T}。 \texttt{push} 的參數是 \texttt{const T \&}，因為 \texttt{T}
可能是 \texttt{string} 這種大型別，照第~2~章的原則用 \texttt{const} 參考。
\item 使用時要在角括號寫出 \texttt{T}：\texttt{Stack<int>}、
\texttt{Stack<string>}。 \texttt{Stack<int>} 和 \texttt{Stack<string>} 是
兩個不同的型別，彼此不能指定。
\end{itemize}

這就是你從第~3~章開始用的 \texttt{vector}：標準函式庫裡有一份
\texttt{vector} 的類別模板，寫 \texttt{vector<int>} 時，編譯器替
\texttt{int} 實例化一份。 \texttt{v.push\_back(x)}、\texttt{v.size()} 都是
那份模板的成員函式。

\subsection{讀懂標準函式庫的宣告}

\texttt{vector} 在標準裡的宣告大致是：

\begin{code}{cpp}{vector 的宣告（簡化）}
template <class T, class Allocator = allocator<T>>
class vector;
\end{code}

它其實有兩個型別參數。 第二個 \texttt{Allocator} 決定記憶體怎麼配置，有
預設值（跟第~2~章的預設引數一樣），平常不用寫。 所以錯誤訊息裡的 \texttt{vector<int>} 常常寫成
\texttt{std::vector<int, std::allocator<int> >}，第~3~章的錯誤訊息就是這
樣。 對 \texttt{vector<int>} 呼叫 \texttt{push\_back("hi")}，錯誤訊息裡有
這樣一行：

\begin{code}{text}{編譯器的錯誤訊息（節錄）}
note: candidate: 'void std::vector<_Tp, _Alloc>::push_back(const value_type&) [with _Tp = int; _Alloc = std::allocator<int>; value_type = int]' (near match)
\end{code}

\texttt{\_Tp} 和 \texttt{\_Alloc} 是函式庫原始碼裡給型別參數取的名字，方
括號裡列出這次實例化時它們各是什麼。 \texttt{value\_type} 是容器裡的元素
型別。 所以這行的意思是「\texttt{push\_back} 要一個 \texttt{const int \&}」，
傳進去的卻是字串常數。 另外，錯誤訊息裡的 \texttt{std::\_\_cxx11::basic\_string<char>}
就是 \texttt{string}。

\subsection{\texttt{pair} 與結構化綁定}

標準函式庫裡最簡單的類別模板是 \texttt{pair}（在 \texttt{<utility>} 裡）：
\texttt{pair<T, U>} 把兩個值綁在一起，分別叫 \texttt{first} 和
\texttt{second}。 需要「名字和分數」「座標 $(x, y)$」這種兩個一組的資料，
又不想為它寫一個 \texttt{struct} 時，就用 \texttt{pair}。 後面的
\texttt{map} 也用它來放「鍵和值」。

\begin{example}[label=ex:pair]{cpp}{(pair\_demo.cpp) pair、結構化綁定與排序}
#include <algorithm>
#include <iostream>
#include <string>
#include <utility>
#include <vector>
using namespace std;

pair<int, int> minMax(const vector<int> &v) {
    int lo = v[0], hi = v[0];
    for (int x : v) {
        lo = min(lo, x);
        hi = max(hi, x);
    }
    return {lo, hi};
}

int main() {
    pair<string, int> p = {"Ada", 95};
    cout << p.first << ' ' << p.second << '\n';
    p.second += 3;

    auto [name, score] = p;
    cout << name << ' ' << score << '\n';

    auto [lo, hi] = minMax({3, 9, -2, 7});
    cout << lo << ' ' << hi << '\n';

    vector<pair<int, string>> v = {{88, "Grace"}, {95, "Ada"}, {88, "Alan"}};
    sort(v.begin(), v.end());
    for (const auto &[s, n] : v) {
        cout << s << ' ' << n << '\n';
    }
    return 0;
}
\end{example}

\iosample{}{
Ada 95 \\
Ada 98 \\
-2 9 \\
88 Alan \\
88 Grace \\
95 Ada \\
}

\begin{itemize}
\item \texttt{pair<string, int> p = \{"Ada", 95\};} 用大括號給兩個值。
\texttt{p.first}、\texttt{p.second} 是一般的成員變數，可以直接讀寫。
\item 函式要回傳兩個值，C 的做法是多傳兩個指標參數；C++ 可以回傳一個
\texttt{pair}，\texttt{return \{lo, hi\};} 會自動建立
\texttt{pair<int, int>}。
\item \texttt{auto [name, score] = p;} 叫\textbf{結構化綁定}（structured
binding，C++17），把 \texttt{pair} 拆成兩個有名字的變數，比
\texttt{p.first}、\texttt{p.second} 好讀。 這樣寫的 \texttt{name}、
\texttt{score} 是複本，跟 \texttt{auto x = p;} 一樣。 要直接改
\texttt{p} 就寫 \texttt{auto \&[name, score] = p;}，只讀不複製就寫
\texttt{const auto \&[...]}，選法跟第~3~章範圍 \texttt{for} 的表格相同。
\item \texttt{pair} 已經定義好比較運算子：先比 \texttt{first}，相同再比
\texttt{second}，跟字串的字典順序一樣。 所以 \texttt{sort} 之後分數由小
到大，同分的照名字排。 想讓 \texttt{sort} 照某個欄位排，把那個欄位放在
\texttt{first} 就好。
\end{itemize}

\subsection{常見錯誤}

\begin{itemize}
\item \textbf{沒寫型別參數}：\texttt{Stack s;} 報 \texttt{class template
argument deduction failed}。 C++17 可以從建構子的引數推斷型別參數，例如
\texttt{vector v = \{1, 2, 3\};} 推斷成 \texttt{vector<int>}，但
\texttt{Stack s;} 沒有引數可以推斷。 本書一律寫出型別參數。
\item \textbf{在類別外面定義成員函式，少了模板的部分}：類別裡只寫
\texttt{void push(const T \&x);}，外面寫
\texttt{void Stack::push(const T \&x) \{ ... \}}，報
\texttt{'template<class T> class Stack' used without template arguments}
和 \texttt{'T' does not name a type}。 要寫成
\texttt{template <typename T> void Stack<T>::push(const T \&x) \{ ... \}}。
模板的成員函式通常直接寫在類別裡面，省掉這一長串。
\item \textbf{改了結構化綁定的複本}：\texttt{auto [name, score] = p;
score += 5;} 不會改到 \texttt{p.second}。 跟第~3~章「想修改卻用複本走訪」
是同一件事，要改就加 \texttt{\&}。
\end{itemize}

\section{\texttt{set} 與 \texttt{map}}
\label{sec:map-set}

解題常要回答兩種問題：「這個數字出現過嗎？」和「每個單字出現幾次？」
用 \texttt{vector} 存，每次都要從頭找一遍，資料有 $n$ 筆，查一次就要看
$n$ 個元素。 C 的話要自己維持一個排好序的陣列再二分搜尋，插入新資料時還
得搬動後面的元素；或者自己寫一棵二元搜尋樹。

\texttt{set}（在 \texttt{<set>} 裡）是一個\textbf{排好序、不重複}的集合。
加入、刪除、查詢都只要 $O(\log n)$ 的時間，大約是 $\log_2 n$ 次比較：
$n$ 是十萬時大約 $17$ 次，$n$ 是一百萬時大約 $20$ 次。 它的內部是一棵會
自己保持平衡的二元搜尋樹，這些細節是資料結構課的內容，這裡只要會用。

\begin{example}[label=ex:set]{cpp}{(set\_demo.cpp) set 的加入、查詢與刪除}
#include <iostream>
#include <set>
using namespace std;

int main() {
    set<int> s;
    int x;
    while (cin >> x) {
        s.insert(x);
    }
    cout << s.size() << '\n';
    for (int y : s) {
        cout << y << ' ';
    }
    cout << '\n';

    cout << s.count(4) << ' ' << s.count(7) << '\n';
    s.erase(4);
    s.erase(7);                  // 沒有 7，什麼都不做
    if (s.find(4) == s.end()) {
        cout << "4 已經刪掉了" << '\n';
    }
    cout << *s.begin() << ' ' << *s.rbegin() << '\n';
    return 0;
}
\end{example}

\iosample{
5 1 4 1 5 9 2 6 \\
}{
6 \\
1 2 4 5 6 9 \\
1 0 \\
4 已經刪掉了 \\
1 9 \\
}

\begin{itemize}
\item \texttt{s.insert(x)}：加入 \texttt{x}。 已經有了就不加，所以讀了
$8$ 個數字，\texttt{size()} 是 $6$。
\item 範圍 \texttt{for} 一定由小到大走訪，跟加入的順序無關。
\item \texttt{s.count(x)}：\texttt{x} 在不在。 \texttt{set} 不重複，所以
只會是 $0$ 或 $1$，可以直接拿來當條件：\texttt{if (s.count(x))}。
\item \texttt{s.erase(x)}：刪掉 \texttt{x}，不存在也沒關係。
\item \texttt{s.find(x)}：找到時回傳指向那個元素的\textbf{迭代器}
（iterator），找不到時回傳 \texttt{s.end()}。 迭代器的用法跟指標很像，
\texttt{*} 取出它指著的元素；\texttt{end()} 是「最後一個元素的下一個」，
不指著任何元素，專門用來表示「沒找到」。 判斷有沒有找到，一律跟
\texttt{end()} 比，就像第~3~章 \texttt{string} 的 \texttt{find} 跟
\texttt{npos} 比。 迭代器第~9~章再詳細講。
\item \texttt{*s.begin()} 是最小的元素，\texttt{*s.rbegin()} 是最大的，
\texttt{set} 不能是空的。
\end{itemize}

\subsection{\texttt{map}：從鍵查到值}

\texttt{map<K, V>}（在 \texttt{<map>} 裡）存的是一對一對的
\texttt{pair<const K, V>}：\textbf{鍵}（key）不重複，每個鍵對應一個
\textbf{值}（value），依鍵的大小排好。 可以把它想成一個下標不必是整數的
陣列：\texttt{cnt["the"]} 是 \texttt{"the"} 對應的值。 下面的程式計算每
個單字出現幾次：

\begin{example}[label=ex:word-count]{cpp}{(word\_count.cpp) 用 map 計算單字出現次數}
#include <iostream>
#include <map>
#include <string>
using namespace std;

int main() {
    map<string, int> cnt;
    string w;
    while (cin >> w) {
        cnt[w]++;
    }

    for (const auto &[word, n] : cnt) {
        cout << word << ' ' << n << '\n';
    }

    auto it = cnt.find("the");
    if (it != cnt.end()) {
        cout << "the 出現 " << it->second << " 次" << '\n';
    }
    cout << cnt.count("fox") << ' ' << cnt.size() << '\n';
    return 0;
}
\end{example}

\iosample{
the cat and the dog \\
and the bird \\
}{
and 2 \\
bird 1 \\
cat 1 \\
dog 1 \\
the 3 \\
the 出現 3 次 \\
0 5 \\
}

\begin{itemize}
\item \texttt{cnt[w]++}：\texttt{w} 第一次出現時，\texttt{map} 裡還沒有這
個鍵，\texttt{[]} 會先\textbf{新增}一個 \texttt{(w, 0)}，再加 $1$。 之後
再遇到 \texttt{w}，就是把原本的值加 $1$。 數數這種事，這個行為剛好省掉
一個 \texttt{if}。
\item 走訪時每個元素是 \texttt{pair<const string, int>}，用結構化綁定拆成
\texttt{word} 和 \texttt{n}。 跟 \texttt{set} 一樣依鍵由小到大，字串就是
字典順序。
\item \texttt{cnt.find("the")} 回傳指向那個 \texttt{pair} 的迭代器。 它像
指向結構的指標，用 \texttt{it->first}、\texttt{it->second} 取出鍵和值。
\item \texttt{count}、\texttt{erase} 跟 \texttt{set} 一樣，以鍵為準。
\end{itemize}

\texttt{map} 的 \texttt{[]}、\texttt{find}、\texttt{count}、
\texttt{insert}、\texttt{erase} 都是 $O(\log n)$，$n$ 是目前有幾個鍵。

需要重複的元素時，有 \texttt{multiset} 和 \texttt{multimap}，用法類似，
\texttt{count} 可能大於 $1$。

\subsection{\texttt{[]} 查詢時會新增鍵}

\texttt{cnt[w]++} 的新增行為，在「只是想查查看」的時候就變成問題：

\begin{code}{cpp}{用 [] 查詢，map 變大了（錯誤示範）}
map<string, int> score = {{"Ada", 95}, {"Alan", 88}};
cout << score.size() << '\n';            // 2
if (score["Grace"] > 90) {               // 只是想查查看
    cout << "Grace 很高分" << '\n';
}
cout << score.size() << '\n';            // 3
\end{code}

\texttt{score["Grace"]} 找不到 \texttt{Grace}，就新增一個分數 $0$ 的
\texttt{Grace}。 之後再走訪，就多出一行 \texttt{Grace 0}；如果分數本來就可能
是 $0$，也分不出「考 $0$ 分」和「沒有這個人」。 只是查詢的時候，不要用
\texttt{[]}：

\begin{itemize}
\item 只問有沒有：\texttt{score.count("Grace")}。
\item 有的話要拿值：\texttt{auto it = score.find("Grace");}，
\texttt{it != score.end()} 時再用 \texttt{it->second}。
\item 確定一定有：\texttt{score.at("Grace")}。 不存在時丟出例外，程式中
止並印出 \texttt{what():  map::at}，跟第~3~章 \texttt{vector} 的
\texttt{at()} 一樣。
\end{itemize}

因為 \texttt{[]} 可能修改 \texttt{map}，\texttt{const map} 不能用
\texttt{[]}。 函式參數寫成 \texttt{const map<string, int> \&m}，裡面寫
\texttt{m["Ada"]} 會編譯失敗：

\begin{code}{text}{編譯器的錯誤訊息}
error: passing 'const std::map<std::__cxx11::basic_string<char>, int>' as 'this' argument discards qualifiers [-fpermissive]
\end{code}

這是第~5~章的訊息：對 \texttt{const} 物件呼叫了非 \texttt{const} 的成員函
式。 修法是改用 \texttt{find} 或 \texttt{at}，不是把 \texttt{const} 拿掉
（本章練習第~2~題）。

\subsection{常見錯誤}

\begin{itemize}
\item \textbf{修改 \texttt{set} 的元素}：\texttt{*s.begin() = 10;} 報
\texttt{assignment of read-only location}；\texttt{for (int \&x : s)} 報
\texttt{binding reference of type 'int\&' to 'const int' discards
qualifiers}。 \texttt{set} 靠元素的大小決定它放在樹的哪裡，直接改值，順序
就亂了，所以元素一律是 \texttt{const}。 要改，就先 \texttt{erase} 舊的值，
再 \texttt{insert} 新的值。 \texttt{map} 的鍵也一樣是 \texttt{const}
（\texttt{pair<const K, V>}），值則可以改：\texttt{for (auto \&[k, v] : m)}
裡可以寫 \texttt{v++}，不能寫 \texttt{k = ...}。
\item \textbf{自己的型別沒有 \texttt{<}}：\texttt{set<Point>} 加入元素時，
錯誤指著標準函式庫的 \texttt{stl\_function.h}：\texttt{no match for
'operator<' (operand types are 'const Point' and 'const Point')}。 跟第~6~
章 \texttt{sort} 的情況一樣，\texttt{set} 和 \texttt{map} 的鍵要能用
\texttt{<} 比較。 替 \texttt{Point} 寫 \texttt{operator<}，或者直接用
\texttt{pair<int, int>} 當鍵。
\end{itemize}

\section{\texttt{unordered\_set} 與 \texttt{unordered\_map}}
\label{sec:unordered}

\texttt{unordered\_set} 和 \texttt{unordered\_map}（在
\texttt{<unordered\_set>}、\texttt{<unordered\_map>} 裡）的用法跟
\texttt{set}、\texttt{map} 幾乎一樣：\texttt{insert}、\texttt{count}、
\texttt{find}、\texttt{erase}、\texttt{[]} 都有，\texttt{[]} 一樣會新增
鍵。 差別在內部用的是\textbf{雜湊表}（hash table）：用一個雜湊函式把鍵算
成一個位置，直接到那個位置找。 所以：

\begin{itemize}
\item 加入、刪除、查詢\textbf{平均}是 $O(1)$，跟資料有多少筆無關。
\item 元素\textbf{沒有順序}。 走訪的順序跟加入的順序、大小都沒有關係。
\end{itemize}

下面的程式先讀入商品的價格表，再讀入一串要買的東西；不賣的商品印出提示，
同一樣東西只買一次：

\begin{example}[label=ex:unordered]{cpp}{(unordered\_demo.cpp) 用 unordered\_map 查價格，unordered\_set 記住買過的}
#include <iostream>
#include <string>
#include <unordered_map>
#include <unordered_set>
using namespace std;

int main() {
    int n;
    cin >> n;
    unordered_map<string, int> price;
    for (int i = 0; i < n; i++) {
        string item;
        int p;
        cin >> item >> p;
        price[item] = p;
    }

    unordered_set<string> bought;
    int total = 0;
    string item;
    while (cin >> item) {
        auto it = price.find(item);
        if (it == price.end()) {
            cout << item << " 沒有賣" << '\n';
        } else if (bought.count(item) > 0) {
            cout << item << " 已經買過了" << '\n';
        } else {
            bought.insert(item);
            total += it->second;
        }
    }
    cout << "買了 " << bought.size() << " 樣，共 " << total << " 元" << '\n';
    return 0;
}
\end{example}

\iosample{
3 \\
tea 30 \\
bread 45 \\
milk 60 \\
tea bread juice tea milk \\
}{
juice 沒有賣 \\
tea 已經買過了 \\
買了 3 樣，共 135 元 \\
}

這支程式只做查詢，沒有走訪 \texttt{price} 或 \texttt{bought}，所以輸出跟
雜湊表內部怎麼排無關。 換成 \texttt{map} 和 \texttt{set}，輸出完全相同。

\subsection{該用哪一個}

\begin{itemize}
\item 需要\textbf{順序}：依大小走訪、要最小或最大的元素、要找「比 $x$ 大
的第一個」，用 \texttt{map}、\texttt{set}。 雜湊表做不到這些。
\item 只需要「有沒有」「對應到什麼」：兩種都可以。 \texttt{unordered\_} 版
本平均比較快。
\item 「平均」$O(1)$ 不是保證。 很多鍵被算到同一個位置時，一次查詢可能要
看過所有元素，變成 $O(n)$。 一般資料很少遇到，但競賽裡有人會故意出這種測
資；\texttt{map}、\texttt{set} 的 $O(\log n)$ 則是每次都成立。
\end{itemize}

APCS 和一般作業的資料量，$O(\log n)$ 幾乎不會是瓶頸。 不確定的時候，先用
\texttt{map}、\texttt{set}；走訪起來有順序，除錯時也比較好看。

\subsection{常見錯誤}

\begin{itemize}
\item \textbf{以為走訪有固定順序}：用範圍 \texttt{for} 印出
\texttt{unordered\_map} 的內容，順序看起來像是隨便排的。 這個順序沒有任何
保證，換一個編譯器、多加幾個元素（雜湊表變大時會重新排列）都可能改變。
某次執行剛好照加入順序印出來，也不代表下次還是。 輸出需要排序就用
\texttt{map}，或者把內容放進 \texttt{vector} 再排序（第~9~章）。
\item \textbf{用 \texttt{pair} 當鍵}：\texttt{unordered\_set<pair<int,
int>>} 編譯失敗，錯誤訊息很長，關鍵是這一句：
\begin{code}{text}{編譯器的錯誤訊息（節錄）}
error: use of deleted function 'std::hash<std::pair<int, int> >::hash()'
\end{code}
標準函式庫沒有替 \texttt{pair} 準備雜湊函式。 簡單的解法是改用
\texttt{set<pair<int, int>>}；座標的話也可以把 $(i, j)$ 換算成一個整數
\texttt{i * c + j} 當鍵。
\end{itemize}

\section{\texttt{stack}、\texttt{queue} 與 \texttt{deque}}
\label{sec:stack-queue}

第~7~章把 \texttt{vector} 包在 \texttt{Stack} 裡面，只開放
\texttt{push}、\texttt{pop}、\texttt{top}。 標準函式庫的 \texttt{stack}
和 \texttt{queue} 就是這樣做的：裡面放一個容器（預設是 \texttt{deque}），
外面只看得到少數幾個操作。 這兩個在標準裡叫\textbf{容器配接器}
（container adapter）。

\begin{center}
\begin{tabular}{lll}
\toprule
& \texttt{stack}（\texttt{<stack>}） & \texttt{queue}（\texttt{<queue>}） \\
\midrule
順序 & 後進先出（LIFO） & 先進先出（FIFO） \\
放入 & \texttt{push(x)}：放到頂端 & \texttt{push(x)}：放到尾端 \\
看 & \texttt{top()} & \texttt{front()}、\texttt{back()} \\
取出 & \texttt{pop()}：拿掉頂端 & \texttt{pop()}：拿掉前端 \\
\bottomrule
\end{tabular}
\end{center}

兩個都有 \texttt{empty()} 和 \texttt{size()}，所有操作都是 $O(1)$。 注意
\texttt{pop()} 不回傳值，要先用 \texttt{top()} 或 \texttt{front()} 拿到元
素，再 \texttt{pop()}。

\texttt{stack} 適合「最近的先處理」，例如括號配對：遇到左括號就放進去，
遇到右括號就跟最近一個還沒配對的左括號比。

\begin{example}[label=ex:brackets]{cpp}{(brackets.cpp) 用 stack 檢查括號配對}
#include <iostream>
#include <stack>
#include <string>
using namespace std;

bool balanced(const string &s) {
    stack<char> st;
    for (char c : s) {
        if (c == '(' || c == '[' || c == '{') {
            st.push(c);
        } else if (c == ')' || c == ']' || c == '}') {
            if (st.empty()) {
                return false;            // 右括號多了
            }
            char open = st.top();
            st.pop();
            if ((c == ')' && open != '(') || (c == ']' && open != '[') ||
                (c == '}' && open != '{')) {
                return false;            // 種類對不上
            }
        }
    }
    return st.empty();                   // 左括號沒有全部配完也不行
}

int main() {
    string line;
    while (getline(cin, line)) {
        cout << (balanced(line) ? "yes" : "no") << '\n';
    }
    return 0;
}
\end{example}

\iosample{
([]\{\}) \\
([)] \\
((x) \\
a{[}i] = f(b{[}j]); \\
}{
yes \\
no \\
no \\
yes \\
}

三種不合法的情況各有一個 \texttt{return false}：右括號出現時沒有左括號可
以配、種類對不上、讀完之後還剩左括號。 呼叫 \texttt{st.top()} 之前先檢查
\texttt{st.empty()}，原因見本節的常見錯誤。

\subsection{\texttt{queue} 與廣度優先搜尋}

\texttt{queue} 適合「先來的先處理」。 它最常見的用途是\textbf{廣度優先搜
尋}（breadth-first search，BFS）：從起點出發，先處理距離 $1$ 的點，再處理
距離 $2$ 的點……，求出每個點離起點最少要走幾步。 BFS 為什麼正確是資料結
構課的內容，這裡看 \texttt{queue} 在裡面怎麼用。

下面的程式讀入一張 $n$ 個點、$m$ 條邊的無向圖，印出每個點到點 $1$ 的最
短距離，走不到的印 $-1$。 每個點的鄰居存成 \texttt{vector<vector<int>>}，
\texttt{adj[u]} 是和 \texttt{u} 相連的點：

\begin{example}[label=ex:bfs]{cpp}{(bfs.cpp) 用 queue 求每個點到點 1 的距離}
#include <iostream>
#include <queue>
#include <vector>
using namespace std;

int main() {
    int n, m;
    cin >> n >> m;
    vector<vector<int>> adj(n + 1);
    for (int i = 0; i < m; i++) {
        int a, b;
        cin >> a >> b;
        adj[a].push_back(b);
        adj[b].push_back(a);
    }

    vector<int> dist(n + 1, -1);         // -1 表示還沒到過
    queue<int> q;
    dist[1] = 0;
    q.push(1);
    while (!q.empty()) {
        int u = q.front();
        q.pop();
        for (int v : adj[u]) {
            if (dist[v] == -1) {
                dist[v] = dist[u] + 1;
                q.push(v);
            }
        }
    }

    for (int i = 1; i <= n; i++) {
        cout << i << ": " << dist[i] << '\n';
    }
    return 0;
}
\end{example}

\iosample{
6 5 \\
1 2 \\
1 3 \\
2 4 \\
3 4 \\
4 5 \\
}{
1: 0 \\
2: 1 \\
3: 1 \\
4: 2 \\
5: 3 \\
6: -1 \\
}

\begin{itemize}
\item 點從 $1$ 編號到 $n$，所以 \texttt{adj} 和 \texttt{dist} 開 $n + 1$
個，第 $0$ 個不用。
\item 迴圈每一輪從前端拿出一個點 \texttt{u}，把它還沒到過的鄰居設好距離、
放到尾端。 距離 $1$ 的點都排在距離 $2$ 的點前面，所以點是依距離由近到遠
被處理的。
\item \texttt{dist[v] == -1} 同時當作「還沒到過」的記號。 每個點只會放進
\texttt{queue} 一次，整個 BFS 的時間跟點數加邊數成正比。
\end{itemize}

這個「\texttt{front}、\texttt{pop}、看鄰居、\texttt{push}」的迴圈是 BFS
的固定寫法，在格子地圖上也一樣，只是點變成 \texttt{pair<int, int>} 座標
（本章練習第~5~題）。

\subsection{\texttt{deque}}

\texttt{deque}（在 \texttt{<deque>} 裡，念作 deck）是兩端都能放入和取出
的序列，兩端的操作都是 $O(1)$，也能用 \texttt{[]} 存取：

\begin{code}{cpp}{deque 的兩端操作}
deque<int> d = {2, 3};
d.push_front(1);           // 1 2 3
d.push_back(4);            // 1 2 3 4
d.pop_front();             // 2 3 4
cout << d.front() << ' ' << d.back() << ' ' << d[1] << '\n';   // 2 4 3
\end{code}

\texttt{vector} 也能在前端插入，但要把後面所有元素往後搬一格，是 $O(n)$。
需要在兩端增刪時用 \texttt{deque}，例如滑動視窗的題目；只用到一端或只需
要先進先出時，用 \texttt{stack}、\texttt{queue} 把意圖寫清楚。

\subsection{常見錯誤}

\begin{itemize}
\item \textbf{以為 \texttt{pop()} 會回傳元素}：\texttt{int x = s.pop();}
報 \texttt{void value not ignored as it ought to be}。 先
\texttt{int x = s.top();}，再 \texttt{s.pop();}。
\item \textbf{對空的容器用 \texttt{top()}、\texttt{front()}、\texttt{pop()}}：
這三個都不檢查是不是空的，跟 \texttt{vector} 的 \texttt{back()} 一樣是未定
義行為。 在本書測試的環境下，對空的 \texttt{stack} 呼叫 \texttt{top()}，
程式直接當掉（Segmentation fault）。 只要有可能是空的，先檢查
\texttt{empty()}，像 \texttt{balanced} 那樣。
\item \textbf{想走訪 \texttt{stack} 或 \texttt{queue}}：用範圍 \texttt{for}
走訪 \texttt{queue} 編譯失敗，報 \texttt{'class std::queue<int>' has no member named
'begin'}。 配接器故意只開放兩端，需要走訪的話，一開始就該用
\texttt{vector} 或 \texttt{deque}。
\end{itemize}

\section{\texttt{priority\_queue}}
\label{sec:pq}

有些問題每次都要拿出「目前最大（或最小）的那一個」，而且中途還會再放新
的東西進來：一直排序太慢，每次從頭找最大值是 $O(n)$。
\texttt{priority\_queue}（優先佇列，在 \texttt{<queue>} 裡）專門做這件事：

\begin{itemize}
\item \texttt{push(x)}：放入，$O(\log n)$。
\item \texttt{top()}：看目前最大的元素，$O(1)$。
\item \texttt{pop()}：拿掉目前最大的元素，$O(\log n)$。
\end{itemize}

它的內部是一個\textbf{堆積}（heap），存在一個 \texttt{vector} 裡，只保證最
大的在最上面，其他元素沒有排好。

\begin{example}[label=ex:pq]{cpp}{(pq\_demo.cpp) 最大堆積、最小堆積與 pair}
#include <functional>
#include <iostream>
#include <queue>
#include <string>
#include <vector>
using namespace std;

int main() {
    priority_queue<int> big;
    priority_queue<int, vector<int>, greater<int>> small;
    for (int x : {5, 1, 4, 1, 5, 9, 2, 6}) {
        big.push(x);
        small.push(x);
    }
    cout << big.top() << ' ' << small.top() << ' ' << big.size() << '\n';

    while (!small.empty()) {
        cout << small.top() << ' ';
        small.pop();
    }
    cout << '\n';

    // (時間, 名字)：時間小的先做，時間一樣就照名字
    priority_queue<pair<int, string>, vector<pair<int, string>>,
                   greater<pair<int, string>>> tasks;
    tasks.push({30, "Grace"});
    tasks.push({10, "Alan"});
    tasks.push({30, "Ada"});
    tasks.push({20, "Linus"});
    while (!tasks.empty()) {
        auto [t, name] = tasks.top();
        tasks.pop();
        cout << t << ' ' << name << '\n';
    }
    return 0;
}
\end{example}

\iosample{}{
9 1 8 \\
1 1 2 4 5 5 6 9 \\
10 Alan \\
20 Linus \\
30 Ada \\
30 Grace \\
}

\begin{itemize}
\item \texttt{priority\_queue<int>} 預設是\textbf{最大堆積}，\texttt{top()}
是最大的。
\item 要最小的在上面，得寫出三個型別參數：元素型別、裡面用的容器
（\texttt{vector<int>}）、比較方式 \texttt{greater<int>}（在
\texttt{<functional>} 裡）。 第二個參數有預設值，但要給第三個就得把第二
個也寫出來，跟函式的預設引數一樣不能跳過。
\item 一直 \texttt{top()}、\texttt{pop()} 到空為止，就依序拿出所有元素，
最小堆積拿出來是由小到大。
\item 元素是 \texttt{pair} 時用 \texttt{pair} 的比較：先比
\texttt{first}，相同再比 \texttt{second}。 把要優先比較的東西放在
\texttt{first}，例如 \texttt{(時間, 名字)}；之後學最短路徑時，會在堆積裡放
\texttt{(距離, 點)}。 \texttt{greater} 的型別要跟元素型別一致，寫成
\texttt{greater<pair<int, string>>}。
\end{itemize}

下面是一個典型的用法：有 $n$ 堆東西，每次可以把兩堆合成一堆，代價是兩堆
重量的和，問全部合成一堆的最小總代價。 每次都合併目前最輕的兩堆，就能得
到最小值（為什麼這樣是對的，屬於演算法課的貪婪法）。 每一輪要拿出最小的
兩個、再放回合併後的那一個，正好是最小堆積的工作：

\begin{example}[label=ex:merge]{cpp}{(merge.cpp) 每次合併最輕的兩堆}
#include <functional>
#include <iostream>
#include <queue>
#include <vector>
using namespace std;

int main() {
    int n;
    cin >> n;
    priority_queue<long long, vector<long long>, greater<long long>> pq;
    for (int i = 0; i < n; i++) {
        long long w;
        cin >> w;
        pq.push(w);
    }

    long long cost = 0;
    while (pq.size() > 1) {
        long long a = pq.top();
        pq.pop();
        long long b = pq.top();
        pq.pop();
        cost += a + b;
        pq.push(a + b);
    }
    cout << cost << '\n';
    return 0;
}
\end{example}

\iosample{
4 \\
4 1 3 2 \\
}{
19 \\
}

合併的過程是 $1 + 2 = 3$、$3 + 3 = 6$、$4 + 6 = 10$，總代價
$3 + 6 + 10 = 19$。 每一輪 \texttt{pop} 兩次、\texttt{push} 一次，共
$n - 1$ 輪，總時間是 $O(n \log n)$。 代價累加起來可能很大，所以用
\texttt{long long}。

\subsection{常見錯誤}

\begin{itemize}
\item \textbf{以為 \texttt{top()} 是最小的}：\texttt{priority\_queue<int>}
是最大堆積。 預設的比較方式是 \texttt{less}（小於），結果卻是最大的在上
面，很容易記反。 寫完先拿兩三個數字試一下。
\item \textbf{只寫兩個型別參數}：\texttt{priority\_queue<int,
greater<int>>} 把 \texttt{greater<int>} 放到了「裡面用的容器」的位置，報
\texttt{no type named 'value\_type' in 'struct std::greater<int>'}。 要寫
\texttt{priority\_queue<int, vector<int>, greater<int>>}。
\item \textbf{對空的 \texttt{priority\_queue} 用 \texttt{top()} 或
\texttt{pop()}}：未定義行為，在本書測試的環境下直接當掉。 編譯時加上
\texttt{-D\_GLIBCXX\_ASSERTIONS}，標準函式庫會替這類操作加上檢查，執行時
印出 \texttt{Assertion '!this->empty()' failed.} 並中止，比當掉好查。
迴圈條件寫 \texttt{while (!pq.empty())}，或像 \texttt{merge.cpp} 那樣先確
定裡面還有足夠的元素。 另外，\texttt{priority\_queue} 跟 \texttt{stack} 一
樣不能走訪。
\end{itemize}

\subsection{怎麼選容器}

最後把本章的容器整理在一起。 $n$ 是容器裡的元素個數：

\begin{center}
\begin{tabular}{lll}
\toprule
需要 & 容器 & 主要操作的時間 \\
\midrule
依下標存取、從尾端增減 & \texttt{vector} & $O(1)$ \\
兩端都要增減 & \texttt{deque} & $O(1)$ \\
後進先出 & \texttt{stack} & $O(1)$ \\
先進先出（BFS） & \texttt{queue} & $O(1)$ \\
有沒有出現過，要有順序 & \texttt{set} & $O(\log n)$ \\
鍵對應到值，要有順序 & \texttt{map} & $O(\log n)$ \\
有沒有出現過，不需要順序 & \texttt{unordered\_set} & 平均 $O(1)$ \\
鍵對應到值，不需要順序 & \texttt{unordered\_map} & 平均 $O(1)$ \\
反覆拿出最大或最小的 & \texttt{priority\_queue} & $O(\log n)$，\texttt{top} 是 $O(1)$ \\
\bottomrule
\end{tabular}
\end{center}

先想清楚程式要對資料做哪些操作、做幾次，再挑一個讓最常做的操作最快的容
器。 都不確定時，從 \texttt{vector} 開始。

\section{本章練習}
\label{sec:ch8-exercises}

\begin{drills}[label=drills:ch8]
\item 預測下面程式的輸出。
\begin{code}{cpp}{練習：預測輸出}
#include <iostream>
#include <map>
#include <queue>
#include <set>
#include <stack>
#include <string>
using namespace std;

int main() {
    map<string, int> m;
    m["b"] = 2;
    m["a"] = 1;
    cout << m["c"] << ' ' << m.size() << '\n';
    for (const auto &[k, v] : m) {
        cout << k << v << ' ';
    }
    cout << '\n';

    set<int> s = {3, 1, 3, 2};
    s.insert(2);
    cout << s.size() << ' ' << s.count(3) << ' ' << *s.begin() << '\n';

    stack<int> st;
    queue<int> q;
    priority_queue<int> pq;
    for (int x : {3, 1, 4}) {
        st.push(x);
        q.push(x);
        pq.push(x);
    }
    pq.pop();
    cout << st.top() << ' ' << q.front() << ' ' << pq.top() << '\n';
    return 0;
}
\end{code}
\item 下面的程式先讀入 $n$ 個人的姓名和分數，再讀入一串要查詢的姓名。 它
編譯失敗。 有人把 \texttt{lookup} 參數的 \texttt{const} 拿掉，編譯就過了，
但輸入
\begin{code}{text}{練習 2 的輸入}
3
Ada 95
Alan 88
Linus 0
Ada Grace Linus
\end{code}
卻輸出 \texttt{Ada 95}、\texttt{Grace 查無此人}、\texttt{Linus 查無此人}、
\texttt{共 4 人}。 說明編譯錯誤和兩個錯誤輸出的原因，並正確地改好它，讓
\texttt{Linus} 印出 \texttt{Linus 0}，最後印出 \texttt{共 3 人}。
\begin{code}{cpp}{練習：成績查詢}
#include <iostream>
#include <map>
#include <string>
using namespace std;

int lookup(const map<string, int> &score, const string &name) {
    if (score[name] == 0) {
        return -1;                       // 查無此人
    }
    return score[name];
}

int main() {
    map<string, int> score;
    int n;
    cin >> n;
    for (int i = 0; i < n; i++) {
        string name;
        int s;
        cin >> name >> s;
        score[name] = s;
    }

    string q;
    while (cin >> q) {
        int s = lookup(score, q);
        if (s == -1) {
            cout << q << " 查無此人" << '\n';
        } else {
            cout << q << ' ' << s << '\n';
        }
    }
    cout << "共 " << score.size() << " 人" << '\n';
    return 0;
}
\end{code}
\item 寫一支程式，讀入一段英文直到輸入結束（單字之間以空白或換行隔開，不
用處理標點和大小寫），依字典順序印出每個單字和它出現的次數，最後印出出現
最多次的單字；次數相同時，取字典順序最前面的。 輸入
\texttt{to be or not to be} 和 \texttt{that is the question} 兩行時，最後
一行應該是 \texttt{最多的是 be 2}。
\item 讀入 $n$ 和 $n$ 個整數。 第一行依「第一次出現的順序」印出所有不同的
數字，第二行由小到大印出所有不同的數字，數字之間以空白隔開。 輸入
\texttt{8} 和 \texttt{5 3 5 1 3 9 1 5} 時應輸出 \texttt{5 3 1 9} 和
\texttt{1 3 5 9}。 只用一個 \texttt{set} 為什麼做不到第一行？
\item 迷宮是 $r$ 列 $c$ 行的格子，\texttt{S} 是起點、\texttt{E} 是終點、
\texttt{\#} 是牆、\texttt{.} 是空地，每一步可以往上下左右走一格，不能走進
牆或走出邊界。 讀入 $r$、$c$ 和地圖，印出從 \texttt{S} 走到 \texttt{E} 最
少要幾步，走不到就印 $-1$。 用 \texttt{queue<pair<int, int>>} 做 BFS。 輸入
\begin{code}{text}{練習 5 的輸入}
4 5
S.#..
.##.#
...#E
.#...
\end{code}
應該輸出 \texttt{8}。
\item 讀入 $n$、$k$ 和 $n$ 個整數，由大到小印出最大的 $k$ 個（$k \le n$）。
輸入 \texttt{8 3} 和 \texttt{5 1 4 1 5 9 2 6} 時應輸出 \texttt{9 6 5}。 用
\texttt{priority\_queue} 做，而且堆積裡最多只放 $k + 1$ 個元素。 提示：用
最小堆積保存「目前最大的 $k$ 個」。
\end{drills}

\addcontentsline{toc}{section}{練習解答}
\section*{練習解答}

\begin{enumerate}
\item 輸出四行：\texttt{0 3}、\texttt{a1 b2 c0 }、\texttt{3 1 1}、
\texttt{4 3 3}。 \texttt{m["c"]} 新增了鍵 \texttt{c}，值是 $0$，所以接著印
出的 \texttt{size()} 是 $3$。 C++17 規定 \texttt{<<} 串起來的運算元由左到
右求值，\texttt{m["c"]} 一定在 \texttt{m.size()} 之前。 \texttt{map} 依鍵
走訪，所以是 \texttt{a}、\texttt{b}、\texttt{c}，跟放進去的順序無關。
\texttt{set} 去掉重複後剩 \texttt{1 2 3}，再加入 \texttt{2} 也不會變。
\texttt{stack} 的頂端是最後放的 $4$，\texttt{queue} 的前端是最先放的 $3$，
\texttt{priority\_queue} 拿掉最大的 $4$ 之後頂端是 $3$。
\item 編譯錯誤：\texttt{score} 是 \texttt{const} 參考，\texttt{[]} 可能新
增鍵，不能用在 \texttt{const map} 上，報 \texttt{passing 'const
std::map<...>' as 'this' argument discards qualifiers}。 拿掉
\texttt{const} 只是讓錯誤藏起來：
(1) \texttt{score["Grace"]} 找不到 \texttt{Grace}，新增了一個值為 $0$ 的
\texttt{Grace}，所以最後是 \texttt{共 4 人}；(2) 用「值是 $0$」表示「查無
此人」，分不出真的考 $0$ 分的 \texttt{Linus}。 正確的改法是保留
\texttt{const}，改用 \texttt{find}：
\begin{code}{cpp}{練習 2：lookup 改用 find}
int lookup(const map<string, int> &score, const string &name) {
    auto it = score.find(name);
    if (it == score.end()) {
        return -1;                       // 查無此人
    }
    return it->second;
}
\end{code}
改完輸出 \texttt{Ada 95}、\texttt{Grace 查無此人}、\texttt{Linus 0}、
\texttt{共 3 人}。 \texttt{find} 只查不改，也順便只找了一次，原本的程式
找到的話要找兩次。 這裡仍然用 $-1$ 表示查無此人，前提是分數不會是負的。
\item 一種寫法：
\begin{code}{cpp}{練習 3 參考解答（讀完之後）}
// 讀入的部分同範例 word_count.cpp：while (cin >> w) cnt[w]++;
string best;
int bestCount = 0;
for (const auto &[word, n] : cnt) {
    cout << word << ' ' << n << '\n';
    if (n > bestCount) {
        best = word;
        bestCount = n;
    }
}
cout << "最多的是 " << best << ' ' << bestCount << '\n';
\end{code}
\texttt{map} 本來就依字典順序走訪，印出來就是要求的順序。 找最多次時用
\texttt{>} 而不是 \texttt{>=}：同樣次數的單字，只有字典順序最前面、最早遇
到的那個會被記下來。 \texttt{be} 和 \texttt{to} 都出現 $2$ 次，所以答案是
\texttt{be}。 這裡用 \texttt{map} 而不是 \texttt{unordered\_map}，因為輸出
需要順序。
\item 一種寫法：
\begin{code}{cpp}{練習 4 參考解答（讀入的迴圈）}
unordered_set<int> seen;
vector<int> order;
set<int> sorted;
for (int i = 0; i < n; i++) {
    int x;
    cin >> x;
    if (seen.count(x) == 0) {
        seen.insert(x);
        order.push_back(x);
    }
    sorted.insert(x);
}
// 接著用範圍 for 依序印出 order，再印出 sorted
\end{code}
\texttt{set} 只記得「有哪些數字」，走訪時一定由小到大，加入的順序已經丟
掉了。 要保留第一次出現的順序，就另外用 \texttt{vector} 記下來，
\texttt{seen} 只負責回答「出現過沒有」。 \texttt{seen} 從來不走訪，所以
用 \texttt{unordered\_set} 沒有順序的問題；換成 \texttt{set} 也可以。 其
實 \texttt{sorted} 就能兼做 \texttt{seen} 的工作（\texttt{sorted.count(x)}），
這裡分開寫比較清楚。 行尾會多一個空白，不想要的話照第~3~章練習 4 的寫法
處理。
\item 一種寫法：
\begin{code}{cpp}{練習 5 參考解答（BFS 的部分）}
// g 是 vector<string>，讀法同第 3 章的踩地雷範例
// 先找到 S 的位置 (si, sj)
vector<vector<int>> dist(r, vector<int>(c, -1));
queue<pair<int, int>> q;
dist[si][sj] = 0;
q.push({si, sj});

int di[] = {-1, 1, 0, 0};
int dj[] = {0, 0, -1, 1};
while (!q.empty()) {
    auto [i, j] = q.front();
    q.pop();
    if (g[i][j] == 'E') {
        cout << dist[i][j] << '\n';
        return 0;
    }
    for (int d = 0; d < 4; d++) {
        int ni = i + di[d], nj = j + dj[d];
        if (ni >= 0 && ni < r && nj >= 0 && nj < c && g[ni][nj] != '#' &&
            dist[ni][nj] == -1) {
            dist[ni][nj] = dist[i][j] + 1;
            q.push({ni, nj});
        }
    }
}
cout << -1 << '\n';
\end{code}
跟範例~\ref{ex:bfs} 是同一個迴圈，點換成座標 \texttt{(i, j)}，鄰居換成上
下左右四格，用 \texttt{di}、\texttt{dj} 兩個陣列列出四個方向。 第~3~章踩
地雷範例的邊界檢查在這裡也用得上，另外多檢查「不是牆」和「還沒到過」。
\texttt{auto [i, j] = q.front();} 是複本，下一行 \texttt{pop()} 之後仍然可
以用；寫成 \texttt{auto \&[i, j]} 的話，\texttt{pop()} 之後就是懸置參考了。
\item 一種寫法：
\begin{code}{cpp}{練習 6 參考解答（主要部分）}
priority_queue<int, vector<int>, greater<int>> pq;   // 目前最大的 k 個
for (int i = 0; i < n; i++) {
    int x;
    cin >> x;
    pq.push(x);
    if ((int)pq.size() > k) {
        pq.pop();                                     // 丟掉最小的
    }
}
vector<int> top;
while (!pq.empty()) {
    top.push_back(pq.top());
    pq.pop();
}
// top 是由小到大，倒著印出來就是 9 6 5
\end{code}
最小堆積的頂端是「目前最大的 $k$ 個」裡最小的那一個。 新的數字放進去後
如果超過 $k$ 個，就把最小的丟掉；丟掉的那個一定不會是最後的答案，因為已經
有 $k$ 個比它大或一樣大的數字。 讀完之後，從最小堆積拿出來是由小到大，
所以先存進 \texttt{vector} 再倒著印（第~3~章範例的倒序迴圈）。 堆積裡最多只有 $k + 1$ 個元素，每
次操作 $O(\log k)$，總時間 $O(n \log k)$；把全部 $n$ 個放進最大堆積再拿
$k$ 次也可以，但堆積會有 $n$ 個元素。
\end{enumerate}

\chapter{迭代器與演算法}
\label{ch:iter-algo}

第~3~章開始就一直在寫 \texttt{sort(v.begin(), v.end())}，第~8~章的
\texttt{find} 也回傳一個「跟 \texttt{end()} 比」的東西，當時都說留到這一
章。 這一章先把迭代器講清楚：它是把指標推廣到各種容器上的東西，
\texttt{begin()}、\texttt{end()} 圍出一個範圍，演算法就在這個範圍上工作。

後半章介紹 \texttt{<algorithm>} 裡解題最常用的幾個函式：自訂順序的
\texttt{sort}、寫比較函式很方便的 lambda、搜尋與計數、去掉重複、二分搜尋，
以及枚舉所有排列的 \texttt{next\_permutation}。 最後談什麼時候自己寫迴圈，
什麼時候用現成的演算法。

\section{迭代器：一般化的指標}
\label{sec:iterators}

\begin{crecap}[用指標走訪陣列]
C 走訪陣列除了用下標，也可以用指標：

\begin{code}{c}{用指標走訪陣列}
int a[] = {3, 1, 4, 1, 5};
int n = 5;
for (int *p = a; p != a + n; p++) {
    printf("%d ", *p);
}
\end{code}

\texttt{p} 從 \texttt{a}（第一個元素的位址）開始，\texttt{*p} 取出目前的
元素，\texttt{p++} 移到下一個。 \texttt{a + n} 是「最後一個元素的下一個」
位置，那裡沒有元素，不能用 \texttt{*} 去讀，但可以拿來比較：\texttt{p}
走到那裡，迴圈就結束。 C 允許指標指到陣列結尾的下一格，就是為了這種寫法。
兩個指向同一個陣列的指標還可以相減，\texttt{(a + n) - a} 是 $n$。
\end{crecap}

\textbf{迭代器}（iterator）是同一個想法用在容器上：\texttt{v.begin()} 對
應 \texttt{a}，\texttt{v.end()} 對應 \texttt{a + n}，\texttt{*it} 取出元
素，\texttt{++it} 移到下一個。 寫法跟指標一模一樣，所以上面的迴圈換成
\texttt{vector}、\texttt{set}、\texttt{map} 都能照寫：

\begin{example}[label=ex:iter-basics]{cpp}{(iter\_basics.cpp) 用迭代器走訪 vector 和 set}
#include <iostream>
#include <set>
#include <string>
#include <vector>
using namespace std;

int main() {
    vector<int> v = {3, 1, 4, 1, 5};
    for (vector<int>::iterator it = v.begin(); it != v.end();
         ++it) {
        *it *= 10;
    }
    for (auto it = v.begin(); it != v.end(); ++it) {
        cout << *it << ' ';
    }
    cout << '\n';

    auto mid = v.begin() + 2;
    cout << *mid << ' ' << mid - v.begin() << ' ';
    cout << v.end() - v.begin() << '\n';

    set<string> s = {"pear", "apple", "fig"};
    for (auto it = s.begin(); it != s.end(); ++it) {
        cout << *it << ' ' << it->size() << '\n';
    }
    return 0;
}
\end{example}

\iosample{}{
30 10 40 10 50 \\
40 2 5 \\
apple 5 \\
fig 3 \\
pear 4 \\
}

\begin{itemize}
\item 迭代器的型別寫出來是 \texttt{vector<int>::iterator}，意思是「屬於
\texttt{vector<int>} 的迭代器型別」。 \texttt{set<string>} 的迭代器是
\texttt{set<string>::iterator}，名字更長，所以平常都寫 \texttt{auto}，這
就是第~2~章說 \texttt{auto} 最有用的地方。
\item \texttt{*it *= 10} 透過迭代器修改元素，跟 \texttt{*p *= 10} 一樣。
\item \texttt{vector} 的迭代器能像指標一樣加減整數：\texttt{v.begin() + 2}
指向第 $2$ 個元素，兩個迭代器相減得到它們差幾格，所以
\texttt{mid - v.begin()} 是下標 $2$，\texttt{v.end() - v.begin()} 是元素
個數 $5$。 演算法回傳迭代器時，常用這個減法換回下標。
\item 元素是結構或物件時，\texttt{it->size()} 等於 \texttt{(*it).size()}，
跟指向結構的指標用法相同。 \texttt{map} 的元素是 \texttt{pair}，所以寫
\texttt{it->first}、\texttt{it->second}，第~8~章已經用過。
\end{itemize}

第~3~章的範圍 \texttt{for} 就是這種迴圈的簡寫：\texttt{for (int x : v)}
會被編譯器展開成從 \texttt{v.begin()} 走到 \texttt{v.end()}、每一輪
\texttt{x = *it} 的迴圈。 只是要從頭到尾看一遍時，用範圍 \texttt{for}
就好；需要知道「現在走到哪裡」、要把位置交給別的函式時，才用迭代器。

\subsection{半開區間}

\texttt{begin()} 指著第一個元素，\texttt{end()} 指著最後一個元素的下一
格。 這樣圍出的範圍包含開頭、不包含結尾，數學上寫成 $[\text{first},
\text{last})$，叫\textbf{半開區間}。 標準函式庫的演算法一律收兩個迭代器
\texttt{first}、\texttt{last}，表示 $[\text{first}, \text{last})$ 這一段。
這樣寫，元素個數就是 \texttt{last - first}，不用加減 $1$；\texttt{first
== last} 就是空範圍，空的 \texttt{vector} 迴圈一次也不跑，不必特別處理。

 \texttt{sort(v.begin(), v.begin() + 3)} 只排前 $3$ 個元素，
\texttt{sort(v.begin() + 1, v.end())} 排第 $1$ 個以後的元素。 C 的陣列也
能交給演算法，指標本身就是迭代器：\texttt{int a[100]} 的前 $n$ 個元素用
\texttt{sort(a, a + n)} 排序。 建立 \texttt{vector} 時也可以給一個範圍，
\texttt{vector<int> w(v.begin(), v.begin() + 3);} 複製 \texttt{v} 的前
$3$ 個元素。

\subsection{迭代器的種類}

指標能做的事，不是每一種迭代器都能做。 \texttt{vector} 的元素一個接一個
放在連續的記憶體裡，所以「往後跳 $k$ 格」就是位址加上 $k$ 個元素的大小，
跟指標一樣是 $O(1)$。 \texttt{set} 和 \texttt{map} 是一棵樹，「下一個元
素」要沿著樹走過去，沒辦法一步跳 $k$ 格。 標準函式庫依照能做的操作，把
迭代器分成幾類，這裡只需要記三種：隨機存取（random access）、雙向
（bidirectional）和單向（forward）。

\begin{center}
\begin{tabular}{lll}
\toprule
種類 & 能做的事 & 容器 \\
\midrule
隨機存取 & 跟指標一樣：\texttt{++}、\texttt{--}、\texttt{+ k}、相減 & \texttt{vector}、\texttt{deque}、\texttt{string}、陣列 \\
雙向 & \texttt{++}、\texttt{--} & \texttt{set}、\texttt{map}、\texttt{list} \\
單向 & 只有 \texttt{++} & \texttt{unordered\_set}、\texttt{unordered\_map} \\
\bottomrule
\end{tabular}
\end{center}

每一種都能用 \texttt{*}、\texttt{==}、\texttt{!=}。 演算法也在宣告裡說明
它需要哪一種迭代器。 \texttt{sort} 排序時要一直跳到範圍的中間、算兩個位置
差幾格，所以需要隨機存取迭代器。 \texttt{set} 不用排序（本來就排好了）；
\texttt{list}（雙向鏈結串列，本書沒有介紹）要排序就用它自己的成員函式
\texttt{a.sort()}。 交給 \texttt{std::sort} 會編譯失敗，見本節的常見錯誤。

\subsection{迭代器失效}

迭代器跟指標一樣，指著某一塊記憶體裡的元素。 第~4~章說過，
\texttt{push\_back} 讓 \texttt{size()} 超過 \texttt{capacity()} 時，
\texttt{vector} 會換一塊記憶體，指向舊元素的參考和指標都懸置了；迭代器也
一樣，這叫\textbf{迭代器失效}（iterator invalidation）。 下面的程式邊走
訪邊 \texttt{push\_back}：

\begin{code}{cpp}{(iter\_dangling.cpp) 走訪時 push\_back（錯誤示範，main 的內容）}
vector<int> v = {1, 2, 3};
for (auto it = v.begin(); it != v.end(); ++it) {
    if (*it % 2 == 1) {
        v.push_back(*it * 10);   // 走訪時加元素
    }
}
cout << v.size() << '\n';
\end{code}

\begin{code}{text}{用 ASan 編譯後執行（節錄）}
==3486867==ERROR: AddressSanitizer: heap-use-after-free on address 0x502000000014
READ of size 4 at 0x502000000014 thread T0
    #0 0x5f7d71da0682 in main iter_dangling.cpp:8

freed by thread T0 here:
    #6 0x5f7d71da0fd9 in std::vector<int, std::allocator<int> >::push_back(int&&) /usr/include/c++/13/bits/stl_vector.h:1299
    #7 0x5f7d71da075c in main iter_dangling.cpp:9
\end{code}

第一輪 \texttt{*it} 是 $1$，\texttt{push\_back(10)} 超過容量，換了記憶
體；\texttt{it} 還指著舊的那塊，第二輪讀 \texttt{*it}（報告裡的第 $8$ 行）時就讀到
已經釋放的記憶體。 不開 ASan 時，在本書測試的環境下程式照樣跑完，印出
\texttt{11}，看不出哪裡出了錯；換一個環境，結果可能不同，也可能當掉。 修法跟第~3~章一樣：改用下
標迴圈，或者把要加的元素先放進另一個 \texttt{vector}。

刪除也會讓迭代器失效：\texttt{v.erase(it)} 之後，\texttt{it} 和它後面的
迭代器都不能再用。 要邊走訪邊刪除，用 \texttt{erase} 的回傳值，它指向被
刪掉的元素的下一個：

\begin{code}{cpp}{邊走訪邊刪除偶數}
vector<int> v = {4, 1, 6, 6, 3, 8};
for (auto it = v.begin(); it != v.end();) {
    if (*it % 2 == 0) {
        it = v.erase(it);      // erase 回傳下一個元素的迭代器
    } else {
        ++it;
    }
}
// v 剩下 1 3
\end{code}

\texttt{for} 的第三個部分空著，因為刪掉元素時 \texttt{it} 已經是下一個，
不能再 \texttt{++}。 \texttt{set} 和 \texttt{map} 的 \texttt{erase(it)}
也回傳下一個位置，同樣的迴圈可以照用。 \texttt{set}、\texttt{map} 加入元
素時，其他元素的迭代器不會失效；刪除時只有指向被刪元素的那一個失效。

\subsection{常見錯誤}

\begin{itemize}
\item \textbf{讀 \texttt{*v.end()}}：\texttt{end()} 不指著任何元素，跟
C 的 \texttt{a[n]} 一樣是越界，屬於未定義行為。 最後一個元素是
\texttt{*(v.end() - 1)} 或 \texttt{v.back()}，而且 \texttt{v} 不能是空的。
空容器的 \texttt{*v.begin()} 也是同樣的錯誤。
\item \textbf{對 \texttt{set} 的迭代器加整數}：\texttt{s.begin() + 1} 編譯
失敗，錯誤訊息是 \texttt{no match for 'operator+'}。 往前、往後一格用
\texttt{++it}、\texttt{--it}；不想改動 \texttt{it} 時用 \texttt{next(it)}、
\texttt{prev(it)}。 \texttt{next(it, k)} 往後 $k$ 格，但它是一格一格走的，
要花 $O(k)$。
\item \textbf{對 \texttt{list} 或 \texttt{set} 用 \texttt{sort}}：
\texttt{sort(a.begin(), a.end())} 的錯誤指向 \texttt{stl\_algo.h}，從上往
下找到這兩行：
\begin{code}{text}{編譯器的錯誤訊息（節錄）}
required from 'void std::sort(_RAIter, _RAIter) [with _RAIter = _List_iterator<int>]'
error: no match for 'operator-' (operand types are 'std::_List_iterator<int>' and 'std::_List_iterator<int>')
\end{code}
\texttt{\_RAIter} 是 random access iterator 的縮寫，表示 \texttt{sort}
要隨機存取迭代器；錯誤的原因是兩個 \texttt{list} 迭代器不能相減。
\texttt{set} 的訊息相同，只是型別換成 \texttt{\_Rb\_tree\_const\_iterator}。
\item \textbf{兩個迭代器來自不同的容器}：\texttt{sort(a.begin(), b.end())}
編譯得過，執行時是未定義行為。 \texttt{sort(get().begin(), get().end())}
也是：\texttt{get()} 回傳複本的話，兩次呼叫是兩個不同的暫時物件。
\end{itemize}

\section{\texttt{sort} 與比較函式}
\label{sec:sort}

\texttt{sort(first, last)} 把範圍內的元素用 \texttt{<} 由小到大排好，時
間是 $O(n \log n)$。 要用別的順序，就多傳第三個引數：一個\textbf{比較函
式}，收兩個元素 \texttt{a}、\texttt{b}，當 \texttt{a} 應該排在 \texttt{b}
前面時回傳 \texttt{true}。 最簡單的是第~8~章用過的 \texttt{greater}：
\texttt{sort(v.begin(), v.end(), greater<int>())} 用 \texttt{>} 比較，排
成由大到小，後面的 \texttt{()} 是建立一個 \texttt{greater<int>} 物件傳進
去。 比較自訂的結構，就自己寫一個函式：

\begin{example}[label=ex:sort-cmp]{cpp}{(sort\_cmp.cpp) 依分數和名字排序}
#include <algorithm>
#include <iostream>
#include <string>
#include <vector>
using namespace std;

struct Student {
    string name;
    int score;
};

bool byScoreDesc(const Student &a, const Student &b) {
    if (a.score != b.score) {
        return a.score > b.score;     // 分數高的在前
    }
    return a.name < b.name;           // 同分照名字
}

int main() {
    vector<Student> s = {
        {"Linus", 88}, {"Ada", 95}, {"Grace", 88}, {"Alan", 70}};
    sort(s.begin(), s.end(), byScoreDesc);
    for (const auto &st : s) {
        cout << st.name << ' ' << st.score << '\n';
    }
    return 0;
}
\end{example}

\iosample{}{
Ada 95 \\
Grace 88 \\
Linus 88 \\
Alan 70 \\
}

\begin{itemize}
\item \texttt{byScoreDesc} 是一般的函式，傳的時候只寫名字，不加括號。
\texttt{sort} 排序時會反覆呼叫它，問「\texttt{a} 要不要排在 \texttt{b}
前面」。
\item 比較好幾個欄位的寫法是固定的：第一個欄位不同就比第一個欄位，相同才往
下比。 第~6~章替型別寫 \texttt{operator<} 只能有一種順序，比較函式則讓同
一個型別可以有好幾種排法。
\item C 的 \texttt{qsort} 的比較函式回傳負數、$0$、正數；\texttt{sort}
的只回傳 \texttt{bool}，意思是「\texttt{a} 是否嚴格排在 \texttt{b} 前
面」。 這個差別是本節常見錯誤的來源。
\end{itemize}

\subsection{\texttt{stable\_sort}}

\texttt{sort} 不保證相等的元素維持原來的先後。 上面 \texttt{Grace} 和
\texttt{Linus} 同分，是靠比較函式比了名字才有固定的順序；只比分數的話，
誰在前面沒有規定。 \texttt{stable\_sort} 用法跟 \texttt{sort} 一樣，但
保證相等的元素照原來的順序。 所以下面兩步的結果跟 \texttt{byScoreDesc}
相同：先依名字排，再依分數\textbf{穩定}排序，同分的人仍然照名字。

\begin{code}{cpp}{先排次要欄位，再穩定排序主要欄位}
sort(s.begin(), s.end(), [](const Student &a, const Student &b) {
    return a.name < b.name;
});
stable_sort(s.begin(), s.end(),
            [](const Student &a, const Student &b) {
                return a.score > b.score;
            });
// Ada 95、Grace 88、Linus 88、Alan 70
\end{code}

比較函式直接寫在呼叫的地方，叫 lambda，下一節介紹。 解題時多半一次比完所
有欄位；原本的順序有意義（例如報名先後）、只想依某個欄位重排時，才需要
\texttt{stable\_sort}。

\subsection{常見錯誤}

\begin{itemize}
\item \textbf{用 \texttt{<=} 當比較函式}：比較函式必須是嚴格的「小於」，
第~6~章列過條件：\texttt{cmp(a, a)} 一定是假，\texttt{cmp(a, b)} 和
\texttt{cmp(b, c)} 為真時 \texttt{cmp(a, c)} 也要為真。 標準裡叫\textbf{嚴
格弱序}（strict weak ordering）。 \texttt{a <= b} 對相等的元素回傳真，不符合。
下面的程式排序 $20$ 個 $5$：
\begin{code}{cpp}{用 <= 排序（錯誤示範）}
vector<int> v(20, 5);            // 20 個 5
sort(v.begin(), v.end(), [](int a, int b) { return a <= b; });
\end{code}
在本書測試的環境下，不開檢查時程式照常結束，看不出問題。 用 ASan 跑，
\texttt{sort} 讀到了 \texttt{vector} 的範圍外面：
\begin{code}{text}{用 ASan 編譯後執行（節錄）}
==3521734==ERROR: AddressSanitizer: heap-buffer-overflow on address 0x5070000000e0
    #1 0x5e6119314035 in __unguarded_partition<...> /usr/include/c++/13/bits/stl_algo.h:1877
\end{code}
\texttt{sort} 為了快，分割時假設「一定會遇到一個不小於基準的元素而停下
來」，所以不檢查邊界；比較函式說每個元素都「小於等於」基準，它就一路走出
去了。 編譯時加 \texttt{-D\_GLIBCXX\_DEBUG}，標準函式庫會檢查這個條件：\texttt{Error: comparison doesn't meet irreflexive
requirements, assert(!(a < a)).} 由大到小要寫 \texttt{>}，不是
\texttt{>=}。
\item \textbf{照 \texttt{qsort} 的習慣寫 \texttt{return a - b;}}：回傳值
轉成 \texttt{bool}，只要 \texttt{a} 和 \texttt{b} 不相等就是 \texttt{true}，
等於說「\texttt{a} 在 \texttt{b} 前面，\texttt{b} 也在 \texttt{a} 前面」。
編譯沒有警告；在本書測試的環境下，\texttt{\{3, 1, 2\}} 排完是
\texttt{2 1 3}。 要寫 \texttt{return a < b;}。
\end{itemize}

\section{lambda}
\label{sec:lambda}

\texttt{byScoreDesc} 只在一個地方用到，卻要寫在 \texttt{main} 外面、另外
取名字，讀程式的人得跳過去看。 \textbf{lambda}（lambda expression）讓你
在需要函式的地方直接寫一個沒有名字的函式：

\begin{example}[label=ex:lambda]{cpp}{(lambda.cpp) lambda 與捕獲}
#include <algorithm>
#include <cstdlib>
#include <iostream>
#include <vector>
using namespace std;

int main() {
    auto add = [](int a, int b) { return a + b; };
    cout << add(2, 3) << '\n';

    int limit = 10;
    auto byValue = [limit](int x) { return x > limit; };
    auto byRef = [&limit](int x) { return x > limit; };
    limit = 100;
    cout << byValue(50) << ' ' << byRef(50) << '\n';

    int sum = 0;
    auto addTo = [&sum](int x) { sum += x; };
    addTo(3);
    addTo(4);
    cout << sum << '\n';

    int target = 10;
    vector<int> v = {3, 14, 8, 20, 11};
    sort(v.begin(), v.end(), [target](int a, int b) {
        return abs(a - target) < abs(b - target);
    });
    for (int x : v) {
        cout << x << ' ';
    }
    cout << '\n';
    return 0;
}
\end{example}

\iosample{}{
5 \\
1 0 \\
7 \\
11 8 14 3 20 \\
}

一個 lambda 由三個部分組成：方括號 \texttt{[...]} 是\textbf{捕獲列表}
（capture list），列出函式本體要用到的外面的區域變數，\texttt{add} 不需要，
寫空的 \texttt{[]}；接著是參數 \texttt{(int a, int b)}，跟一般函式一樣；
最後是函式本體，回傳型別由 \texttt{return} 的運算式推斷，不用寫出來。
lambda 是一個運算式，結果是一個可以呼叫的物件。 它的型別沒有名字，所以存
成變數時寫 \texttt{auto}；更常見的用法是像 \texttt{sort} 那行一樣，直
接當成引數傳進去。

\subsection{傳值捕獲與傳參考捕獲}

lambda 本體裡不能直接用外面的區域變數，要先捕獲。 捕獲的方式跟傳參數一樣
有兩種：

\begin{itemize}
\item \texttt{[limit]}：\textbf{傳值捕獲}。 建立 lambda 的那一刻，把
\texttt{limit} 的值複製一份存在 lambda 裡。 之後 \texttt{limit} 改成
$100$，\texttt{byValue} 裡的仍然是 $10$，所以 \texttt{byValue(50)} 是真。
複製進來的值是唯讀的，只能讀不能改。
\item \texttt{[\&limit]}：\textbf{傳參考捕獲}。 lambda 裡的
\texttt{limit} 是外面那個變數的參考，看得到之後的修改，所以
\texttt{byRef(50)} 用的是 $100$，結果是假。 要在 lambda 裡修改外面的變數，
例如 \texttt{addTo} 累加 \texttt{sum}，就要傳參考捕獲。
\end{itemize}

選法跟第~2~章選參數一樣：小型別只讀就傳值，要修改或型別很大（例如一個
\texttt{vector}）就傳參考。 好幾個變數用逗號隔開：\texttt{[target,
\&cnt]}。 簡寫 \texttt{[=]}、\texttt{[\&]} 表示用到的全部傳值、全部傳
參考，但看不出捕獲了什麼，本書不用。

範例最後的 \texttt{sort} 依「離 \texttt{target} 多近」排序。
\texttt{sort} 只傳兩個元素給比較函式，\texttt{target} 沒有地方放；一般函
式只能把它改成全域變數，lambda 捕獲進來就好。 這是 lambda 比一般函式好用
的主要原因。

\subsection{常見錯誤}

\begin{itemize}
\item \textbf{沒有捕獲就使用}：\texttt{[](int x) \{ return x > limit; \}}
報 \texttt{error: 'limit' is not captured}。 把 \texttt{limit} 寫進方括號。
\item \textbf{修改傳值捕獲的變數}：\texttt{[sum](int x) \{ sum += x; \}}
報 \texttt{error: assignment of read-only variable 'sum'}。 想改的是外面的
\texttt{sum}，就要寫 \texttt{[\&sum]}。
\item \textbf{傳參考捕獲的變數比 lambda 先消失}：下面的函式回傳一個
lambda，它用參考捕獲了參數 \texttt{limit}：
\begin{code}{cpp}{lambda 捕獲了將要消失的參數（錯誤示範）}
auto makeAbove(int limit) {
    return [&limit](int x) { return x > limit; };
}

int main() {
    auto above10 = makeAbove(10);
    cout << above10(50) << '\n';
    return 0;
}
\end{code}
\texttt{makeAbove} 一回傳，參數 \texttt{limit} 的生命週期就結束了，
\texttt{above10} 裡的參考跟著懸置，跟第~4~章「回傳區域變數的參考」是同
一件事。 \texttt{-Wall -Wextra} 不會警告；在本書測試的環境下，程式印出
\texttt{0}，不是預期的 \texttt{1}。 ASan 報 \texttt{stack-use-after-return}：
\begin{code}{text}{用 ASan 編譯後執行（節錄）}
==3527040==ERROR: AddressSanitizer: stack-use-after-return on address 0x7febbeb00060
READ of size 4 at 0x7febbeb00060 thread T0
    #0 0x59ac85a0d2f7 in operator() lambda_dangling.cpp:5
    #1 0x59ac85a0d4a5 in main lambda_dangling.cpp:10
\end{code}
改成傳值捕獲 \texttt{[limit]}，值就跟著 lambda 走。 規則是：lambda 只在
當下用掉（例如傳給 \texttt{sort}），傳參考捕獲沒問題；lambda 會被存起來、
回傳出去，就用傳值捕獲。
\end{itemize}

\section{常用的演算法}
\label{sec:algorithms}

\texttt{<algorithm>} 裡有上百個函式，解題常用的不多。 下面的範例一次示範
大部分，每一個都收一個範圍 $[\text{first}, \text{last})$：

\begin{example}[label=ex:algos]{cpp}{(algos.cpp) 搜尋、計數、最大最小、反轉、加總、去重複}
#include <algorithm>
#include <iostream>
#include <numeric>
#include <vector>
using namespace std;

void print(const vector<int> &v) {
    for (int x : v) {
        cout << x << ' ';
    }
    cout << '\n';
}

int main() {
    vector<int> v = {7, 3, 9, 3, 12, 5, 3};

    auto it = find(v.begin(), v.end(), 9);
    if (it != v.end()) {
        cout << "9 在下標 " << it - v.begin() << '\n';
    }
    auto ev = find_if(v.begin(), v.end(),
                      [](int x) { return x % 2 == 0; });
    cout << "第一個偶數 " << *ev << '\n';

    cout << count(v.begin(), v.end(), 3) << ' ';
    cout << count_if(v.begin(), v.end(),
                     [](int x) { return x > 4; }) << '\n';

    auto lo = min_element(v.begin(), v.end());
    auto hi = max_element(v.begin(), v.end());
    cout << *lo << " 在 " << lo - v.begin() << "，"
         << *hi << " 在 " << hi - v.begin() << '\n';

    cout << accumulate(v.begin(), v.end(), 0) << '\n';

    reverse(v.begin(), v.end());
    print(v);

    sort(v.begin(), v.end());
    v.erase(unique(v.begin(), v.end()), v.end());
    print(v);
    return 0;
}
\end{example}

\iosample{}{
9 在下標 2 \\
第一個偶數 12 \\
3 4 \\
3 在 1，12 在 4 \\
42 \\
3 5 12 3 9 3 7 \\
3 5 7 9 12 \\
}

\begin{center}
\begin{tabular}{ll}
\toprule
呼叫 & 結果 \\
\midrule
\texttt{find(first, last, x)} & 第一個等於 \texttt{x} 的位置，沒有就是 \texttt{last} \\
\texttt{find\_if(first, last, pred)} & 第一個讓 \texttt{pred} 為真的位置，沒有就是 \texttt{last} \\
\texttt{count(first, last, x)} & 等於 \texttt{x} 的個數 \\
\texttt{count\_if(first, last, pred)} & 讓 \texttt{pred} 為真的個數 \\
\texttt{min\_element}、\texttt{max\_element} & 最小、最大元素的位置（有好幾個時取第一個） \\
\texttt{reverse(first, last)} & 把範圍反轉，沒有回傳值 \\
\texttt{accumulate(first, last, init)} & \texttt{init} 加上所有元素（在 \texttt{<numeric>} 裡） \\
\bottomrule
\end{tabular}
\end{center}

回傳位置的演算法都回傳迭代器，找不到時是 \texttt{last}，跟第~8~章
\texttt{set} 的 \texttt{find} 回傳 \texttt{end()} 同一個慣例。 要下標就
減掉 \texttt{v.begin()}，要值就用 \texttt{*}。 範例裡 \texttt{ev} 沒有跟
\texttt{v.end()} 比就取值，是因為確定有偶數。 \texttt{pred} 是回傳
\texttt{bool} 的函式，叫\textbf{述詞}（predicate），通常寫成 lambda。 這
些演算法都只是一個迴圈，時間是 $O(n)$。

\subsection{\texttt{unique} 與 \texttt{erase}}

範例最後兩行是去掉重複元素的固定寫法。 \texttt{unique} 把\textbf{相鄰}的
重複元素只留一個，往前搬，回傳「不重複部分的結尾」。 它不會改變
\texttt{vector} 的大小：演算法只拿到兩個迭代器，碰得到元素，碰不到容器本
身，沒辦法刪除元素。 所以要再用 \texttt{erase} 把後面那段刪掉：

\begin{code}{cpp}{排序、unique、erase 去掉重複}
sort(v.begin(), v.end());                       // 相同的排在一起
v.erase(unique(v.begin(), v.end()), v.end());   // 刪掉後面那段
\end{code}

\texttt{erase} 收兩個迭代器時，刪掉這個半開區間裡的元素。 做完之後
\texttt{v} 排好序又不重複，可以直接用下一節的二分搜尋。

\subsection{常見錯誤}

\begin{itemize}
\item \textbf{\texttt{unique} 沒有接 \texttt{erase}}：
\begin{code}{cpp}{只呼叫 unique（錯誤示範）}
vector<int> v = {1, 1, 2, 3, 3, 3, 4};
auto newEnd = unique(v.begin(), v.end());
cout << newEnd - v.begin() << ' ' << v.size() << '\n';   // 4 7
\end{code}
前 $4$ 個是 \texttt{1 2 3 4}，但 \texttt{v.size()} 還是 $7$。 後面 $3$ 個
元素的值沒有規定，在本書測試的環境下整個 \texttt{v} 印出來是
\texttt{1 2 3 4 3 3 4}。 之後如果拿 \texttt{v.size()} 或範圍 \texttt{for}
去用，就會多出這幾個舊值。
\item \textbf{\texttt{unique} 之前沒有排序}：\texttt{\{3, 1, 3, 1\}} 沒有
相鄰的重複，\texttt{unique} 什麼都不做，結果還是 \texttt{3 1 3 1}。
\item \textbf{\texttt{accumulate} 的初始值型別}：回傳值和加總過程都用
\texttt{init} 的型別。 \texttt{vector<double>} 寫
\texttt{accumulate(d.begin(), d.end(), 0)}，\texttt{0} 是 \texttt{int}，
每一步都截掉小數：\texttt{\{0.5, 1.5, 2.25\}} 加出來是 \texttt{3}，不是
\texttt{4.25}。 要寫 \texttt{0.0}。 加總可能超過 \texttt{int} 時寫
\texttt{0LL}：兩個 \texttt{2000000000} 用 \texttt{0} 加是有號整數溢位
（未定義行為），本書測試的環境下印出 \texttt{-294967296}；加
\texttt{-fsanitize=undefined} 編譯，執行時會報 \texttt{runtime error: signed
integer overflow: 2000000000 + 2000000000 cannot be represented in type
'int'}。
\item \textbf{對空範圍取 \texttt{*max\_element}}：空的 \texttt{vector}，
\texttt{max\_element} 回傳 \texttt{end()}，再用 \texttt{*} 就是讀
\texttt{*v.end()}。 在本書測試的環境下用 ASan 跑，報
\texttt{SEGV on unknown address 0x000000000000}。 先確定範圍不是空的。
\end{itemize}

\section{二分搜尋}
\label{sec:binary-search}

在排好序的資料裡找東西，可以每次看中間那個，丟掉不可能的一半，$n$ 個元素
只要看大約 $\log_2 n$ 個，這是\textbf{二分搜尋}（binary search）。 自己寫
時，\texttt{lo}、\texttt{hi} 是否包含、迴圈條件是 \texttt{<} 還是
\texttt{<=}，很容易差一。 標準函式庫有三個現成的版本，用在\textbf{已經由小
到大排好}的範圍上：

\begin{itemize}
\item \texttt{binary\_search(first, last, x)}：有沒有 \texttt{x}，回傳
\texttt{bool}。
\item \texttt{lower\_bound(first, last, x)}：第一個 $\ge x$ 的位置。
\item \texttt{upper\_bound(first, last, x)}：第一個 $> x$ 的位置。
\end{itemize}

兩個 \texttt{bound} 都回傳迭代器，找不到時是 \texttt{last}。 三個都是
$O(\log n)$。

\begin{example}[label=ex:bounds]{cpp}{(bounds.cpp) binary\_search、lower\_bound、upper\_bound}
#include <algorithm>
#include <iostream>
#include <set>
#include <vector>
using namespace std;

int main() {
    vector<int> v = {1, 3, 3, 3, 5, 8};
    cout << binary_search(v.begin(), v.end(), 3) << ' '
         << binary_search(v.begin(), v.end(), 4) << '\n';

    auto lo = lower_bound(v.begin(), v.end(), 3);
    auto hi = upper_bound(v.begin(), v.end(), 3);
    cout << lo - v.begin() << ' ' << hi - v.begin() << ' ';
    cout << hi - lo << '\n';

    cout << lower_bound(v.begin(), v.end(), 4) - v.begin() << ' '
         << upper_bound(v.begin(), v.end(), 4) - v.begin() << '\n';
    if (lower_bound(v.begin(), v.end(), 9) == v.end()) {
        cout << "沒有 >= 9 的元素" << '\n';
    }

    set<int> s = {10, 20, 30};
    cout << *s.lower_bound(15) << ' ' << *s.lower_bound(20) << ' '
         << *s.upper_bound(20) << '\n';
    return 0;
}
\end{example}

\iosample{}{
1 0 \\
1 4 3 \\
4 4 \\
沒有 >= 9 的元素 \\
20 20 30 \\
}

\begin{itemize}
\item \texttt{v} 裡的 $3$ 在下標 $1$ 到 $3$。 \texttt{lower\_bound} 指著
第一個 $3$（下標 $1$），\texttt{upper\_bound} 指著最後一個 $3$ 的下一格
（下標 $4$）。 兩者圍出的 $[\text{lo}, \text{hi})$ 正好是所有的 $3$，
\texttt{hi - lo} 是 $3$ 出現的次數。 又是一個半開區間。
\item 找不存在的 $4$ 時，兩個都指著下標 $4$，也就是 $4$ 應該插入的位置，
\texttt{hi - lo} 是 $0$。 \texttt{lower\_bound} 不只回答有沒有，還告訴你
「比 $x$ 小的有幾個」：就是它減掉 \texttt{v.begin()}。
\item 比所有元素都大的 $9$，\texttt{lower\_bound} 回傳 \texttt{end()}。 要
用 \texttt{*} 之前，先跟 \texttt{end()} 比。
\item \texttt{set} 和 \texttt{map} 有同名的成員函式 \texttt{s.lower\_bound(x)}、
\texttt{s.upper\_bound(x)}，意思相同，也是 $O(\log n)$；\texttt{map} 的
版本以鍵比較。 第~8~章說「要找比 $x$ 大的第一個就用 \texttt{set}」，指的
就是這兩個函式。 要「最後一個 $\le x$」，就先找第一個 $> x$ 再
\texttt{--it}，前提是它不是 \texttt{begin()}。
\end{itemize}

\texttt{vector} 和 \texttt{set} 怎麼選：資料讀進來之後就不再變動，排序一次再
用 \texttt{lower\_bound} 查，最簡單也最快；查詢之間還要插入或刪除，排序過的
\texttt{vector} 每次插入要搬動 $O(n)$ 個元素，就改用 \texttt{set} 的成員函式。

\subsection{例題：最近的商店}

APCS 和解題網站常見的題型是「先給 $n$ 個數字，再問 $q$ 次」。 每次查詢都
從頭找一遍是 $O(nq)$，$n$、$q$ 都是 $2 \times 10^5$ 時要做 $4 \times 10^{10}$
次，太慢；先排序，每次查詢用二分搜尋，總時間是 $O((n + q) \log n)$。

題目：一條街上有 $n$ 家商店，位置是整數 $p_1, \dots, p_n$（沒有排序）。
接著有 $q$ 次查詢，每次給一個位置 $x$，輸出 $x$ 到最近一家商店的距離。

\begin{example}[label=ex:nearest]{cpp}{(nearest.cpp) 用 lower\_bound 找最近的商店}
#include <algorithm>
#include <iostream>
#include <vector>
using namespace std;

int main() {
    ios::sync_with_stdio(false);
    cin.tie(nullptr);
    int n, q;
    cin >> n >> q;
    vector<int> p(n);
    for (int &x : p) {
        cin >> x;
    }
    sort(p.begin(), p.end());

    for (int i = 0; i < q; i++) {
        int x;
        cin >> x;
        auto it = lower_bound(p.begin(), p.end(), x);  // >= x
        int best = -1;
        if (it != p.end()) {
            best = *it - x;              // 右邊最近的
        }
        if (it != p.begin()) {
            int left = x - *(it - 1);    // 左邊最近的
            if (best == -1 || left < best) {
                best = left;
            }
        }
        cout << best << '\n';
    }
    return 0;
}
\end{example}

\iosample{
5 5 \\
30 2 17 9 40 \\
10 1 35 17 50 \\
}{
1 \\
1 \\
5 \\
0 \\
10 \\
}

排好序之後商店是 \texttt{2 9 17 30 40}。 最近的商店一定是「第一個 $\ge x$
的」或「它前面那一個」，所以只要看 \texttt{it} 和 \texttt{it - 1} 兩家。
兩個 \texttt{if} 分別處理兩端：$x = 50$ 比所有商店都大，\texttt{it} 是
\texttt{end()}，只能看左邊；$x = 1$ 比所有商店都小，\texttt{it} 是
\texttt{begin()}，前面沒有商店。 輸入加速是第~1~章的寫法，查詢多的題目都
該加上。

\subsection{常見錯誤}

\begin{itemize}
\item \textbf{範圍沒有排序}：這三個函式不檢查範圍有沒有排好，只照二分搜尋
的步驟走。 下面的 \texttt{v} 裡明明有 $1$、$8$、$5$：
\begin{code}{cpp}{對沒有排序的 vector 二分搜尋（錯誤示範）}
vector<int> v = {8, 1, 5, 3, 9, 2};      // 沒有排序
cout << binary_search(v.begin(), v.end(), 1) << ' '
     << binary_search(v.begin(), v.end(), 8) << ' '
     << lower_bound(v.begin(), v.end(), 5) - v.begin() << '\n';
\end{code}
在本書測試的環境下印出 \texttt{0 0 6}：兩個都說找不到，\texttt{lower\_bound}
回傳了 \texttt{end()}。 結果沒有保證，換一組資料可能剛好對，所以小測資
過了不代表沒問題。 加 \texttt{-D\_GLIBCXX\_DEBUG} 編譯，執行時會報
\texttt{Error: elements in iterator range [first, last) are not partitioned
by the value \_\_val.} 記得先 \texttt{sort}。
\item \textbf{對 \texttt{set} 用 \texttt{std::lower\_bound}}：
\texttt{lower\_bound(s.begin(), s.end(), x)} 編譯得過，結果也對，但
\texttt{set} 的迭代器不能跳到中間，演算法只好一格一格走，變成 $O(n)$。 查
詢多的時候會超時。 用成員函式 \texttt{s.lower\_bound(x)}。
\texttt{std::find}、\texttt{std::count} 用在 \texttt{set} 上也一樣是
$O(n)$，要用成員函式 \texttt{s.find(x)}、\texttt{s.count(x)}。
\item \textbf{沒有檢查 \texttt{end()} 或 \texttt{begin()} 就取值}：
\texttt{*lower\_bound(...)} 在 $x$ 比所有元素都大時是讀
\texttt{*v.end()}；\texttt{*(it - 1)} 在 \texttt{it} 是 \texttt{begin()} 時
讀到範圍前面。 像 \texttt{nearest.cpp} 一樣，兩端都要檢查。
\end{itemize}

\section{\texttt{next\_permutation} 與暴力枚舉}
\label{sec:permutation}

有些題目的 $n$ 很小，最直接的做法是試遍所有排列：$n$ 個人排隊的每一種順序、
$n$ 件工作分給 $n$ 個人的每一種分法。 \texttt{next\_permutation(first, last)}
把範圍改成字典順序的下一個排列，回傳 \texttt{true}；已經是最後一個（由大到
小）時，把範圍改回第一個（由小到大），回傳 \texttt{false}。 從排好序的範圍開
始，用 \texttt{do}-\texttt{while} 呼叫到它回傳 \texttt{false}，就恰好走過
每一個排列一次。

下面的程式先印出 \texttt{1 2 3} 的所有排列，再解一個指派問題：$n$ 個人做
$n$ 件工作，每人做一件，第 $i$ 個人做第 $j$ 件的成本是 $c_{ij}$，求總成本
最小的分法。

\begin{example}[label=ex:perm]{cpp}{(perm.cpp) 列出所有排列，暴力解指派問題}
#include <algorithm>
#include <iostream>
#include <vector>
using namespace std;

int main() {
    vector<int> a = {1, 2, 3};
    do {
        cout << a[0] << a[1] << a[2] << ' ';
    } while (next_permutation(a.begin(), a.end()));
    cout << '\n';

    int n;
    cin >> n;
    vector<vector<int>> cost(n, vector<int>(n));
    for (auto &row : cost) {
        for (int &c : row) {
            cin >> c;
        }
    }

    vector<int> job(n);              // job[i]：第 i 個人做哪一件工作
    for (int i = 0; i < n; i++) {
        job[i] = i;
    }
    int best = -1;
    vector<int> bestJob;
    do {
        int total = 0;
        for (int i = 0; i < n; i++) {
            total += cost[i][job[i]];
        }
        if (best == -1 || total < best) {
            best = total;
            bestJob = job;
        }
    } while (next_permutation(job.begin(), job.end()));

    cout << best << '\n';
    for (int i = 0; i < n; i++) {
        cout << "第 " << i + 1 << " 人做工作 " << bestJob[i] + 1
             << '\n';
    }
    return 0;
}
\end{example}

\iosample{
3 \\
9 2 7 \\
6 4 3 \\
5 8 1 \\
}{
123 132 213 231 312 321 \\
9 \\
第 1 人做工作 2 \\
第 2 人做工作 1 \\
第 3 人做工作 3 \\
}

\begin{itemize}
\item 用 \texttt{do}-\texttt{while} 是因為一開始的 \texttt{1 2 3} 本身就是
第一個排列，要先處理它，再問有沒有下一個。
\item 一個排列 \texttt{job} 代表一種分法：第 $i$ 個人做第 \texttt{job[i]}
件。 \texttt{job} 從 $0, 1, \dots, n - 1$ 開始，已經由小到大。
\item 答案是 $2 + 6 + 1 = 9$。 \texttt{bestJob = job} 複製一份當時的排列，
因為 \texttt{job} 接下來還會繼續變。
\end{itemize}

$n$ 個元素有 $n!$ 個排列，每個排列再花 $O(n)$ 計算，總時間是 $O(n \cdot n!)$。
$8! = 40320$，$10! = 3628800$，到 $n = 10$ 還跑得完；$12!$ 約是 $4.8$ 億，
再乘上 $n$ 就太慢了。 題目的 $n$ 不超過 $8$ 到 $10$，常常就是在暗示可以
暴力枚舉。

\subsection{常見錯誤}

\begin{itemize}
\item \textbf{沒有從排好序的狀態開始}：從 \texttt{\{2, 1, 3\}} 開始，只會
走過 $213$ 之後的排列，到 $321$ 就停了，共 $4$ 個，漏掉 $123$ 和 $132$。
迴圈前先 \texttt{sort}。
\item \textbf{寫成 \texttt{while}}：\texttt{while (next\_permutation(...))
\{ ... \}} 第一次進迴圈前就已經換到下一個排列，少處理了最初的那一個，
\texttt{1 2 3} 只會處理 $5$ 個排列。
\item \textbf{以為有重複元素時也是 $n!$ 個}：\texttt{\{1, 1, 2\}} 只會走過
\texttt{112}、\texttt{121}、\texttt{211} 三個。 \texttt{next\_permutation}
依字典順序找下一個「不同的」排列，重複的不會出現兩次。 這通常正是題目要的，
但不要拿 $n!$ 去算答案的個數。
\end{itemize}

\subsection{自己寫迴圈，還是用演算法}

本章的演算法，每一個都能用幾行迴圈寫出來。 用哪一個，可以照下面的原則：

\begin{itemize}
\item \textbf{排序和二分搜尋一律用標準函式庫。} 自己寫的排序很難做到
$O(n \log n)$ 又沒有錯；二分搜尋的邊界最容易差一，\texttt{lower\_bound}
和 \texttt{upper\_bound} 已經把它定義清楚了。
\item \textbf{名字比迴圈清楚時用演算法。} \texttt{count\_if}、
\texttt{max\_element}、\texttt{reverse} 一看就知道在做什麼，同樣的迴圈要
讀完三四行才知道。
\item \textbf{一次走訪要做好幾件事時，寫迴圈。} 同時要算總和、最大值、
偶數個數，呼叫三個演算法就走了三遍，也不比一個迴圈好讀。 lambda 要寫十幾
行時，也直接寫迴圈。
\item \textbf{邊走訪邊增刪元素時，寫迴圈}，並且照第~\ref{sec:iterators}~節
處理失效：刪除用 \texttt{it = v.erase(it)}，增加用下標或另一個
\texttt{vector}。
\end{itemize}

標準函式庫的演算法通常跟自己寫的迴圈一樣快，所以選擇的標準是哪一種比較好
讀、比較不容易寫錯。

\section{本章練習}
\label{sec:ch9-exercises}

\begin{drills}[label=drills:ch9]
\item 預測下面程式的輸出。
\begin{code}{cpp}{練習：預測輸出}
#include <algorithm>
#include <iostream>
#include <numeric>
#include <vector>
using namespace std;

int main() {
    vector<int> v = {4, 2, 7, 2, 9, 2};
    cout << count(v.begin(), v.end(), 2) << ' '
         << *max_element(v.begin(), v.end()) << ' '
         << find(v.begin(), v.end(), 7) - v.begin() << '\n';

    sort(v.begin(), v.end());
    auto b = v.begin(), e = v.end();
    cout << lower_bound(b, e, 2) - b << ' ';
    cout << upper_bound(b, e, 2) - b << ' ';
    cout << lower_bound(b, e, 5) - b << '\n';

    v.erase(unique(v.begin(), v.end()), v.end());
    cout << v.size() << ' ';
    cout << accumulate(v.begin(), v.end(), 0) << '\n';

    int k = 3;
    auto f = [k](int x) { return x * k; };
    auto g = [&k](int x) { return x * k; };
    k = 10;
    cout << f(2) << ' ' << g(2) << '\n';
    return 0;
}
\end{code}
\item 下面的程式讀入 $n$ 個整數，印出有幾個不同的數字，再對之後讀到的每個
數字回答「有」或「沒有」。 輸入
\begin{code}{text}{練習 2 的輸入}
6
5 3 5 1 3 9
9 4 1
\end{code}
時它輸出 \texttt{不同的數字有 6 個}、\texttt{9 沒有}、\texttt{4 沒有}、
\texttt{1 沒有}。 找出程式裡的三個錯誤，並改好它，讓輸出變成
\texttt{不同的數字有 4 個}、\texttt{9 有}、\texttt{4 沒有}、\texttt{1 有}。
\begin{code}{cpp}{練習：去重複與查詢}
#include <algorithm>
#include <functional>
#include <iostream>
#include <vector>
using namespace std;

int main() {
    int n;
    cin >> n;
    vector<int> v(n);
    for (int &x : v) {
        cin >> x;
    }
    unique(v.begin(), v.end());                  // 去掉重複
    sort(v.begin(), v.end(), greater<int>());    // 由大到小
    cout << "不同的數字有 " << v.size() << " 個" << '\n';

    int x;
    while (cin >> x) {
        bool found = binary_search(v.begin(), v.end(), x);
        cout << x << (found ? " 有" : " 沒有") << '\n';
    }
    return 0;
}
\end{code}
\item （區間計數）讀入 $n$、$q$ 和 $n$ 個學生的分數，接著有 $q$ 次查詢，每
次給 $l$、$r$，輸出分數在 $l$ 到 $r$ 之間（包含兩端）的學生有幾個。
$n, q \le 2 \times 10^5$，每次查詢要在 $O(\log n)$ 內完成。 輸入
\begin{code}{text}{練習 3 的輸入}
8 4
72 45 90 60 88 60 100 33
60 89
0 59
91 99
100 100
\end{code}
應該輸出 \texttt{4}、\texttt{2}、\texttt{0}、\texttt{1}。
\item 讀入一串整數直到輸入結束，依「出現次數多的在前；次數相同時，數字小的
在前」印出每個不同的數字和它的出現次數，一行一個。 輸入
\texttt{4 6 2 6 4 9 6 2 7} 時應輸出 \texttt{6 3}、\texttt{2 2}、
\texttt{4 2}、\texttt{7 1}、\texttt{9 1}。 先用第~8~章的 \texttt{map} 計
數，再搬進 \texttt{vector} 用 lambda 排序。
\item （旅行推銷員）有 $n$ 個城市（$2 \le n \le 8$），讀入 $n \times n$
的距離矩陣，第 $i$ 列第 $j$ 個數字是城市 $i$ 到城市 $j$ 的距離。 從城市
$1$ 出發，每個城市恰好經過一次，最後回到城市 $1$，求最短的總距離。 用
\texttt{next\_permutation} 枚舉。 輸入
\begin{code}{text}{練習 5 的輸入}
4
0 10 15 20
10 0 35 25
15 35 0 30
20 25 30 0
\end{code}
應該輸出 \texttt{80}。
\end{drills}

\addcontentsline{toc}{section}{練習解答}
\section*{練習解答}

\begin{enumerate}
\item 輸出四行：\texttt{3 9 2}、\texttt{0 3 4}、\texttt{4 22}、
\texttt{6 20}。 排序後是 \texttt{2 2 2 4 7 9}：第一個 $\ge 2$ 在下標 $0$，
第一個 $> 2$ 在下標 $3$，$5$ 不存在，第一個 $\ge 5$ 是下標 $4$ 的 $7$。
去掉重複剩 \texttt{2 4 7 9}，和是 $22$。 \texttt{f} 在建立時複製了
\texttt{k} 的值 $3$，\texttt{g} 綁著 \texttt{k} 本身，呼叫時已經是 $10$。
\item 三個錯誤：(1) \texttt{unique} 的回傳值沒有拿去 \texttt{erase}，大小
還是 $6$；(2) \texttt{unique} 只去掉\textbf{相鄰}的重複，必須在排序之後；
(3) \texttt{binary\_search} 預設假設範圍由小到大，\texttt{v} 卻是由大到小，
前提不成立，結果沒有保證，在本書測試的環境下 $9$ 和 $1$ 都找不到。 改法是
排成由小到大，再 \texttt{unique} 加 \texttt{erase}：
\begin{code}{cpp}{練習 2：改好的去重複}
sort(v.begin(), v.end());
v.erase(unique(v.begin(), v.end()), v.end());
\end{code}
\texttt{binary\_search} 那一行不用改。 也可以保留由大到小，但
\texttt{binary\_search} 要傳同一個比較函式：
\texttt{binary\_search(v.begin(), v.end(), x, greater<int>())}。
\item 一種寫法：
\begin{code}{cpp}{練習 3 參考解答（讀入分數 a 之後）}
sort(a.begin(), a.end());
for (int i = 0; i < q; i++) {
    int l, r;
    cin >> l >> r;
    auto lo = lower_bound(a.begin(), a.end(), l);  // >= l
    auto hi = upper_bound(a.begin(), a.end(), r);  // > r
    cout << hi - lo << '\n';
}
\end{code}
排序後是 \texttt{33 45 60 60 72 88 90 100}。 分數在 $[l, r]$ 的學生，就是
從第一個 $\ge l$ 到第一個 $> r$ 之前的那一段，半開區間
$[\text{lo}, \text{hi})$ 的長度就是答案。 沒有人在範圍內時，\texttt{lo}
和 \texttt{hi} 指著同一格，相減是 $0$，不用另外處理。 排序 $O(n \log n)$，
每次查詢 $O(\log n)$。
\item 一種寫法：
\begin{code}{cpp}{練習 4 參考解答（main 的內容）}
map<int, int> cnt;
int x;
while (cin >> x) {
    cnt[x]++;
}
// v 的元素是 (數字, 次數)
vector<pair<int, int>> v(cnt.begin(), cnt.end());
sort(v.begin(), v.end(),
     [](const pair<int, int> &a, const pair<int, int> &b) {
         if (a.second != b.second) {
             return a.second > b.second;   // 次數多的在前
         }
         return a.first < b.first;         // 數字小的在前
     });
for (const auto &[num, c] : v) {
    cout << num << ' ' << c << '\n';
}
\end{code}
\texttt{map} 只能依鍵排序，所以把內容搬進 \texttt{vector}。
\texttt{vector<pair<int, int>> v(cnt.begin(), cnt.end())} 用第~\ref{sec:iterators}~節
「用一個範圍建立容器」的寫法，\texttt{map} 的元素是
\texttt{pair<const int, int>}，複製時會轉成 \texttt{pair<int, int>}。 比
較函式先比次數（由大到小，用 \texttt{>}），相同再比數字（用 \texttt{<}），
兩邊都是嚴格的比較。
\item 一種寫法：
\begin{code}{cpp}{練習 5 參考解答（讀完距離矩陣 d 之後）}
vector<int> order;              // 城市 1 以外的城市
for (int i = 1; i < n; i++) {
    order.push_back(i);
}
int best = -1;
do {
    int total = 0, cur = 0;         // 從城市 0 出發
    for (int c : order) {
        total += d[cur][c];
        cur = c;
    }
    total += d[cur][0];             // 回到出發點
    if (best == -1 || total < best) {
        best = total;
    }
} while (next_permutation(order.begin(), order.end()));
cout << best << '\n';
\end{code}
起點固定是城市 $1$（程式裡的下標 $0$），只要枚舉其餘 $n - 1$ 個城市的順序，
共 $(n - 1)!$ 種，$n = 8$ 時是 $5040$ 種。 \texttt{order} 一開始就由小到
大，不用再排序。 範例的最短路線
是 $1 \to 2 \to 4 \to 3 \to 1$，距離 $10 + 25 + 30 + 15 = 80$。
\end{enumerate}

\chapter{銜接 APCS 與資料結構}
\label{ch:bridge}

前九章把 C++ 的語法和標準函式庫講完了。 這一章幾乎不再介紹新語法，而是
把學過的東西用在接下來會遇到的題目上：APCS 這類解題考試，以及大學的資料
結構與演算法課。

前半章講解題：一支解題程式的基本寫法、怎麼從題目的範圍估計演算法夠不夠快，
以及最常見的兩種解法，遞迴搜尋和動態規劃。 後半章用指標和類別自己寫一個鏈
結串列，這是資料結構課的第一個作業常見的形式。 最後整理讀完這本書之後可
以往哪裡走。

\section{解題程式的基本功}
\label{sec:contest-basics}

解題程式從標準輸入讀資料、把答案印到標準輸出，評分系統拿你的輸出跟正確答
案逐字比對。 本書的範例已經是這個形式，固定的部分整理起來是這樣：

\begin{code}{cpp}{解題程式的骨架}
#include <algorithm>
#include <iostream>
#include <vector>
using namespace std;

int main() {
    ios::sync_with_stdio(false);
    cin.tie(nullptr);
    // 讀入、計算、輸出
    return 0;
}
\end{code}

標頭檔用到什麼就 include 什麼。 網路上的解題程式常寫
\texttt{\#include <bits/stdc++.h>}，一行包含所有標準標頭檔。 它是
\texttt{g++} 特有的檔案，不是標準的一部分，換成別的編譯器可能找不到，編譯
也比較慢。 本書一律列出實際用到的標頭檔，讀程式的人也能從開頭看出這支程
式用了哪些東西。

輸入加速的兩行是第~1~章的內容，資料量大的題目一定要加，加了之後就不要再
混用 \texttt{scanf}／\texttt{printf}。

\subsection{從資料範圍估計時間}

題目會寫資料範圍，例如 $n \le 2 \times 10^5$。 這不只是告訴你陣列要開多
大，也在告訴你需要多快的演算法。 粗略的估計是：評分系統一秒大約能做
$10^8$ 次簡單運算（加減、比較、存取陣列）。 把 $n$ 的上限代進演算法的時間
複雜度，算出來在 $10^8$ 左右或更少，通常就來得及：

\begin{center}
\begin{tabular}{lll}
\toprule
$n$ 的上限 & 可以接受的複雜度 & 例子 \\
\midrule
$10$ 左右 & $O(n!)$、$O(n \cdot n!)$ & \texttt{next\_permutation} 枚舉 \\
$20$ 左右 & $O(2^n)$ & 枚舉所有子集合 \\
$500$ & $O(n^3)$ & 三層迴圈 \\
$5000$ & $O(n^2)$ & 雙層迴圈、$O(n^2)$ 的動態規劃 \\
$2 \times 10^5$ & $O(n \log n)$ & 排序、二分搜尋、\texttt{set}、\texttt{map}、\texttt{priority\_queue} \\
$10^7$ & $O(n)$ & 一個迴圈、\texttt{unordered\_map} \\
\bottomrule
\end{tabular}
\end{center}

這張表只是大概：常數大的演算法要打折，時間限制不是一秒也要換算。 但它足
以在動手之前排除不可能的做法。 下面是一個例子。

題目：給 $n$ 個整數和一個整數 $k$，問有幾對 $i < j$ 滿足
$a_i + a_j = k$。 $n \le 2 \times 10^5$，$|a_i|, |k| \le 10^9$。

最直接的做法是雙層迴圈檢查每一對，$O(n^2)$。 $n = 2 \times 10^5$ 時大約是
$2 \times 10^{10}$ 次比較，照上表一定來不及；在本書測試的環境下，加
\texttt{-O2} 編譯，跑一組 $n = 2 \times 10^5$ 的資料要 $6.7$ 秒左右。 換個
想法：由左往右讀，讀到 $a_j$ 時，它能和前面每一個等於 $k - a_j$ 的數字配
成一對。 只要用第~8~章的 \texttt{unordered\_map} 記住前面每個數字出現幾次，
每個數字就只要 $O(1)$：

\begin{example}[label=ex:pairs]{cpp}{(pairs.cpp) 數有幾對的和等於 k}
#include <iostream>
#include <unordered_map>
using namespace std;

int main() {
    ios::sync_with_stdio(false);
    cin.tie(nullptr);
    int n;
    long long k;
    cin >> n >> k;
    unordered_map<long long, int> seen;   // 前面每個數字出現幾次
    long long ans = 0;
    for (int i = 0; i < n; i++) {
        long long x;
        cin >> x;
        auto it = seen.find(k - x);
        if (it != seen.end()) {
            ans += it->second;            // 跟前面的 k - x 配對
        }
        seen[x]++;
    }
    cout << ans << '\n';
    return 0;
}
\end{example}

\iosample{
6 7 \\
3 4 2 5 4 3 \\
}{
5 \\
}

和是 $7$ 的五對是兩個 $3$ 各配兩個 $4$，再加上 $2$ 和 $5$。 同一組
$n = 2 \times 10^5$ 的資料，這個版本在本書測試的環境下不到 $0.1$ 秒。 查
詢用 \texttt{find} 而不是 \texttt{seen[k - x]}，避免第~8~章說的「查詢時新
增鍵」。 \texttt{k - x} 最大可能到 $2 \times 10^9$，超過 \texttt{int}，所以
數字都用 \texttt{long long}。

\subsection{用檔案測試}

範例輸入通常不只一組，每次都手打很累，也容易打錯。 把輸入存成檔案，執行時
用 \texttt{<} 讓程式從檔案讀，用 \texttt{>} 把輸出寫進檔案，再跟正確答案比：

\begin{code}{bash}{用檔案測試}
g++ -std=c++17 -Wall -Wextra pairs.cpp -o pairs
./pairs < in1.txt > out1.txt
diff out1.txt ans1.txt
\end{code}

這是命令列的功能，程式本身照樣用 \texttt{cin}、\texttt{cout}，不用改。
\texttt{diff} 沒有印出任何東西，就表示兩個檔案完全一樣；Windows 的命令提
示字元用 \texttt{fc out1.txt ans1.txt}。 自己再補幾組邊界的測資：$n = 1$、
所有數字都一樣、數字都是最大值。 很多錯誤只在這些情況出現。

\subsection{常見錯誤}

\begin{itemize}
\item \textbf{答案超過 \texttt{int}}：把 \texttt{pairs.cpp} 的 \texttt{ans}
改成 \texttt{int}，輸入 $n = 10^5$ 個 $1$、$k = 2$，正確答案是
$\binom{10^5}{2} = 4999950000$。 在本書測試的環境下程式印出
\texttt{704982704}，沒有任何警告。 這是有號整數溢位，屬於未定義行為；加
\texttt{-fsanitize=undefined} 編譯，執行時會指出是哪一行：
\begin{code}{text}{用 -fsanitize=undefined 編譯後執行}
pairs_int.cpp:18:17: runtime error: signed integer overflow: 2147450880 + 65536 cannot be represented in type 'int'
\end{code}
寫程式之前先估答案最大可能多大：「幾對」最多是 $n^2 / 2$，總和最多是 $n$
乘上最大值。 超過 $2 \times 10^9$ 就用 \texttt{long long}。
\item \textbf{只用範例測資測試}：範例通常很小，$O(n^2)$ 的程式和溢位的程
式都會通過。 自己用最大的範圍產生一組資料跑跑看，時間和答案都要檢查。
\item \textbf{評分系統和自己的電腦結果不同}：未定義行為（越界、沒有初始化
的變數、溢位）在自己的電腦上「剛好對」，到評分系統上就錯了。 交出去之前用
\texttt{-fsanitize=address,undefined} 跑一次範例測資。
\end{itemize}

\section{遞迴與深度優先搜尋}
\label{sec:dfs}

第~8~章用 \texttt{queue} 做了廣度優先搜尋（BFS）。 另一種走遍整張圖的方法
是\textbf{深度優先搜尋}（depth-first search，DFS）：從一個點出發，能往前走
就一直往前走，走不下去才退回上一步換個方向。 用遞迴寫最短，「往前走」就是
呼叫自己，「退回」就是函式回傳。

典型的題目是數地圖上有幾塊陸地。 地圖是 $r$ 列 $c$ 行的格子，\texttt{\#}
是陸地，\texttt{.} 是水，上下左右相連的陸地算同一塊。 輸出有幾塊陸地，以
及最大的一塊有幾格：

\begin{example}[label=ex:islands]{cpp}{(islands.cpp) 用遞迴數陸地}
#include <algorithm>
#include <iostream>
#include <string>
#include <vector>
using namespace std;

// 把 (i, j) 所在的那塊陸地全部改成水，回傳它有幾格
int sink(vector<string> &g, int i, int j) {
    int r = g.size(), c = g[0].size();
    if (i < 0 || i >= r || j < 0 || j >= c || g[i][j] != '#') {
        return 0;
    }
    g[i][j] = '.';                         // 先標記，再往四個方向走
    return 1 + sink(g, i - 1, j) + sink(g, i + 1, j) +
           sink(g, i, j - 1) + sink(g, i, j + 1);
}

int main() {
    int r, c;
    cin >> r >> c;
    vector<string> g(r);
    for (auto &row : g) {
        cin >> row;
    }
    int count = 0, biggest = 0;
    for (int i = 0; i < r; i++) {
        for (int j = 0; j < c; j++) {
            if (g[i][j] == '#') {
                count++;
                biggest = max(biggest, sink(g, i, j));
            }
        }
    }
    cout << count << ' ' << biggest << '\n';
    return 0;
}
\end{example}

\iosample{
4 6 \\
\#\#..\#. \\
\#..\#\#. \\
..\#... \\
\#\#..\#\# \\
}{
5 3 \\
}

\begin{itemize}
\item \texttt{sink} 的第一個 \texttt{if} 是遞迴的終止條件：出界、是水、或
已經走過（已經改成水），就回傳 $0$。 邊界檢查一定要在 \texttt{g[i][j]} 之
前，\texttt{||} 由左到右求值，出界時不會去讀 \texttt{g[i][j]}。
\item 進到一格陸地，先把它改成水，再往四個方向遞迴。 改成水同時當作「走過了」
的記號，這塊陸地之後就不會再被數到，也不會繞回來。
\item \texttt{main} 掃過每一格，遇到還沒被淹掉的陸地就是新的一塊，呼叫一次
\texttt{sink} 把整塊淹掉。 每一格最多被處理一次，總時間是 $O(rc)$。
\item \texttt{g} 用參考傳進去。 所有遞迴呼叫改的都是同一張地圖，這正是這
個演算法需要的。
\end{itemize}

寫遞迴時，先寫終止條件，再寫「做一點事，把剩下的交給自己」。 數陸地、走迷
宮、枚舉組合、樹的走訪，都是這個形狀。

\subsection{常見錯誤}

\begin{itemize}
\item \textbf{沒有標記走過的格子}：拿掉 \texttt{g[i][j] = '.';}，兩格相鄰
的陸地會互相呼叫對方，永遠不會結束。 每次呼叫都佔用一些 stack 空間，用完
就當掉，在本書測試的環境下是 \texttt{Segmentation fault}。 ASan 會直接說
是 \texttt{stack-overflow}：
\begin{code}{text}{用 ASan 編譯後執行（節錄）}
==508446==ERROR: AddressSanitizer: stack-overflow on address 0x7ffc9fb2fff8
    #1 0x64bbc8bc44cb in sink(...) nomark.cpp:9
\end{code}
\item \textbf{地圖用傳值}：\texttt{int sink(vector<string> g, int i, int j)}
編譯沒有警告。 每一層遞迴都複製一整張地圖，標記只改到自己的複本，回到上一
層就不見了。 範例的輸入在本書測試的環境下輸出 \texttt{11 3}，陸地被重複數
了；地圖大一點時還會慢到跑不完。 照第~2~章的原則，要修改的參數用參考。
\item \textbf{遞迴太深}：$1000 \times 1000$ 全是陸地的地圖，遞迴最深可能有
$10^6$ 層。 本書測試的環境下 stack 是 $8$ MB，範例程式直接
\texttt{Segmentation fault}，加 \texttt{-O2} 也一樣。 有些評分系統會把
stack 設得比較大，但不能指望。 格子很多時改用第~8~章的 BFS，把待處理的格子
放在 \texttt{queue} 裡，它在 heap 上，不受 stack 大小限制（本章練習第~2~
題）。
\end{itemize}

\section{動態規劃}
\label{sec:dp}

\textbf{動態規劃}（dynamic programming，DP）把問題拆成較小的子問題，把每個
子問題的答案存在表格裡，大問題直接查表組合出來，不必重算。 C 的話表格是陣
列，大小要事先定好；C++ 用 \texttt{vector}，大小可以依輸入決定，越界時還能
用 \texttt{at()} 檢查。

這一節的例子是\textbf{最長遞增子序列}（longest increasing subsequence，
LIS）：從數列裡挑出一些數字，保持原來的先後，而且嚴格遞增，最長可以挑幾
個？ \texttt{3 1 4 1 5 9 2 6} 的答案是 $4$，例如 \texttt{1 4 5 9}。 下面
用兩種方法算：

\begin{example}[label=ex:lis]{cpp}{(lis.cpp) 最長遞增子序列的兩種算法}
#include <algorithm>
#include <iostream>
#include <vector>
using namespace std;

int lisSlow(const vector<int> &a) {        // O(n^2)
    int n = a.size(), best = 0;
    vector<int> dp(n, 1);   // dp[i]：以 a[i] 結尾的最長長度
    for (int i = 0; i < n; i++) {
        for (int j = 0; j < i; j++) {
            if (a[j] < a[i]) {
                dp[i] = max(dp[i], dp[j] + 1);
            }
        }
        best = max(best, dp[i]);
    }
    return best;
}

int lisFast(const vector<int> &a) {        // O(n log n)
    vector<int> tail;   // tail[k]：長度 k + 1 的最小結尾
    for (int x : a) {
        auto it = lower_bound(tail.begin(), tail.end(), x);
        if (it == tail.end()) {
            tail.push_back(x);   // 比所有結尾都大：最長的再加一
        } else {
            *it = x;             // 讓同樣長度的結尾變小
        }
    }
    return tail.size();
}

int main() {
    int n;
    cin >> n;
    vector<int> a(n);
    for (int &x : a) {
        cin >> x;
    }
    cout << lisSlow(a) << ' ' << lisFast(a) << '\n';
    return 0;
}
\end{example}

\iosample{
8 \\
3 1 4 1 5 9 2 6 \\
}{
4 4 \\
}

\texttt{lisSlow} 是標準的 DP 寫法，設計時回答三個問題：

\begin{itemize}
\item \textbf{狀態}：\texttt{dp[i]} 是「最後一個數字是 \texttt{a[i]}」的遞增
子序列最長有多長。
\item \textbf{轉移}：\texttt{a[i]} 可以接在任何一個 \texttt{a[j] < a[i]}
（$j < i$）的子序列後面，所以 \texttt{dp[i]} 是這些 \texttt{dp[j] + 1} 的最
大值；一個都沒有的話，\texttt{a[i]} 自己就是長度 $1$。
\item \textbf{順序}：算 \texttt{dp[i]} 時要用到所有 $j < i$ 的
\texttt{dp[j]}，所以 $i$ 由小到大算。
\end{itemize}

答案是所有 \texttt{dp[i]} 裡最大的那個，不一定是最後一個。 雙層迴圈，時間
是 $O(n^2)$，照第~\ref{sec:contest-basics}~節的表，$n$ 到 $5000$ 還可以。

$n = 2 \times 10^5$ 時就要用 \texttt{lisFast}。 它換了一個狀態：
\texttt{tail[k]} 是「長度 $k + 1$ 的遞增子序列，結尾最小可以是多少」。 結
尾越小，後面越容易接上新的數字，所以每種長度只要記最小的結尾。
\texttt{tail} 一定是遞增的，可以用第~9~章的 \texttt{lower\_bound} 找到第一
個 $\ge x$ 的結尾：找不到，表示 $x$ 比所有結尾都大，可以接在最長的後面；找
到了，就用 $x$ 取代它，同樣的長度有了更小的結尾。 範例的 \texttt{tail} 依序
變成 \texttt{3}、\texttt{1}、\texttt{1 4}、\texttt{1 4}、\texttt{1 4 5}、
\texttt{1 4 5 9}、\texttt{1 2 5 9}、\texttt{1 2 5 6}，長度是 $4$。 每個數字一
次二分搜尋，總時間 $O(n \log n)$。

\subsection{常見錯誤}

\begin{itemize}
\item \textbf{以為 \texttt{tail} 就是那個子序列}：範例最後的
\texttt{tail} 是 \texttt{1 2 5 6}，但 $2$ 在原數列裡排在 $5$ 後面，它不是一
個合法的子序列。 \texttt{tail} 只保證\textbf{長度}正確。 要印出子序列本身，
得另外記下每個數字是接在誰後面。
\item \textbf{「遞增」和「不下降」用錯函式}：題目允許相等（\texttt{1 3 3 4}
也算）時，要用 \texttt{upper\_bound}，讓相等的數字可以接在後面。 用錯的話，
\texttt{2 2 2} 會算出 $1$ 而不是 $3$（本章練習第~4~題）。
\item \textbf{忘了初始值}：\texttt{vector<int> dp(n);} 的元素都是 $0$，少
了「自己就是長度 $1$」，每個 \texttt{dp[i]} 都少 $1$。 DP 的表格要想清楚
最小的子問題是什麼，先填好。
\end{itemize}

\section{自己實作資料結構：鏈結串列}
\label{sec:linked-list}

第~8~章的容器都是現成的。 資料結構課則要你自己寫：鏈結串列、堆疊、二元搜
尋樹，了解它們裡面怎麼運作。 這一節用鏈結串列示範怎麼把 C 的寫法搬到
C++：節點還是結構加指標，但包進一個類別，用建構子、解構子和模板管理。

\begin{crecap}[鏈結串列]
鏈結串列的每個節點存一個值和一個指向下一個節點的指標，最後一個節點的
\texttt{next} 是 \texttt{NULL}：

\begin{code}{c}{C 的鏈結串列}
struct Node {
    int value;
    struct Node *next;
};

struct Node *pushFront(struct Node *head, int x) {
    struct Node *node = malloc(sizeof(struct Node));
    node->value = x;
    node->next = head;
    return node;               /* 新的節點變成開頭 */
}

void freeList(struct Node *head) {
    while (head != NULL) {
        struct Node *next = head->next;
        free(head);
        head = next;
    }
}
\end{code}

在開頭加入是 $O(1)$，不用像陣列那樣搬動元素。 代價是不能用下標直接跳到第
$k$ 個，只能從 \texttt{head} 沿著 \texttt{next} 一個一個走。 用完一定要
呼叫 \texttt{freeList}，忘了就是記憶體洩漏；只寫 \texttt{int} 的版本，要存
字串就得整份重寫。
\end{crecap}

C++ 版本把 \texttt{head} 和節點都藏在類別裡，釋放交給解構子：

\begin{example}[label=ex:linked-list]{cpp}{(linked\_list.cpp) 鏈結串列的類別模板}
#include <iostream>
#include <string>
using namespace std;

template <typename T>
class LinkedList {
public:
    LinkedList() = default;
    LinkedList(const LinkedList &) = delete;      // 不允許複製
    LinkedList &operator=(const LinkedList &) = delete;
    ~LinkedList() {
        while (head != nullptr) {
            Node *next = head->next;   // 先記下下一個
            delete head;
            head = next;
        }
    }

    void pushFront(const T &x) {
        head = new Node{x, head};
        size_++;
    }
    void reverse() {
        Node *prev = nullptr;
        while (head != nullptr) {
            Node *next = head->next;
            head->next = prev;         // 把箭頭轉向
            prev = head;
            head = next;
        }
        head = prev;
    }
    void print() const {
        for (Node *p = head; p != nullptr; p = p->next) {
            cout << p->value << ' ';
        }
        cout << '\n';
    }
    int size() const { return size_; }

private:
    struct Node {
        T value;
        Node *next;
    };
    Node *head = nullptr;
    int size_ = 0;
};

int main() {
    LinkedList<int> a;
    for (int x : {3, 1, 4, 1, 5}) {
        a.pushFront(x);
    }
    a.print();
    a.reverse();
    a.print();

    LinkedList<string> words;
    words.pushFront("world");
    words.pushFront("hello");
    words.print();
    cout << a.size() << ' ' << words.size() << '\n';
    return 0;
}
\end{example}

\iosample{}{
5 1 4 1 3 \\
3 1 4 1 5 \\
hello world \\
5 2 \\
}

\begin{itemize}
\item \texttt{Node} 宣告在類別的 \texttt{private} 裡，外面的程式碰不到節點和
指標，只能用 \texttt{pushFront}、\texttt{reverse} 這些成員函式。 第~5~章
說的不變量（例如 \texttt{size\_} 一定等於節點個數）就由類別自己維持。
\item \texttt{new Node\{x, head\}} 配置一個節點，同時用大括號依序給兩個成員
初值，取代 C 的 \texttt{malloc} 再逐一指定。
\item 解構子做的就是 C 的 \texttt{freeList}，\texttt{a} 和 \texttt{words}
離開 \texttt{main} 時自動執行（第~4~章的 RAII）。 用 ASan 跑，沒有洩漏。
\item 類別自己管理 \texttt{new} 來的記憶體，照第~6~章的規則，複製建構子和
複製指定要嘛自己寫成深複製，要嘛 \texttt{= delete}。 這裡選擇禁止複製。
\item \texttt{reverse} 是資料結構課的經典題：不配置新節點，一個一個把
\texttt{next} 轉向。 \texttt{prev} 是已經轉好的部分的開頭，\texttt{head}
是還沒處理的部分的開頭。
\item 改成模板之後，同一份程式可以存 \texttt{int} 也可以存 \texttt{string}。
\end{itemize}

標準函式庫其實有雙向鏈結串列 \texttt{list}（第~9~章提過），解題時直接用
它或 \texttt{vector} 就好。 自己寫一次的目的是看懂指標怎麼串起資料；之後
的二元搜尋樹也是同樣的寫法，只是每個節點有 \texttt{left}、\texttt{right}
兩個指標。

\subsection{常見錯誤}

\begin{itemize}
\item \textbf{刪掉節點之後再讀它的 \texttt{next}}：解構子寫成
\begin{code}{cpp}{先 delete 再讀 next（錯誤示範）}
for (Node *p = head; p != nullptr; p = p->next) {
    delete p;        // 迴圈的 p = p->next 讀的是已經刪掉的節點
}
\end{code}
\texttt{delete p} 之後，\texttt{p->next} 是已經釋放的記憶體。 在本書測試
的環境下，程式在結束時當掉（Segmentation fault）；ASan 報
\texttt{heap-use-after-free}，位置在解構子裡。 範例先用 \texttt{next} 記下
下一個節點，再刪掉目前這個，順序不能反過來。
\item \textbf{只刪了 \texttt{head}}：解構子只寫 \texttt{delete head;}，其他
節點沒有指標記得了。 程式照常結束，看不出問題，ASan 在結束時報告洩漏：
\begin{code}{text}{ASan 的洩漏報告（節錄）}
==552332==ERROR: LeakSanitizer: detected memory leaks
SUMMARY: AddressSanitizer: 104 byte(s) leaked in 5 allocation(s).
\end{code}
$5$ 個是 \texttt{a} 剩下的 $4$ 個節點和 \texttt{words} 剩下的 $1$ 個。
\item \textbf{複製串列}：\texttt{LinkedList<int> b = a;} 報
\texttt{error: use of deleted function 'LinkedList<T>::LinkedList(const
LinkedList<T>\&) [with T = int]'}。 這正是 \texttt{= delete} 的作用：沒有它
的話，編譯器產生的複製只複製 \texttt{head} 指標，兩個串列共用同一串節點，
結束時各刪一次，就是第~5~章的重複釋放。
\item \textbf{忘了處理空串列}：對空的串列取 \texttt{head->value}，就是透過
\texttt{nullptr} 讀取，程式當掉。 每個會用到 \texttt{head} 的成員函式，都
要想一下 \texttt{head} 是 \texttt{nullptr} 時怎麼辦。
\end{itemize}

\section{讀完這本書之後}
\label{sec:next}

本書刻意寫短，只講從 C 走到 C++ 一定要補的部分。 接下來可以往三個方向走。

\textbf{解題}：APCS 實作題和各個解題網站常見的主題，像排序與二分搜尋、遞迴與
枚舉、貪心、動態規劃、BFS 與 DFS、樹與圖，本書已經把需要的語言工具都介紹
過了，剩下的是演算法本身和大量練習。 練習時照第~\ref{sec:contest-basics}~節
的習慣：先看資料範圍決定複雜度，估答案會不會超過 \texttt{int}，交出去之前
自己補邊界測資。

\textbf{資料結構與演算法課}：課本上的鏈結串列、堆疊、佇列、堆積、雜湊表、二
元搜尋樹、圖，就是第~8~章那些容器的內部。 自己實作時用本章的寫法：節點是結
構加指標，整個資料結構包成類別模板，用解構子釋放記憶體，禁止或自己寫複製。
寫完用 ASan 跑一次，確定沒有洩漏和懸置指標。

\textbf{C++ 本身}：本書沒有談的主題還有不少，例如例外的處理（\texttt{try} 和
\texttt{catch}）、檔案輸入輸出（\texttt{<fstream>}）、更完整的移動語意，以
及 C++20 之後的新功能。 這些在寫大型程式時才比較常用，需要時再查
cppreference.com 這類參考資料就好。 遇到看不懂的錯誤訊息，照第~8~章的方法，
先找 \texttt{error:} 那一行和 \texttt{required from here}。

\section{本章練習}
\label{sec:ch10-exercises}

\begin{drills}[label=drills:ch10]
\item 時間限制一秒，估計下面每種做法來不來得及，說明理由。
\begin{enumerate}
\item $n \le 2 \times 10^5$，對每個 $i$ 用一個迴圈檢查所有 $j > i$。
\item $n, q \le 2 \times 10^5$，先排序，每次查詢用 \texttt{lower\_bound}。
\item $n \le 10$，用 \texttt{next\_permutation} 枚舉所有排列，每個排列花
$O(n)$ 計算。
\item $n \le 5000$，$O(n^2)$ 的動態規劃。
\item $n \le 10^5$，每個 $a_i \le 10^6$，用 \texttt{int} 存所有 $a_i$ 的總和。
\end{enumerate}
\item 把範例~\ref{ex:islands} 的 \texttt{sink} 改成用 \texttt{queue} 做 BFS，
不用遞迴，\texttt{main} 不變。 範例的輸入要輸出 \texttt{5 3}；$1000 \times
1000$ 全是陸地的地圖要輸出 \texttt{1 1000000}，不能當掉。
\item （0/1 背包）有 $n$ 件物品，第 $i$ 件的重量是 $w_i$、價值是 $v_i$，
背包最多能裝重量 $W$。 每件物品最多拿一次，求能裝的最大總價值。
$n \le 100$，$W \le 10^5$。 輸入
\begin{code}{text}{練習 3 的輸入}
4 7
3 4
4 5
2 3
5 6
\end{code}
（第一行是 $n$ 和 $W$，接著每行是一件物品的重量和價值）應該輸出 \texttt{9}。
提示：\texttt{dp[w]} 是重量不超過 \texttt{w} 時的最大價值。
\item 讀入 $n$ 和 $n$ 個整數，輸出最長嚴格遞增子序列和最長不下降子序列（允許
相等）的長度，要求 $O(n \log n)$。 輸入 \texttt{7} 和
\texttt{1 3 3 2 2 4 1} 時應輸出 \texttt{3 4}。
\item 替範例~\ref{ex:linked-list} 的 \texttt{LinkedList} 加上
\texttt{pushBack(x)}（加在尾端，要 $O(1)$）和 \texttt{popFront()}（刪掉第一
個節點，串列不是空的）。 提示：多記一個指向最後一個節點的 \texttt{tail}。
下面的程式應該輸出 \texttt{0 1 2}、\texttt{7 8}、\texttt{2} 三行，用 ASan
跑沒有錯誤：
\begin{code}{cpp}{練習 5 的 main}
LinkedList<int> q;
q.pushBack(1);
q.pushBack(2);
q.pushFront(0);
q.print();
q.popFront();
q.popFront();
q.popFront();
q.pushBack(7);
q.pushBack(8);
q.print();
cout << q.size() << '\n';
\end{code}
\end{drills}

\addcontentsline{toc}{section}{練習解答}
\section*{練習解答}

\begin{enumerate}
\item (a) 來不及。 大約 $n^2 / 2 = 2 \times 10^{10}$ 次，是 $10^8$ 的兩百倍。
(b) 來得及。 排序和查詢都是 $(2 \times 10^5) \times \log_2(2 \times 10^5)
\approx 3.6 \times 10^6$ 次左右。 (c) 來得及。 $10! \times 10 \approx 3.6
\times 10^7$。 (d) 來得及。 $5000^2 = 2.5 \times 10^7$。 (e) 時間沒問題，但
總和最大是 $10^5 \times 10^6 = 10^{11}$，超過 \texttt{int} 的上限（大約
$2.1 \times 10^9$），要用 \texttt{long long}。
\item 一種寫法：
\begin{code}{cpp}{練習 2 參考解答}
int sink(vector<string> &g, int si, int sj) {
    int r = g.size(), c = g[0].size();
    int di[] = {-1, 1, 0, 0};
    int dj[] = {0, 0, -1, 1};
    queue<pair<int, int>> q;
    g[si][sj] = '.';                   // 放進 queue 時就標記
    q.push({si, sj});
    int cells = 0;
    while (!q.empty()) {
        auto [i, j] = q.front();
        q.pop();
        cells++;
        for (int d = 0; d < 4; d++) {
            int ni = i + di[d], nj = j + dj[d];
            if (ni >= 0 && ni < r && nj >= 0 && nj < c &&
                g[ni][nj] == '#') {
                g[ni][nj] = '.';
                q.push({ni, nj});
            }
        }
    }
    return cells;
}
\end{code}
要另外 include \texttt{<queue>} 和 \texttt{<utility>}。 跟第~8~章迷宮的
BFS 是同一個迴圈，差別是這裡不需要距離，只要數有幾格。 格子在放進
\texttt{queue} 的時候就標記成水，同一格才不會被放進去兩次。 \texttt{queue}
的元素放在 heap 上，$10^6$ 格也不會用完 stack。
\item 一種寫法：
\begin{code}{cpp}{練習 3 參考解答}
#include <algorithm>
#include <iostream>
#include <vector>
using namespace std;

int main() {
    int n, cap;
    cin >> n >> cap;
    vector<long long> dp(cap + 1, 0);
    for (int i = 0; i < n; i++) {
        int w;
        long long v;
        cin >> w >> v;
        for (int j = cap; j >= w; j--) {   // 由大到小
            dp[j] = max(dp[j], dp[j - w] + v);
        }
    }
    cout << dp[cap] << '\n';
    return 0;
}
\end{code}
處理第 $i$ 件物品時，\texttt{dp[j]} 有兩種選擇：不拿，維持原值；拿，就是
\texttt{dp[j - w] + v}。 \texttt{j} 一定要由大到小：這樣讀到的
\texttt{dp[j - w]} 還是「沒考慮第 $i$ 件」時的值，每件只會被拿一次。 寫成
由小到大，\texttt{dp[j - w]} 可能已經拿過這件物品，同一件會被拿好幾次，範
例會輸出 \texttt{10}（重量 $2$ 的物品拿兩次再加重量 $3$ 的）。 時間是
$O(nW) = 10^7$。
\item 一種寫法：
\begin{code}{cpp}{練習 4 參考解答（讀入 a 之後）}
vector<int> strict, nondec;
for (int x : a) {
    auto it = lower_bound(strict.begin(), strict.end(), x);
    if (it == strict.end()) {
        strict.push_back(x);
    } else {
        *it = x;
    }
    auto jt = upper_bound(nondec.begin(), nondec.end(), x);
    if (jt == nondec.end()) {
        nondec.push_back(x);
    } else {
        *jt = x;
    }
}
cout << strict.size() << ' ' << nondec.size() << '\n';
\end{code}
不下降時，等於結尾的數字也可以接上去，所以要找第一個\textbf{大於} $x$ 的結
尾來取代，用 \texttt{upper\_bound}。 範例的嚴格遞增是 \texttt{1 3 4}（或
\texttt{1 2 4}），不下降是 \texttt{1 3 3 4}（或 \texttt{1 2 2 4}）。
\item 一種寫法，新增的成員：
\begin{code}{cpp}{練習 5 參考解答（新增的部分）}
// private 裡多一個成員：Node *tail = nullptr;
void pushBack(const T &x) {
    Node *node = new Node{x, nullptr};
    if (tail == nullptr) {
        head = tail = node;            // 原本是空的
    } else {
        tail->next = node;
        tail = node;
    }
    size_++;
}
void popFront() {
    Node *old = head;
    head = head->next;
    if (head == nullptr) {
        tail = nullptr;                // 刪完變空的
    }
    delete old;
    size_--;
}
\end{code}
多了 \texttt{tail} 之後，原本的成員函式也要維持「\texttt{tail} 指著最後一個節
點」這個不變量：\texttt{pushFront} 在串列原本是空的時候要把 \texttt{tail}
設成新節點；\texttt{reverse} 之後，原本的第一個節點變成最後一個，要在迴圈開
始前記下 \texttt{head}，最後設給 \texttt{tail}。 最容易漏掉的是兩個邊界：空
串列加入第一個節點時 \texttt{head} 和 \texttt{tail} 都要設；刪掉最後一個節
點時 \texttt{tail} 要改回 \texttt{nullptr}，不然它會指著已經刪掉的節點，下
一次 \texttt{pushBack} 就是寫入已釋放的記憶體。 練習的 \texttt{main} 正好
經過這兩種情況：三次 \texttt{popFront} 把串列刪空，再 \texttt{pushBack}。
\end{enumerate}


\appendix
% 附錄的章標題寫成「附錄 A」，不是「第 A 章」。
\titleformat{\chapter}{\normalfont\huge\bfseries}{附錄~\thechapter}{1em}{}

\chapter{補充：例外、\texttt{std::array} 與計時}
\label{app:supplement}

這篇附錄補三樣正文沒有講、但讀完之後很快會碰到的東西：例外的接法、固定大小
的 \texttt{std::array}，以及用 \texttt{<chrono>} 量程式跑多久。 最後列出本
書沒有涵蓋的主題，方便之後找進階的書。

\section{例外}
\label{sec:exceptions}

\begin{crecap}[用回傳值回報錯誤]
C 的函式出錯時，通常用回傳值告訴呼叫者：\texttt{fopen} 回傳 \texttt{NULL}，
\texttt{malloc} 回傳 \texttt{NULL}，\texttt{scanf} 回傳成功讀到的個數。 呼
叫者要自己寫 \texttt{if (fp == NULL)} 這樣的檢查，忘了檢查，程式就帶著錯誤
的值繼續跑。 而且有些函式沒有多餘的回傳值可以表
示錯誤：\texttt{atoi("abc")} 回傳 $0$，跟 \texttt{atoi("0")} 分不出來。
\end{crecap}

C++ 多了一種方式：\textbf{例外}（exception）。 函式遇到沒辦法處理的情況時，
用 \texttt{throw} 丟出一個例外物件，函式立刻結束，一路往外層傳，直到某個
\texttt{try} 區塊後面有對應型別的 \texttt{catch} 接住它。 第~3~章的
\texttt{stoi}、\texttt{at()}、\texttt{substr} 出錯時就是丟出例外。

\begin{example}[label=ex:exceptions]{cpp}{(exceptions.cpp) try、catch 與 throw}
#include <iostream>
#include <stdexcept>
#include <string>
#include <vector>
using namespace std;

int average(const vector<int> &v) {
    if (v.empty()) {
        throw invalid_argument("average: 空的 vector");
    }
    long long sum = 0;
    for (int x : v) {
        sum += x;
    }
    return sum / (long long)v.size();
}

int main() {
    vector<string> words = {"42", "abc", "99999999999", "7"};
    vector<int> nums;
    for (const string &w : words) {
        try {
            nums.push_back(stoi(w));
            cout << w << " 轉換成功" << '\n';
        } catch (const invalid_argument &e) {
            cout << w << " 不是數字（" << e.what() << "）" << '\n';
        } catch (const out_of_range &e) {
            cout << w << " 太大（" << e.what() << "）" << '\n';
        }
    }

    try {
        cout << average(nums) << '\n';
        cout << average({}) << '\n';
        cout << "這一行不會執行" << '\n';
    } catch (const exception &e) {
        cout << "錯誤：" << e.what() << '\n';
    }
    return 0;
}
\end{example}

\iosample{}{
42 轉換成功 \\
abc 不是數字（stoi） \\
99999999999 太大（stoi） \\
7 轉換成功 \\
24 \\
錯誤：average: 空的 vector \\
}

\begin{itemize}
\item \texttt{try \{ ... \}} 裡的程式照常執行。 其中某一行丟出例外時，
\texttt{try} 裡剩下的程式都不執行，直接跳到型別相符的 \texttt{catch}。
\texttt{"abc"} 那一輪，\texttt{push\_back} 和印出「轉換成功」都被跳過了。
\item 一個 \texttt{try} 後面可以接好幾個 \texttt{catch}，由上往下比對型別。
\texttt{stoi} 遇到不是數字的字串丟 \texttt{invalid\_argument}，數字太大丟
\texttt{out\_of\_range}。 接住之後程式從整個 \texttt{try}-\texttt{catch} 的後
面繼續，迴圈照樣處理下一個字串。
\item 標準函式庫的例外都繼承自 \texttt{exception}（在 \texttt{<stdexcept>}
裡），\texttt{what()} 回傳說明文字。 \texttt{catch (const exception \&e)}
可以接住全部，這是第~7~章的多型：用基底類別的參考接住衍生類別的物件。
\item \texttt{throw invalid\_argument("...")} 丟出自己的例外，括號裡的字串
就是 \texttt{what()} 的內容。 \texttt{average} 不能計算時就丟出例外，不必
回傳特殊值，呼叫者也沒辦法忽略它。 第二個 \texttt{try} 裡，
\texttt{average(\{\})} 丟出例外之後，下一行沒有執行。
\end{itemize}

\subsection{沒接住會怎樣}

沒有任何 \texttt{catch} 接住的例外，會讓程式呼叫 \texttt{std::terminate}
中止。 下面的程式對三個元素的 \texttt{vector} 呼叫 \texttt{v.at(3)}：

\begin{code}{cpp}{(uncaught.cpp) 沒有接住的例外}
vector<int> v = {1, 2, 3};
cout << "開始" << '\n';
cout << v.at(3) << '\n';
cout << "結束" << '\n';
\end{code}

\begin{code}{text}{在終端機執行}
開始
terminate called after throwing an instance of 'std::out_of_range'
  what():  vector::_M_range_check: __n (which is 3) >= this->size() (which is 3)
Aborted (core dumped)
\end{code}

前三行是程式印的，最後一行是 shell 印的，表示程式被中止，結束狀態是
$134$。 「結束」沒有印出來。 在本書測試的環境下，把輸出導向檔案
（\texttt{./uncaught > out.txt}）時，連「開始」都沒有寫進檔案：\texttt{cout}
的內容還在緩衝區裡，程式就被中止了。 評分系統看到的是非正常結束（通常標成
Runtime Error），輸出是什麼都不重要了。

\subsection{解題時通常不用例外}

例外適合「發生錯誤的地方不知道怎麼處理，要交給外層決定」的情況，例如讀設
定檔的函式不知道檔案壞了該重試還是結束，就丟給呼叫者。 解題程式通常用不到：
題目保證輸入格式正確，\texttt{stoi} 不會遇到 \texttt{"abc"}；程式出錯就是答案
錯，接住例外也救不回來。 而且例外只處理「有人丟出」的錯誤，\texttt{v[10]}
越界、整數溢位、懸置指標都是未定義行為，不會丟出例外，\texttt{try} 接不到。

所以解題時看到 \texttt{terminate called after throwing ...}，要做的是看
\texttt{what()} 找出是哪個函式、哪個下標，修好程式。 寫大一點的程式（例如
資料結構課的作業要求對不合法的操作回報錯誤）時，再用 \texttt{throw} 和
\texttt{try}-\texttt{catch}。

\subsection{常見錯誤}

\begin{itemize}
\item \textbf{用傳值接例外}：\texttt{catch (exception e)} 會把例外物件複製
成一個 \texttt{exception}，衍生類別的部分被切掉，就是第~7~章的物件切割。
\texttt{-Wall} 會警告 \texttt{catching polymorphic type 'class
std::exception' by value}；接住 \texttt{stoi} 的例外後，\texttt{e.what()}
在本書測試的環境下印出 \texttt{std::exception}，原本的 \texttt{stoi} 不見了。
一律寫 \texttt{catch (const exception \&e)}。
\item \textbf{基底類別寫在前面}：接 \texttt{exception} 的 \texttt{catch}
寫在接 \texttt{invalid\_argument} 的前面，後者永遠輪不到，因為
\texttt{catch} 由上往下比對，前面的已經接住了。 \texttt{g++} 會警告
\texttt{exception of type 'std::invalid\_argument' will be caught by earlier
handler}。 具體的型別寫在前面，\texttt{exception} 放最後。
\end{itemize}

\section{\texttt{std::array}}
\label{sec:std-array}

\begin{crecap}[陣列不能整個指定]
C 的陣列有幾個不方便的地方：

\begin{code}{c}{C 陣列的限制}
int a[5] = {3, 1, 4, 1, 5}, b[5];
b = a;                  /* 編譯失敗 */
void f(int a[5]) {
    printf("%zu\n", sizeof(a));   /* 8，是指標的大小 */
}
\end{code}

\texttt{b = a} 報 \texttt{assignment to expression with array type}，要用
迴圈或 \texttt{memcpy} 複製。 陣列傳進函式時會變成指向第一個元素的指標，
參數上的 \texttt{[5]} 沒有作用，\texttt{sizeof(a)} 是指標的大小
（\texttt{gcc -Wall} 會警告），函式不知道陣列有多長，只能另外傳 \texttt{n}。
\texttt{==} 比的也是兩個位址，不是內容。
\end{crecap}

\texttt{std::array<T, N>}（在 \texttt{<array>} 裡）是 $N$ 個 \texttt{T} 的固
定大小陣列。 元素跟 C 陣列一樣直接放在物件裡，不另外配置記憶體，但它是一個
真正的型別，上面這些限制都沒有了：

\begin{example}[label=ex:array]{cpp}{(array\_demo.cpp) 複製、比較、傳給函式、放進 vector}
#include <algorithm>
#include <array>
#include <iostream>
#include <vector>
using namespace std;

void print(const array<int, 5> &a) {
    for (int x : a) {
        cout << x << ' ';
    }
    cout << '(' << a.size() << " 個)" << '\n';
}

int main() {
    array<int, 5> a = {3, 1, 4, 1, 5};
    array<int, 5> b = a;               // 整個複製
    b[0] = 9;
    sort(b.begin(), b.end());
    print(a);
    print(b);
    cout << (a == b) << ' ' << (a < b) << '\n';

    // 每條邊是 {重量, 起點, 終點}，依重量排序
    vector<array<int, 3>> edges = {{7, 1, 2}, {3, 2, 3}, {5, 1, 3}};
    sort(edges.begin(), edges.end());
    for (const auto &[w, u, v] : edges) {
        cout << u << '-' << v << ' ' << w << '\n';
    }
    return 0;
}
\end{example}

\iosample{}{
3 1 4 1 5 (5 個) \\
1 1 4 5 9 (5 個) \\
0 0 \\
2-3 3 \\
1-3 5 \\
1-2 7 \\
}

\begin{itemize}
\item \texttt{b = a} 複製全部元素，之後改 \texttt{b} 不影響 \texttt{a}。
\texttt{==}、\texttt{<} 比的是內容，\texttt{<} 是字典順序，跟 \texttt{vector}
一樣。
\item 傳給函式時大小是型別的一部分，\texttt{print} 用 \texttt{a.size()} 就知
道有幾個元素。 照第~2~章的原則用 \texttt{const} 參考傳。
\item 有 \texttt{begin()}、\texttt{end()}，第~9~章的演算法都能用。
\texttt{[]} 跟 \texttt{vector} 一樣不檢查越界，要檢查就用 \texttt{at()}。
\item \texttt{array<int, 5> z\{\};} 的大括號讓所有元素都是 $0$。 沒有大括號
的區域變數 \texttt{array<int, 5> z;} 跟 C 的區域陣列一樣不會初始化，裡面是
垃圾值。
\item \texttt{vector<array<int, 3>>} 常用來存「好幾個欄位都是整數」的資料，
例如圖的邊。 比 \texttt{pair} 多一個欄位，又不用為它寫一個結構，排序時一樣
是先比第一個欄位；用結構化綁定拆開。
\end{itemize}

什麼時候用哪一個：大小在寫程式時就固定（例如四個方向、三個欄位、$26$ 個英文
字母的計數），用 \texttt{array}；大小要等執行時讀了輸入才知道，或者會增減，
用 \texttt{vector}。 C 陣列仍然能用，但上面的限制都還在。

\subsection{常見錯誤}

\begin{itemize}
\item \textbf{大小用變數}：讀入 \texttt{n} 之後寫 \texttt{array<int, n> a;}，
報 \texttt{error: the value of 'n' is not usable in a constant
expression}。 \texttt{array} 的大小是型別的一部分，必須在編譯時就確定。
要依輸入決定大小，用 \texttt{vector<int> a(n);}。
\item \textbf{大小不同就是不同型別}：\texttt{array<int, 5>} 和
\texttt{array<int, 6>} 是兩個型別，不能互相指定，也不能傳給同一個
\texttt{print}。 這是固定大小的代價。
\end{itemize}

\section{用 \texttt{<chrono>} 計時}
\label{sec:chrono}

第~10~章用資料範圍估計演算法夠不夠快。 估計之外，有時也想實際量一段程式花
了多久，例如比較兩種寫法。 \texttt{<chrono>} 的 \texttt{steady\_clock} 是一
個只會往前走的時鐘，在要量的程式前後各讀一次，相減就是經過的時間：

\begin{example}[label=ex:timing]{cpp}{(timing.cpp) 量 sort 花了多久}
#include <algorithm>
#include <chrono>
#include <iostream>
#include <vector>
using namespace std;

int main() {
    const int n = 1000000;
    vector<int> v(n);
    for (int i = 0; i < n; i++) {
        v[i] = (long long)i * 7919 % n;   // 打亂順序的 0 到 n - 1
    }

    auto start = chrono::steady_clock::now();
    sort(v.begin(), v.end());
    auto end = chrono::steady_clock::now();

    auto ms = chrono::duration_cast<chrono::milliseconds>(
        end - start);
    cout << "sort " << n << " 個：" << ms.count() << " ms" << '\n';
    cout << v[0] << ' ' << v[n - 1] << '\n';
    return 0;
}
\end{example}

\begin{code}{text}{本書測試的環境下，兩種編譯方式各執行一次}
$ g++ -std=c++17 -Wall -Wextra timing.cpp -o timing && ./timing
sort 1000000 個：775 ms
0 999999
$ g++ -std=c++17 -Wall -Wextra -O2 timing.cpp -o timing && ./timing
sort 1000000 個：113 ms
0 999999
\end{code}

\begin{itemize}
\item \texttt{chrono::steady\_clock::now()} 回傳目前的時間點，兩個時間點相減
得到一段時間。 \texttt{duration\_cast<chrono::milliseconds>} 把它換成毫秒，
\texttt{.count()} 取出數字。 型別很長，所以都寫 \texttt{auto}。
\item 毫秒數每次執行都不一樣。 在本書測試的環境下，不加 \texttt{-O2} 連跑三
次是 $775$、$940$、$1990$ 毫秒，要比較時多跑幾次。 第二行是排序的結果，用來
確認程式真的做完了，每次都一樣。
\item 有沒有 \texttt{-O2} 差了好幾倍。 評分系統幾乎都開最佳化編譯，自己量時
間也要加 \texttt{-O2}，結果才有參考價值。
\end{itemize}

\subsection{自己量的時間和評分系統的時間}

評分系統顯示的時間跟上面量到的數字通常對不上。 評分系統量的是\textbf{整支
程式}，包括讀輸入和印輸出，上面只量了中間一段；輸入有幾十萬個數字時，讀入本
身就可能佔掉不少時間，所以第~1~章的輸入加速很重要。 很多評分系統量的是 CPU
時間，\texttt{steady\_clock} 量的是實際經過的時間，電腦同時在跑別的程式時，
後者會變長。 機器和編譯選項不同，速度也不同，自己電腦上的秒數不能直接拿去
跟時間限制比。

所以自己量的時間只適合拿來比較「同一台電腦上，兩種寫法哪個快、快多少」。
程式能不能在時間限制內跑完，主要還是靠第~10~章的複雜度估計，計時只是確認。
要量整支程式，Linux 和 macOS 上可以用 \texttt{time ./timing < in.txt}。

\subsection{常見錯誤}

\begin{itemize}
\item \textbf{量的計算結果沒用到}：下面的迴圈算了 $10^9$ 次，但
\texttt{sum} 從來沒有印出來：
\begin{code}{cpp}{結果沒用到的迴圈（錯誤示範）}
auto start = chrono::steady_clock::now();
long long sum = 0;
for (long long i = 0; i < 1000000000; i++) {
    sum += i % 7;
}
auto end = chrono::steady_clock::now();
\end{code}
不加 \texttt{-O2} 時，在本書測試的環境下量到 \texttt{4525 ms}；加了
\texttt{-O2}，量到 \texttt{0 ms}。 編譯器發現 \texttt{sum} 沒有被用到，整個
迴圈直接刪掉了。 要量的計算，結果一定要用到，例如最後把 \texttt{sum} 印出
來。
\item \textbf{換成秒數，小數不見了}：
\texttt{duration\_cast<chrono::seconds>} 會無條件捨去，跑了 $0.57$ 秒的程式
印出 \texttt{0}。 用 \texttt{milliseconds}，或者寫
\texttt{chrono::duration<double>(end - start).count()}，得到以秒為單位的小
數（本書測試的環境下是 \texttt{0.570522}）。
\item \textbf{用 \texttt{system\_clock} 計時}：\texttt{system\_clock} 是「現
在幾點」的時鐘，系統校正時間時可能往回調，算出來的時間差可能是負的。 計時用
\texttt{steady\_clock}。
\end{itemize}

\section{本書沒有涵蓋的主題}
\label{sec:not-covered}

下面這些也是標準 C++ 的一部分，本書沒有講。 要學的話可以讀 Bjarne
Stroustrup 的《A Tour of C++》，它用不長的篇幅走過整個語言和標準函式庫；查單
一函式的細節用 cppreference.com。

\begin{itemize}
\item \textbf{並行程式}（\texttt{<thread>}、\texttt{<mutex>}、\texttt{<future>}）：
\texttt{thread} 讓程式同時跑好幾段程式碼，\texttt{mutex} 防止它們同時改同一
個變數，\texttt{async} 把一個函式丟到背景執行、之後再拿結果。
\item \textbf{正規表達式}（\texttt{<regex>}）：用一個樣式描述「長得像什麼的
字串」，例如電子郵件地址，再拿來搜尋或取代。
\item \textbf{concepts}（C++20）：替模板的型別參數寫出要求，例如「\texttt{T}
要能用 \texttt{<} 比較」，不符合時錯誤訊息直接說哪個要求沒做到，不會像第~8~
章那樣長。
\item \textbf{可變參數模板}（variadic template）：參數個數不固定的模板，
\texttt{template <typename... Args>}。 標準函式庫的 \texttt{make\_unique}、
\texttt{emplace\_back} 就是這樣寫的。
\item \textbf{type traits}（\texttt{<type\_traits>}）：在編譯時查詢型別的性
質，例如「\texttt{T} 是不是整數型別」，讓模板依型別選擇不同的寫法。
\item \textbf{\texttt{function} 與 \texttt{bind}}（\texttt{<functional>}）：
\texttt{function<int(int)>} 可以存任何「收一個 \texttt{int}、回傳
\texttt{int}」的東西，一般函式、lambda 都行，用來把函式存進變數或容器。
\texttt{bind} 先固定函式的部分參數，現在多半改用 lambda。
\end{itemize}


\end{document}
